% 高一71_中国中学_10月月考.tex
%   —— 高一上数学“2026-2027 年中国中学高一上 10 月月考”（讲评版）
% 试卷版：xelatex 高一71_中国中学_10月月考.tex
% 答案版：xelatex "\def\isanswer{1}% 高一71_中国中学_10月月考.tex
%   —— 高一上数学“2026-2027 年中国中学高一上 10 月月考”（讲评版）
% 试卷版：xelatex 高一71_中国中学_10月月考.tex
% 答案版：xelatex "\def\isanswer{1}% 高一71_中国中学_10月月考.tex
%   —— 高一上数学“2026-2027 年中国中学高一上 10 月月考”（讲评版）
% 试卷版：xelatex 高一71_中国中学_10月月考.tex
% 答案版：xelatex "\def\isanswer{1}% 高一71_中国中学_10月月考.tex
%   —— 高一上数学“2026-2027 年中国中学高一上 10 月月考”（讲评版）
% 试卷版：xelatex 高一71_中国中学_10月月考.tex
% 答案版：xelatex "\def\isanswer{1}\input{高一71_中国中学_10月月考.tex}"
% 紧凑版：xelatex "\def\iscompact{1}\input{高一71_中国中学_10月月考.tex}"
%
% 原始材料：微信公众号“魔都数学加油站”文章，连载“27学年高一71”；
% 原文标题“27学年高一71——2026-2027年上海市中国中学高一上10月月考”；
% 原文链接：https://mp.weixin.qq.com/s/UZrP53l3GIin040L1EkljA。
% 页面标注发布时间 2026-10-19 22:47（上海）（与 高一69、高一70 页面读出同一 ct 值）。
% 正文为 7 张 PNG 页（`原图/高一71_中国中学_10月月考/01..07.png`），全部 4762×6735，
% 均按页序读取。7 页全为试卷正文页（卷首标题在第 1 页页顶，卷末解答题第 21 题 (3)
% 解析止于第 7 页），无封面页、目录页、二维码或宣传页。粉色斜排“魔都数学加油站”
% 为卷面水印，与公众号页脚一并未录入。
%
% 卷首标题：照卷面印作“2026-2027 年中国中学高一上 10 月月考”（黑体加粗一行）。
% 原卷年份区间用半角连字符，本稿按全稿惯例排 en dash（同 高一68、高一69、高一70）；
% 原卷只有一行标题、无副标题。
%
% 篇幅与题号：原卷自第 3 题起编号（第 1、2 题不在本卷内）。填空题 3—12（10 题）、
% 选择题 13—16（4 题）、解答题 17—21（5 题），共 19 题。本稿沿用原节与原题号，
% 只改排为全稿统一的大节样式（原卷小节标题为“一、填空题”“二、选择题”“三、解答题”）。
%
% 讲评版：原卷每题下直接印【答案】或【解析】。**只印【答案】的 3 题**（第 3、4、5 题），
% 本稿只给 \ans{}；**第 13 题既无【答案】也无【解析】段，答案“B”直接印在题干括号内
% （“(  B  )”）**，按 高一37 第 13 题、高一63 第 14 题、高一67 第 15 题成例把括号留空、
% 答案另提为 \ans{$B$}。其余 15 题全部印【解析】（按全稿口径这类题不另提 \ans{}——
% 第 14—16 题的“故选 D／B／B”写在解析末句，亦照录在解析内）。
% 除第 13 题外，全卷选择题作答括号内均为空白（第 14、15、16 题）。
%
% 副标题：原卷未印副标题。本稿按“知识模块”口径标作“集合与常用逻辑用语、不等式”：
% 19 题中集合与常用逻辑用语 10 题（第 3、4、5、6、9、11、12、13、16、21 题；其中第 4 题
% 属反证法，一并计入常用逻辑用语），不等式 9 题（第 7、8、10、14、15、17、18、19、20 题；
% 第 19 题 (1) 属补集、(2) 属分式不等式，按主干计入不等式），两类均确有分量，故并标。
% 口径同 高一63、高一68、高一69、高一70。
%
% 与原卷的排版差异：
%   ① 原卷每题下直接印【答案】/【解析】，本稿按全稿惯例将答案与解析置于答案版，试卷版
%      保持干净；第 13 题印在题干括号内的答案“B”同样移入 \ans{}，见上；
%   ② 原卷未印副标题，本稿按“知识模块”口径补出；
%   ③ 原卷填空横线约四个字宽，本稿按全稿惯例统一排 \blankline[4.4em]；
%   ④ 原卷数集记号不统一：实数集排斜体 $R$（第 19 题 (1)），自然数集排斜体 $N$
%      （第 5 题题干、第 12 题题干与解析共三处），第 16 题解析又排正体 $\mathrm{N}^*$；
%      本稿按全稿惯例统一改用 $\R$、$\N$、$\N^*$（同 高一68／69／70 排版差异④）；
%      空集记号原卷作 $\varnothing$，照录；
%   ⑤ 原卷句末标点混用：第 15 题题干“（3）$x^2+2+\dfrac1{x^2+2}$**．**其中”、
%      第 17 题题干末“$\dfrac{x-2}{x-1}\geqslant0$**．**”两处放大确认（第 15 题为全角
%      句点、第 17 题为西文句点），其余各题末为西文句点，已逐处照录；题干内冒号——
%      第 5 题“$\{y\mid\cdots\}$**：**”、第 12 题““×”**：**”、第 17 题“不等式组**：**”、
%      第 21 题“定义**：**”——均为全角冒号，照录；数学式内的逗号（如第 12 题
%      “$x\in A,y\in B$”、第 16 题“$A=\{a_1,\cdots\},B=\{b_1,\cdots\},n>2$”）按原卷作
%      半角，中文行文里的逗号按原卷作全角，均未强改；
%   ⑥ 第 17—21 题的子问原卷印作“（1）…（2）…”（另起一段），本稿按 高一69／70 成例用
%      行内“（1）（2）（3）”段，解析中的子问标号一律用 \stepm{(1)}{…}；
%      第 21 题 (2) 解析内层的“①”“②”“③”照原卷作 \stepm{①}{…} 等；
%   ⑦ 每道解答题原卷未留作答空白，本稿按全稿惯例加 \ansspace（末题用 \ansspace*）；
%   ⑧ 选择题的作答括号原卷排半角“(　　)”（第 14、15、16 题括号内空白、第 13 题括号内
%      印 B），本稿按全稿惯例排 \mbox{（\quad）}；选项行排照原卷——第 13 题两行两条
%      （A B 一行、C D 一行），第 14、15、16 题一行四条；
%   ⑨ 第 8 题解析、第 9 题解析、第 17 题题干与解析中的多行不等式组，原卷均排作“{”
%      连接的多行式，本稿按全稿惯例排 cases；
%   ⑩ 第 10 题解析、第 17 题解析、第 19 题 (1) 解析的并集符号原卷排大号 $\bigcup$
%      （显示体），本稿按全稿惯例排行内 $\cup$；
%   ⑪ 第 8 题解析的判别式照原卷作空心“$\triangle$”（与第 7 题解析的“$\Delta$”不同形），
%      照录未改；
%   ⑫ 第 16 题解析第 4 行的“$\forall$”（全称量词）照原卷符号排，未改写为中文“任意”。
%
% 原卷笔误/存疑（均照录未改，正文对应处另有 % 注）：
%   ① 第 19 题 (1) 连用两个“若”——“若 $a=4$，若全集为实数集 $R$”，后一“若”应作
%      “当”或删；原卷如此，照录未改。
%   ② 第 20 题题干“可以售出 8 万本**、**据市场调查”——“万本”后用顿号，按文意应为
%      逗号或句号；原卷如此，照录未改。
%   ③ 第 15 题解析对代数式 (1) 只得出“$y/x+x/y\geqslant2$”（同号时），未讨论取等是否
%      可得——取 $x=y$ 时确能取到 2——却总结“最小值为 2 的代数式有 1 个，故选 $B$”；
%      按解析自身的推理 (1) 也应计入，结论“1 个”与推理矛盾（若计 (1) 应为 2 个、选 C）。
%      原卷如此，照录未改。
%   ④ 第 21 题题干“定义：对于集合 $A$，则称集合 $B=\cdots$”——“对于……”与“则称……”
%      句式杂糅（应作“对于集合 $A$，称……”或“若集合 $A$，则称……”）；原卷如此，照录未改。
%   ⑤ 第 13 题题干末“……的（　）条件”后无任何标点（其余各题题干末均有句点）；
%      原卷如此，照录未改。
%
% 行数核对：逐题数 \step/\stepm 条数与原卷【解析】印刷行数，15 题逐一相符，
% 合计 86 条 / 86 行（Q6：1；Q7：1；Q8：7；Q9：2；Q10：4；Q11：7；Q12：5；Q14：7；
% Q15：8；Q16：7；Q17：3；Q18：7；Q19：8；Q20：6；Q21：13）；另有 3 题只印【答案】
% （第 3、4、5 题）与第 13 题括号内答案，各提 \ans{}，不计入 \step 账。
%
% 插图：无题目配图或数学图形（7 页题面均无“如图/右图/图中”字样；页面另有粉色斜排水印，
% 见上）。
%
\documentclass[a4paper,11pt]{article}
\usepackage{common}

\begin{document}

\examtitlest[3]{2026--2027 年中国中学高一上 10 月月考}{集合与常用逻辑用语、不等式}

\bigsec{一、填空题}

\begin{enumerate}
\setcounter{enumi}{2}

\item 已知集合 $A=[1,+\infty)$，$B=[-1,3)$，则 $A\cap B=$\blankline[4.4em].

\ifanswer
\ans{$[1,3)$}
\fi

\item 用反证法证明命题“已知 $a$、$b$ 为正数，$n$ 为正整数，求证：如果 $a^n>b^n$，那么 $a>b$”时，应当作出的假设为：“假设\blankline[4.4em]”.

\ifanswer
\ans{$a\leqslant b$}
\fi

\item 用列举法表示集合 $\{y\mid y=x+1,x\in\N,x\leqslant1\}$：\blankline[4.4em].

\ifanswer
\ans{$\{1,2\}$}
\fi

\item 某班在一次测验中，有 36 人数学成绩不低于 80 分，有 20 人物理成绩不低于 80 分，有 15 人的数学、物理成绩都不低于 80 分，则这两科成绩中至少有一科不低于 80 分的人数为\blankline[4.4em]人.

\ifanswer
\solbegin
\step{至少有 $36+20-15=41$ 人.}
\solend
\fi

\item 若关于 $x$ 的不等式 $x^2-x+m<0$ 的解集是 $\varnothing$，则实数 $m$ 的取值范围是\blankline[4.4em].

\ifanswer
\solbegin
\step{由题意得 $\Delta=1-4m\leqslant0$，所以 $m\geqslant\dfrac14$.}
\solend
\fi

% 第 8 题解析的判别式原卷作空心“△”（与第 7 题解析的“Δ”不同形），照录（见文件头“排版差异”⑪）
\item 已知关于 $x$ 的方程 $x^2+mx+2m=0$ 的两个实数根分别为 $\alpha$，$\beta$，若 $\alpha^2+\beta^2=5$，则实数 $m=$\blankline[4.4em].

\ifanswer
\solbegin
\step{方程 $x^2+mx+2m=0$ 的两个实数根为 $\alpha$、$\beta$，}
\step{由根与系数的关系得 $\begin{cases}\alpha+\beta=-m\\ \alpha\beta=2m\end{cases}$，又因为 $\alpha^2+\beta^2=5$，}
\step{所以 $\alpha^2+\beta^2=(\alpha+\beta)^2-2\alpha\beta=m^2-4m=5$，整理得 $m^2-4m-5=0$，}
\step{解得 $m=5$ 或 $m=-1$，}
\step{当 $m=5$ 时，$\triangle=25-40=-15<0$，无实根，舍去；}
\step{当 $m=-1$ 时，$\triangle=1+8=9>0$，符合题意，}
\step{所以实数 $m=-1$.}
\solend
\fi

\item 已知集合 $A=\{x\mid6x+a>0\}$，若 $1\in A$，$0\notin A$，则实数 $a$ 的取值范围为\blankline[4.4em].

\ifanswer
\solbegin
\step{由 $A=\{x\mid6x+a>0\}$，$1\in A$，$0\notin A$，得 $\begin{cases}6+a>0\\ a\leqslant0\end{cases}$，解得 $-6<a\leqslant0$，}
\step{即实数 $a$ 的取值范围是 $(-6,0]$.}
\solend
\fi

\item 关于 $x$ 的不等式 $ax-b>0$ 的解集为 $(1,+\infty)$，则关于 $x$ 的不等式 $(ax+b)(x-2)>0$ 的解集为\blankline[4.4em].

\ifanswer
\solbegin
\step{由不等式 $ax-b>0$ 的解集为 $(1,+\infty)$，得 $a>0$ 且 $\dfrac ba=1$，即 $b=a$，}
\step{则 $(ax+b)(x-2)>0$，可化为 $a(x+1)(x-2)>0$.}
\step{即 $(x+1)(x-2)>0$，解得 $x<-1$ 或 $x>2$，}
\step{所以不等式的解集为 $(-\infty,-1)\cup(2,+\infty)$.}
\solend
\fi

\item 已知实数 $a_1$、$a_2$、$a_3$、$a_4$、$a_5$ 满足 $\{x\mid(x-a_1)(x-a_2)(x-a_3)(x-a_4)(x-a_5)=0\}=\{1,2,3,4\}$，则 $a_1+a_2+a_3+a_4+a_5$ 的最大值为\blankline[4.4em].

\ifanswer
\solbegin
\step{5 个实数 $a_1$、$a_2$、$a_3$、$a_4$、$a_5$ 的不同取值只有 1，2，3，4，}
\step{即这 5 个数均从 1，2，3，4 中选取.}
\step{集合有 4 个元素，说明 1，2，3，4 每个数至少出现一次，}
\step{因此恰好有一个数重复出现两次，其余三个数各出现一次.}
\step{要让 5 个数的总和最大，应让重复出现两次的数取集合中的最大值 4，}
\step{此时 5 个数为 1，2，3，4，4，总和为 $1+2+3+4+4=14$，}
\step{即所求最大值为 14.}
\solend
\fi

\item 定义集合运算“$\times$”：$A\times B=\{(x,y)\mid x\in A,y\in B\}$，称为 $A$、$B$ 两个集合的“卡氏积”，若 $A=\{x\mid x^2-|x|\leqslant0,x\in\N\}$，$B=\{1,2,3\}$，则 $(A\times B)\cap(B\times A)=$\blankline[4.4em].

\ifanswer
\solbegin
\step{因为 $x^2-|x|\leqslant0$，所以 $0\leqslant|x|\leqslant1$，又因为 $x\in\N$，}
\step{所以 $A=\{x\mid x^2-|x|\leqslant0,x\in\N\}=\{0,1\}$，又 $B=\{1,2,3\}$，}
\step{由“卡氏积”的定义得 $A\times B=\{(0,1),(0,2),(0,3),(1,1),(1,2),(1,3)\}$，}
\step{$B\times A=\{(1,0),(1,1),(2,0),(2,1),(3,0),(3,1)\}$，}
\step{所以 $(A\times B)\cap(B\times A)=\{(1,1)\}$.}
\solend
\fi

\end{enumerate}

\bigsec{二、选择题}

\begin{enumerate}
\setcounter{enumi}{12}

% 第 13 题答案“B”原卷印在题干括号内，且无【答案】/【解析】段；
% 按 高一37 第 13 题、高一63 第 14 题、高一67 第 15 题成例括号留空、答案另提 \ans{$B$}
% 原卷题干末“条件”后无标点，照录（见文件头“原卷笔误/存疑”⑤）
\item “$A\subseteq B$”是“$A\subset B$”的\mbox{（\quad）}条件

\keepwithprev
A. 充分不必要条件\qquad B. 必要不充分条件

C. 充要条件\qquad D. 既不充分也不必要条件

\ifanswer
\ans{$B$}
\fi

\item 已知 $a$、$b$、$c$ 均为实数，若 $a>b$ 且 $ab\neq0$，则下列不等关系中成立的是\mbox{（\quad）}

\keepwithprev
A. $ac^2>bc^2$\qquad B. $\dfrac1b>\dfrac1a$\qquad C. $|a|>|b|$\qquad D. $c-b>c-a$

\ifanswer
\solbegin
\step{对于 $A$，若 $a>b$，当 $c=0$ 时，有 $ac^2=bc^2=0$，故 $A$ 错误；}
\step{对于 $B$，$\dfrac1b-\dfrac1a=\dfrac{a-b}{ab}$，因为 $a>b$，所以 $a-b>0$，}
\step{当 $a$、$b$ 同号时，$ab>0$，则 $\dfrac1b-\dfrac1a>0$，即 $\dfrac1b>\dfrac1a$，}
\step{当 $a$、$b$ 异号时，$ab<0$，则 $\dfrac1b-\dfrac1a<0$，即 $\dfrac1b<\dfrac1a$，故 $B$ 错误；}
\step{对于 $C$，当 $a=1$，$b=-3$ 时，$|a|<|b|$，故 $C$ 错误；}
\step{对于 $D$，若 $a>b$，则 $-b>-a$，$c-b>c-a$，故 $D$ 正确.}
\step{故选 $D$.}
\solend
\fi

% 原卷解析对 (1) 未讨论取等是否可得（x=y 时确能取 2），却总结“有 1 个、故选 B”，
% 与推理矛盾，照录（见文件头“原卷笔误/存疑”③）
\item 以下三个代数式：（1）$\dfrac yx+\dfrac xy$；（2）$|x|+\dfrac1{|x|}$；（3）$x^2+2+\dfrac1{x^2+2}$．其中的字母均表示实数，且代数式均有意义，那么代数式的最小值为 2 的个数是\mbox{（\quad）}

\keepwithprev
A. $0$\qquad B. $1$\qquad C. $2$\qquad D. $3$

\ifanswer
\solbegin
\step{（1）仅当 $x,y$ 同号时，$\dfrac yx,\dfrac xy>0$，由基本不等式得 $\dfrac yx+\dfrac xy\geqslant2$；}
\step{若 $x,y$ 异号，$\dfrac yx,\dfrac xy<0$，故 $\dfrac yx+\dfrac xy<0$，此时最小值不是 2；}
\step{（2）$|x|+\dfrac1{|x|}$ 中，$|x|>0$，由基本不等式 $|x|+\dfrac1{|x|}\geqslant2\sqrt{|x|\cdot\dfrac1{|x|}}=2$，}
\step{当且仅当 $|x|=1$，即 $x=\pm1$ 时取等号，故最小值为 2；}
\step{（3）令 $t=x^2+2$，则 $t\geqslant2$，函数化为 $t+\dfrac1t$，}
\step{对勾函数 $t+\dfrac1t$ 在 $[2,+\infty)$ 上严格增，最小值为 $2+\dfrac12=\dfrac52>2$，}
\step{等号成立条件 $t=1$ 无法取到，故最小值不是 2；}
\step{综上，最小值为 2 的代数式有 1 个，故选 $B$.}
\solend
\fi

\item 已知 $A$、$B$ 为非空实数集，$\Omega$ 为平面直角坐标系中的一些点构成的集合，集合 $C=\{y\mid$对任意 $x\in A$，有 $(x,y)\in\Omega\}$，集合 $D=\{x\mid$对任意 $y\in B$，有 $(x,y)\in\Omega\}$，对于下列两个命题：①若 $C\subseteq B$，则 $D\subseteq A$；②若 $C\supseteq B$，则 $D\subseteq A$；其中判断正确的是\mbox{（\quad）}

\keepwithprev
A. ①②都正确\qquad B. ①②都错误\qquad C. ①正确②错误\qquad D. ①错误②正确

\ifanswer
\solbegin
\step{设 $A=\{a_1,a_2,a_3,\cdots,a_n\},B=\{b_1,b_2,b_3,\cdots,b_n\},n>2$，}
\step{若 $\Omega=\{(c,b_1),(c,b_2),\cdots,(c,b_n)\},c\notin A$，}
\step{此时 $C=\varnothing\subseteq B,D=\{c\},D\subseteq A$ 错误，故①错误；}
% 原卷此处 N* 排正体（\mathrm{N}*），本稿统一 \N^*（见文件头“排版差异”④）；
% 全称量词“∀”照原卷符号排（见文件头“排版差异”⑫）
\step{若 $C\supseteq B$，则 $\forall(a_i,b_j)(1\leqslant i,j\leqslant n,i,j\in\N^*)\in\Omega$，则 $D\supseteq A$，}
\step{且 $A=\{a_1\},B=\{b_1,b_2\}$，}
\step{若 $\Omega=\{(a_1,b_1),(a_2,b_1),(a_1,b_2),(a_2,b_2),(c,b_1),(c,b_2)\},D$ 真包含 $A$，故②错误.}
\step{故选 $B$.}
\solend
\fi

\end{enumerate}

\bigsec{三、解答题}

\begin{enumerate}
\setcounter{enumi}{16}

\item 解下列关于 $x$ 的不等式组：$\begin{cases}|x-1|<2\\ \dfrac{x-2}{x-1}\geqslant0\end{cases}$.

\ifanswer
\solbegin
\step{解 $|x-1|<2$ 得 $-2<x-1<2$，解得 $-1<x<3$.}
\step{分式不等式 $\dfrac{x-2}{x-1}\geqslant0$ 等价于 $\begin{cases}(x-2)(x-1)\geqslant0\\ x-1\neq0\end{cases}$，解得 $x<1$ 或 $x\geqslant2$.}
\step{即 $\begin{cases}|x-1|<2\\ \dfrac{x-2}{x-1}\geqslant0\end{cases}$ 的解集为 $(-1,1)\cup[2,3)$.}
\solend
\fi

\ansspace[5]

\item 比较多项式 $2x^2-3x+1$ 与 $x^2-1$ 的大小.

\ifanswer
\solbegin
\step{作差得 $(2x^2-3x+1)-(x^2-1)=x^2-3x+2=(x-1)(x-2)$，}
\step{当 $x=1$ 或 $x=2$ 时，$(x-1)(x-2)=0$，即 $2x^2-3x+1=x^2-1$，}
\step{当 $x<1$ 或 $x>2$ 时，$(x-1)(x-2)>0$，即 $2x^2-3x+1>x^2-1$，}
\step{当 $1<x<2$ 时，$(x-1)(x-2)<0$，即 $2x^2-3x+1<x^2-1$；}
\step{综上，当 $x=1$ 或 $x=2$ 时，$2x^2-3x+1=x^2-1$；}
\step{当 $x<1$ 或 $x>2$ 时，$2x^2-3x+1>x^2-1$；}
\step{当 $1<x<2$ 时，$2x^2-3x+1<x^2-1$.}
\solend
\fi

\ansspace[6]

% 原卷 (1) 连用两个“若”（“若 a=4，若全集为实数集 R”），照录（见文件头“原卷笔误/存疑”①）
\item 已知集合 $A=\left\{x\mid\dfrac{(a+1)x}{x-1}<a\right\}$，集合 $B=[1,4]$，$a$ 为实常数；

（1）若 $a=4$，若全集为实数集 $\R$，求集合 $\overline{A}$（用区间形式表示）；

（2）若 $A\subseteq B$，求 $a$ 的取值范围.

\ifanswer
\solbegin
\stepm{(1)}{若 $a=4$，集合 $A=\left\{x\mid\dfrac{5x}{x-1}<4\right\}=\left\{x\mid\dfrac{5x}{x-1}-4<0\right\}=\left\{x\mid\dfrac{x+4}{x-1}<0\right\}$}
\step{$=(-4,1)$，因此 $\overline{A}=(-\infty,-4]\cup[1,+\infty)$；}
\stepm{(2)}{当 $-a<1$，即 $a>-1$ 时，$A=(-a,1)$，}
\step{要使 $A\subseteq[1,4]$，需 $-a\geqslant1$，即 $a\leqslant-1$，与 $a>-1$ 矛盾，舍去；}
\step{当 $-a>1$，即 $a<-1$ 时，$A=(1,-a)$，}
\step{要使 $A\subseteq[1,4]$，需 $-a\leqslant4$，即 $a\geqslant-4$，所以 $-4\leqslant a<-1$；}
\step{当 $-a=1$，即 $a=-1$ 时，$A=\varnothing$，满足 $A\subseteq B$；}
\step{综上，实数 $a$ 的取值范围为 $[-4,-1]$.}
\solend
\fi

\ansspace[7]

% 原卷题干“可以售出 8 万本、据市场调查”处用顿号，照录（见文件头“原卷笔误/存疑”②）
\item 某种杂志原以每本 25 元的价格销售，可以售出 8 万本、据市场调查，杂志的单价每提高 1 元，销售量就减少 0.2 万本，假定杂志的成本是每本 10 元（不计其他成本）.

（1）当杂志以每本 30 元定价销售时，求销售该杂志所获利润；

（2）在每本 25 元的价格的基础上提高定价，试确定杂志的定价在什么范围内可使得销售利润不低于 140 万元.

\ifanswer
\solbegin
\stepm{(1)}{当杂志以 30 元定价销售时，}
\step{销售该杂志所获利润为 $[8-(30-25)\times0.2]\times(30-10)=140$ 万元；}
\stepm{(2)}{在每本 25 元的价格的基础上提高定价，设杂志的定价为每本 $x$ 元，}
\step{由题意得 $[8-(x-25)\times0.2]\times(x-10)\geqslant140$，解得 $30\leqslant x\leqslant45$，}
\step{所以杂志的定价在每本不低于 30 元且不高于 45 元的范围内，}
\step{可使得销售利润不低于 140 万元.}
\solend
\fi

\ansspace[7]

% 原卷题干“定义：对于集合 A，则称集合 B=…”句式杂糅，照录（见文件头“原卷笔误/存疑”④）
\item 定义：对于集合 $A$，则称集合 $B=\left\{x\mid x=\dfrac{m+n}{1+mn},m,n\in A,1+mn\neq0\right\}$ 为集合 $A$ 的生成集（其中 $m$、$n$ 可以相等）.

（1）若集合 $A=\{-1,1\}$ 的生成集为 $B$，求集合 $B$；

（2）若集合 $A=\{a,2\}$ 的生成集为 $B$，$B$ 的子集个数为 4 个，求实数 $a$ 的值；

（3）若集合 $A=\{x\mid-1<x<1\}$ 的生成集为 $B$，求证：$B\subseteq A$.

\ifanswer
\solbegin
\stepm{(1)}{因为 $1+mn\neq0$，}
\step{当 $m=n=1$ 时，$x=\dfrac{m+n}{1+mn}=1$；当 $m=n=-1$ 时，$x=\dfrac{m+n}{1+mn}=-1$；}
\step{所以 $B=\{-1,1\}$；}
\stepm{(2)}{$B$ 的子集个数为 4 个，故 $B$ 集合的元素有 2 个，}
\step{当 $m=n=2$ 时，$x=\dfrac{m+n}{1+mn}=\dfrac45$，当 $m=n=a$ 时，$x=\dfrac{m+n}{1+mn}=\dfrac{2a}{1+a^2}$，}
\step{当 $m$ 与 $n$ 分别取 2 和 $a$ 时，$x=\dfrac{a+2}{1+2a}$，}
\stepm{①}{$\dfrac{2a}{1+a^2}=\dfrac45$，解得 $a=2$（舍）或 $a=\dfrac12$，}
\stepm{②}{$\dfrac{a+2}{1+2a}=\dfrac45$，解得 $a=2$（舍），}
\stepm{③}{$\dfrac{2a}{1+a^2}=\dfrac{a+2}{1+2a}$，解得 $a=2$（舍）或 $a=1$ 或 $a=-1$，}
\step{综上所述，$a=1$ 或 $a=-1$ 或 $a=\dfrac12$，即 $a\in\left\{-1,1,\dfrac12\right\}$；}
\stepm{(3)}{$\forall m$、$n\in A$，则 $\dfrac{m+n}{1+mn}+1=\dfrac{(m+1)(n+1)}{1+mn}>0$，故 $\dfrac{m+n}{1+mn}>-1$；}
\step{$\dfrac{m+n}{1+mn}-1=-\dfrac{(m-1)(n-1)}{1+mn}<0$，故 $\dfrac{m+n}{1+mn}<1$；}
\step{故 $B\subseteq\{x\mid-1<x<1\}$，即 $B\subseteq A$.}
\solend
\fi

\ansspace*[10]

\end{enumerate}

\end{document}
"
% 紧凑版：xelatex "\def\iscompact{1}% 高一71_中国中学_10月月考.tex
%   —— 高一上数学“2026-2027 年中国中学高一上 10 月月考”（讲评版）
% 试卷版：xelatex 高一71_中国中学_10月月考.tex
% 答案版：xelatex "\def\isanswer{1}\input{高一71_中国中学_10月月考.tex}"
% 紧凑版：xelatex "\def\iscompact{1}\input{高一71_中国中学_10月月考.tex}"
%
% 原始材料：微信公众号“魔都数学加油站”文章，连载“27学年高一71”；
% 原文标题“27学年高一71——2026-2027年上海市中国中学高一上10月月考”；
% 原文链接：https://mp.weixin.qq.com/s/UZrP53l3GIin040L1EkljA。
% 页面标注发布时间 2026-10-19 22:47（上海）（与 高一69、高一70 页面读出同一 ct 值）。
% 正文为 7 张 PNG 页（`原图/高一71_中国中学_10月月考/01..07.png`），全部 4762×6735，
% 均按页序读取。7 页全为试卷正文页（卷首标题在第 1 页页顶，卷末解答题第 21 题 (3)
% 解析止于第 7 页），无封面页、目录页、二维码或宣传页。粉色斜排“魔都数学加油站”
% 为卷面水印，与公众号页脚一并未录入。
%
% 卷首标题：照卷面印作“2026-2027 年中国中学高一上 10 月月考”（黑体加粗一行）。
% 原卷年份区间用半角连字符，本稿按全稿惯例排 en dash（同 高一68、高一69、高一70）；
% 原卷只有一行标题、无副标题。
%
% 篇幅与题号：原卷自第 3 题起编号（第 1、2 题不在本卷内）。填空题 3—12（10 题）、
% 选择题 13—16（4 题）、解答题 17—21（5 题），共 19 题。本稿沿用原节与原题号，
% 只改排为全稿统一的大节样式（原卷小节标题为“一、填空题”“二、选择题”“三、解答题”）。
%
% 讲评版：原卷每题下直接印【答案】或【解析】。**只印【答案】的 3 题**（第 3、4、5 题），
% 本稿只给 \ans{}；**第 13 题既无【答案】也无【解析】段，答案“B”直接印在题干括号内
% （“(  B  )”）**，按 高一37 第 13 题、高一63 第 14 题、高一67 第 15 题成例把括号留空、
% 答案另提为 \ans{$B$}。其余 15 题全部印【解析】（按全稿口径这类题不另提 \ans{}——
% 第 14—16 题的“故选 D／B／B”写在解析末句，亦照录在解析内）。
% 除第 13 题外，全卷选择题作答括号内均为空白（第 14、15、16 题）。
%
% 副标题：原卷未印副标题。本稿按“知识模块”口径标作“集合与常用逻辑用语、不等式”：
% 19 题中集合与常用逻辑用语 10 题（第 3、4、5、6、9、11、12、13、16、21 题；其中第 4 题
% 属反证法，一并计入常用逻辑用语），不等式 9 题（第 7、8、10、14、15、17、18、19、20 题；
% 第 19 题 (1) 属补集、(2) 属分式不等式，按主干计入不等式），两类均确有分量，故并标。
% 口径同 高一63、高一68、高一69、高一70。
%
% 与原卷的排版差异：
%   ① 原卷每题下直接印【答案】/【解析】，本稿按全稿惯例将答案与解析置于答案版，试卷版
%      保持干净；第 13 题印在题干括号内的答案“B”同样移入 \ans{}，见上；
%   ② 原卷未印副标题，本稿按“知识模块”口径补出；
%   ③ 原卷填空横线约四个字宽，本稿按全稿惯例统一排 \blankline[4.4em]；
%   ④ 原卷数集记号不统一：实数集排斜体 $R$（第 19 题 (1)），自然数集排斜体 $N$
%      （第 5 题题干、第 12 题题干与解析共三处），第 16 题解析又排正体 $\mathrm{N}^*$；
%      本稿按全稿惯例统一改用 $\R$、$\N$、$\N^*$（同 高一68／69／70 排版差异④）；
%      空集记号原卷作 $\varnothing$，照录；
%   ⑤ 原卷句末标点混用：第 15 题题干“（3）$x^2+2+\dfrac1{x^2+2}$**．**其中”、
%      第 17 题题干末“$\dfrac{x-2}{x-1}\geqslant0$**．**”两处放大确认（第 15 题为全角
%      句点、第 17 题为西文句点），其余各题末为西文句点，已逐处照录；题干内冒号——
%      第 5 题“$\{y\mid\cdots\}$**：**”、第 12 题““×”**：**”、第 17 题“不等式组**：**”、
%      第 21 题“定义**：**”——均为全角冒号，照录；数学式内的逗号（如第 12 题
%      “$x\in A,y\in B$”、第 16 题“$A=\{a_1,\cdots\},B=\{b_1,\cdots\},n>2$”）按原卷作
%      半角，中文行文里的逗号按原卷作全角，均未强改；
%   ⑥ 第 17—21 题的子问原卷印作“（1）…（2）…”（另起一段），本稿按 高一69／70 成例用
%      行内“（1）（2）（3）”段，解析中的子问标号一律用 \stepm{(1)}{…}；
%      第 21 题 (2) 解析内层的“①”“②”“③”照原卷作 \stepm{①}{…} 等；
%   ⑦ 每道解答题原卷未留作答空白，本稿按全稿惯例加 \ansspace（末题用 \ansspace*）；
%   ⑧ 选择题的作答括号原卷排半角“(　　)”（第 14、15、16 题括号内空白、第 13 题括号内
%      印 B），本稿按全稿惯例排 \mbox{（\quad）}；选项行排照原卷——第 13 题两行两条
%      （A B 一行、C D 一行），第 14、15、16 题一行四条；
%   ⑨ 第 8 题解析、第 9 题解析、第 17 题题干与解析中的多行不等式组，原卷均排作“{”
%      连接的多行式，本稿按全稿惯例排 cases；
%   ⑩ 第 10 题解析、第 17 题解析、第 19 题 (1) 解析的并集符号原卷排大号 $\bigcup$
%      （显示体），本稿按全稿惯例排行内 $\cup$；
%   ⑪ 第 8 题解析的判别式照原卷作空心“$\triangle$”（与第 7 题解析的“$\Delta$”不同形），
%      照录未改；
%   ⑫ 第 16 题解析第 4 行的“$\forall$”（全称量词）照原卷符号排，未改写为中文“任意”。
%
% 原卷笔误/存疑（均照录未改，正文对应处另有 % 注）：
%   ① 第 19 题 (1) 连用两个“若”——“若 $a=4$，若全集为实数集 $R$”，后一“若”应作
%      “当”或删；原卷如此，照录未改。
%   ② 第 20 题题干“可以售出 8 万本**、**据市场调查”——“万本”后用顿号，按文意应为
%      逗号或句号；原卷如此，照录未改。
%   ③ 第 15 题解析对代数式 (1) 只得出“$y/x+x/y\geqslant2$”（同号时），未讨论取等是否
%      可得——取 $x=y$ 时确能取到 2——却总结“最小值为 2 的代数式有 1 个，故选 $B$”；
%      按解析自身的推理 (1) 也应计入，结论“1 个”与推理矛盾（若计 (1) 应为 2 个、选 C）。
%      原卷如此，照录未改。
%   ④ 第 21 题题干“定义：对于集合 $A$，则称集合 $B=\cdots$”——“对于……”与“则称……”
%      句式杂糅（应作“对于集合 $A$，称……”或“若集合 $A$，则称……”）；原卷如此，照录未改。
%   ⑤ 第 13 题题干末“……的（　）条件”后无任何标点（其余各题题干末均有句点）；
%      原卷如此，照录未改。
%
% 行数核对：逐题数 \step/\stepm 条数与原卷【解析】印刷行数，15 题逐一相符，
% 合计 86 条 / 86 行（Q6：1；Q7：1；Q8：7；Q9：2；Q10：4；Q11：7；Q12：5；Q14：7；
% Q15：8；Q16：7；Q17：3；Q18：7；Q19：8；Q20：6；Q21：13）；另有 3 题只印【答案】
% （第 3、4、5 题）与第 13 题括号内答案，各提 \ans{}，不计入 \step 账。
%
% 插图：无题目配图或数学图形（7 页题面均无“如图/右图/图中”字样；页面另有粉色斜排水印，
% 见上）。
%
\documentclass[a4paper,11pt]{article}
\usepackage{common}

\begin{document}

\examtitlest[3]{2026--2027 年中国中学高一上 10 月月考}{集合与常用逻辑用语、不等式}

\bigsec{一、填空题}

\begin{enumerate}
\setcounter{enumi}{2}

\item 已知集合 $A=[1,+\infty)$，$B=[-1,3)$，则 $A\cap B=$\blankline[4.4em].

\ifanswer
\ans{$[1,3)$}
\fi

\item 用反证法证明命题“已知 $a$、$b$ 为正数，$n$ 为正整数，求证：如果 $a^n>b^n$，那么 $a>b$”时，应当作出的假设为：“假设\blankline[4.4em]”.

\ifanswer
\ans{$a\leqslant b$}
\fi

\item 用列举法表示集合 $\{y\mid y=x+1,x\in\N,x\leqslant1\}$：\blankline[4.4em].

\ifanswer
\ans{$\{1,2\}$}
\fi

\item 某班在一次测验中，有 36 人数学成绩不低于 80 分，有 20 人物理成绩不低于 80 分，有 15 人的数学、物理成绩都不低于 80 分，则这两科成绩中至少有一科不低于 80 分的人数为\blankline[4.4em]人.

\ifanswer
\solbegin
\step{至少有 $36+20-15=41$ 人.}
\solend
\fi

\item 若关于 $x$ 的不等式 $x^2-x+m<0$ 的解集是 $\varnothing$，则实数 $m$ 的取值范围是\blankline[4.4em].

\ifanswer
\solbegin
\step{由题意得 $\Delta=1-4m\leqslant0$，所以 $m\geqslant\dfrac14$.}
\solend
\fi

% 第 8 题解析的判别式原卷作空心“△”（与第 7 题解析的“Δ”不同形），照录（见文件头“排版差异”⑪）
\item 已知关于 $x$ 的方程 $x^2+mx+2m=0$ 的两个实数根分别为 $\alpha$，$\beta$，若 $\alpha^2+\beta^2=5$，则实数 $m=$\blankline[4.4em].

\ifanswer
\solbegin
\step{方程 $x^2+mx+2m=0$ 的两个实数根为 $\alpha$、$\beta$，}
\step{由根与系数的关系得 $\begin{cases}\alpha+\beta=-m\\ \alpha\beta=2m\end{cases}$，又因为 $\alpha^2+\beta^2=5$，}
\step{所以 $\alpha^2+\beta^2=(\alpha+\beta)^2-2\alpha\beta=m^2-4m=5$，整理得 $m^2-4m-5=0$，}
\step{解得 $m=5$ 或 $m=-1$，}
\step{当 $m=5$ 时，$\triangle=25-40=-15<0$，无实根，舍去；}
\step{当 $m=-1$ 时，$\triangle=1+8=9>0$，符合题意，}
\step{所以实数 $m=-1$.}
\solend
\fi

\item 已知集合 $A=\{x\mid6x+a>0\}$，若 $1\in A$，$0\notin A$，则实数 $a$ 的取值范围为\blankline[4.4em].

\ifanswer
\solbegin
\step{由 $A=\{x\mid6x+a>0\}$，$1\in A$，$0\notin A$，得 $\begin{cases}6+a>0\\ a\leqslant0\end{cases}$，解得 $-6<a\leqslant0$，}
\step{即实数 $a$ 的取值范围是 $(-6,0]$.}
\solend
\fi

\item 关于 $x$ 的不等式 $ax-b>0$ 的解集为 $(1,+\infty)$，则关于 $x$ 的不等式 $(ax+b)(x-2)>0$ 的解集为\blankline[4.4em].

\ifanswer
\solbegin
\step{由不等式 $ax-b>0$ 的解集为 $(1,+\infty)$，得 $a>0$ 且 $\dfrac ba=1$，即 $b=a$，}
\step{则 $(ax+b)(x-2)>0$，可化为 $a(x+1)(x-2)>0$.}
\step{即 $(x+1)(x-2)>0$，解得 $x<-1$ 或 $x>2$，}
\step{所以不等式的解集为 $(-\infty,-1)\cup(2,+\infty)$.}
\solend
\fi

\item 已知实数 $a_1$、$a_2$、$a_3$、$a_4$、$a_5$ 满足 $\{x\mid(x-a_1)(x-a_2)(x-a_3)(x-a_4)(x-a_5)=0\}=\{1,2,3,4\}$，则 $a_1+a_2+a_3+a_4+a_5$ 的最大值为\blankline[4.4em].

\ifanswer
\solbegin
\step{5 个实数 $a_1$、$a_2$、$a_3$、$a_4$、$a_5$ 的不同取值只有 1，2，3，4，}
\step{即这 5 个数均从 1，2，3，4 中选取.}
\step{集合有 4 个元素，说明 1，2，3，4 每个数至少出现一次，}
\step{因此恰好有一个数重复出现两次，其余三个数各出现一次.}
\step{要让 5 个数的总和最大，应让重复出现两次的数取集合中的最大值 4，}
\step{此时 5 个数为 1，2，3，4，4，总和为 $1+2+3+4+4=14$，}
\step{即所求最大值为 14.}
\solend
\fi

\item 定义集合运算“$\times$”：$A\times B=\{(x,y)\mid x\in A,y\in B\}$，称为 $A$、$B$ 两个集合的“卡氏积”，若 $A=\{x\mid x^2-|x|\leqslant0,x\in\N\}$，$B=\{1,2,3\}$，则 $(A\times B)\cap(B\times A)=$\blankline[4.4em].

\ifanswer
\solbegin
\step{因为 $x^2-|x|\leqslant0$，所以 $0\leqslant|x|\leqslant1$，又因为 $x\in\N$，}
\step{所以 $A=\{x\mid x^2-|x|\leqslant0,x\in\N\}=\{0,1\}$，又 $B=\{1,2,3\}$，}
\step{由“卡氏积”的定义得 $A\times B=\{(0,1),(0,2),(0,3),(1,1),(1,2),(1,3)\}$，}
\step{$B\times A=\{(1,0),(1,1),(2,0),(2,1),(3,0),(3,1)\}$，}
\step{所以 $(A\times B)\cap(B\times A)=\{(1,1)\}$.}
\solend
\fi

\end{enumerate}

\bigsec{二、选择题}

\begin{enumerate}
\setcounter{enumi}{12}

% 第 13 题答案“B”原卷印在题干括号内，且无【答案】/【解析】段；
% 按 高一37 第 13 题、高一63 第 14 题、高一67 第 15 题成例括号留空、答案另提 \ans{$B$}
% 原卷题干末“条件”后无标点，照录（见文件头“原卷笔误/存疑”⑤）
\item “$A\subseteq B$”是“$A\subset B$”的\mbox{（\quad）}条件

\keepwithprev
A. 充分不必要条件\qquad B. 必要不充分条件

C. 充要条件\qquad D. 既不充分也不必要条件

\ifanswer
\ans{$B$}
\fi

\item 已知 $a$、$b$、$c$ 均为实数，若 $a>b$ 且 $ab\neq0$，则下列不等关系中成立的是\mbox{（\quad）}

\keepwithprev
A. $ac^2>bc^2$\qquad B. $\dfrac1b>\dfrac1a$\qquad C. $|a|>|b|$\qquad D. $c-b>c-a$

\ifanswer
\solbegin
\step{对于 $A$，若 $a>b$，当 $c=0$ 时，有 $ac^2=bc^2=0$，故 $A$ 错误；}
\step{对于 $B$，$\dfrac1b-\dfrac1a=\dfrac{a-b}{ab}$，因为 $a>b$，所以 $a-b>0$，}
\step{当 $a$、$b$ 同号时，$ab>0$，则 $\dfrac1b-\dfrac1a>0$，即 $\dfrac1b>\dfrac1a$，}
\step{当 $a$、$b$ 异号时，$ab<0$，则 $\dfrac1b-\dfrac1a<0$，即 $\dfrac1b<\dfrac1a$，故 $B$ 错误；}
\step{对于 $C$，当 $a=1$，$b=-3$ 时，$|a|<|b|$，故 $C$ 错误；}
\step{对于 $D$，若 $a>b$，则 $-b>-a$，$c-b>c-a$，故 $D$ 正确.}
\step{故选 $D$.}
\solend
\fi

% 原卷解析对 (1) 未讨论取等是否可得（x=y 时确能取 2），却总结“有 1 个、故选 B”，
% 与推理矛盾，照录（见文件头“原卷笔误/存疑”③）
\item 以下三个代数式：（1）$\dfrac yx+\dfrac xy$；（2）$|x|+\dfrac1{|x|}$；（3）$x^2+2+\dfrac1{x^2+2}$．其中的字母均表示实数，且代数式均有意义，那么代数式的最小值为 2 的个数是\mbox{（\quad）}

\keepwithprev
A. $0$\qquad B. $1$\qquad C. $2$\qquad D. $3$

\ifanswer
\solbegin
\step{（1）仅当 $x,y$ 同号时，$\dfrac yx,\dfrac xy>0$，由基本不等式得 $\dfrac yx+\dfrac xy\geqslant2$；}
\step{若 $x,y$ 异号，$\dfrac yx,\dfrac xy<0$，故 $\dfrac yx+\dfrac xy<0$，此时最小值不是 2；}
\step{（2）$|x|+\dfrac1{|x|}$ 中，$|x|>0$，由基本不等式 $|x|+\dfrac1{|x|}\geqslant2\sqrt{|x|\cdot\dfrac1{|x|}}=2$，}
\step{当且仅当 $|x|=1$，即 $x=\pm1$ 时取等号，故最小值为 2；}
\step{（3）令 $t=x^2+2$，则 $t\geqslant2$，函数化为 $t+\dfrac1t$，}
\step{对勾函数 $t+\dfrac1t$ 在 $[2,+\infty)$ 上严格增，最小值为 $2+\dfrac12=\dfrac52>2$，}
\step{等号成立条件 $t=1$ 无法取到，故最小值不是 2；}
\step{综上，最小值为 2 的代数式有 1 个，故选 $B$.}
\solend
\fi

\item 已知 $A$、$B$ 为非空实数集，$\Omega$ 为平面直角坐标系中的一些点构成的集合，集合 $C=\{y\mid$对任意 $x\in A$，有 $(x,y)\in\Omega\}$，集合 $D=\{x\mid$对任意 $y\in B$，有 $(x,y)\in\Omega\}$，对于下列两个命题：①若 $C\subseteq B$，则 $D\subseteq A$；②若 $C\supseteq B$，则 $D\subseteq A$；其中判断正确的是\mbox{（\quad）}

\keepwithprev
A. ①②都正确\qquad B. ①②都错误\qquad C. ①正确②错误\qquad D. ①错误②正确

\ifanswer
\solbegin
\step{设 $A=\{a_1,a_2,a_3,\cdots,a_n\},B=\{b_1,b_2,b_3,\cdots,b_n\},n>2$，}
\step{若 $\Omega=\{(c,b_1),(c,b_2),\cdots,(c,b_n)\},c\notin A$，}
\step{此时 $C=\varnothing\subseteq B,D=\{c\},D\subseteq A$ 错误，故①错误；}
% 原卷此处 N* 排正体（\mathrm{N}*），本稿统一 \N^*（见文件头“排版差异”④）；
% 全称量词“∀”照原卷符号排（见文件头“排版差异”⑫）
\step{若 $C\supseteq B$，则 $\forall(a_i,b_j)(1\leqslant i,j\leqslant n,i,j\in\N^*)\in\Omega$，则 $D\supseteq A$，}
\step{且 $A=\{a_1\},B=\{b_1,b_2\}$，}
\step{若 $\Omega=\{(a_1,b_1),(a_2,b_1),(a_1,b_2),(a_2,b_2),(c,b_1),(c,b_2)\},D$ 真包含 $A$，故②错误.}
\step{故选 $B$.}
\solend
\fi

\end{enumerate}

\bigsec{三、解答题}

\begin{enumerate}
\setcounter{enumi}{16}

\item 解下列关于 $x$ 的不等式组：$\begin{cases}|x-1|<2\\ \dfrac{x-2}{x-1}\geqslant0\end{cases}$.

\ifanswer
\solbegin
\step{解 $|x-1|<2$ 得 $-2<x-1<2$，解得 $-1<x<3$.}
\step{分式不等式 $\dfrac{x-2}{x-1}\geqslant0$ 等价于 $\begin{cases}(x-2)(x-1)\geqslant0\\ x-1\neq0\end{cases}$，解得 $x<1$ 或 $x\geqslant2$.}
\step{即 $\begin{cases}|x-1|<2\\ \dfrac{x-2}{x-1}\geqslant0\end{cases}$ 的解集为 $(-1,1)\cup[2,3)$.}
\solend
\fi

\ansspace[5]

\item 比较多项式 $2x^2-3x+1$ 与 $x^2-1$ 的大小.

\ifanswer
\solbegin
\step{作差得 $(2x^2-3x+1)-(x^2-1)=x^2-3x+2=(x-1)(x-2)$，}
\step{当 $x=1$ 或 $x=2$ 时，$(x-1)(x-2)=0$，即 $2x^2-3x+1=x^2-1$，}
\step{当 $x<1$ 或 $x>2$ 时，$(x-1)(x-2)>0$，即 $2x^2-3x+1>x^2-1$，}
\step{当 $1<x<2$ 时，$(x-1)(x-2)<0$，即 $2x^2-3x+1<x^2-1$；}
\step{综上，当 $x=1$ 或 $x=2$ 时，$2x^2-3x+1=x^2-1$；}
\step{当 $x<1$ 或 $x>2$ 时，$2x^2-3x+1>x^2-1$；}
\step{当 $1<x<2$ 时，$2x^2-3x+1<x^2-1$.}
\solend
\fi

\ansspace[6]

% 原卷 (1) 连用两个“若”（“若 a=4，若全集为实数集 R”），照录（见文件头“原卷笔误/存疑”①）
\item 已知集合 $A=\left\{x\mid\dfrac{(a+1)x}{x-1}<a\right\}$，集合 $B=[1,4]$，$a$ 为实常数；

（1）若 $a=4$，若全集为实数集 $\R$，求集合 $\overline{A}$（用区间形式表示）；

（2）若 $A\subseteq B$，求 $a$ 的取值范围.

\ifanswer
\solbegin
\stepm{(1)}{若 $a=4$，集合 $A=\left\{x\mid\dfrac{5x}{x-1}<4\right\}=\left\{x\mid\dfrac{5x}{x-1}-4<0\right\}=\left\{x\mid\dfrac{x+4}{x-1}<0\right\}$}
\step{$=(-4,1)$，因此 $\overline{A}=(-\infty,-4]\cup[1,+\infty)$；}
\stepm{(2)}{当 $-a<1$，即 $a>-1$ 时，$A=(-a,1)$，}
\step{要使 $A\subseteq[1,4]$，需 $-a\geqslant1$，即 $a\leqslant-1$，与 $a>-1$ 矛盾，舍去；}
\step{当 $-a>1$，即 $a<-1$ 时，$A=(1,-a)$，}
\step{要使 $A\subseteq[1,4]$，需 $-a\leqslant4$，即 $a\geqslant-4$，所以 $-4\leqslant a<-1$；}
\step{当 $-a=1$，即 $a=-1$ 时，$A=\varnothing$，满足 $A\subseteq B$；}
\step{综上，实数 $a$ 的取值范围为 $[-4,-1]$.}
\solend
\fi

\ansspace[7]

% 原卷题干“可以售出 8 万本、据市场调查”处用顿号，照录（见文件头“原卷笔误/存疑”②）
\item 某种杂志原以每本 25 元的价格销售，可以售出 8 万本、据市场调查，杂志的单价每提高 1 元，销售量就减少 0.2 万本，假定杂志的成本是每本 10 元（不计其他成本）.

（1）当杂志以每本 30 元定价销售时，求销售该杂志所获利润；

（2）在每本 25 元的价格的基础上提高定价，试确定杂志的定价在什么范围内可使得销售利润不低于 140 万元.

\ifanswer
\solbegin
\stepm{(1)}{当杂志以 30 元定价销售时，}
\step{销售该杂志所获利润为 $[8-(30-25)\times0.2]\times(30-10)=140$ 万元；}
\stepm{(2)}{在每本 25 元的价格的基础上提高定价，设杂志的定价为每本 $x$ 元，}
\step{由题意得 $[8-(x-25)\times0.2]\times(x-10)\geqslant140$，解得 $30\leqslant x\leqslant45$，}
\step{所以杂志的定价在每本不低于 30 元且不高于 45 元的范围内，}
\step{可使得销售利润不低于 140 万元.}
\solend
\fi

\ansspace[7]

% 原卷题干“定义：对于集合 A，则称集合 B=…”句式杂糅，照录（见文件头“原卷笔误/存疑”④）
\item 定义：对于集合 $A$，则称集合 $B=\left\{x\mid x=\dfrac{m+n}{1+mn},m,n\in A,1+mn\neq0\right\}$ 为集合 $A$ 的生成集（其中 $m$、$n$ 可以相等）.

（1）若集合 $A=\{-1,1\}$ 的生成集为 $B$，求集合 $B$；

（2）若集合 $A=\{a,2\}$ 的生成集为 $B$，$B$ 的子集个数为 4 个，求实数 $a$ 的值；

（3）若集合 $A=\{x\mid-1<x<1\}$ 的生成集为 $B$，求证：$B\subseteq A$.

\ifanswer
\solbegin
\stepm{(1)}{因为 $1+mn\neq0$，}
\step{当 $m=n=1$ 时，$x=\dfrac{m+n}{1+mn}=1$；当 $m=n=-1$ 时，$x=\dfrac{m+n}{1+mn}=-1$；}
\step{所以 $B=\{-1,1\}$；}
\stepm{(2)}{$B$ 的子集个数为 4 个，故 $B$ 集合的元素有 2 个，}
\step{当 $m=n=2$ 时，$x=\dfrac{m+n}{1+mn}=\dfrac45$，当 $m=n=a$ 时，$x=\dfrac{m+n}{1+mn}=\dfrac{2a}{1+a^2}$，}
\step{当 $m$ 与 $n$ 分别取 2 和 $a$ 时，$x=\dfrac{a+2}{1+2a}$，}
\stepm{①}{$\dfrac{2a}{1+a^2}=\dfrac45$，解得 $a=2$（舍）或 $a=\dfrac12$，}
\stepm{②}{$\dfrac{a+2}{1+2a}=\dfrac45$，解得 $a=2$（舍），}
\stepm{③}{$\dfrac{2a}{1+a^2}=\dfrac{a+2}{1+2a}$，解得 $a=2$（舍）或 $a=1$ 或 $a=-1$，}
\step{综上所述，$a=1$ 或 $a=-1$ 或 $a=\dfrac12$，即 $a\in\left\{-1,1,\dfrac12\right\}$；}
\stepm{(3)}{$\forall m$、$n\in A$，则 $\dfrac{m+n}{1+mn}+1=\dfrac{(m+1)(n+1)}{1+mn}>0$，故 $\dfrac{m+n}{1+mn}>-1$；}
\step{$\dfrac{m+n}{1+mn}-1=-\dfrac{(m-1)(n-1)}{1+mn}<0$，故 $\dfrac{m+n}{1+mn}<1$；}
\step{故 $B\subseteq\{x\mid-1<x<1\}$，即 $B\subseteq A$.}
\solend
\fi

\ansspace*[10]

\end{enumerate}

\end{document}
"
%
% 原始材料：微信公众号“魔都数学加油站”文章，连载“27学年高一71”；
% 原文标题“27学年高一71——2026-2027年上海市中国中学高一上10月月考”；
% 原文链接：https://mp.weixin.qq.com/s/UZrP53l3GIin040L1EkljA。
% 页面标注发布时间 2026-10-19 22:47（上海）（与 高一69、高一70 页面读出同一 ct 值）。
% 正文为 7 张 PNG 页（`原图/高一71_中国中学_10月月考/01..07.png`），全部 4762×6735，
% 均按页序读取。7 页全为试卷正文页（卷首标题在第 1 页页顶，卷末解答题第 21 题 (3)
% 解析止于第 7 页），无封面页、目录页、二维码或宣传页。粉色斜排“魔都数学加油站”
% 为卷面水印，与公众号页脚一并未录入。
%
% 卷首标题：照卷面印作“2026-2027 年中国中学高一上 10 月月考”（黑体加粗一行）。
% 原卷年份区间用半角连字符，本稿按全稿惯例排 en dash（同 高一68、高一69、高一70）；
% 原卷只有一行标题、无副标题。
%
% 篇幅与题号：原卷自第 3 题起编号（第 1、2 题不在本卷内）。填空题 3—12（10 题）、
% 选择题 13—16（4 题）、解答题 17—21（5 题），共 19 题。本稿沿用原节与原题号，
% 只改排为全稿统一的大节样式（原卷小节标题为“一、填空题”“二、选择题”“三、解答题”）。
%
% 讲评版：原卷每题下直接印【答案】或【解析】。**只印【答案】的 3 题**（第 3、4、5 题），
% 本稿只给 \ans{}；**第 13 题既无【答案】也无【解析】段，答案“B”直接印在题干括号内
% （“(  B  )”）**，按 高一37 第 13 题、高一63 第 14 题、高一67 第 15 题成例把括号留空、
% 答案另提为 \ans{$B$}。其余 15 题全部印【解析】（按全稿口径这类题不另提 \ans{}——
% 第 14—16 题的“故选 D／B／B”写在解析末句，亦照录在解析内）。
% 除第 13 题外，全卷选择题作答括号内均为空白（第 14、15、16 题）。
%
% 副标题：原卷未印副标题。本稿按“知识模块”口径标作“集合与常用逻辑用语、不等式”：
% 19 题中集合与常用逻辑用语 10 题（第 3、4、5、6、9、11、12、13、16、21 题；其中第 4 题
% 属反证法，一并计入常用逻辑用语），不等式 9 题（第 7、8、10、14、15、17、18、19、20 题；
% 第 19 题 (1) 属补集、(2) 属分式不等式，按主干计入不等式），两类均确有分量，故并标。
% 口径同 高一63、高一68、高一69、高一70。
%
% 与原卷的排版差异：
%   ① 原卷每题下直接印【答案】/【解析】，本稿按全稿惯例将答案与解析置于答案版，试卷版
%      保持干净；第 13 题印在题干括号内的答案“B”同样移入 \ans{}，见上；
%   ② 原卷未印副标题，本稿按“知识模块”口径补出；
%   ③ 原卷填空横线约四个字宽，本稿按全稿惯例统一排 \blankline[4.4em]；
%   ④ 原卷数集记号不统一：实数集排斜体 $R$（第 19 题 (1)），自然数集排斜体 $N$
%      （第 5 题题干、第 12 题题干与解析共三处），第 16 题解析又排正体 $\mathrm{N}^*$；
%      本稿按全稿惯例统一改用 $\R$、$\N$、$\N^*$（同 高一68／69／70 排版差异④）；
%      空集记号原卷作 $\varnothing$，照录；
%   ⑤ 原卷句末标点混用：第 15 题题干“（3）$x^2+2+\dfrac1{x^2+2}$**．**其中”、
%      第 17 题题干末“$\dfrac{x-2}{x-1}\geqslant0$**．**”两处放大确认（第 15 题为全角
%      句点、第 17 题为西文句点），其余各题末为西文句点，已逐处照录；题干内冒号——
%      第 5 题“$\{y\mid\cdots\}$**：**”、第 12 题““×”**：**”、第 17 题“不等式组**：**”、
%      第 21 题“定义**：**”——均为全角冒号，照录；数学式内的逗号（如第 12 题
%      “$x\in A,y\in B$”、第 16 题“$A=\{a_1,\cdots\},B=\{b_1,\cdots\},n>2$”）按原卷作
%      半角，中文行文里的逗号按原卷作全角，均未强改；
%   ⑥ 第 17—21 题的子问原卷印作“（1）…（2）…”（另起一段），本稿按 高一69／70 成例用
%      行内“（1）（2）（3）”段，解析中的子问标号一律用 \stepm{(1)}{…}；
%      第 21 题 (2) 解析内层的“①”“②”“③”照原卷作 \stepm{①}{…} 等；
%   ⑦ 每道解答题原卷未留作答空白，本稿按全稿惯例加 \ansspace（末题用 \ansspace*）；
%   ⑧ 选择题的作答括号原卷排半角“(　　)”（第 14、15、16 题括号内空白、第 13 题括号内
%      印 B），本稿按全稿惯例排 \mbox{（\quad）}；选项行排照原卷——第 13 题两行两条
%      （A B 一行、C D 一行），第 14、15、16 题一行四条；
%   ⑨ 第 8 题解析、第 9 题解析、第 17 题题干与解析中的多行不等式组，原卷均排作“{”
%      连接的多行式，本稿按全稿惯例排 cases；
%   ⑩ 第 10 题解析、第 17 题解析、第 19 题 (1) 解析的并集符号原卷排大号 $\bigcup$
%      （显示体），本稿按全稿惯例排行内 $\cup$；
%   ⑪ 第 8 题解析的判别式照原卷作空心“$\triangle$”（与第 7 题解析的“$\Delta$”不同形），
%      照录未改；
%   ⑫ 第 16 题解析第 4 行的“$\forall$”（全称量词）照原卷符号排，未改写为中文“任意”。
%
% 原卷笔误/存疑（均照录未改，正文对应处另有 % 注）：
%   ① 第 19 题 (1) 连用两个“若”——“若 $a=4$，若全集为实数集 $R$”，后一“若”应作
%      “当”或删；原卷如此，照录未改。
%   ② 第 20 题题干“可以售出 8 万本**、**据市场调查”——“万本”后用顿号，按文意应为
%      逗号或句号；原卷如此，照录未改。
%   ③ 第 15 题解析对代数式 (1) 只得出“$y/x+x/y\geqslant2$”（同号时），未讨论取等是否
%      可得——取 $x=y$ 时确能取到 2——却总结“最小值为 2 的代数式有 1 个，故选 $B$”；
%      按解析自身的推理 (1) 也应计入，结论“1 个”与推理矛盾（若计 (1) 应为 2 个、选 C）。
%      原卷如此，照录未改。
%   ④ 第 21 题题干“定义：对于集合 $A$，则称集合 $B=\cdots$”——“对于……”与“则称……”
%      句式杂糅（应作“对于集合 $A$，称……”或“若集合 $A$，则称……”）；原卷如此，照录未改。
%   ⑤ 第 13 题题干末“……的（　）条件”后无任何标点（其余各题题干末均有句点）；
%      原卷如此，照录未改。
%
% 行数核对：逐题数 \step/\stepm 条数与原卷【解析】印刷行数，15 题逐一相符，
% 合计 86 条 / 86 行（Q6：1；Q7：1；Q8：7；Q9：2；Q10：4；Q11：7；Q12：5；Q14：7；
% Q15：8；Q16：7；Q17：3；Q18：7；Q19：8；Q20：6；Q21：13）；另有 3 题只印【答案】
% （第 3、4、5 题）与第 13 题括号内答案，各提 \ans{}，不计入 \step 账。
%
% 插图：无题目配图或数学图形（7 页题面均无“如图/右图/图中”字样；页面另有粉色斜排水印，
% 见上）。
%
\documentclass[a4paper,11pt]{article}
\usepackage{common}

\begin{document}

\examtitlest[3]{2026--2027 年中国中学高一上 10 月月考}{集合与常用逻辑用语、不等式}

\bigsec{一、填空题}

\begin{enumerate}
\setcounter{enumi}{2}

\item 已知集合 $A=[1,+\infty)$，$B=[-1,3)$，则 $A\cap B=$\blankline[4.4em].

\ifanswer
\ans{$[1,3)$}
\fi

\item 用反证法证明命题“已知 $a$、$b$ 为正数，$n$ 为正整数，求证：如果 $a^n>b^n$，那么 $a>b$”时，应当作出的假设为：“假设\blankline[4.4em]”.

\ifanswer
\ans{$a\leqslant b$}
\fi

\item 用列举法表示集合 $\{y\mid y=x+1,x\in\N,x\leqslant1\}$：\blankline[4.4em].

\ifanswer
\ans{$\{1,2\}$}
\fi

\item 某班在一次测验中，有 36 人数学成绩不低于 80 分，有 20 人物理成绩不低于 80 分，有 15 人的数学、物理成绩都不低于 80 分，则这两科成绩中至少有一科不低于 80 分的人数为\blankline[4.4em]人.

\ifanswer
\solbegin
\step{至少有 $36+20-15=41$ 人.}
\solend
\fi

\item 若关于 $x$ 的不等式 $x^2-x+m<0$ 的解集是 $\varnothing$，则实数 $m$ 的取值范围是\blankline[4.4em].

\ifanswer
\solbegin
\step{由题意得 $\Delta=1-4m\leqslant0$，所以 $m\geqslant\dfrac14$.}
\solend
\fi

% 第 8 题解析的判别式原卷作空心“△”（与第 7 题解析的“Δ”不同形），照录（见文件头“排版差异”⑪）
\item 已知关于 $x$ 的方程 $x^2+mx+2m=0$ 的两个实数根分别为 $\alpha$，$\beta$，若 $\alpha^2+\beta^2=5$，则实数 $m=$\blankline[4.4em].

\ifanswer
\solbegin
\step{方程 $x^2+mx+2m=0$ 的两个实数根为 $\alpha$、$\beta$，}
\step{由根与系数的关系得 $\begin{cases}\alpha+\beta=-m\\ \alpha\beta=2m\end{cases}$，又因为 $\alpha^2+\beta^2=5$，}
\step{所以 $\alpha^2+\beta^2=(\alpha+\beta)^2-2\alpha\beta=m^2-4m=5$，整理得 $m^2-4m-5=0$，}
\step{解得 $m=5$ 或 $m=-1$，}
\step{当 $m=5$ 时，$\triangle=25-40=-15<0$，无实根，舍去；}
\step{当 $m=-1$ 时，$\triangle=1+8=9>0$，符合题意，}
\step{所以实数 $m=-1$.}
\solend
\fi

\item 已知集合 $A=\{x\mid6x+a>0\}$，若 $1\in A$，$0\notin A$，则实数 $a$ 的取值范围为\blankline[4.4em].

\ifanswer
\solbegin
\step{由 $A=\{x\mid6x+a>0\}$，$1\in A$，$0\notin A$，得 $\begin{cases}6+a>0\\ a\leqslant0\end{cases}$，解得 $-6<a\leqslant0$，}
\step{即实数 $a$ 的取值范围是 $(-6,0]$.}
\solend
\fi

\item 关于 $x$ 的不等式 $ax-b>0$ 的解集为 $(1,+\infty)$，则关于 $x$ 的不等式 $(ax+b)(x-2)>0$ 的解集为\blankline[4.4em].

\ifanswer
\solbegin
\step{由不等式 $ax-b>0$ 的解集为 $(1,+\infty)$，得 $a>0$ 且 $\dfrac ba=1$，即 $b=a$，}
\step{则 $(ax+b)(x-2)>0$，可化为 $a(x+1)(x-2)>0$.}
\step{即 $(x+1)(x-2)>0$，解得 $x<-1$ 或 $x>2$，}
\step{所以不等式的解集为 $(-\infty,-1)\cup(2,+\infty)$.}
\solend
\fi

\item 已知实数 $a_1$、$a_2$、$a_3$、$a_4$、$a_5$ 满足 $\{x\mid(x-a_1)(x-a_2)(x-a_3)(x-a_4)(x-a_5)=0\}=\{1,2,3,4\}$，则 $a_1+a_2+a_3+a_4+a_5$ 的最大值为\blankline[4.4em].

\ifanswer
\solbegin
\step{5 个实数 $a_1$、$a_2$、$a_3$、$a_4$、$a_5$ 的不同取值只有 1，2，3，4，}
\step{即这 5 个数均从 1，2，3，4 中选取.}
\step{集合有 4 个元素，说明 1，2，3，4 每个数至少出现一次，}
\step{因此恰好有一个数重复出现两次，其余三个数各出现一次.}
\step{要让 5 个数的总和最大，应让重复出现两次的数取集合中的最大值 4，}
\step{此时 5 个数为 1，2，3，4，4，总和为 $1+2+3+4+4=14$，}
\step{即所求最大值为 14.}
\solend
\fi

\item 定义集合运算“$\times$”：$A\times B=\{(x,y)\mid x\in A,y\in B\}$，称为 $A$、$B$ 两个集合的“卡氏积”，若 $A=\{x\mid x^2-|x|\leqslant0,x\in\N\}$，$B=\{1,2,3\}$，则 $(A\times B)\cap(B\times A)=$\blankline[4.4em].

\ifanswer
\solbegin
\step{因为 $x^2-|x|\leqslant0$，所以 $0\leqslant|x|\leqslant1$，又因为 $x\in\N$，}
\step{所以 $A=\{x\mid x^2-|x|\leqslant0,x\in\N\}=\{0,1\}$，又 $B=\{1,2,3\}$，}
\step{由“卡氏积”的定义得 $A\times B=\{(0,1),(0,2),(0,3),(1,1),(1,2),(1,3)\}$，}
\step{$B\times A=\{(1,0),(1,1),(2,0),(2,1),(3,0),(3,1)\}$，}
\step{所以 $(A\times B)\cap(B\times A)=\{(1,1)\}$.}
\solend
\fi

\end{enumerate}

\bigsec{二、选择题}

\begin{enumerate}
\setcounter{enumi}{12}

% 第 13 题答案“B”原卷印在题干括号内，且无【答案】/【解析】段；
% 按 高一37 第 13 题、高一63 第 14 题、高一67 第 15 题成例括号留空、答案另提 \ans{$B$}
% 原卷题干末“条件”后无标点，照录（见文件头“原卷笔误/存疑”⑤）
\item “$A\subseteq B$”是“$A\subset B$”的\mbox{（\quad）}条件

\keepwithprev
A. 充分不必要条件\qquad B. 必要不充分条件

C. 充要条件\qquad D. 既不充分也不必要条件

\ifanswer
\ans{$B$}
\fi

\item 已知 $a$、$b$、$c$ 均为实数，若 $a>b$ 且 $ab\neq0$，则下列不等关系中成立的是\mbox{（\quad）}

\keepwithprev
A. $ac^2>bc^2$\qquad B. $\dfrac1b>\dfrac1a$\qquad C. $|a|>|b|$\qquad D. $c-b>c-a$

\ifanswer
\solbegin
\step{对于 $A$，若 $a>b$，当 $c=0$ 时，有 $ac^2=bc^2=0$，故 $A$ 错误；}
\step{对于 $B$，$\dfrac1b-\dfrac1a=\dfrac{a-b}{ab}$，因为 $a>b$，所以 $a-b>0$，}
\step{当 $a$、$b$ 同号时，$ab>0$，则 $\dfrac1b-\dfrac1a>0$，即 $\dfrac1b>\dfrac1a$，}
\step{当 $a$、$b$ 异号时，$ab<0$，则 $\dfrac1b-\dfrac1a<0$，即 $\dfrac1b<\dfrac1a$，故 $B$ 错误；}
\step{对于 $C$，当 $a=1$，$b=-3$ 时，$|a|<|b|$，故 $C$ 错误；}
\step{对于 $D$，若 $a>b$，则 $-b>-a$，$c-b>c-a$，故 $D$ 正确.}
\step{故选 $D$.}
\solend
\fi

% 原卷解析对 (1) 未讨论取等是否可得（x=y 时确能取 2），却总结“有 1 个、故选 B”，
% 与推理矛盾，照录（见文件头“原卷笔误/存疑”③）
\item 以下三个代数式：（1）$\dfrac yx+\dfrac xy$；（2）$|x|+\dfrac1{|x|}$；（3）$x^2+2+\dfrac1{x^2+2}$．其中的字母均表示实数，且代数式均有意义，那么代数式的最小值为 2 的个数是\mbox{（\quad）}

\keepwithprev
A. $0$\qquad B. $1$\qquad C. $2$\qquad D. $3$

\ifanswer
\solbegin
\step{（1）仅当 $x,y$ 同号时，$\dfrac yx,\dfrac xy>0$，由基本不等式得 $\dfrac yx+\dfrac xy\geqslant2$；}
\step{若 $x,y$ 异号，$\dfrac yx,\dfrac xy<0$，故 $\dfrac yx+\dfrac xy<0$，此时最小值不是 2；}
\step{（2）$|x|+\dfrac1{|x|}$ 中，$|x|>0$，由基本不等式 $|x|+\dfrac1{|x|}\geqslant2\sqrt{|x|\cdot\dfrac1{|x|}}=2$，}
\step{当且仅当 $|x|=1$，即 $x=\pm1$ 时取等号，故最小值为 2；}
\step{（3）令 $t=x^2+2$，则 $t\geqslant2$，函数化为 $t+\dfrac1t$，}
\step{对勾函数 $t+\dfrac1t$ 在 $[2,+\infty)$ 上严格增，最小值为 $2+\dfrac12=\dfrac52>2$，}
\step{等号成立条件 $t=1$ 无法取到，故最小值不是 2；}
\step{综上，最小值为 2 的代数式有 1 个，故选 $B$.}
\solend
\fi

\item 已知 $A$、$B$ 为非空实数集，$\Omega$ 为平面直角坐标系中的一些点构成的集合，集合 $C=\{y\mid$对任意 $x\in A$，有 $(x,y)\in\Omega\}$，集合 $D=\{x\mid$对任意 $y\in B$，有 $(x,y)\in\Omega\}$，对于下列两个命题：①若 $C\subseteq B$，则 $D\subseteq A$；②若 $C\supseteq B$，则 $D\subseteq A$；其中判断正确的是\mbox{（\quad）}

\keepwithprev
A. ①②都正确\qquad B. ①②都错误\qquad C. ①正确②错误\qquad D. ①错误②正确

\ifanswer
\solbegin
\step{设 $A=\{a_1,a_2,a_3,\cdots,a_n\},B=\{b_1,b_2,b_3,\cdots,b_n\},n>2$，}
\step{若 $\Omega=\{(c,b_1),(c,b_2),\cdots,(c,b_n)\},c\notin A$，}
\step{此时 $C=\varnothing\subseteq B,D=\{c\},D\subseteq A$ 错误，故①错误；}
% 原卷此处 N* 排正体（\mathrm{N}*），本稿统一 \N^*（见文件头“排版差异”④）；
% 全称量词“∀”照原卷符号排（见文件头“排版差异”⑫）
\step{若 $C\supseteq B$，则 $\forall(a_i,b_j)(1\leqslant i,j\leqslant n,i,j\in\N^*)\in\Omega$，则 $D\supseteq A$，}
\step{且 $A=\{a_1\},B=\{b_1,b_2\}$，}
\step{若 $\Omega=\{(a_1,b_1),(a_2,b_1),(a_1,b_2),(a_2,b_2),(c,b_1),(c,b_2)\},D$ 真包含 $A$，故②错误.}
\step{故选 $B$.}
\solend
\fi

\end{enumerate}

\bigsec{三、解答题}

\begin{enumerate}
\setcounter{enumi}{16}

\item 解下列关于 $x$ 的不等式组：$\begin{cases}|x-1|<2\\ \dfrac{x-2}{x-1}\geqslant0\end{cases}$.

\ifanswer
\solbegin
\step{解 $|x-1|<2$ 得 $-2<x-1<2$，解得 $-1<x<3$.}
\step{分式不等式 $\dfrac{x-2}{x-1}\geqslant0$ 等价于 $\begin{cases}(x-2)(x-1)\geqslant0\\ x-1\neq0\end{cases}$，解得 $x<1$ 或 $x\geqslant2$.}
\step{即 $\begin{cases}|x-1|<2\\ \dfrac{x-2}{x-1}\geqslant0\end{cases}$ 的解集为 $(-1,1)\cup[2,3)$.}
\solend
\fi

\ansspace[5]

\item 比较多项式 $2x^2-3x+1$ 与 $x^2-1$ 的大小.

\ifanswer
\solbegin
\step{作差得 $(2x^2-3x+1)-(x^2-1)=x^2-3x+2=(x-1)(x-2)$，}
\step{当 $x=1$ 或 $x=2$ 时，$(x-1)(x-2)=0$，即 $2x^2-3x+1=x^2-1$，}
\step{当 $x<1$ 或 $x>2$ 时，$(x-1)(x-2)>0$，即 $2x^2-3x+1>x^2-1$，}
\step{当 $1<x<2$ 时，$(x-1)(x-2)<0$，即 $2x^2-3x+1<x^2-1$；}
\step{综上，当 $x=1$ 或 $x=2$ 时，$2x^2-3x+1=x^2-1$；}
\step{当 $x<1$ 或 $x>2$ 时，$2x^2-3x+1>x^2-1$；}
\step{当 $1<x<2$ 时，$2x^2-3x+1<x^2-1$.}
\solend
\fi

\ansspace[6]

% 原卷 (1) 连用两个“若”（“若 a=4，若全集为实数集 R”），照录（见文件头“原卷笔误/存疑”①）
\item 已知集合 $A=\left\{x\mid\dfrac{(a+1)x}{x-1}<a\right\}$，集合 $B=[1,4]$，$a$ 为实常数；

（1）若 $a=4$，若全集为实数集 $\R$，求集合 $\overline{A}$（用区间形式表示）；

（2）若 $A\subseteq B$，求 $a$ 的取值范围.

\ifanswer
\solbegin
\stepm{(1)}{若 $a=4$，集合 $A=\left\{x\mid\dfrac{5x}{x-1}<4\right\}=\left\{x\mid\dfrac{5x}{x-1}-4<0\right\}=\left\{x\mid\dfrac{x+4}{x-1}<0\right\}$}
\step{$=(-4,1)$，因此 $\overline{A}=(-\infty,-4]\cup[1,+\infty)$；}
\stepm{(2)}{当 $-a<1$，即 $a>-1$ 时，$A=(-a,1)$，}
\step{要使 $A\subseteq[1,4]$，需 $-a\geqslant1$，即 $a\leqslant-1$，与 $a>-1$ 矛盾，舍去；}
\step{当 $-a>1$，即 $a<-1$ 时，$A=(1,-a)$，}
\step{要使 $A\subseteq[1,4]$，需 $-a\leqslant4$，即 $a\geqslant-4$，所以 $-4\leqslant a<-1$；}
\step{当 $-a=1$，即 $a=-1$ 时，$A=\varnothing$，满足 $A\subseteq B$；}
\step{综上，实数 $a$ 的取值范围为 $[-4,-1]$.}
\solend
\fi

\ansspace[7]

% 原卷题干“可以售出 8 万本、据市场调查”处用顿号，照录（见文件头“原卷笔误/存疑”②）
\item 某种杂志原以每本 25 元的价格销售，可以售出 8 万本、据市场调查，杂志的单价每提高 1 元，销售量就减少 0.2 万本，假定杂志的成本是每本 10 元（不计其他成本）.

（1）当杂志以每本 30 元定价销售时，求销售该杂志所获利润；

（2）在每本 25 元的价格的基础上提高定价，试确定杂志的定价在什么范围内可使得销售利润不低于 140 万元.

\ifanswer
\solbegin
\stepm{(1)}{当杂志以 30 元定价销售时，}
\step{销售该杂志所获利润为 $[8-(30-25)\times0.2]\times(30-10)=140$ 万元；}
\stepm{(2)}{在每本 25 元的价格的基础上提高定价，设杂志的定价为每本 $x$ 元，}
\step{由题意得 $[8-(x-25)\times0.2]\times(x-10)\geqslant140$，解得 $30\leqslant x\leqslant45$，}
\step{所以杂志的定价在每本不低于 30 元且不高于 45 元的范围内，}
\step{可使得销售利润不低于 140 万元.}
\solend
\fi

\ansspace[7]

% 原卷题干“定义：对于集合 A，则称集合 B=…”句式杂糅，照录（见文件头“原卷笔误/存疑”④）
\item 定义：对于集合 $A$，则称集合 $B=\left\{x\mid x=\dfrac{m+n}{1+mn},m,n\in A,1+mn\neq0\right\}$ 为集合 $A$ 的生成集（其中 $m$、$n$ 可以相等）.

（1）若集合 $A=\{-1,1\}$ 的生成集为 $B$，求集合 $B$；

（2）若集合 $A=\{a,2\}$ 的生成集为 $B$，$B$ 的子集个数为 4 个，求实数 $a$ 的值；

（3）若集合 $A=\{x\mid-1<x<1\}$ 的生成集为 $B$，求证：$B\subseteq A$.

\ifanswer
\solbegin
\stepm{(1)}{因为 $1+mn\neq0$，}
\step{当 $m=n=1$ 时，$x=\dfrac{m+n}{1+mn}=1$；当 $m=n=-1$ 时，$x=\dfrac{m+n}{1+mn}=-1$；}
\step{所以 $B=\{-1,1\}$；}
\stepm{(2)}{$B$ 的子集个数为 4 个，故 $B$ 集合的元素有 2 个，}
\step{当 $m=n=2$ 时，$x=\dfrac{m+n}{1+mn}=\dfrac45$，当 $m=n=a$ 时，$x=\dfrac{m+n}{1+mn}=\dfrac{2a}{1+a^2}$，}
\step{当 $m$ 与 $n$ 分别取 2 和 $a$ 时，$x=\dfrac{a+2}{1+2a}$，}
\stepm{①}{$\dfrac{2a}{1+a^2}=\dfrac45$，解得 $a=2$（舍）或 $a=\dfrac12$，}
\stepm{②}{$\dfrac{a+2}{1+2a}=\dfrac45$，解得 $a=2$（舍），}
\stepm{③}{$\dfrac{2a}{1+a^2}=\dfrac{a+2}{1+2a}$，解得 $a=2$（舍）或 $a=1$ 或 $a=-1$，}
\step{综上所述，$a=1$ 或 $a=-1$ 或 $a=\dfrac12$，即 $a\in\left\{-1,1,\dfrac12\right\}$；}
\stepm{(3)}{$\forall m$、$n\in A$，则 $\dfrac{m+n}{1+mn}+1=\dfrac{(m+1)(n+1)}{1+mn}>0$，故 $\dfrac{m+n}{1+mn}>-1$；}
\step{$\dfrac{m+n}{1+mn}-1=-\dfrac{(m-1)(n-1)}{1+mn}<0$，故 $\dfrac{m+n}{1+mn}<1$；}
\step{故 $B\subseteq\{x\mid-1<x<1\}$，即 $B\subseteq A$.}
\solend
\fi

\ansspace*[10]

\end{enumerate}

\end{document}
"
% 紧凑版：xelatex "\def\iscompact{1}% 高一71_中国中学_10月月考.tex
%   —— 高一上数学“2026-2027 年中国中学高一上 10 月月考”（讲评版）
% 试卷版：xelatex 高一71_中国中学_10月月考.tex
% 答案版：xelatex "\def\isanswer{1}% 高一71_中国中学_10月月考.tex
%   —— 高一上数学“2026-2027 年中国中学高一上 10 月月考”（讲评版）
% 试卷版：xelatex 高一71_中国中学_10月月考.tex
% 答案版：xelatex "\def\isanswer{1}\input{高一71_中国中学_10月月考.tex}"
% 紧凑版：xelatex "\def\iscompact{1}\input{高一71_中国中学_10月月考.tex}"
%
% 原始材料：微信公众号“魔都数学加油站”文章，连载“27学年高一71”；
% 原文标题“27学年高一71——2026-2027年上海市中国中学高一上10月月考”；
% 原文链接：https://mp.weixin.qq.com/s/UZrP53l3GIin040L1EkljA。
% 页面标注发布时间 2026-10-19 22:47（上海）（与 高一69、高一70 页面读出同一 ct 值）。
% 正文为 7 张 PNG 页（`原图/高一71_中国中学_10月月考/01..07.png`），全部 4762×6735，
% 均按页序读取。7 页全为试卷正文页（卷首标题在第 1 页页顶，卷末解答题第 21 题 (3)
% 解析止于第 7 页），无封面页、目录页、二维码或宣传页。粉色斜排“魔都数学加油站”
% 为卷面水印，与公众号页脚一并未录入。
%
% 卷首标题：照卷面印作“2026-2027 年中国中学高一上 10 月月考”（黑体加粗一行）。
% 原卷年份区间用半角连字符，本稿按全稿惯例排 en dash（同 高一68、高一69、高一70）；
% 原卷只有一行标题、无副标题。
%
% 篇幅与题号：原卷自第 3 题起编号（第 1、2 题不在本卷内）。填空题 3—12（10 题）、
% 选择题 13—16（4 题）、解答题 17—21（5 题），共 19 题。本稿沿用原节与原题号，
% 只改排为全稿统一的大节样式（原卷小节标题为“一、填空题”“二、选择题”“三、解答题”）。
%
% 讲评版：原卷每题下直接印【答案】或【解析】。**只印【答案】的 3 题**（第 3、4、5 题），
% 本稿只给 \ans{}；**第 13 题既无【答案】也无【解析】段，答案“B”直接印在题干括号内
% （“(  B  )”）**，按 高一37 第 13 题、高一63 第 14 题、高一67 第 15 题成例把括号留空、
% 答案另提为 \ans{$B$}。其余 15 题全部印【解析】（按全稿口径这类题不另提 \ans{}——
% 第 14—16 题的“故选 D／B／B”写在解析末句，亦照录在解析内）。
% 除第 13 题外，全卷选择题作答括号内均为空白（第 14、15、16 题）。
%
% 副标题：原卷未印副标题。本稿按“知识模块”口径标作“集合与常用逻辑用语、不等式”：
% 19 题中集合与常用逻辑用语 10 题（第 3、4、5、6、9、11、12、13、16、21 题；其中第 4 题
% 属反证法，一并计入常用逻辑用语），不等式 9 题（第 7、8、10、14、15、17、18、19、20 题；
% 第 19 题 (1) 属补集、(2) 属分式不等式，按主干计入不等式），两类均确有分量，故并标。
% 口径同 高一63、高一68、高一69、高一70。
%
% 与原卷的排版差异：
%   ① 原卷每题下直接印【答案】/【解析】，本稿按全稿惯例将答案与解析置于答案版，试卷版
%      保持干净；第 13 题印在题干括号内的答案“B”同样移入 \ans{}，见上；
%   ② 原卷未印副标题，本稿按“知识模块”口径补出；
%   ③ 原卷填空横线约四个字宽，本稿按全稿惯例统一排 \blankline[4.4em]；
%   ④ 原卷数集记号不统一：实数集排斜体 $R$（第 19 题 (1)），自然数集排斜体 $N$
%      （第 5 题题干、第 12 题题干与解析共三处），第 16 题解析又排正体 $\mathrm{N}^*$；
%      本稿按全稿惯例统一改用 $\R$、$\N$、$\N^*$（同 高一68／69／70 排版差异④）；
%      空集记号原卷作 $\varnothing$，照录；
%   ⑤ 原卷句末标点混用：第 15 题题干“（3）$x^2+2+\dfrac1{x^2+2}$**．**其中”、
%      第 17 题题干末“$\dfrac{x-2}{x-1}\geqslant0$**．**”两处放大确认（第 15 题为全角
%      句点、第 17 题为西文句点），其余各题末为西文句点，已逐处照录；题干内冒号——
%      第 5 题“$\{y\mid\cdots\}$**：**”、第 12 题““×”**：**”、第 17 题“不等式组**：**”、
%      第 21 题“定义**：**”——均为全角冒号，照录；数学式内的逗号（如第 12 题
%      “$x\in A,y\in B$”、第 16 题“$A=\{a_1,\cdots\},B=\{b_1,\cdots\},n>2$”）按原卷作
%      半角，中文行文里的逗号按原卷作全角，均未强改；
%   ⑥ 第 17—21 题的子问原卷印作“（1）…（2）…”（另起一段），本稿按 高一69／70 成例用
%      行内“（1）（2）（3）”段，解析中的子问标号一律用 \stepm{(1)}{…}；
%      第 21 题 (2) 解析内层的“①”“②”“③”照原卷作 \stepm{①}{…} 等；
%   ⑦ 每道解答题原卷未留作答空白，本稿按全稿惯例加 \ansspace（末题用 \ansspace*）；
%   ⑧ 选择题的作答括号原卷排半角“(　　)”（第 14、15、16 题括号内空白、第 13 题括号内
%      印 B），本稿按全稿惯例排 \mbox{（\quad）}；选项行排照原卷——第 13 题两行两条
%      （A B 一行、C D 一行），第 14、15、16 题一行四条；
%   ⑨ 第 8 题解析、第 9 题解析、第 17 题题干与解析中的多行不等式组，原卷均排作“{”
%      连接的多行式，本稿按全稿惯例排 cases；
%   ⑩ 第 10 题解析、第 17 题解析、第 19 题 (1) 解析的并集符号原卷排大号 $\bigcup$
%      （显示体），本稿按全稿惯例排行内 $\cup$；
%   ⑪ 第 8 题解析的判别式照原卷作空心“$\triangle$”（与第 7 题解析的“$\Delta$”不同形），
%      照录未改；
%   ⑫ 第 16 题解析第 4 行的“$\forall$”（全称量词）照原卷符号排，未改写为中文“任意”。
%
% 原卷笔误/存疑（均照录未改，正文对应处另有 % 注）：
%   ① 第 19 题 (1) 连用两个“若”——“若 $a=4$，若全集为实数集 $R$”，后一“若”应作
%      “当”或删；原卷如此，照录未改。
%   ② 第 20 题题干“可以售出 8 万本**、**据市场调查”——“万本”后用顿号，按文意应为
%      逗号或句号；原卷如此，照录未改。
%   ③ 第 15 题解析对代数式 (1) 只得出“$y/x+x/y\geqslant2$”（同号时），未讨论取等是否
%      可得——取 $x=y$ 时确能取到 2——却总结“最小值为 2 的代数式有 1 个，故选 $B$”；
%      按解析自身的推理 (1) 也应计入，结论“1 个”与推理矛盾（若计 (1) 应为 2 个、选 C）。
%      原卷如此，照录未改。
%   ④ 第 21 题题干“定义：对于集合 $A$，则称集合 $B=\cdots$”——“对于……”与“则称……”
%      句式杂糅（应作“对于集合 $A$，称……”或“若集合 $A$，则称……”）；原卷如此，照录未改。
%   ⑤ 第 13 题题干末“……的（　）条件”后无任何标点（其余各题题干末均有句点）；
%      原卷如此，照录未改。
%
% 行数核对：逐题数 \step/\stepm 条数与原卷【解析】印刷行数，15 题逐一相符，
% 合计 86 条 / 86 行（Q6：1；Q7：1；Q8：7；Q9：2；Q10：4；Q11：7；Q12：5；Q14：7；
% Q15：8；Q16：7；Q17：3；Q18：7；Q19：8；Q20：6；Q21：13）；另有 3 题只印【答案】
% （第 3、4、5 题）与第 13 题括号内答案，各提 \ans{}，不计入 \step 账。
%
% 插图：无题目配图或数学图形（7 页题面均无“如图/右图/图中”字样；页面另有粉色斜排水印，
% 见上）。
%
\documentclass[a4paper,11pt]{article}
\usepackage{common}

\begin{document}

\examtitlest[3]{2026--2027 年中国中学高一上 10 月月考}{集合与常用逻辑用语、不等式}

\bigsec{一、填空题}

\begin{enumerate}
\setcounter{enumi}{2}

\item 已知集合 $A=[1,+\infty)$，$B=[-1,3)$，则 $A\cap B=$\blankline[4.4em].

\ifanswer
\ans{$[1,3)$}
\fi

\item 用反证法证明命题“已知 $a$、$b$ 为正数，$n$ 为正整数，求证：如果 $a^n>b^n$，那么 $a>b$”时，应当作出的假设为：“假设\blankline[4.4em]”.

\ifanswer
\ans{$a\leqslant b$}
\fi

\item 用列举法表示集合 $\{y\mid y=x+1,x\in\N,x\leqslant1\}$：\blankline[4.4em].

\ifanswer
\ans{$\{1,2\}$}
\fi

\item 某班在一次测验中，有 36 人数学成绩不低于 80 分，有 20 人物理成绩不低于 80 分，有 15 人的数学、物理成绩都不低于 80 分，则这两科成绩中至少有一科不低于 80 分的人数为\blankline[4.4em]人.

\ifanswer
\solbegin
\step{至少有 $36+20-15=41$ 人.}
\solend
\fi

\item 若关于 $x$ 的不等式 $x^2-x+m<0$ 的解集是 $\varnothing$，则实数 $m$ 的取值范围是\blankline[4.4em].

\ifanswer
\solbegin
\step{由题意得 $\Delta=1-4m\leqslant0$，所以 $m\geqslant\dfrac14$.}
\solend
\fi

% 第 8 题解析的判别式原卷作空心“△”（与第 7 题解析的“Δ”不同形），照录（见文件头“排版差异”⑪）
\item 已知关于 $x$ 的方程 $x^2+mx+2m=0$ 的两个实数根分别为 $\alpha$，$\beta$，若 $\alpha^2+\beta^2=5$，则实数 $m=$\blankline[4.4em].

\ifanswer
\solbegin
\step{方程 $x^2+mx+2m=0$ 的两个实数根为 $\alpha$、$\beta$，}
\step{由根与系数的关系得 $\begin{cases}\alpha+\beta=-m\\ \alpha\beta=2m\end{cases}$，又因为 $\alpha^2+\beta^2=5$，}
\step{所以 $\alpha^2+\beta^2=(\alpha+\beta)^2-2\alpha\beta=m^2-4m=5$，整理得 $m^2-4m-5=0$，}
\step{解得 $m=5$ 或 $m=-1$，}
\step{当 $m=5$ 时，$\triangle=25-40=-15<0$，无实根，舍去；}
\step{当 $m=-1$ 时，$\triangle=1+8=9>0$，符合题意，}
\step{所以实数 $m=-1$.}
\solend
\fi

\item 已知集合 $A=\{x\mid6x+a>0\}$，若 $1\in A$，$0\notin A$，则实数 $a$ 的取值范围为\blankline[4.4em].

\ifanswer
\solbegin
\step{由 $A=\{x\mid6x+a>0\}$，$1\in A$，$0\notin A$，得 $\begin{cases}6+a>0\\ a\leqslant0\end{cases}$，解得 $-6<a\leqslant0$，}
\step{即实数 $a$ 的取值范围是 $(-6,0]$.}
\solend
\fi

\item 关于 $x$ 的不等式 $ax-b>0$ 的解集为 $(1,+\infty)$，则关于 $x$ 的不等式 $(ax+b)(x-2)>0$ 的解集为\blankline[4.4em].

\ifanswer
\solbegin
\step{由不等式 $ax-b>0$ 的解集为 $(1,+\infty)$，得 $a>0$ 且 $\dfrac ba=1$，即 $b=a$，}
\step{则 $(ax+b)(x-2)>0$，可化为 $a(x+1)(x-2)>0$.}
\step{即 $(x+1)(x-2)>0$，解得 $x<-1$ 或 $x>2$，}
\step{所以不等式的解集为 $(-\infty,-1)\cup(2,+\infty)$.}
\solend
\fi

\item 已知实数 $a_1$、$a_2$、$a_3$、$a_4$、$a_5$ 满足 $\{x\mid(x-a_1)(x-a_2)(x-a_3)(x-a_4)(x-a_5)=0\}=\{1,2,3,4\}$，则 $a_1+a_2+a_3+a_4+a_5$ 的最大值为\blankline[4.4em].

\ifanswer
\solbegin
\step{5 个实数 $a_1$、$a_2$、$a_3$、$a_4$、$a_5$ 的不同取值只有 1，2，3，4，}
\step{即这 5 个数均从 1，2，3，4 中选取.}
\step{集合有 4 个元素，说明 1，2，3，4 每个数至少出现一次，}
\step{因此恰好有一个数重复出现两次，其余三个数各出现一次.}
\step{要让 5 个数的总和最大，应让重复出现两次的数取集合中的最大值 4，}
\step{此时 5 个数为 1，2，3，4，4，总和为 $1+2+3+4+4=14$，}
\step{即所求最大值为 14.}
\solend
\fi

\item 定义集合运算“$\times$”：$A\times B=\{(x,y)\mid x\in A,y\in B\}$，称为 $A$、$B$ 两个集合的“卡氏积”，若 $A=\{x\mid x^2-|x|\leqslant0,x\in\N\}$，$B=\{1,2,3\}$，则 $(A\times B)\cap(B\times A)=$\blankline[4.4em].

\ifanswer
\solbegin
\step{因为 $x^2-|x|\leqslant0$，所以 $0\leqslant|x|\leqslant1$，又因为 $x\in\N$，}
\step{所以 $A=\{x\mid x^2-|x|\leqslant0,x\in\N\}=\{0,1\}$，又 $B=\{1,2,3\}$，}
\step{由“卡氏积”的定义得 $A\times B=\{(0,1),(0,2),(0,3),(1,1),(1,2),(1,3)\}$，}
\step{$B\times A=\{(1,0),(1,1),(2,0),(2,1),(3,0),(3,1)\}$，}
\step{所以 $(A\times B)\cap(B\times A)=\{(1,1)\}$.}
\solend
\fi

\end{enumerate}

\bigsec{二、选择题}

\begin{enumerate}
\setcounter{enumi}{12}

% 第 13 题答案“B”原卷印在题干括号内，且无【答案】/【解析】段；
% 按 高一37 第 13 题、高一63 第 14 题、高一67 第 15 题成例括号留空、答案另提 \ans{$B$}
% 原卷题干末“条件”后无标点，照录（见文件头“原卷笔误/存疑”⑤）
\item “$A\subseteq B$”是“$A\subset B$”的\mbox{（\quad）}条件

\keepwithprev
A. 充分不必要条件\qquad B. 必要不充分条件

C. 充要条件\qquad D. 既不充分也不必要条件

\ifanswer
\ans{$B$}
\fi

\item 已知 $a$、$b$、$c$ 均为实数，若 $a>b$ 且 $ab\neq0$，则下列不等关系中成立的是\mbox{（\quad）}

\keepwithprev
A. $ac^2>bc^2$\qquad B. $\dfrac1b>\dfrac1a$\qquad C. $|a|>|b|$\qquad D. $c-b>c-a$

\ifanswer
\solbegin
\step{对于 $A$，若 $a>b$，当 $c=0$ 时，有 $ac^2=bc^2=0$，故 $A$ 错误；}
\step{对于 $B$，$\dfrac1b-\dfrac1a=\dfrac{a-b}{ab}$，因为 $a>b$，所以 $a-b>0$，}
\step{当 $a$、$b$ 同号时，$ab>0$，则 $\dfrac1b-\dfrac1a>0$，即 $\dfrac1b>\dfrac1a$，}
\step{当 $a$、$b$ 异号时，$ab<0$，则 $\dfrac1b-\dfrac1a<0$，即 $\dfrac1b<\dfrac1a$，故 $B$ 错误；}
\step{对于 $C$，当 $a=1$，$b=-3$ 时，$|a|<|b|$，故 $C$ 错误；}
\step{对于 $D$，若 $a>b$，则 $-b>-a$，$c-b>c-a$，故 $D$ 正确.}
\step{故选 $D$.}
\solend
\fi

% 原卷解析对 (1) 未讨论取等是否可得（x=y 时确能取 2），却总结“有 1 个、故选 B”，
% 与推理矛盾，照录（见文件头“原卷笔误/存疑”③）
\item 以下三个代数式：（1）$\dfrac yx+\dfrac xy$；（2）$|x|+\dfrac1{|x|}$；（3）$x^2+2+\dfrac1{x^2+2}$．其中的字母均表示实数，且代数式均有意义，那么代数式的最小值为 2 的个数是\mbox{（\quad）}

\keepwithprev
A. $0$\qquad B. $1$\qquad C. $2$\qquad D. $3$

\ifanswer
\solbegin
\step{（1）仅当 $x,y$ 同号时，$\dfrac yx,\dfrac xy>0$，由基本不等式得 $\dfrac yx+\dfrac xy\geqslant2$；}
\step{若 $x,y$ 异号，$\dfrac yx,\dfrac xy<0$，故 $\dfrac yx+\dfrac xy<0$，此时最小值不是 2；}
\step{（2）$|x|+\dfrac1{|x|}$ 中，$|x|>0$，由基本不等式 $|x|+\dfrac1{|x|}\geqslant2\sqrt{|x|\cdot\dfrac1{|x|}}=2$，}
\step{当且仅当 $|x|=1$，即 $x=\pm1$ 时取等号，故最小值为 2；}
\step{（3）令 $t=x^2+2$，则 $t\geqslant2$，函数化为 $t+\dfrac1t$，}
\step{对勾函数 $t+\dfrac1t$ 在 $[2,+\infty)$ 上严格增，最小值为 $2+\dfrac12=\dfrac52>2$，}
\step{等号成立条件 $t=1$ 无法取到，故最小值不是 2；}
\step{综上，最小值为 2 的代数式有 1 个，故选 $B$.}
\solend
\fi

\item 已知 $A$、$B$ 为非空实数集，$\Omega$ 为平面直角坐标系中的一些点构成的集合，集合 $C=\{y\mid$对任意 $x\in A$，有 $(x,y)\in\Omega\}$，集合 $D=\{x\mid$对任意 $y\in B$，有 $(x,y)\in\Omega\}$，对于下列两个命题：①若 $C\subseteq B$，则 $D\subseteq A$；②若 $C\supseteq B$，则 $D\subseteq A$；其中判断正确的是\mbox{（\quad）}

\keepwithprev
A. ①②都正确\qquad B. ①②都错误\qquad C. ①正确②错误\qquad D. ①错误②正确

\ifanswer
\solbegin
\step{设 $A=\{a_1,a_2,a_3,\cdots,a_n\},B=\{b_1,b_2,b_3,\cdots,b_n\},n>2$，}
\step{若 $\Omega=\{(c,b_1),(c,b_2),\cdots,(c,b_n)\},c\notin A$，}
\step{此时 $C=\varnothing\subseteq B,D=\{c\},D\subseteq A$ 错误，故①错误；}
% 原卷此处 N* 排正体（\mathrm{N}*），本稿统一 \N^*（见文件头“排版差异”④）；
% 全称量词“∀”照原卷符号排（见文件头“排版差异”⑫）
\step{若 $C\supseteq B$，则 $\forall(a_i,b_j)(1\leqslant i,j\leqslant n,i,j\in\N^*)\in\Omega$，则 $D\supseteq A$，}
\step{且 $A=\{a_1\},B=\{b_1,b_2\}$，}
\step{若 $\Omega=\{(a_1,b_1),(a_2,b_1),(a_1,b_2),(a_2,b_2),(c,b_1),(c,b_2)\},D$ 真包含 $A$，故②错误.}
\step{故选 $B$.}
\solend
\fi

\end{enumerate}

\bigsec{三、解答题}

\begin{enumerate}
\setcounter{enumi}{16}

\item 解下列关于 $x$ 的不等式组：$\begin{cases}|x-1|<2\\ \dfrac{x-2}{x-1}\geqslant0\end{cases}$.

\ifanswer
\solbegin
\step{解 $|x-1|<2$ 得 $-2<x-1<2$，解得 $-1<x<3$.}
\step{分式不等式 $\dfrac{x-2}{x-1}\geqslant0$ 等价于 $\begin{cases}(x-2)(x-1)\geqslant0\\ x-1\neq0\end{cases}$，解得 $x<1$ 或 $x\geqslant2$.}
\step{即 $\begin{cases}|x-1|<2\\ \dfrac{x-2}{x-1}\geqslant0\end{cases}$ 的解集为 $(-1,1)\cup[2,3)$.}
\solend
\fi

\ansspace[5]

\item 比较多项式 $2x^2-3x+1$ 与 $x^2-1$ 的大小.

\ifanswer
\solbegin
\step{作差得 $(2x^2-3x+1)-(x^2-1)=x^2-3x+2=(x-1)(x-2)$，}
\step{当 $x=1$ 或 $x=2$ 时，$(x-1)(x-2)=0$，即 $2x^2-3x+1=x^2-1$，}
\step{当 $x<1$ 或 $x>2$ 时，$(x-1)(x-2)>0$，即 $2x^2-3x+1>x^2-1$，}
\step{当 $1<x<2$ 时，$(x-1)(x-2)<0$，即 $2x^2-3x+1<x^2-1$；}
\step{综上，当 $x=1$ 或 $x=2$ 时，$2x^2-3x+1=x^2-1$；}
\step{当 $x<1$ 或 $x>2$ 时，$2x^2-3x+1>x^2-1$；}
\step{当 $1<x<2$ 时，$2x^2-3x+1<x^2-1$.}
\solend
\fi

\ansspace[6]

% 原卷 (1) 连用两个“若”（“若 a=4，若全集为实数集 R”），照录（见文件头“原卷笔误/存疑”①）
\item 已知集合 $A=\left\{x\mid\dfrac{(a+1)x}{x-1}<a\right\}$，集合 $B=[1,4]$，$a$ 为实常数；

（1）若 $a=4$，若全集为实数集 $\R$，求集合 $\overline{A}$（用区间形式表示）；

（2）若 $A\subseteq B$，求 $a$ 的取值范围.

\ifanswer
\solbegin
\stepm{(1)}{若 $a=4$，集合 $A=\left\{x\mid\dfrac{5x}{x-1}<4\right\}=\left\{x\mid\dfrac{5x}{x-1}-4<0\right\}=\left\{x\mid\dfrac{x+4}{x-1}<0\right\}$}
\step{$=(-4,1)$，因此 $\overline{A}=(-\infty,-4]\cup[1,+\infty)$；}
\stepm{(2)}{当 $-a<1$，即 $a>-1$ 时，$A=(-a,1)$，}
\step{要使 $A\subseteq[1,4]$，需 $-a\geqslant1$，即 $a\leqslant-1$，与 $a>-1$ 矛盾，舍去；}
\step{当 $-a>1$，即 $a<-1$ 时，$A=(1,-a)$，}
\step{要使 $A\subseteq[1,4]$，需 $-a\leqslant4$，即 $a\geqslant-4$，所以 $-4\leqslant a<-1$；}
\step{当 $-a=1$，即 $a=-1$ 时，$A=\varnothing$，满足 $A\subseteq B$；}
\step{综上，实数 $a$ 的取值范围为 $[-4,-1]$.}
\solend
\fi

\ansspace[7]

% 原卷题干“可以售出 8 万本、据市场调查”处用顿号，照录（见文件头“原卷笔误/存疑”②）
\item 某种杂志原以每本 25 元的价格销售，可以售出 8 万本、据市场调查，杂志的单价每提高 1 元，销售量就减少 0.2 万本，假定杂志的成本是每本 10 元（不计其他成本）.

（1）当杂志以每本 30 元定价销售时，求销售该杂志所获利润；

（2）在每本 25 元的价格的基础上提高定价，试确定杂志的定价在什么范围内可使得销售利润不低于 140 万元.

\ifanswer
\solbegin
\stepm{(1)}{当杂志以 30 元定价销售时，}
\step{销售该杂志所获利润为 $[8-(30-25)\times0.2]\times(30-10)=140$ 万元；}
\stepm{(2)}{在每本 25 元的价格的基础上提高定价，设杂志的定价为每本 $x$ 元，}
\step{由题意得 $[8-(x-25)\times0.2]\times(x-10)\geqslant140$，解得 $30\leqslant x\leqslant45$，}
\step{所以杂志的定价在每本不低于 30 元且不高于 45 元的范围内，}
\step{可使得销售利润不低于 140 万元.}
\solend
\fi

\ansspace[7]

% 原卷题干“定义：对于集合 A，则称集合 B=…”句式杂糅，照录（见文件头“原卷笔误/存疑”④）
\item 定义：对于集合 $A$，则称集合 $B=\left\{x\mid x=\dfrac{m+n}{1+mn},m,n\in A,1+mn\neq0\right\}$ 为集合 $A$ 的生成集（其中 $m$、$n$ 可以相等）.

（1）若集合 $A=\{-1,1\}$ 的生成集为 $B$，求集合 $B$；

（2）若集合 $A=\{a,2\}$ 的生成集为 $B$，$B$ 的子集个数为 4 个，求实数 $a$ 的值；

（3）若集合 $A=\{x\mid-1<x<1\}$ 的生成集为 $B$，求证：$B\subseteq A$.

\ifanswer
\solbegin
\stepm{(1)}{因为 $1+mn\neq0$，}
\step{当 $m=n=1$ 时，$x=\dfrac{m+n}{1+mn}=1$；当 $m=n=-1$ 时，$x=\dfrac{m+n}{1+mn}=-1$；}
\step{所以 $B=\{-1,1\}$；}
\stepm{(2)}{$B$ 的子集个数为 4 个，故 $B$ 集合的元素有 2 个，}
\step{当 $m=n=2$ 时，$x=\dfrac{m+n}{1+mn}=\dfrac45$，当 $m=n=a$ 时，$x=\dfrac{m+n}{1+mn}=\dfrac{2a}{1+a^2}$，}
\step{当 $m$ 与 $n$ 分别取 2 和 $a$ 时，$x=\dfrac{a+2}{1+2a}$，}
\stepm{①}{$\dfrac{2a}{1+a^2}=\dfrac45$，解得 $a=2$（舍）或 $a=\dfrac12$，}
\stepm{②}{$\dfrac{a+2}{1+2a}=\dfrac45$，解得 $a=2$（舍），}
\stepm{③}{$\dfrac{2a}{1+a^2}=\dfrac{a+2}{1+2a}$，解得 $a=2$（舍）或 $a=1$ 或 $a=-1$，}
\step{综上所述，$a=1$ 或 $a=-1$ 或 $a=\dfrac12$，即 $a\in\left\{-1,1,\dfrac12\right\}$；}
\stepm{(3)}{$\forall m$、$n\in A$，则 $\dfrac{m+n}{1+mn}+1=\dfrac{(m+1)(n+1)}{1+mn}>0$，故 $\dfrac{m+n}{1+mn}>-1$；}
\step{$\dfrac{m+n}{1+mn}-1=-\dfrac{(m-1)(n-1)}{1+mn}<0$，故 $\dfrac{m+n}{1+mn}<1$；}
\step{故 $B\subseteq\{x\mid-1<x<1\}$，即 $B\subseteq A$.}
\solend
\fi

\ansspace*[10]

\end{enumerate}

\end{document}
"
% 紧凑版：xelatex "\def\iscompact{1}% 高一71_中国中学_10月月考.tex
%   —— 高一上数学“2026-2027 年中国中学高一上 10 月月考”（讲评版）
% 试卷版：xelatex 高一71_中国中学_10月月考.tex
% 答案版：xelatex "\def\isanswer{1}\input{高一71_中国中学_10月月考.tex}"
% 紧凑版：xelatex "\def\iscompact{1}\input{高一71_中国中学_10月月考.tex}"
%
% 原始材料：微信公众号“魔都数学加油站”文章，连载“27学年高一71”；
% 原文标题“27学年高一71——2026-2027年上海市中国中学高一上10月月考”；
% 原文链接：https://mp.weixin.qq.com/s/UZrP53l3GIin040L1EkljA。
% 页面标注发布时间 2026-10-19 22:47（上海）（与 高一69、高一70 页面读出同一 ct 值）。
% 正文为 7 张 PNG 页（`原图/高一71_中国中学_10月月考/01..07.png`），全部 4762×6735，
% 均按页序读取。7 页全为试卷正文页（卷首标题在第 1 页页顶，卷末解答题第 21 题 (3)
% 解析止于第 7 页），无封面页、目录页、二维码或宣传页。粉色斜排“魔都数学加油站”
% 为卷面水印，与公众号页脚一并未录入。
%
% 卷首标题：照卷面印作“2026-2027 年中国中学高一上 10 月月考”（黑体加粗一行）。
% 原卷年份区间用半角连字符，本稿按全稿惯例排 en dash（同 高一68、高一69、高一70）；
% 原卷只有一行标题、无副标题。
%
% 篇幅与题号：原卷自第 3 题起编号（第 1、2 题不在本卷内）。填空题 3—12（10 题）、
% 选择题 13—16（4 题）、解答题 17—21（5 题），共 19 题。本稿沿用原节与原题号，
% 只改排为全稿统一的大节样式（原卷小节标题为“一、填空题”“二、选择题”“三、解答题”）。
%
% 讲评版：原卷每题下直接印【答案】或【解析】。**只印【答案】的 3 题**（第 3、4、5 题），
% 本稿只给 \ans{}；**第 13 题既无【答案】也无【解析】段，答案“B”直接印在题干括号内
% （“(  B  )”）**，按 高一37 第 13 题、高一63 第 14 题、高一67 第 15 题成例把括号留空、
% 答案另提为 \ans{$B$}。其余 15 题全部印【解析】（按全稿口径这类题不另提 \ans{}——
% 第 14—16 题的“故选 D／B／B”写在解析末句，亦照录在解析内）。
% 除第 13 题外，全卷选择题作答括号内均为空白（第 14、15、16 题）。
%
% 副标题：原卷未印副标题。本稿按“知识模块”口径标作“集合与常用逻辑用语、不等式”：
% 19 题中集合与常用逻辑用语 10 题（第 3、4、5、6、9、11、12、13、16、21 题；其中第 4 题
% 属反证法，一并计入常用逻辑用语），不等式 9 题（第 7、8、10、14、15、17、18、19、20 题；
% 第 19 题 (1) 属补集、(2) 属分式不等式，按主干计入不等式），两类均确有分量，故并标。
% 口径同 高一63、高一68、高一69、高一70。
%
% 与原卷的排版差异：
%   ① 原卷每题下直接印【答案】/【解析】，本稿按全稿惯例将答案与解析置于答案版，试卷版
%      保持干净；第 13 题印在题干括号内的答案“B”同样移入 \ans{}，见上；
%   ② 原卷未印副标题，本稿按“知识模块”口径补出；
%   ③ 原卷填空横线约四个字宽，本稿按全稿惯例统一排 \blankline[4.4em]；
%   ④ 原卷数集记号不统一：实数集排斜体 $R$（第 19 题 (1)），自然数集排斜体 $N$
%      （第 5 题题干、第 12 题题干与解析共三处），第 16 题解析又排正体 $\mathrm{N}^*$；
%      本稿按全稿惯例统一改用 $\R$、$\N$、$\N^*$（同 高一68／69／70 排版差异④）；
%      空集记号原卷作 $\varnothing$，照录；
%   ⑤ 原卷句末标点混用：第 15 题题干“（3）$x^2+2+\dfrac1{x^2+2}$**．**其中”、
%      第 17 题题干末“$\dfrac{x-2}{x-1}\geqslant0$**．**”两处放大确认（第 15 题为全角
%      句点、第 17 题为西文句点），其余各题末为西文句点，已逐处照录；题干内冒号——
%      第 5 题“$\{y\mid\cdots\}$**：**”、第 12 题““×”**：**”、第 17 题“不等式组**：**”、
%      第 21 题“定义**：**”——均为全角冒号，照录；数学式内的逗号（如第 12 题
%      “$x\in A,y\in B$”、第 16 题“$A=\{a_1,\cdots\},B=\{b_1,\cdots\},n>2$”）按原卷作
%      半角，中文行文里的逗号按原卷作全角，均未强改；
%   ⑥ 第 17—21 题的子问原卷印作“（1）…（2）…”（另起一段），本稿按 高一69／70 成例用
%      行内“（1）（2）（3）”段，解析中的子问标号一律用 \stepm{(1)}{…}；
%      第 21 题 (2) 解析内层的“①”“②”“③”照原卷作 \stepm{①}{…} 等；
%   ⑦ 每道解答题原卷未留作答空白，本稿按全稿惯例加 \ansspace（末题用 \ansspace*）；
%   ⑧ 选择题的作答括号原卷排半角“(　　)”（第 14、15、16 题括号内空白、第 13 题括号内
%      印 B），本稿按全稿惯例排 \mbox{（\quad）}；选项行排照原卷——第 13 题两行两条
%      （A B 一行、C D 一行），第 14、15、16 题一行四条；
%   ⑨ 第 8 题解析、第 9 题解析、第 17 题题干与解析中的多行不等式组，原卷均排作“{”
%      连接的多行式，本稿按全稿惯例排 cases；
%   ⑩ 第 10 题解析、第 17 题解析、第 19 题 (1) 解析的并集符号原卷排大号 $\bigcup$
%      （显示体），本稿按全稿惯例排行内 $\cup$；
%   ⑪ 第 8 题解析的判别式照原卷作空心“$\triangle$”（与第 7 题解析的“$\Delta$”不同形），
%      照录未改；
%   ⑫ 第 16 题解析第 4 行的“$\forall$”（全称量词）照原卷符号排，未改写为中文“任意”。
%
% 原卷笔误/存疑（均照录未改，正文对应处另有 % 注）：
%   ① 第 19 题 (1) 连用两个“若”——“若 $a=4$，若全集为实数集 $R$”，后一“若”应作
%      “当”或删；原卷如此，照录未改。
%   ② 第 20 题题干“可以售出 8 万本**、**据市场调查”——“万本”后用顿号，按文意应为
%      逗号或句号；原卷如此，照录未改。
%   ③ 第 15 题解析对代数式 (1) 只得出“$y/x+x/y\geqslant2$”（同号时），未讨论取等是否
%      可得——取 $x=y$ 时确能取到 2——却总结“最小值为 2 的代数式有 1 个，故选 $B$”；
%      按解析自身的推理 (1) 也应计入，结论“1 个”与推理矛盾（若计 (1) 应为 2 个、选 C）。
%      原卷如此，照录未改。
%   ④ 第 21 题题干“定义：对于集合 $A$，则称集合 $B=\cdots$”——“对于……”与“则称……”
%      句式杂糅（应作“对于集合 $A$，称……”或“若集合 $A$，则称……”）；原卷如此，照录未改。
%   ⑤ 第 13 题题干末“……的（　）条件”后无任何标点（其余各题题干末均有句点）；
%      原卷如此，照录未改。
%
% 行数核对：逐题数 \step/\stepm 条数与原卷【解析】印刷行数，15 题逐一相符，
% 合计 86 条 / 86 行（Q6：1；Q7：1；Q8：7；Q9：2；Q10：4；Q11：7；Q12：5；Q14：7；
% Q15：8；Q16：7；Q17：3；Q18：7；Q19：8；Q20：6；Q21：13）；另有 3 题只印【答案】
% （第 3、4、5 题）与第 13 题括号内答案，各提 \ans{}，不计入 \step 账。
%
% 插图：无题目配图或数学图形（7 页题面均无“如图/右图/图中”字样；页面另有粉色斜排水印，
% 见上）。
%
\documentclass[a4paper,11pt]{article}
\usepackage{common}

\begin{document}

\examtitlest[3]{2026--2027 年中国中学高一上 10 月月考}{集合与常用逻辑用语、不等式}

\bigsec{一、填空题}

\begin{enumerate}
\setcounter{enumi}{2}

\item 已知集合 $A=[1,+\infty)$，$B=[-1,3)$，则 $A\cap B=$\blankline[4.4em].

\ifanswer
\ans{$[1,3)$}
\fi

\item 用反证法证明命题“已知 $a$、$b$ 为正数，$n$ 为正整数，求证：如果 $a^n>b^n$，那么 $a>b$”时，应当作出的假设为：“假设\blankline[4.4em]”.

\ifanswer
\ans{$a\leqslant b$}
\fi

\item 用列举法表示集合 $\{y\mid y=x+1,x\in\N,x\leqslant1\}$：\blankline[4.4em].

\ifanswer
\ans{$\{1,2\}$}
\fi

\item 某班在一次测验中，有 36 人数学成绩不低于 80 分，有 20 人物理成绩不低于 80 分，有 15 人的数学、物理成绩都不低于 80 分，则这两科成绩中至少有一科不低于 80 分的人数为\blankline[4.4em]人.

\ifanswer
\solbegin
\step{至少有 $36+20-15=41$ 人.}
\solend
\fi

\item 若关于 $x$ 的不等式 $x^2-x+m<0$ 的解集是 $\varnothing$，则实数 $m$ 的取值范围是\blankline[4.4em].

\ifanswer
\solbegin
\step{由题意得 $\Delta=1-4m\leqslant0$，所以 $m\geqslant\dfrac14$.}
\solend
\fi

% 第 8 题解析的判别式原卷作空心“△”（与第 7 题解析的“Δ”不同形），照录（见文件头“排版差异”⑪）
\item 已知关于 $x$ 的方程 $x^2+mx+2m=0$ 的两个实数根分别为 $\alpha$，$\beta$，若 $\alpha^2+\beta^2=5$，则实数 $m=$\blankline[4.4em].

\ifanswer
\solbegin
\step{方程 $x^2+mx+2m=0$ 的两个实数根为 $\alpha$、$\beta$，}
\step{由根与系数的关系得 $\begin{cases}\alpha+\beta=-m\\ \alpha\beta=2m\end{cases}$，又因为 $\alpha^2+\beta^2=5$，}
\step{所以 $\alpha^2+\beta^2=(\alpha+\beta)^2-2\alpha\beta=m^2-4m=5$，整理得 $m^2-4m-5=0$，}
\step{解得 $m=5$ 或 $m=-1$，}
\step{当 $m=5$ 时，$\triangle=25-40=-15<0$，无实根，舍去；}
\step{当 $m=-1$ 时，$\triangle=1+8=9>0$，符合题意，}
\step{所以实数 $m=-1$.}
\solend
\fi

\item 已知集合 $A=\{x\mid6x+a>0\}$，若 $1\in A$，$0\notin A$，则实数 $a$ 的取值范围为\blankline[4.4em].

\ifanswer
\solbegin
\step{由 $A=\{x\mid6x+a>0\}$，$1\in A$，$0\notin A$，得 $\begin{cases}6+a>0\\ a\leqslant0\end{cases}$，解得 $-6<a\leqslant0$，}
\step{即实数 $a$ 的取值范围是 $(-6,0]$.}
\solend
\fi

\item 关于 $x$ 的不等式 $ax-b>0$ 的解集为 $(1,+\infty)$，则关于 $x$ 的不等式 $(ax+b)(x-2)>0$ 的解集为\blankline[4.4em].

\ifanswer
\solbegin
\step{由不等式 $ax-b>0$ 的解集为 $(1,+\infty)$，得 $a>0$ 且 $\dfrac ba=1$，即 $b=a$，}
\step{则 $(ax+b)(x-2)>0$，可化为 $a(x+1)(x-2)>0$.}
\step{即 $(x+1)(x-2)>0$，解得 $x<-1$ 或 $x>2$，}
\step{所以不等式的解集为 $(-\infty,-1)\cup(2,+\infty)$.}
\solend
\fi

\item 已知实数 $a_1$、$a_2$、$a_3$、$a_4$、$a_5$ 满足 $\{x\mid(x-a_1)(x-a_2)(x-a_3)(x-a_4)(x-a_5)=0\}=\{1,2,3,4\}$，则 $a_1+a_2+a_3+a_4+a_5$ 的最大值为\blankline[4.4em].

\ifanswer
\solbegin
\step{5 个实数 $a_1$、$a_2$、$a_3$、$a_4$、$a_5$ 的不同取值只有 1，2，3，4，}
\step{即这 5 个数均从 1，2，3，4 中选取.}
\step{集合有 4 个元素，说明 1，2，3，4 每个数至少出现一次，}
\step{因此恰好有一个数重复出现两次，其余三个数各出现一次.}
\step{要让 5 个数的总和最大，应让重复出现两次的数取集合中的最大值 4，}
\step{此时 5 个数为 1，2，3，4，4，总和为 $1+2+3+4+4=14$，}
\step{即所求最大值为 14.}
\solend
\fi

\item 定义集合运算“$\times$”：$A\times B=\{(x,y)\mid x\in A,y\in B\}$，称为 $A$、$B$ 两个集合的“卡氏积”，若 $A=\{x\mid x^2-|x|\leqslant0,x\in\N\}$，$B=\{1,2,3\}$，则 $(A\times B)\cap(B\times A)=$\blankline[4.4em].

\ifanswer
\solbegin
\step{因为 $x^2-|x|\leqslant0$，所以 $0\leqslant|x|\leqslant1$，又因为 $x\in\N$，}
\step{所以 $A=\{x\mid x^2-|x|\leqslant0,x\in\N\}=\{0,1\}$，又 $B=\{1,2,3\}$，}
\step{由“卡氏积”的定义得 $A\times B=\{(0,1),(0,2),(0,3),(1,1),(1,2),(1,3)\}$，}
\step{$B\times A=\{(1,0),(1,1),(2,0),(2,1),(3,0),(3,1)\}$，}
\step{所以 $(A\times B)\cap(B\times A)=\{(1,1)\}$.}
\solend
\fi

\end{enumerate}

\bigsec{二、选择题}

\begin{enumerate}
\setcounter{enumi}{12}

% 第 13 题答案“B”原卷印在题干括号内，且无【答案】/【解析】段；
% 按 高一37 第 13 题、高一63 第 14 题、高一67 第 15 题成例括号留空、答案另提 \ans{$B$}
% 原卷题干末“条件”后无标点，照录（见文件头“原卷笔误/存疑”⑤）
\item “$A\subseteq B$”是“$A\subset B$”的\mbox{（\quad）}条件

\keepwithprev
A. 充分不必要条件\qquad B. 必要不充分条件

C. 充要条件\qquad D. 既不充分也不必要条件

\ifanswer
\ans{$B$}
\fi

\item 已知 $a$、$b$、$c$ 均为实数，若 $a>b$ 且 $ab\neq0$，则下列不等关系中成立的是\mbox{（\quad）}

\keepwithprev
A. $ac^2>bc^2$\qquad B. $\dfrac1b>\dfrac1a$\qquad C. $|a|>|b|$\qquad D. $c-b>c-a$

\ifanswer
\solbegin
\step{对于 $A$，若 $a>b$，当 $c=0$ 时，有 $ac^2=bc^2=0$，故 $A$ 错误；}
\step{对于 $B$，$\dfrac1b-\dfrac1a=\dfrac{a-b}{ab}$，因为 $a>b$，所以 $a-b>0$，}
\step{当 $a$、$b$ 同号时，$ab>0$，则 $\dfrac1b-\dfrac1a>0$，即 $\dfrac1b>\dfrac1a$，}
\step{当 $a$、$b$ 异号时，$ab<0$，则 $\dfrac1b-\dfrac1a<0$，即 $\dfrac1b<\dfrac1a$，故 $B$ 错误；}
\step{对于 $C$，当 $a=1$，$b=-3$ 时，$|a|<|b|$，故 $C$ 错误；}
\step{对于 $D$，若 $a>b$，则 $-b>-a$，$c-b>c-a$，故 $D$ 正确.}
\step{故选 $D$.}
\solend
\fi

% 原卷解析对 (1) 未讨论取等是否可得（x=y 时确能取 2），却总结“有 1 个、故选 B”，
% 与推理矛盾，照录（见文件头“原卷笔误/存疑”③）
\item 以下三个代数式：（1）$\dfrac yx+\dfrac xy$；（2）$|x|+\dfrac1{|x|}$；（3）$x^2+2+\dfrac1{x^2+2}$．其中的字母均表示实数，且代数式均有意义，那么代数式的最小值为 2 的个数是\mbox{（\quad）}

\keepwithprev
A. $0$\qquad B. $1$\qquad C. $2$\qquad D. $3$

\ifanswer
\solbegin
\step{（1）仅当 $x,y$ 同号时，$\dfrac yx,\dfrac xy>0$，由基本不等式得 $\dfrac yx+\dfrac xy\geqslant2$；}
\step{若 $x,y$ 异号，$\dfrac yx,\dfrac xy<0$，故 $\dfrac yx+\dfrac xy<0$，此时最小值不是 2；}
\step{（2）$|x|+\dfrac1{|x|}$ 中，$|x|>0$，由基本不等式 $|x|+\dfrac1{|x|}\geqslant2\sqrt{|x|\cdot\dfrac1{|x|}}=2$，}
\step{当且仅当 $|x|=1$，即 $x=\pm1$ 时取等号，故最小值为 2；}
\step{（3）令 $t=x^2+2$，则 $t\geqslant2$，函数化为 $t+\dfrac1t$，}
\step{对勾函数 $t+\dfrac1t$ 在 $[2,+\infty)$ 上严格增，最小值为 $2+\dfrac12=\dfrac52>2$，}
\step{等号成立条件 $t=1$ 无法取到，故最小值不是 2；}
\step{综上，最小值为 2 的代数式有 1 个，故选 $B$.}
\solend
\fi

\item 已知 $A$、$B$ 为非空实数集，$\Omega$ 为平面直角坐标系中的一些点构成的集合，集合 $C=\{y\mid$对任意 $x\in A$，有 $(x,y)\in\Omega\}$，集合 $D=\{x\mid$对任意 $y\in B$，有 $(x,y)\in\Omega\}$，对于下列两个命题：①若 $C\subseteq B$，则 $D\subseteq A$；②若 $C\supseteq B$，则 $D\subseteq A$；其中判断正确的是\mbox{（\quad）}

\keepwithprev
A. ①②都正确\qquad B. ①②都错误\qquad C. ①正确②错误\qquad D. ①错误②正确

\ifanswer
\solbegin
\step{设 $A=\{a_1,a_2,a_3,\cdots,a_n\},B=\{b_1,b_2,b_3,\cdots,b_n\},n>2$，}
\step{若 $\Omega=\{(c,b_1),(c,b_2),\cdots,(c,b_n)\},c\notin A$，}
\step{此时 $C=\varnothing\subseteq B,D=\{c\},D\subseteq A$ 错误，故①错误；}
% 原卷此处 N* 排正体（\mathrm{N}*），本稿统一 \N^*（见文件头“排版差异”④）；
% 全称量词“∀”照原卷符号排（见文件头“排版差异”⑫）
\step{若 $C\supseteq B$，则 $\forall(a_i,b_j)(1\leqslant i,j\leqslant n,i,j\in\N^*)\in\Omega$，则 $D\supseteq A$，}
\step{且 $A=\{a_1\},B=\{b_1,b_2\}$，}
\step{若 $\Omega=\{(a_1,b_1),(a_2,b_1),(a_1,b_2),(a_2,b_2),(c,b_1),(c,b_2)\},D$ 真包含 $A$，故②错误.}
\step{故选 $B$.}
\solend
\fi

\end{enumerate}

\bigsec{三、解答题}

\begin{enumerate}
\setcounter{enumi}{16}

\item 解下列关于 $x$ 的不等式组：$\begin{cases}|x-1|<2\\ \dfrac{x-2}{x-1}\geqslant0\end{cases}$.

\ifanswer
\solbegin
\step{解 $|x-1|<2$ 得 $-2<x-1<2$，解得 $-1<x<3$.}
\step{分式不等式 $\dfrac{x-2}{x-1}\geqslant0$ 等价于 $\begin{cases}(x-2)(x-1)\geqslant0\\ x-1\neq0\end{cases}$，解得 $x<1$ 或 $x\geqslant2$.}
\step{即 $\begin{cases}|x-1|<2\\ \dfrac{x-2}{x-1}\geqslant0\end{cases}$ 的解集为 $(-1,1)\cup[2,3)$.}
\solend
\fi

\ansspace[5]

\item 比较多项式 $2x^2-3x+1$ 与 $x^2-1$ 的大小.

\ifanswer
\solbegin
\step{作差得 $(2x^2-3x+1)-(x^2-1)=x^2-3x+2=(x-1)(x-2)$，}
\step{当 $x=1$ 或 $x=2$ 时，$(x-1)(x-2)=0$，即 $2x^2-3x+1=x^2-1$，}
\step{当 $x<1$ 或 $x>2$ 时，$(x-1)(x-2)>0$，即 $2x^2-3x+1>x^2-1$，}
\step{当 $1<x<2$ 时，$(x-1)(x-2)<0$，即 $2x^2-3x+1<x^2-1$；}
\step{综上，当 $x=1$ 或 $x=2$ 时，$2x^2-3x+1=x^2-1$；}
\step{当 $x<1$ 或 $x>2$ 时，$2x^2-3x+1>x^2-1$；}
\step{当 $1<x<2$ 时，$2x^2-3x+1<x^2-1$.}
\solend
\fi

\ansspace[6]

% 原卷 (1) 连用两个“若”（“若 a=4，若全集为实数集 R”），照录（见文件头“原卷笔误/存疑”①）
\item 已知集合 $A=\left\{x\mid\dfrac{(a+1)x}{x-1}<a\right\}$，集合 $B=[1,4]$，$a$ 为实常数；

（1）若 $a=4$，若全集为实数集 $\R$，求集合 $\overline{A}$（用区间形式表示）；

（2）若 $A\subseteq B$，求 $a$ 的取值范围.

\ifanswer
\solbegin
\stepm{(1)}{若 $a=4$，集合 $A=\left\{x\mid\dfrac{5x}{x-1}<4\right\}=\left\{x\mid\dfrac{5x}{x-1}-4<0\right\}=\left\{x\mid\dfrac{x+4}{x-1}<0\right\}$}
\step{$=(-4,1)$，因此 $\overline{A}=(-\infty,-4]\cup[1,+\infty)$；}
\stepm{(2)}{当 $-a<1$，即 $a>-1$ 时，$A=(-a,1)$，}
\step{要使 $A\subseteq[1,4]$，需 $-a\geqslant1$，即 $a\leqslant-1$，与 $a>-1$ 矛盾，舍去；}
\step{当 $-a>1$，即 $a<-1$ 时，$A=(1,-a)$，}
\step{要使 $A\subseteq[1,4]$，需 $-a\leqslant4$，即 $a\geqslant-4$，所以 $-4\leqslant a<-1$；}
\step{当 $-a=1$，即 $a=-1$ 时，$A=\varnothing$，满足 $A\subseteq B$；}
\step{综上，实数 $a$ 的取值范围为 $[-4,-1]$.}
\solend
\fi

\ansspace[7]

% 原卷题干“可以售出 8 万本、据市场调查”处用顿号，照录（见文件头“原卷笔误/存疑”②）
\item 某种杂志原以每本 25 元的价格销售，可以售出 8 万本、据市场调查，杂志的单价每提高 1 元，销售量就减少 0.2 万本，假定杂志的成本是每本 10 元（不计其他成本）.

（1）当杂志以每本 30 元定价销售时，求销售该杂志所获利润；

（2）在每本 25 元的价格的基础上提高定价，试确定杂志的定价在什么范围内可使得销售利润不低于 140 万元.

\ifanswer
\solbegin
\stepm{(1)}{当杂志以 30 元定价销售时，}
\step{销售该杂志所获利润为 $[8-(30-25)\times0.2]\times(30-10)=140$ 万元；}
\stepm{(2)}{在每本 25 元的价格的基础上提高定价，设杂志的定价为每本 $x$ 元，}
\step{由题意得 $[8-(x-25)\times0.2]\times(x-10)\geqslant140$，解得 $30\leqslant x\leqslant45$，}
\step{所以杂志的定价在每本不低于 30 元且不高于 45 元的范围内，}
\step{可使得销售利润不低于 140 万元.}
\solend
\fi

\ansspace[7]

% 原卷题干“定义：对于集合 A，则称集合 B=…”句式杂糅，照录（见文件头“原卷笔误/存疑”④）
\item 定义：对于集合 $A$，则称集合 $B=\left\{x\mid x=\dfrac{m+n}{1+mn},m,n\in A,1+mn\neq0\right\}$ 为集合 $A$ 的生成集（其中 $m$、$n$ 可以相等）.

（1）若集合 $A=\{-1,1\}$ 的生成集为 $B$，求集合 $B$；

（2）若集合 $A=\{a,2\}$ 的生成集为 $B$，$B$ 的子集个数为 4 个，求实数 $a$ 的值；

（3）若集合 $A=\{x\mid-1<x<1\}$ 的生成集为 $B$，求证：$B\subseteq A$.

\ifanswer
\solbegin
\stepm{(1)}{因为 $1+mn\neq0$，}
\step{当 $m=n=1$ 时，$x=\dfrac{m+n}{1+mn}=1$；当 $m=n=-1$ 时，$x=\dfrac{m+n}{1+mn}=-1$；}
\step{所以 $B=\{-1,1\}$；}
\stepm{(2)}{$B$ 的子集个数为 4 个，故 $B$ 集合的元素有 2 个，}
\step{当 $m=n=2$ 时，$x=\dfrac{m+n}{1+mn}=\dfrac45$，当 $m=n=a$ 时，$x=\dfrac{m+n}{1+mn}=\dfrac{2a}{1+a^2}$，}
\step{当 $m$ 与 $n$ 分别取 2 和 $a$ 时，$x=\dfrac{a+2}{1+2a}$，}
\stepm{①}{$\dfrac{2a}{1+a^2}=\dfrac45$，解得 $a=2$（舍）或 $a=\dfrac12$，}
\stepm{②}{$\dfrac{a+2}{1+2a}=\dfrac45$，解得 $a=2$（舍），}
\stepm{③}{$\dfrac{2a}{1+a^2}=\dfrac{a+2}{1+2a}$，解得 $a=2$（舍）或 $a=1$ 或 $a=-1$，}
\step{综上所述，$a=1$ 或 $a=-1$ 或 $a=\dfrac12$，即 $a\in\left\{-1,1,\dfrac12\right\}$；}
\stepm{(3)}{$\forall m$、$n\in A$，则 $\dfrac{m+n}{1+mn}+1=\dfrac{(m+1)(n+1)}{1+mn}>0$，故 $\dfrac{m+n}{1+mn}>-1$；}
\step{$\dfrac{m+n}{1+mn}-1=-\dfrac{(m-1)(n-1)}{1+mn}<0$，故 $\dfrac{m+n}{1+mn}<1$；}
\step{故 $B\subseteq\{x\mid-1<x<1\}$，即 $B\subseteq A$.}
\solend
\fi

\ansspace*[10]

\end{enumerate}

\end{document}
"
%
% 原始材料：微信公众号“魔都数学加油站”文章，连载“27学年高一71”；
% 原文标题“27学年高一71——2026-2027年上海市中国中学高一上10月月考”；
% 原文链接：https://mp.weixin.qq.com/s/UZrP53l3GIin040L1EkljA。
% 页面标注发布时间 2026-10-19 22:47（上海）（与 高一69、高一70 页面读出同一 ct 值）。
% 正文为 7 张 PNG 页（`原图/高一71_中国中学_10月月考/01..07.png`），全部 4762×6735，
% 均按页序读取。7 页全为试卷正文页（卷首标题在第 1 页页顶，卷末解答题第 21 题 (3)
% 解析止于第 7 页），无封面页、目录页、二维码或宣传页。粉色斜排“魔都数学加油站”
% 为卷面水印，与公众号页脚一并未录入。
%
% 卷首标题：照卷面印作“2026-2027 年中国中学高一上 10 月月考”（黑体加粗一行）。
% 原卷年份区间用半角连字符，本稿按全稿惯例排 en dash（同 高一68、高一69、高一70）；
% 原卷只有一行标题、无副标题。
%
% 篇幅与题号：原卷自第 3 题起编号（第 1、2 题不在本卷内）。填空题 3—12（10 题）、
% 选择题 13—16（4 题）、解答题 17—21（5 题），共 19 题。本稿沿用原节与原题号，
% 只改排为全稿统一的大节样式（原卷小节标题为“一、填空题”“二、选择题”“三、解答题”）。
%
% 讲评版：原卷每题下直接印【答案】或【解析】。**只印【答案】的 3 题**（第 3、4、5 题），
% 本稿只给 \ans{}；**第 13 题既无【答案】也无【解析】段，答案“B”直接印在题干括号内
% （“(  B  )”）**，按 高一37 第 13 题、高一63 第 14 题、高一67 第 15 题成例把括号留空、
% 答案另提为 \ans{$B$}。其余 15 题全部印【解析】（按全稿口径这类题不另提 \ans{}——
% 第 14—16 题的“故选 D／B／B”写在解析末句，亦照录在解析内）。
% 除第 13 题外，全卷选择题作答括号内均为空白（第 14、15、16 题）。
%
% 副标题：原卷未印副标题。本稿按“知识模块”口径标作“集合与常用逻辑用语、不等式”：
% 19 题中集合与常用逻辑用语 10 题（第 3、4、5、6、9、11、12、13、16、21 题；其中第 4 题
% 属反证法，一并计入常用逻辑用语），不等式 9 题（第 7、8、10、14、15、17、18、19、20 题；
% 第 19 题 (1) 属补集、(2) 属分式不等式，按主干计入不等式），两类均确有分量，故并标。
% 口径同 高一63、高一68、高一69、高一70。
%
% 与原卷的排版差异：
%   ① 原卷每题下直接印【答案】/【解析】，本稿按全稿惯例将答案与解析置于答案版，试卷版
%      保持干净；第 13 题印在题干括号内的答案“B”同样移入 \ans{}，见上；
%   ② 原卷未印副标题，本稿按“知识模块”口径补出；
%   ③ 原卷填空横线约四个字宽，本稿按全稿惯例统一排 \blankline[4.4em]；
%   ④ 原卷数集记号不统一：实数集排斜体 $R$（第 19 题 (1)），自然数集排斜体 $N$
%      （第 5 题题干、第 12 题题干与解析共三处），第 16 题解析又排正体 $\mathrm{N}^*$；
%      本稿按全稿惯例统一改用 $\R$、$\N$、$\N^*$（同 高一68／69／70 排版差异④）；
%      空集记号原卷作 $\varnothing$，照录；
%   ⑤ 原卷句末标点混用：第 15 题题干“（3）$x^2+2+\dfrac1{x^2+2}$**．**其中”、
%      第 17 题题干末“$\dfrac{x-2}{x-1}\geqslant0$**．**”两处放大确认（第 15 题为全角
%      句点、第 17 题为西文句点），其余各题末为西文句点，已逐处照录；题干内冒号——
%      第 5 题“$\{y\mid\cdots\}$**：**”、第 12 题““×”**：**”、第 17 题“不等式组**：**”、
%      第 21 题“定义**：**”——均为全角冒号，照录；数学式内的逗号（如第 12 题
%      “$x\in A,y\in B$”、第 16 题“$A=\{a_1,\cdots\},B=\{b_1,\cdots\},n>2$”）按原卷作
%      半角，中文行文里的逗号按原卷作全角，均未强改；
%   ⑥ 第 17—21 题的子问原卷印作“（1）…（2）…”（另起一段），本稿按 高一69／70 成例用
%      行内“（1）（2）（3）”段，解析中的子问标号一律用 \stepm{(1)}{…}；
%      第 21 题 (2) 解析内层的“①”“②”“③”照原卷作 \stepm{①}{…} 等；
%   ⑦ 每道解答题原卷未留作答空白，本稿按全稿惯例加 \ansspace（末题用 \ansspace*）；
%   ⑧ 选择题的作答括号原卷排半角“(　　)”（第 14、15、16 题括号内空白、第 13 题括号内
%      印 B），本稿按全稿惯例排 \mbox{（\quad）}；选项行排照原卷——第 13 题两行两条
%      （A B 一行、C D 一行），第 14、15、16 题一行四条；
%   ⑨ 第 8 题解析、第 9 题解析、第 17 题题干与解析中的多行不等式组，原卷均排作“{”
%      连接的多行式，本稿按全稿惯例排 cases；
%   ⑩ 第 10 题解析、第 17 题解析、第 19 题 (1) 解析的并集符号原卷排大号 $\bigcup$
%      （显示体），本稿按全稿惯例排行内 $\cup$；
%   ⑪ 第 8 题解析的判别式照原卷作空心“$\triangle$”（与第 7 题解析的“$\Delta$”不同形），
%      照录未改；
%   ⑫ 第 16 题解析第 4 行的“$\forall$”（全称量词）照原卷符号排，未改写为中文“任意”。
%
% 原卷笔误/存疑（均照录未改，正文对应处另有 % 注）：
%   ① 第 19 题 (1) 连用两个“若”——“若 $a=4$，若全集为实数集 $R$”，后一“若”应作
%      “当”或删；原卷如此，照录未改。
%   ② 第 20 题题干“可以售出 8 万本**、**据市场调查”——“万本”后用顿号，按文意应为
%      逗号或句号；原卷如此，照录未改。
%   ③ 第 15 题解析对代数式 (1) 只得出“$y/x+x/y\geqslant2$”（同号时），未讨论取等是否
%      可得——取 $x=y$ 时确能取到 2——却总结“最小值为 2 的代数式有 1 个，故选 $B$”；
%      按解析自身的推理 (1) 也应计入，结论“1 个”与推理矛盾（若计 (1) 应为 2 个、选 C）。
%      原卷如此，照录未改。
%   ④ 第 21 题题干“定义：对于集合 $A$，则称集合 $B=\cdots$”——“对于……”与“则称……”
%      句式杂糅（应作“对于集合 $A$，称……”或“若集合 $A$，则称……”）；原卷如此，照录未改。
%   ⑤ 第 13 题题干末“……的（　）条件”后无任何标点（其余各题题干末均有句点）；
%      原卷如此，照录未改。
%
% 行数核对：逐题数 \step/\stepm 条数与原卷【解析】印刷行数，15 题逐一相符，
% 合计 86 条 / 86 行（Q6：1；Q7：1；Q8：7；Q9：2；Q10：4；Q11：7；Q12：5；Q14：7；
% Q15：8；Q16：7；Q17：3；Q18：7；Q19：8；Q20：6；Q21：13）；另有 3 题只印【答案】
% （第 3、4、5 题）与第 13 题括号内答案，各提 \ans{}，不计入 \step 账。
%
% 插图：无题目配图或数学图形（7 页题面均无“如图/右图/图中”字样；页面另有粉色斜排水印，
% 见上）。
%
\documentclass[a4paper,11pt]{article}
\usepackage{common}

\begin{document}

\examtitlest[3]{2026--2027 年中国中学高一上 10 月月考}{集合与常用逻辑用语、不等式}

\bigsec{一、填空题}

\begin{enumerate}
\setcounter{enumi}{2}

\item 已知集合 $A=[1,+\infty)$，$B=[-1,3)$，则 $A\cap B=$\blankline[4.4em].

\ifanswer
\ans{$[1,3)$}
\fi

\item 用反证法证明命题“已知 $a$、$b$ 为正数，$n$ 为正整数，求证：如果 $a^n>b^n$，那么 $a>b$”时，应当作出的假设为：“假设\blankline[4.4em]”.

\ifanswer
\ans{$a\leqslant b$}
\fi

\item 用列举法表示集合 $\{y\mid y=x+1,x\in\N,x\leqslant1\}$：\blankline[4.4em].

\ifanswer
\ans{$\{1,2\}$}
\fi

\item 某班在一次测验中，有 36 人数学成绩不低于 80 分，有 20 人物理成绩不低于 80 分，有 15 人的数学、物理成绩都不低于 80 分，则这两科成绩中至少有一科不低于 80 分的人数为\blankline[4.4em]人.

\ifanswer
\solbegin
\step{至少有 $36+20-15=41$ 人.}
\solend
\fi

\item 若关于 $x$ 的不等式 $x^2-x+m<0$ 的解集是 $\varnothing$，则实数 $m$ 的取值范围是\blankline[4.4em].

\ifanswer
\solbegin
\step{由题意得 $\Delta=1-4m\leqslant0$，所以 $m\geqslant\dfrac14$.}
\solend
\fi

% 第 8 题解析的判别式原卷作空心“△”（与第 7 题解析的“Δ”不同形），照录（见文件头“排版差异”⑪）
\item 已知关于 $x$ 的方程 $x^2+mx+2m=0$ 的两个实数根分别为 $\alpha$，$\beta$，若 $\alpha^2+\beta^2=5$，则实数 $m=$\blankline[4.4em].

\ifanswer
\solbegin
\step{方程 $x^2+mx+2m=0$ 的两个实数根为 $\alpha$、$\beta$，}
\step{由根与系数的关系得 $\begin{cases}\alpha+\beta=-m\\ \alpha\beta=2m\end{cases}$，又因为 $\alpha^2+\beta^2=5$，}
\step{所以 $\alpha^2+\beta^2=(\alpha+\beta)^2-2\alpha\beta=m^2-4m=5$，整理得 $m^2-4m-5=0$，}
\step{解得 $m=5$ 或 $m=-1$，}
\step{当 $m=5$ 时，$\triangle=25-40=-15<0$，无实根，舍去；}
\step{当 $m=-1$ 时，$\triangle=1+8=9>0$，符合题意，}
\step{所以实数 $m=-1$.}
\solend
\fi

\item 已知集合 $A=\{x\mid6x+a>0\}$，若 $1\in A$，$0\notin A$，则实数 $a$ 的取值范围为\blankline[4.4em].

\ifanswer
\solbegin
\step{由 $A=\{x\mid6x+a>0\}$，$1\in A$，$0\notin A$，得 $\begin{cases}6+a>0\\ a\leqslant0\end{cases}$，解得 $-6<a\leqslant0$，}
\step{即实数 $a$ 的取值范围是 $(-6,0]$.}
\solend
\fi

\item 关于 $x$ 的不等式 $ax-b>0$ 的解集为 $(1,+\infty)$，则关于 $x$ 的不等式 $(ax+b)(x-2)>0$ 的解集为\blankline[4.4em].

\ifanswer
\solbegin
\step{由不等式 $ax-b>0$ 的解集为 $(1,+\infty)$，得 $a>0$ 且 $\dfrac ba=1$，即 $b=a$，}
\step{则 $(ax+b)(x-2)>0$，可化为 $a(x+1)(x-2)>0$.}
\step{即 $(x+1)(x-2)>0$，解得 $x<-1$ 或 $x>2$，}
\step{所以不等式的解集为 $(-\infty,-1)\cup(2,+\infty)$.}
\solend
\fi

\item 已知实数 $a_1$、$a_2$、$a_3$、$a_4$、$a_5$ 满足 $\{x\mid(x-a_1)(x-a_2)(x-a_3)(x-a_4)(x-a_5)=0\}=\{1,2,3,4\}$，则 $a_1+a_2+a_3+a_4+a_5$ 的最大值为\blankline[4.4em].

\ifanswer
\solbegin
\step{5 个实数 $a_1$、$a_2$、$a_3$、$a_4$、$a_5$ 的不同取值只有 1，2，3，4，}
\step{即这 5 个数均从 1，2，3，4 中选取.}
\step{集合有 4 个元素，说明 1，2，3，4 每个数至少出现一次，}
\step{因此恰好有一个数重复出现两次，其余三个数各出现一次.}
\step{要让 5 个数的总和最大，应让重复出现两次的数取集合中的最大值 4，}
\step{此时 5 个数为 1，2，3，4，4，总和为 $1+2+3+4+4=14$，}
\step{即所求最大值为 14.}
\solend
\fi

\item 定义集合运算“$\times$”：$A\times B=\{(x,y)\mid x\in A,y\in B\}$，称为 $A$、$B$ 两个集合的“卡氏积”，若 $A=\{x\mid x^2-|x|\leqslant0,x\in\N\}$，$B=\{1,2,3\}$，则 $(A\times B)\cap(B\times A)=$\blankline[4.4em].

\ifanswer
\solbegin
\step{因为 $x^2-|x|\leqslant0$，所以 $0\leqslant|x|\leqslant1$，又因为 $x\in\N$，}
\step{所以 $A=\{x\mid x^2-|x|\leqslant0,x\in\N\}=\{0,1\}$，又 $B=\{1,2,3\}$，}
\step{由“卡氏积”的定义得 $A\times B=\{(0,1),(0,2),(0,3),(1,1),(1,2),(1,3)\}$，}
\step{$B\times A=\{(1,0),(1,1),(2,0),(2,1),(3,0),(3,1)\}$，}
\step{所以 $(A\times B)\cap(B\times A)=\{(1,1)\}$.}
\solend
\fi

\end{enumerate}

\bigsec{二、选择题}

\begin{enumerate}
\setcounter{enumi}{12}

% 第 13 题答案“B”原卷印在题干括号内，且无【答案】/【解析】段；
% 按 高一37 第 13 题、高一63 第 14 题、高一67 第 15 题成例括号留空、答案另提 \ans{$B$}
% 原卷题干末“条件”后无标点，照录（见文件头“原卷笔误/存疑”⑤）
\item “$A\subseteq B$”是“$A\subset B$”的\mbox{（\quad）}条件

\keepwithprev
A. 充分不必要条件\qquad B. 必要不充分条件

C. 充要条件\qquad D. 既不充分也不必要条件

\ifanswer
\ans{$B$}
\fi

\item 已知 $a$、$b$、$c$ 均为实数，若 $a>b$ 且 $ab\neq0$，则下列不等关系中成立的是\mbox{（\quad）}

\keepwithprev
A. $ac^2>bc^2$\qquad B. $\dfrac1b>\dfrac1a$\qquad C. $|a|>|b|$\qquad D. $c-b>c-a$

\ifanswer
\solbegin
\step{对于 $A$，若 $a>b$，当 $c=0$ 时，有 $ac^2=bc^2=0$，故 $A$ 错误；}
\step{对于 $B$，$\dfrac1b-\dfrac1a=\dfrac{a-b}{ab}$，因为 $a>b$，所以 $a-b>0$，}
\step{当 $a$、$b$ 同号时，$ab>0$，则 $\dfrac1b-\dfrac1a>0$，即 $\dfrac1b>\dfrac1a$，}
\step{当 $a$、$b$ 异号时，$ab<0$，则 $\dfrac1b-\dfrac1a<0$，即 $\dfrac1b<\dfrac1a$，故 $B$ 错误；}
\step{对于 $C$，当 $a=1$，$b=-3$ 时，$|a|<|b|$，故 $C$ 错误；}
\step{对于 $D$，若 $a>b$，则 $-b>-a$，$c-b>c-a$，故 $D$ 正确.}
\step{故选 $D$.}
\solend
\fi

% 原卷解析对 (1) 未讨论取等是否可得（x=y 时确能取 2），却总结“有 1 个、故选 B”，
% 与推理矛盾，照录（见文件头“原卷笔误/存疑”③）
\item 以下三个代数式：（1）$\dfrac yx+\dfrac xy$；（2）$|x|+\dfrac1{|x|}$；（3）$x^2+2+\dfrac1{x^2+2}$．其中的字母均表示实数，且代数式均有意义，那么代数式的最小值为 2 的个数是\mbox{（\quad）}

\keepwithprev
A. $0$\qquad B. $1$\qquad C. $2$\qquad D. $3$

\ifanswer
\solbegin
\step{（1）仅当 $x,y$ 同号时，$\dfrac yx,\dfrac xy>0$，由基本不等式得 $\dfrac yx+\dfrac xy\geqslant2$；}
\step{若 $x,y$ 异号，$\dfrac yx,\dfrac xy<0$，故 $\dfrac yx+\dfrac xy<0$，此时最小值不是 2；}
\step{（2）$|x|+\dfrac1{|x|}$ 中，$|x|>0$，由基本不等式 $|x|+\dfrac1{|x|}\geqslant2\sqrt{|x|\cdot\dfrac1{|x|}}=2$，}
\step{当且仅当 $|x|=1$，即 $x=\pm1$ 时取等号，故最小值为 2；}
\step{（3）令 $t=x^2+2$，则 $t\geqslant2$，函数化为 $t+\dfrac1t$，}
\step{对勾函数 $t+\dfrac1t$ 在 $[2,+\infty)$ 上严格增，最小值为 $2+\dfrac12=\dfrac52>2$，}
\step{等号成立条件 $t=1$ 无法取到，故最小值不是 2；}
\step{综上，最小值为 2 的代数式有 1 个，故选 $B$.}
\solend
\fi

\item 已知 $A$、$B$ 为非空实数集，$\Omega$ 为平面直角坐标系中的一些点构成的集合，集合 $C=\{y\mid$对任意 $x\in A$，有 $(x,y)\in\Omega\}$，集合 $D=\{x\mid$对任意 $y\in B$，有 $(x,y)\in\Omega\}$，对于下列两个命题：①若 $C\subseteq B$，则 $D\subseteq A$；②若 $C\supseteq B$，则 $D\subseteq A$；其中判断正确的是\mbox{（\quad）}

\keepwithprev
A. ①②都正确\qquad B. ①②都错误\qquad C. ①正确②错误\qquad D. ①错误②正确

\ifanswer
\solbegin
\step{设 $A=\{a_1,a_2,a_3,\cdots,a_n\},B=\{b_1,b_2,b_3,\cdots,b_n\},n>2$，}
\step{若 $\Omega=\{(c,b_1),(c,b_2),\cdots,(c,b_n)\},c\notin A$，}
\step{此时 $C=\varnothing\subseteq B,D=\{c\},D\subseteq A$ 错误，故①错误；}
% 原卷此处 N* 排正体（\mathrm{N}*），本稿统一 \N^*（见文件头“排版差异”④）；
% 全称量词“∀”照原卷符号排（见文件头“排版差异”⑫）
\step{若 $C\supseteq B$，则 $\forall(a_i,b_j)(1\leqslant i,j\leqslant n,i,j\in\N^*)\in\Omega$，则 $D\supseteq A$，}
\step{且 $A=\{a_1\},B=\{b_1,b_2\}$，}
\step{若 $\Omega=\{(a_1,b_1),(a_2,b_1),(a_1,b_2),(a_2,b_2),(c,b_1),(c,b_2)\},D$ 真包含 $A$，故②错误.}
\step{故选 $B$.}
\solend
\fi

\end{enumerate}

\bigsec{三、解答题}

\begin{enumerate}
\setcounter{enumi}{16}

\item 解下列关于 $x$ 的不等式组：$\begin{cases}|x-1|<2\\ \dfrac{x-2}{x-1}\geqslant0\end{cases}$.

\ifanswer
\solbegin
\step{解 $|x-1|<2$ 得 $-2<x-1<2$，解得 $-1<x<3$.}
\step{分式不等式 $\dfrac{x-2}{x-1}\geqslant0$ 等价于 $\begin{cases}(x-2)(x-1)\geqslant0\\ x-1\neq0\end{cases}$，解得 $x<1$ 或 $x\geqslant2$.}
\step{即 $\begin{cases}|x-1|<2\\ \dfrac{x-2}{x-1}\geqslant0\end{cases}$ 的解集为 $(-1,1)\cup[2,3)$.}
\solend
\fi

\ansspace[5]

\item 比较多项式 $2x^2-3x+1$ 与 $x^2-1$ 的大小.

\ifanswer
\solbegin
\step{作差得 $(2x^2-3x+1)-(x^2-1)=x^2-3x+2=(x-1)(x-2)$，}
\step{当 $x=1$ 或 $x=2$ 时，$(x-1)(x-2)=0$，即 $2x^2-3x+1=x^2-1$，}
\step{当 $x<1$ 或 $x>2$ 时，$(x-1)(x-2)>0$，即 $2x^2-3x+1>x^2-1$，}
\step{当 $1<x<2$ 时，$(x-1)(x-2)<0$，即 $2x^2-3x+1<x^2-1$；}
\step{综上，当 $x=1$ 或 $x=2$ 时，$2x^2-3x+1=x^2-1$；}
\step{当 $x<1$ 或 $x>2$ 时，$2x^2-3x+1>x^2-1$；}
\step{当 $1<x<2$ 时，$2x^2-3x+1<x^2-1$.}
\solend
\fi

\ansspace[6]

% 原卷 (1) 连用两个“若”（“若 a=4，若全集为实数集 R”），照录（见文件头“原卷笔误/存疑”①）
\item 已知集合 $A=\left\{x\mid\dfrac{(a+1)x}{x-1}<a\right\}$，集合 $B=[1,4]$，$a$ 为实常数；

（1）若 $a=4$，若全集为实数集 $\R$，求集合 $\overline{A}$（用区间形式表示）；

（2）若 $A\subseteq B$，求 $a$ 的取值范围.

\ifanswer
\solbegin
\stepm{(1)}{若 $a=4$，集合 $A=\left\{x\mid\dfrac{5x}{x-1}<4\right\}=\left\{x\mid\dfrac{5x}{x-1}-4<0\right\}=\left\{x\mid\dfrac{x+4}{x-1}<0\right\}$}
\step{$=(-4,1)$，因此 $\overline{A}=(-\infty,-4]\cup[1,+\infty)$；}
\stepm{(2)}{当 $-a<1$，即 $a>-1$ 时，$A=(-a,1)$，}
\step{要使 $A\subseteq[1,4]$，需 $-a\geqslant1$，即 $a\leqslant-1$，与 $a>-1$ 矛盾，舍去；}
\step{当 $-a>1$，即 $a<-1$ 时，$A=(1,-a)$，}
\step{要使 $A\subseteq[1,4]$，需 $-a\leqslant4$，即 $a\geqslant-4$，所以 $-4\leqslant a<-1$；}
\step{当 $-a=1$，即 $a=-1$ 时，$A=\varnothing$，满足 $A\subseteq B$；}
\step{综上，实数 $a$ 的取值范围为 $[-4,-1]$.}
\solend
\fi

\ansspace[7]

% 原卷题干“可以售出 8 万本、据市场调查”处用顿号，照录（见文件头“原卷笔误/存疑”②）
\item 某种杂志原以每本 25 元的价格销售，可以售出 8 万本、据市场调查，杂志的单价每提高 1 元，销售量就减少 0.2 万本，假定杂志的成本是每本 10 元（不计其他成本）.

（1）当杂志以每本 30 元定价销售时，求销售该杂志所获利润；

（2）在每本 25 元的价格的基础上提高定价，试确定杂志的定价在什么范围内可使得销售利润不低于 140 万元.

\ifanswer
\solbegin
\stepm{(1)}{当杂志以 30 元定价销售时，}
\step{销售该杂志所获利润为 $[8-(30-25)\times0.2]\times(30-10)=140$ 万元；}
\stepm{(2)}{在每本 25 元的价格的基础上提高定价，设杂志的定价为每本 $x$ 元，}
\step{由题意得 $[8-(x-25)\times0.2]\times(x-10)\geqslant140$，解得 $30\leqslant x\leqslant45$，}
\step{所以杂志的定价在每本不低于 30 元且不高于 45 元的范围内，}
\step{可使得销售利润不低于 140 万元.}
\solend
\fi

\ansspace[7]

% 原卷题干“定义：对于集合 A，则称集合 B=…”句式杂糅，照录（见文件头“原卷笔误/存疑”④）
\item 定义：对于集合 $A$，则称集合 $B=\left\{x\mid x=\dfrac{m+n}{1+mn},m,n\in A,1+mn\neq0\right\}$ 为集合 $A$ 的生成集（其中 $m$、$n$ 可以相等）.

（1）若集合 $A=\{-1,1\}$ 的生成集为 $B$，求集合 $B$；

（2）若集合 $A=\{a,2\}$ 的生成集为 $B$，$B$ 的子集个数为 4 个，求实数 $a$ 的值；

（3）若集合 $A=\{x\mid-1<x<1\}$ 的生成集为 $B$，求证：$B\subseteq A$.

\ifanswer
\solbegin
\stepm{(1)}{因为 $1+mn\neq0$，}
\step{当 $m=n=1$ 时，$x=\dfrac{m+n}{1+mn}=1$；当 $m=n=-1$ 时，$x=\dfrac{m+n}{1+mn}=-1$；}
\step{所以 $B=\{-1,1\}$；}
\stepm{(2)}{$B$ 的子集个数为 4 个，故 $B$ 集合的元素有 2 个，}
\step{当 $m=n=2$ 时，$x=\dfrac{m+n}{1+mn}=\dfrac45$，当 $m=n=a$ 时，$x=\dfrac{m+n}{1+mn}=\dfrac{2a}{1+a^2}$，}
\step{当 $m$ 与 $n$ 分别取 2 和 $a$ 时，$x=\dfrac{a+2}{1+2a}$，}
\stepm{①}{$\dfrac{2a}{1+a^2}=\dfrac45$，解得 $a=2$（舍）或 $a=\dfrac12$，}
\stepm{②}{$\dfrac{a+2}{1+2a}=\dfrac45$，解得 $a=2$（舍），}
\stepm{③}{$\dfrac{2a}{1+a^2}=\dfrac{a+2}{1+2a}$，解得 $a=2$（舍）或 $a=1$ 或 $a=-1$，}
\step{综上所述，$a=1$ 或 $a=-1$ 或 $a=\dfrac12$，即 $a\in\left\{-1,1,\dfrac12\right\}$；}
\stepm{(3)}{$\forall m$、$n\in A$，则 $\dfrac{m+n}{1+mn}+1=\dfrac{(m+1)(n+1)}{1+mn}>0$，故 $\dfrac{m+n}{1+mn}>-1$；}
\step{$\dfrac{m+n}{1+mn}-1=-\dfrac{(m-1)(n-1)}{1+mn}<0$，故 $\dfrac{m+n}{1+mn}<1$；}
\step{故 $B\subseteq\{x\mid-1<x<1\}$，即 $B\subseteq A$.}
\solend
\fi

\ansspace*[10]

\end{enumerate}

\end{document}
"
%
% 原始材料：微信公众号“魔都数学加油站”文章，连载“27学年高一71”；
% 原文标题“27学年高一71——2026-2027年上海市中国中学高一上10月月考”；
% 原文链接：https://mp.weixin.qq.com/s/UZrP53l3GIin040L1EkljA。
% 页面标注发布时间 2026-10-19 22:47（上海）（与 高一69、高一70 页面读出同一 ct 值）。
% 正文为 7 张 PNG 页（`原图/高一71_中国中学_10月月考/01..07.png`），全部 4762×6735，
% 均按页序读取。7 页全为试卷正文页（卷首标题在第 1 页页顶，卷末解答题第 21 题 (3)
% 解析止于第 7 页），无封面页、目录页、二维码或宣传页。粉色斜排“魔都数学加油站”
% 为卷面水印，与公众号页脚一并未录入。
%
% 卷首标题：照卷面印作“2026-2027 年中国中学高一上 10 月月考”（黑体加粗一行）。
% 原卷年份区间用半角连字符，本稿按全稿惯例排 en dash（同 高一68、高一69、高一70）；
% 原卷只有一行标题、无副标题。
%
% 篇幅与题号：原卷自第 3 题起编号（第 1、2 题不在本卷内）。填空题 3—12（10 题）、
% 选择题 13—16（4 题）、解答题 17—21（5 题），共 19 题。本稿沿用原节与原题号，
% 只改排为全稿统一的大节样式（原卷小节标题为“一、填空题”“二、选择题”“三、解答题”）。
%
% 讲评版：原卷每题下直接印【答案】或【解析】。**只印【答案】的 3 题**（第 3、4、5 题），
% 本稿只给 \ans{}；**第 13 题既无【答案】也无【解析】段，答案“B”直接印在题干括号内
% （“(  B  )”）**，按 高一37 第 13 题、高一63 第 14 题、高一67 第 15 题成例把括号留空、
% 答案另提为 \ans{$B$}。其余 15 题全部印【解析】（按全稿口径这类题不另提 \ans{}——
% 第 14—16 题的“故选 D／B／B”写在解析末句，亦照录在解析内）。
% 除第 13 题外，全卷选择题作答括号内均为空白（第 14、15、16 题）。
%
% 副标题：原卷未印副标题。本稿按“知识模块”口径标作“集合与常用逻辑用语、不等式”：
% 19 题中集合与常用逻辑用语 10 题（第 3、4、5、6、9、11、12、13、16、21 题；其中第 4 题
% 属反证法，一并计入常用逻辑用语），不等式 9 题（第 7、8、10、14、15、17、18、19、20 题；
% 第 19 题 (1) 属补集、(2) 属分式不等式，按主干计入不等式），两类均确有分量，故并标。
% 口径同 高一63、高一68、高一69、高一70。
%
% 与原卷的排版差异：
%   ① 原卷每题下直接印【答案】/【解析】，本稿按全稿惯例将答案与解析置于答案版，试卷版
%      保持干净；第 13 题印在题干括号内的答案“B”同样移入 \ans{}，见上；
%   ② 原卷未印副标题，本稿按“知识模块”口径补出；
%   ③ 原卷填空横线约四个字宽，本稿按全稿惯例统一排 \blankline[4.4em]；
%   ④ 原卷数集记号不统一：实数集排斜体 $R$（第 19 题 (1)），自然数集排斜体 $N$
%      （第 5 题题干、第 12 题题干与解析共三处），第 16 题解析又排正体 $\mathrm{N}^*$；
%      本稿按全稿惯例统一改用 $\R$、$\N$、$\N^*$（同 高一68／69／70 排版差异④）；
%      空集记号原卷作 $\varnothing$，照录；
%   ⑤ 原卷句末标点混用：第 15 题题干“（3）$x^2+2+\dfrac1{x^2+2}$**．**其中”、
%      第 17 题题干末“$\dfrac{x-2}{x-1}\geqslant0$**．**”两处放大确认（第 15 题为全角
%      句点、第 17 题为西文句点），其余各题末为西文句点，已逐处照录；题干内冒号——
%      第 5 题“$\{y\mid\cdots\}$**：**”、第 12 题““×”**：**”、第 17 题“不等式组**：**”、
%      第 21 题“定义**：**”——均为全角冒号，照录；数学式内的逗号（如第 12 题
%      “$x\in A,y\in B$”、第 16 题“$A=\{a_1,\cdots\},B=\{b_1,\cdots\},n>2$”）按原卷作
%      半角，中文行文里的逗号按原卷作全角，均未强改；
%   ⑥ 第 17—21 题的子问原卷印作“（1）…（2）…”（另起一段），本稿按 高一69／70 成例用
%      行内“（1）（2）（3）”段，解析中的子问标号一律用 \stepm{(1)}{…}；
%      第 21 题 (2) 解析内层的“①”“②”“③”照原卷作 \stepm{①}{…} 等；
%   ⑦ 每道解答题原卷未留作答空白，本稿按全稿惯例加 \ansspace（末题用 \ansspace*）；
%   ⑧ 选择题的作答括号原卷排半角“(　　)”（第 14、15、16 题括号内空白、第 13 题括号内
%      印 B），本稿按全稿惯例排 \mbox{（\quad）}；选项行排照原卷——第 13 题两行两条
%      （A B 一行、C D 一行），第 14、15、16 题一行四条；
%   ⑨ 第 8 题解析、第 9 题解析、第 17 题题干与解析中的多行不等式组，原卷均排作“{”
%      连接的多行式，本稿按全稿惯例排 cases；
%   ⑩ 第 10 题解析、第 17 题解析、第 19 题 (1) 解析的并集符号原卷排大号 $\bigcup$
%      （显示体），本稿按全稿惯例排行内 $\cup$；
%   ⑪ 第 8 题解析的判别式照原卷作空心“$\triangle$”（与第 7 题解析的“$\Delta$”不同形），
%      照录未改；
%   ⑫ 第 16 题解析第 4 行的“$\forall$”（全称量词）照原卷符号排，未改写为中文“任意”。
%
% 原卷笔误/存疑（均照录未改，正文对应处另有 % 注）：
%   ① 第 19 题 (1) 连用两个“若”——“若 $a=4$，若全集为实数集 $R$”，后一“若”应作
%      “当”或删；原卷如此，照录未改。
%   ② 第 20 题题干“可以售出 8 万本**、**据市场调查”——“万本”后用顿号，按文意应为
%      逗号或句号；原卷如此，照录未改。
%   ③ 第 15 题解析对代数式 (1) 只得出“$y/x+x/y\geqslant2$”（同号时），未讨论取等是否
%      可得——取 $x=y$ 时确能取到 2——却总结“最小值为 2 的代数式有 1 个，故选 $B$”；
%      按解析自身的推理 (1) 也应计入，结论“1 个”与推理矛盾（若计 (1) 应为 2 个、选 C）。
%      原卷如此，照录未改。
%   ④ 第 21 题题干“定义：对于集合 $A$，则称集合 $B=\cdots$”——“对于……”与“则称……”
%      句式杂糅（应作“对于集合 $A$，称……”或“若集合 $A$，则称……”）；原卷如此，照录未改。
%   ⑤ 第 13 题题干末“……的（　）条件”后无任何标点（其余各题题干末均有句点）；
%      原卷如此，照录未改。
%
% 行数核对：逐题数 \step/\stepm 条数与原卷【解析】印刷行数，15 题逐一相符，
% 合计 86 条 / 86 行（Q6：1；Q7：1；Q8：7；Q9：2；Q10：4；Q11：7；Q12：5；Q14：7；
% Q15：8；Q16：7；Q17：3；Q18：7；Q19：8；Q20：6；Q21：13）；另有 3 题只印【答案】
% （第 3、4、5 题）与第 13 题括号内答案，各提 \ans{}，不计入 \step 账。
%
% 插图：无题目配图或数学图形（7 页题面均无“如图/右图/图中”字样；页面另有粉色斜排水印，
% 见上）。
%
\documentclass[a4paper,11pt]{article}
\usepackage{common}

\begin{document}

\examtitlest[3]{2026--2027 年中国中学高一上 10 月月考}{集合与常用逻辑用语、不等式}

\bigsec{一、填空题}

\begin{enumerate}
\setcounter{enumi}{2}

\item 已知集合 $A=[1,+\infty)$，$B=[-1,3)$，则 $A\cap B=$\blankline[4.4em].

\ifanswer
\ans{$[1,3)$}
\fi

\item 用反证法证明命题“已知 $a$、$b$ 为正数，$n$ 为正整数，求证：如果 $a^n>b^n$，那么 $a>b$”时，应当作出的假设为：“假设\blankline[4.4em]”.

\ifanswer
\ans{$a\leqslant b$}
\fi

\item 用列举法表示集合 $\{y\mid y=x+1,x\in\N,x\leqslant1\}$：\blankline[4.4em].

\ifanswer
\ans{$\{1,2\}$}
\fi

\item 某班在一次测验中，有 36 人数学成绩不低于 80 分，有 20 人物理成绩不低于 80 分，有 15 人的数学、物理成绩都不低于 80 分，则这两科成绩中至少有一科不低于 80 分的人数为\blankline[4.4em]人.

\ifanswer
\solbegin
\step{至少有 $36+20-15=41$ 人.}
\solend
\fi

\item 若关于 $x$ 的不等式 $x^2-x+m<0$ 的解集是 $\varnothing$，则实数 $m$ 的取值范围是\blankline[4.4em].

\ifanswer
\solbegin
\step{由题意得 $\Delta=1-4m\leqslant0$，所以 $m\geqslant\dfrac14$.}
\solend
\fi

% 第 8 题解析的判别式原卷作空心“△”（与第 7 题解析的“Δ”不同形），照录（见文件头“排版差异”⑪）
\item 已知关于 $x$ 的方程 $x^2+mx+2m=0$ 的两个实数根分别为 $\alpha$，$\beta$，若 $\alpha^2+\beta^2=5$，则实数 $m=$\blankline[4.4em].

\ifanswer
\solbegin
\step{方程 $x^2+mx+2m=0$ 的两个实数根为 $\alpha$、$\beta$，}
\step{由根与系数的关系得 $\begin{cases}\alpha+\beta=-m\\ \alpha\beta=2m\end{cases}$，又因为 $\alpha^2+\beta^2=5$，}
\step{所以 $\alpha^2+\beta^2=(\alpha+\beta)^2-2\alpha\beta=m^2-4m=5$，整理得 $m^2-4m-5=0$，}
\step{解得 $m=5$ 或 $m=-1$，}
\step{当 $m=5$ 时，$\triangle=25-40=-15<0$，无实根，舍去；}
\step{当 $m=-1$ 时，$\triangle=1+8=9>0$，符合题意，}
\step{所以实数 $m=-1$.}
\solend
\fi

\item 已知集合 $A=\{x\mid6x+a>0\}$，若 $1\in A$，$0\notin A$，则实数 $a$ 的取值范围为\blankline[4.4em].

\ifanswer
\solbegin
\step{由 $A=\{x\mid6x+a>0\}$，$1\in A$，$0\notin A$，得 $\begin{cases}6+a>0\\ a\leqslant0\end{cases}$，解得 $-6<a\leqslant0$，}
\step{即实数 $a$ 的取值范围是 $(-6,0]$.}
\solend
\fi

\item 关于 $x$ 的不等式 $ax-b>0$ 的解集为 $(1,+\infty)$，则关于 $x$ 的不等式 $(ax+b)(x-2)>0$ 的解集为\blankline[4.4em].

\ifanswer
\solbegin
\step{由不等式 $ax-b>0$ 的解集为 $(1,+\infty)$，得 $a>0$ 且 $\dfrac ba=1$，即 $b=a$，}
\step{则 $(ax+b)(x-2)>0$，可化为 $a(x+1)(x-2)>0$.}
\step{即 $(x+1)(x-2)>0$，解得 $x<-1$ 或 $x>2$，}
\step{所以不等式的解集为 $(-\infty,-1)\cup(2,+\infty)$.}
\solend
\fi

\item 已知实数 $a_1$、$a_2$、$a_3$、$a_4$、$a_5$ 满足 $\{x\mid(x-a_1)(x-a_2)(x-a_3)(x-a_4)(x-a_5)=0\}=\{1,2,3,4\}$，则 $a_1+a_2+a_3+a_4+a_5$ 的最大值为\blankline[4.4em].

\ifanswer
\solbegin
\step{5 个实数 $a_1$、$a_2$、$a_3$、$a_4$、$a_5$ 的不同取值只有 1，2，3，4，}
\step{即这 5 个数均从 1，2，3，4 中选取.}
\step{集合有 4 个元素，说明 1，2，3，4 每个数至少出现一次，}
\step{因此恰好有一个数重复出现两次，其余三个数各出现一次.}
\step{要让 5 个数的总和最大，应让重复出现两次的数取集合中的最大值 4，}
\step{此时 5 个数为 1，2，3，4，4，总和为 $1+2+3+4+4=14$，}
\step{即所求最大值为 14.}
\solend
\fi

\item 定义集合运算“$\times$”：$A\times B=\{(x,y)\mid x\in A,y\in B\}$，称为 $A$、$B$ 两个集合的“卡氏积”，若 $A=\{x\mid x^2-|x|\leqslant0,x\in\N\}$，$B=\{1,2,3\}$，则 $(A\times B)\cap(B\times A)=$\blankline[4.4em].

\ifanswer
\solbegin
\step{因为 $x^2-|x|\leqslant0$，所以 $0\leqslant|x|\leqslant1$，又因为 $x\in\N$，}
\step{所以 $A=\{x\mid x^2-|x|\leqslant0,x\in\N\}=\{0,1\}$，又 $B=\{1,2,3\}$，}
\step{由“卡氏积”的定义得 $A\times B=\{(0,1),(0,2),(0,3),(1,1),(1,2),(1,3)\}$，}
\step{$B\times A=\{(1,0),(1,1),(2,0),(2,1),(3,0),(3,1)\}$，}
\step{所以 $(A\times B)\cap(B\times A)=\{(1,1)\}$.}
\solend
\fi

\end{enumerate}

\bigsec{二、选择题}

\begin{enumerate}
\setcounter{enumi}{12}

% 第 13 题答案“B”原卷印在题干括号内，且无【答案】/【解析】段；
% 按 高一37 第 13 题、高一63 第 14 题、高一67 第 15 题成例括号留空、答案另提 \ans{$B$}
% 原卷题干末“条件”后无标点，照录（见文件头“原卷笔误/存疑”⑤）
\item “$A\subseteq B$”是“$A\subset B$”的\mbox{（\quad）}条件

\keepwithprev
A. 充分不必要条件\qquad B. 必要不充分条件

C. 充要条件\qquad D. 既不充分也不必要条件

\ifanswer
\ans{$B$}
\fi

\item 已知 $a$、$b$、$c$ 均为实数，若 $a>b$ 且 $ab\neq0$，则下列不等关系中成立的是\mbox{（\quad）}

\keepwithprev
A. $ac^2>bc^2$\qquad B. $\dfrac1b>\dfrac1a$\qquad C. $|a|>|b|$\qquad D. $c-b>c-a$

\ifanswer
\solbegin
\step{对于 $A$，若 $a>b$，当 $c=0$ 时，有 $ac^2=bc^2=0$，故 $A$ 错误；}
\step{对于 $B$，$\dfrac1b-\dfrac1a=\dfrac{a-b}{ab}$，因为 $a>b$，所以 $a-b>0$，}
\step{当 $a$、$b$ 同号时，$ab>0$，则 $\dfrac1b-\dfrac1a>0$，即 $\dfrac1b>\dfrac1a$，}
\step{当 $a$、$b$ 异号时，$ab<0$，则 $\dfrac1b-\dfrac1a<0$，即 $\dfrac1b<\dfrac1a$，故 $B$ 错误；}
\step{对于 $C$，当 $a=1$，$b=-3$ 时，$|a|<|b|$，故 $C$ 错误；}
\step{对于 $D$，若 $a>b$，则 $-b>-a$，$c-b>c-a$，故 $D$ 正确.}
\step{故选 $D$.}
\solend
\fi

% 原卷解析对 (1) 未讨论取等是否可得（x=y 时确能取 2），却总结“有 1 个、故选 B”，
% 与推理矛盾，照录（见文件头“原卷笔误/存疑”③）
\item 以下三个代数式：（1）$\dfrac yx+\dfrac xy$；（2）$|x|+\dfrac1{|x|}$；（3）$x^2+2+\dfrac1{x^2+2}$．其中的字母均表示实数，且代数式均有意义，那么代数式的最小值为 2 的个数是\mbox{（\quad）}

\keepwithprev
A. $0$\qquad B. $1$\qquad C. $2$\qquad D. $3$

\ifanswer
\solbegin
\step{（1）仅当 $x,y$ 同号时，$\dfrac yx,\dfrac xy>0$，由基本不等式得 $\dfrac yx+\dfrac xy\geqslant2$；}
\step{若 $x,y$ 异号，$\dfrac yx,\dfrac xy<0$，故 $\dfrac yx+\dfrac xy<0$，此时最小值不是 2；}
\step{（2）$|x|+\dfrac1{|x|}$ 中，$|x|>0$，由基本不等式 $|x|+\dfrac1{|x|}\geqslant2\sqrt{|x|\cdot\dfrac1{|x|}}=2$，}
\step{当且仅当 $|x|=1$，即 $x=\pm1$ 时取等号，故最小值为 2；}
\step{（3）令 $t=x^2+2$，则 $t\geqslant2$，函数化为 $t+\dfrac1t$，}
\step{对勾函数 $t+\dfrac1t$ 在 $[2,+\infty)$ 上严格增，最小值为 $2+\dfrac12=\dfrac52>2$，}
\step{等号成立条件 $t=1$ 无法取到，故最小值不是 2；}
\step{综上，最小值为 2 的代数式有 1 个，故选 $B$.}
\solend
\fi

\item 已知 $A$、$B$ 为非空实数集，$\Omega$ 为平面直角坐标系中的一些点构成的集合，集合 $C=\{y\mid$对任意 $x\in A$，有 $(x,y)\in\Omega\}$，集合 $D=\{x\mid$对任意 $y\in B$，有 $(x,y)\in\Omega\}$，对于下列两个命题：①若 $C\subseteq B$，则 $D\subseteq A$；②若 $C\supseteq B$，则 $D\subseteq A$；其中判断正确的是\mbox{（\quad）}

\keepwithprev
A. ①②都正确\qquad B. ①②都错误\qquad C. ①正确②错误\qquad D. ①错误②正确

\ifanswer
\solbegin
\step{设 $A=\{a_1,a_2,a_3,\cdots,a_n\},B=\{b_1,b_2,b_3,\cdots,b_n\},n>2$，}
\step{若 $\Omega=\{(c,b_1),(c,b_2),\cdots,(c,b_n)\},c\notin A$，}
\step{此时 $C=\varnothing\subseteq B,D=\{c\},D\subseteq A$ 错误，故①错误；}
% 原卷此处 N* 排正体（\mathrm{N}*），本稿统一 \N^*（见文件头“排版差异”④）；
% 全称量词“∀”照原卷符号排（见文件头“排版差异”⑫）
\step{若 $C\supseteq B$，则 $\forall(a_i,b_j)(1\leqslant i,j\leqslant n,i,j\in\N^*)\in\Omega$，则 $D\supseteq A$，}
\step{且 $A=\{a_1\},B=\{b_1,b_2\}$，}
\step{若 $\Omega=\{(a_1,b_1),(a_2,b_1),(a_1,b_2),(a_2,b_2),(c,b_1),(c,b_2)\},D$ 真包含 $A$，故②错误.}
\step{故选 $B$.}
\solend
\fi

\end{enumerate}

\bigsec{三、解答题}

\begin{enumerate}
\setcounter{enumi}{16}

\item 解下列关于 $x$ 的不等式组：$\begin{cases}|x-1|<2\\ \dfrac{x-2}{x-1}\geqslant0\end{cases}$.

\ifanswer
\solbegin
\step{解 $|x-1|<2$ 得 $-2<x-1<2$，解得 $-1<x<3$.}
\step{分式不等式 $\dfrac{x-2}{x-1}\geqslant0$ 等价于 $\begin{cases}(x-2)(x-1)\geqslant0\\ x-1\neq0\end{cases}$，解得 $x<1$ 或 $x\geqslant2$.}
\step{即 $\begin{cases}|x-1|<2\\ \dfrac{x-2}{x-1}\geqslant0\end{cases}$ 的解集为 $(-1,1)\cup[2,3)$.}
\solend
\fi

\ansspace[5]

\item 比较多项式 $2x^2-3x+1$ 与 $x^2-1$ 的大小.

\ifanswer
\solbegin
\step{作差得 $(2x^2-3x+1)-(x^2-1)=x^2-3x+2=(x-1)(x-2)$，}
\step{当 $x=1$ 或 $x=2$ 时，$(x-1)(x-2)=0$，即 $2x^2-3x+1=x^2-1$，}
\step{当 $x<1$ 或 $x>2$ 时，$(x-1)(x-2)>0$，即 $2x^2-3x+1>x^2-1$，}
\step{当 $1<x<2$ 时，$(x-1)(x-2)<0$，即 $2x^2-3x+1<x^2-1$；}
\step{综上，当 $x=1$ 或 $x=2$ 时，$2x^2-3x+1=x^2-1$；}
\step{当 $x<1$ 或 $x>2$ 时，$2x^2-3x+1>x^2-1$；}
\step{当 $1<x<2$ 时，$2x^2-3x+1<x^2-1$.}
\solend
\fi

\ansspace[6]

% 原卷 (1) 连用两个“若”（“若 a=4，若全集为实数集 R”），照录（见文件头“原卷笔误/存疑”①）
\item 已知集合 $A=\left\{x\mid\dfrac{(a+1)x}{x-1}<a\right\}$，集合 $B=[1,4]$，$a$ 为实常数；

（1）若 $a=4$，若全集为实数集 $\R$，求集合 $\overline{A}$（用区间形式表示）；

（2）若 $A\subseteq B$，求 $a$ 的取值范围.

\ifanswer
\solbegin
\stepm{(1)}{若 $a=4$，集合 $A=\left\{x\mid\dfrac{5x}{x-1}<4\right\}=\left\{x\mid\dfrac{5x}{x-1}-4<0\right\}=\left\{x\mid\dfrac{x+4}{x-1}<0\right\}$}
\step{$=(-4,1)$，因此 $\overline{A}=(-\infty,-4]\cup[1,+\infty)$；}
\stepm{(2)}{当 $-a<1$，即 $a>-1$ 时，$A=(-a,1)$，}
\step{要使 $A\subseteq[1,4]$，需 $-a\geqslant1$，即 $a\leqslant-1$，与 $a>-1$ 矛盾，舍去；}
\step{当 $-a>1$，即 $a<-1$ 时，$A=(1,-a)$，}
\step{要使 $A\subseteq[1,4]$，需 $-a\leqslant4$，即 $a\geqslant-4$，所以 $-4\leqslant a<-1$；}
\step{当 $-a=1$，即 $a=-1$ 时，$A=\varnothing$，满足 $A\subseteq B$；}
\step{综上，实数 $a$ 的取值范围为 $[-4,-1]$.}
\solend
\fi

\ansspace[7]

% 原卷题干“可以售出 8 万本、据市场调查”处用顿号，照录（见文件头“原卷笔误/存疑”②）
\item 某种杂志原以每本 25 元的价格销售，可以售出 8 万本、据市场调查，杂志的单价每提高 1 元，销售量就减少 0.2 万本，假定杂志的成本是每本 10 元（不计其他成本）.

（1）当杂志以每本 30 元定价销售时，求销售该杂志所获利润；

（2）在每本 25 元的价格的基础上提高定价，试确定杂志的定价在什么范围内可使得销售利润不低于 140 万元.

\ifanswer
\solbegin
\stepm{(1)}{当杂志以 30 元定价销售时，}
\step{销售该杂志所获利润为 $[8-(30-25)\times0.2]\times(30-10)=140$ 万元；}
\stepm{(2)}{在每本 25 元的价格的基础上提高定价，设杂志的定价为每本 $x$ 元，}
\step{由题意得 $[8-(x-25)\times0.2]\times(x-10)\geqslant140$，解得 $30\leqslant x\leqslant45$，}
\step{所以杂志的定价在每本不低于 30 元且不高于 45 元的范围内，}
\step{可使得销售利润不低于 140 万元.}
\solend
\fi

\ansspace[7]

% 原卷题干“定义：对于集合 A，则称集合 B=…”句式杂糅，照录（见文件头“原卷笔误/存疑”④）
\item 定义：对于集合 $A$，则称集合 $B=\left\{x\mid x=\dfrac{m+n}{1+mn},m,n\in A,1+mn\neq0\right\}$ 为集合 $A$ 的生成集（其中 $m$、$n$ 可以相等）.

（1）若集合 $A=\{-1,1\}$ 的生成集为 $B$，求集合 $B$；

（2）若集合 $A=\{a,2\}$ 的生成集为 $B$，$B$ 的子集个数为 4 个，求实数 $a$ 的值；

（3）若集合 $A=\{x\mid-1<x<1\}$ 的生成集为 $B$，求证：$B\subseteq A$.

\ifanswer
\solbegin
\stepm{(1)}{因为 $1+mn\neq0$，}
\step{当 $m=n=1$ 时，$x=\dfrac{m+n}{1+mn}=1$；当 $m=n=-1$ 时，$x=\dfrac{m+n}{1+mn}=-1$；}
\step{所以 $B=\{-1,1\}$；}
\stepm{(2)}{$B$ 的子集个数为 4 个，故 $B$ 集合的元素有 2 个，}
\step{当 $m=n=2$ 时，$x=\dfrac{m+n}{1+mn}=\dfrac45$，当 $m=n=a$ 时，$x=\dfrac{m+n}{1+mn}=\dfrac{2a}{1+a^2}$，}
\step{当 $m$ 与 $n$ 分别取 2 和 $a$ 时，$x=\dfrac{a+2}{1+2a}$，}
\stepm{①}{$\dfrac{2a}{1+a^2}=\dfrac45$，解得 $a=2$（舍）或 $a=\dfrac12$，}
\stepm{②}{$\dfrac{a+2}{1+2a}=\dfrac45$，解得 $a=2$（舍），}
\stepm{③}{$\dfrac{2a}{1+a^2}=\dfrac{a+2}{1+2a}$，解得 $a=2$（舍）或 $a=1$ 或 $a=-1$，}
\step{综上所述，$a=1$ 或 $a=-1$ 或 $a=\dfrac12$，即 $a\in\left\{-1,1,\dfrac12\right\}$；}
\stepm{(3)}{$\forall m$、$n\in A$，则 $\dfrac{m+n}{1+mn}+1=\dfrac{(m+1)(n+1)}{1+mn}>0$，故 $\dfrac{m+n}{1+mn}>-1$；}
\step{$\dfrac{m+n}{1+mn}-1=-\dfrac{(m-1)(n-1)}{1+mn}<0$，故 $\dfrac{m+n}{1+mn}<1$；}
\step{故 $B\subseteq\{x\mid-1<x<1\}$，即 $B\subseteq A$.}
\solend
\fi

\ansspace*[10]

\end{enumerate}

\end{document}
"
% 紧凑版：xelatex "\def\iscompact{1}% 高一71_中国中学_10月月考.tex
%   —— 高一上数学“2026-2027 年中国中学高一上 10 月月考”（讲评版）
% 试卷版：xelatex 高一71_中国中学_10月月考.tex
% 答案版：xelatex "\def\isanswer{1}% 高一71_中国中学_10月月考.tex
%   —— 高一上数学“2026-2027 年中国中学高一上 10 月月考”（讲评版）
% 试卷版：xelatex 高一71_中国中学_10月月考.tex
% 答案版：xelatex "\def\isanswer{1}% 高一71_中国中学_10月月考.tex
%   —— 高一上数学“2026-2027 年中国中学高一上 10 月月考”（讲评版）
% 试卷版：xelatex 高一71_中国中学_10月月考.tex
% 答案版：xelatex "\def\isanswer{1}\input{高一71_中国中学_10月月考.tex}"
% 紧凑版：xelatex "\def\iscompact{1}\input{高一71_中国中学_10月月考.tex}"
%
% 原始材料：微信公众号“魔都数学加油站”文章，连载“27学年高一71”；
% 原文标题“27学年高一71——2026-2027年上海市中国中学高一上10月月考”；
% 原文链接：https://mp.weixin.qq.com/s/UZrP53l3GIin040L1EkljA。
% 页面标注发布时间 2026-10-19 22:47（上海）（与 高一69、高一70 页面读出同一 ct 值）。
% 正文为 7 张 PNG 页（`原图/高一71_中国中学_10月月考/01..07.png`），全部 4762×6735，
% 均按页序读取。7 页全为试卷正文页（卷首标题在第 1 页页顶，卷末解答题第 21 题 (3)
% 解析止于第 7 页），无封面页、目录页、二维码或宣传页。粉色斜排“魔都数学加油站”
% 为卷面水印，与公众号页脚一并未录入。
%
% 卷首标题：照卷面印作“2026-2027 年中国中学高一上 10 月月考”（黑体加粗一行）。
% 原卷年份区间用半角连字符，本稿按全稿惯例排 en dash（同 高一68、高一69、高一70）；
% 原卷只有一行标题、无副标题。
%
% 篇幅与题号：原卷自第 3 题起编号（第 1、2 题不在本卷内）。填空题 3—12（10 题）、
% 选择题 13—16（4 题）、解答题 17—21（5 题），共 19 题。本稿沿用原节与原题号，
% 只改排为全稿统一的大节样式（原卷小节标题为“一、填空题”“二、选择题”“三、解答题”）。
%
% 讲评版：原卷每题下直接印【答案】或【解析】。**只印【答案】的 3 题**（第 3、4、5 题），
% 本稿只给 \ans{}；**第 13 题既无【答案】也无【解析】段，答案“B”直接印在题干括号内
% （“(  B  )”）**，按 高一37 第 13 题、高一63 第 14 题、高一67 第 15 题成例把括号留空、
% 答案另提为 \ans{$B$}。其余 15 题全部印【解析】（按全稿口径这类题不另提 \ans{}——
% 第 14—16 题的“故选 D／B／B”写在解析末句，亦照录在解析内）。
% 除第 13 题外，全卷选择题作答括号内均为空白（第 14、15、16 题）。
%
% 副标题：原卷未印副标题。本稿按“知识模块”口径标作“集合与常用逻辑用语、不等式”：
% 19 题中集合与常用逻辑用语 10 题（第 3、4、5、6、9、11、12、13、16、21 题；其中第 4 题
% 属反证法，一并计入常用逻辑用语），不等式 9 题（第 7、8、10、14、15、17、18、19、20 题；
% 第 19 题 (1) 属补集、(2) 属分式不等式，按主干计入不等式），两类均确有分量，故并标。
% 口径同 高一63、高一68、高一69、高一70。
%
% 与原卷的排版差异：
%   ① 原卷每题下直接印【答案】/【解析】，本稿按全稿惯例将答案与解析置于答案版，试卷版
%      保持干净；第 13 题印在题干括号内的答案“B”同样移入 \ans{}，见上；
%   ② 原卷未印副标题，本稿按“知识模块”口径补出；
%   ③ 原卷填空横线约四个字宽，本稿按全稿惯例统一排 \blankline[4.4em]；
%   ④ 原卷数集记号不统一：实数集排斜体 $R$（第 19 题 (1)），自然数集排斜体 $N$
%      （第 5 题题干、第 12 题题干与解析共三处），第 16 题解析又排正体 $\mathrm{N}^*$；
%      本稿按全稿惯例统一改用 $\R$、$\N$、$\N^*$（同 高一68／69／70 排版差异④）；
%      空集记号原卷作 $\varnothing$，照录；
%   ⑤ 原卷句末标点混用：第 15 题题干“（3）$x^2+2+\dfrac1{x^2+2}$**．**其中”、
%      第 17 题题干末“$\dfrac{x-2}{x-1}\geqslant0$**．**”两处放大确认（第 15 题为全角
%      句点、第 17 题为西文句点），其余各题末为西文句点，已逐处照录；题干内冒号——
%      第 5 题“$\{y\mid\cdots\}$**：**”、第 12 题““×”**：**”、第 17 题“不等式组**：**”、
%      第 21 题“定义**：**”——均为全角冒号，照录；数学式内的逗号（如第 12 题
%      “$x\in A,y\in B$”、第 16 题“$A=\{a_1,\cdots\},B=\{b_1,\cdots\},n>2$”）按原卷作
%      半角，中文行文里的逗号按原卷作全角，均未强改；
%   ⑥ 第 17—21 题的子问原卷印作“（1）…（2）…”（另起一段），本稿按 高一69／70 成例用
%      行内“（1）（2）（3）”段，解析中的子问标号一律用 \stepm{(1)}{…}；
%      第 21 题 (2) 解析内层的“①”“②”“③”照原卷作 \stepm{①}{…} 等；
%   ⑦ 每道解答题原卷未留作答空白，本稿按全稿惯例加 \ansspace（末题用 \ansspace*）；
%   ⑧ 选择题的作答括号原卷排半角“(　　)”（第 14、15、16 题括号内空白、第 13 题括号内
%      印 B），本稿按全稿惯例排 \mbox{（\quad）}；选项行排照原卷——第 13 题两行两条
%      （A B 一行、C D 一行），第 14、15、16 题一行四条；
%   ⑨ 第 8 题解析、第 9 题解析、第 17 题题干与解析中的多行不等式组，原卷均排作“{”
%      连接的多行式，本稿按全稿惯例排 cases；
%   ⑩ 第 10 题解析、第 17 题解析、第 19 题 (1) 解析的并集符号原卷排大号 $\bigcup$
%      （显示体），本稿按全稿惯例排行内 $\cup$；
%   ⑪ 第 8 题解析的判别式照原卷作空心“$\triangle$”（与第 7 题解析的“$\Delta$”不同形），
%      照录未改；
%   ⑫ 第 16 题解析第 4 行的“$\forall$”（全称量词）照原卷符号排，未改写为中文“任意”。
%
% 原卷笔误/存疑（均照录未改，正文对应处另有 % 注）：
%   ① 第 19 题 (1) 连用两个“若”——“若 $a=4$，若全集为实数集 $R$”，后一“若”应作
%      “当”或删；原卷如此，照录未改。
%   ② 第 20 题题干“可以售出 8 万本**、**据市场调查”——“万本”后用顿号，按文意应为
%      逗号或句号；原卷如此，照录未改。
%   ③ 第 15 题解析对代数式 (1) 只得出“$y/x+x/y\geqslant2$”（同号时），未讨论取等是否
%      可得——取 $x=y$ 时确能取到 2——却总结“最小值为 2 的代数式有 1 个，故选 $B$”；
%      按解析自身的推理 (1) 也应计入，结论“1 个”与推理矛盾（若计 (1) 应为 2 个、选 C）。
%      原卷如此，照录未改。
%   ④ 第 21 题题干“定义：对于集合 $A$，则称集合 $B=\cdots$”——“对于……”与“则称……”
%      句式杂糅（应作“对于集合 $A$，称……”或“若集合 $A$，则称……”）；原卷如此，照录未改。
%   ⑤ 第 13 题题干末“……的（　）条件”后无任何标点（其余各题题干末均有句点）；
%      原卷如此，照录未改。
%
% 行数核对：逐题数 \step/\stepm 条数与原卷【解析】印刷行数，15 题逐一相符，
% 合计 86 条 / 86 行（Q6：1；Q7：1；Q8：7；Q9：2；Q10：4；Q11：7；Q12：5；Q14：7；
% Q15：8；Q16：7；Q17：3；Q18：7；Q19：8；Q20：6；Q21：13）；另有 3 题只印【答案】
% （第 3、4、5 题）与第 13 题括号内答案，各提 \ans{}，不计入 \step 账。
%
% 插图：无题目配图或数学图形（7 页题面均无“如图/右图/图中”字样；页面另有粉色斜排水印，
% 见上）。
%
\documentclass[a4paper,11pt]{article}
\usepackage{common}

\begin{document}

\examtitlest[3]{2026--2027 年中国中学高一上 10 月月考}{集合与常用逻辑用语、不等式}

\bigsec{一、填空题}

\begin{enumerate}
\setcounter{enumi}{2}

\item 已知集合 $A=[1,+\infty)$，$B=[-1,3)$，则 $A\cap B=$\blankline[4.4em].

\ifanswer
\ans{$[1,3)$}
\fi

\item 用反证法证明命题“已知 $a$、$b$ 为正数，$n$ 为正整数，求证：如果 $a^n>b^n$，那么 $a>b$”时，应当作出的假设为：“假设\blankline[4.4em]”.

\ifanswer
\ans{$a\leqslant b$}
\fi

\item 用列举法表示集合 $\{y\mid y=x+1,x\in\N,x\leqslant1\}$：\blankline[4.4em].

\ifanswer
\ans{$\{1,2\}$}
\fi

\item 某班在一次测验中，有 36 人数学成绩不低于 80 分，有 20 人物理成绩不低于 80 分，有 15 人的数学、物理成绩都不低于 80 分，则这两科成绩中至少有一科不低于 80 分的人数为\blankline[4.4em]人.

\ifanswer
\solbegin
\step{至少有 $36+20-15=41$ 人.}
\solend
\fi

\item 若关于 $x$ 的不等式 $x^2-x+m<0$ 的解集是 $\varnothing$，则实数 $m$ 的取值范围是\blankline[4.4em].

\ifanswer
\solbegin
\step{由题意得 $\Delta=1-4m\leqslant0$，所以 $m\geqslant\dfrac14$.}
\solend
\fi

% 第 8 题解析的判别式原卷作空心“△”（与第 7 题解析的“Δ”不同形），照录（见文件头“排版差异”⑪）
\item 已知关于 $x$ 的方程 $x^2+mx+2m=0$ 的两个实数根分别为 $\alpha$，$\beta$，若 $\alpha^2+\beta^2=5$，则实数 $m=$\blankline[4.4em].

\ifanswer
\solbegin
\step{方程 $x^2+mx+2m=0$ 的两个实数根为 $\alpha$、$\beta$，}
\step{由根与系数的关系得 $\begin{cases}\alpha+\beta=-m\\ \alpha\beta=2m\end{cases}$，又因为 $\alpha^2+\beta^2=5$，}
\step{所以 $\alpha^2+\beta^2=(\alpha+\beta)^2-2\alpha\beta=m^2-4m=5$，整理得 $m^2-4m-5=0$，}
\step{解得 $m=5$ 或 $m=-1$，}
\step{当 $m=5$ 时，$\triangle=25-40=-15<0$，无实根，舍去；}
\step{当 $m=-1$ 时，$\triangle=1+8=9>0$，符合题意，}
\step{所以实数 $m=-1$.}
\solend
\fi

\item 已知集合 $A=\{x\mid6x+a>0\}$，若 $1\in A$，$0\notin A$，则实数 $a$ 的取值范围为\blankline[4.4em].

\ifanswer
\solbegin
\step{由 $A=\{x\mid6x+a>0\}$，$1\in A$，$0\notin A$，得 $\begin{cases}6+a>0\\ a\leqslant0\end{cases}$，解得 $-6<a\leqslant0$，}
\step{即实数 $a$ 的取值范围是 $(-6,0]$.}
\solend
\fi

\item 关于 $x$ 的不等式 $ax-b>0$ 的解集为 $(1,+\infty)$，则关于 $x$ 的不等式 $(ax+b)(x-2)>0$ 的解集为\blankline[4.4em].

\ifanswer
\solbegin
\step{由不等式 $ax-b>0$ 的解集为 $(1,+\infty)$，得 $a>0$ 且 $\dfrac ba=1$，即 $b=a$，}
\step{则 $(ax+b)(x-2)>0$，可化为 $a(x+1)(x-2)>0$.}
\step{即 $(x+1)(x-2)>0$，解得 $x<-1$ 或 $x>2$，}
\step{所以不等式的解集为 $(-\infty,-1)\cup(2,+\infty)$.}
\solend
\fi

\item 已知实数 $a_1$、$a_2$、$a_3$、$a_4$、$a_5$ 满足 $\{x\mid(x-a_1)(x-a_2)(x-a_3)(x-a_4)(x-a_5)=0\}=\{1,2,3,4\}$，则 $a_1+a_2+a_3+a_4+a_5$ 的最大值为\blankline[4.4em].

\ifanswer
\solbegin
\step{5 个实数 $a_1$、$a_2$、$a_3$、$a_4$、$a_5$ 的不同取值只有 1，2，3，4，}
\step{即这 5 个数均从 1，2，3，4 中选取.}
\step{集合有 4 个元素，说明 1，2，3，4 每个数至少出现一次，}
\step{因此恰好有一个数重复出现两次，其余三个数各出现一次.}
\step{要让 5 个数的总和最大，应让重复出现两次的数取集合中的最大值 4，}
\step{此时 5 个数为 1，2，3，4，4，总和为 $1+2+3+4+4=14$，}
\step{即所求最大值为 14.}
\solend
\fi

\item 定义集合运算“$\times$”：$A\times B=\{(x,y)\mid x\in A,y\in B\}$，称为 $A$、$B$ 两个集合的“卡氏积”，若 $A=\{x\mid x^2-|x|\leqslant0,x\in\N\}$，$B=\{1,2,3\}$，则 $(A\times B)\cap(B\times A)=$\blankline[4.4em].

\ifanswer
\solbegin
\step{因为 $x^2-|x|\leqslant0$，所以 $0\leqslant|x|\leqslant1$，又因为 $x\in\N$，}
\step{所以 $A=\{x\mid x^2-|x|\leqslant0,x\in\N\}=\{0,1\}$，又 $B=\{1,2,3\}$，}
\step{由“卡氏积”的定义得 $A\times B=\{(0,1),(0,2),(0,3),(1,1),(1,2),(1,3)\}$，}
\step{$B\times A=\{(1,0),(1,1),(2,0),(2,1),(3,0),(3,1)\}$，}
\step{所以 $(A\times B)\cap(B\times A)=\{(1,1)\}$.}
\solend
\fi

\end{enumerate}

\bigsec{二、选择题}

\begin{enumerate}
\setcounter{enumi}{12}

% 第 13 题答案“B”原卷印在题干括号内，且无【答案】/【解析】段；
% 按 高一37 第 13 题、高一63 第 14 题、高一67 第 15 题成例括号留空、答案另提 \ans{$B$}
% 原卷题干末“条件”后无标点，照录（见文件头“原卷笔误/存疑”⑤）
\item “$A\subseteq B$”是“$A\subset B$”的\mbox{（\quad）}条件

\keepwithprev
A. 充分不必要条件\qquad B. 必要不充分条件

C. 充要条件\qquad D. 既不充分也不必要条件

\ifanswer
\ans{$B$}
\fi

\item 已知 $a$、$b$、$c$ 均为实数，若 $a>b$ 且 $ab\neq0$，则下列不等关系中成立的是\mbox{（\quad）}

\keepwithprev
A. $ac^2>bc^2$\qquad B. $\dfrac1b>\dfrac1a$\qquad C. $|a|>|b|$\qquad D. $c-b>c-a$

\ifanswer
\solbegin
\step{对于 $A$，若 $a>b$，当 $c=0$ 时，有 $ac^2=bc^2=0$，故 $A$ 错误；}
\step{对于 $B$，$\dfrac1b-\dfrac1a=\dfrac{a-b}{ab}$，因为 $a>b$，所以 $a-b>0$，}
\step{当 $a$、$b$ 同号时，$ab>0$，则 $\dfrac1b-\dfrac1a>0$，即 $\dfrac1b>\dfrac1a$，}
\step{当 $a$、$b$ 异号时，$ab<0$，则 $\dfrac1b-\dfrac1a<0$，即 $\dfrac1b<\dfrac1a$，故 $B$ 错误；}
\step{对于 $C$，当 $a=1$，$b=-3$ 时，$|a|<|b|$，故 $C$ 错误；}
\step{对于 $D$，若 $a>b$，则 $-b>-a$，$c-b>c-a$，故 $D$ 正确.}
\step{故选 $D$.}
\solend
\fi

% 原卷解析对 (1) 未讨论取等是否可得（x=y 时确能取 2），却总结“有 1 个、故选 B”，
% 与推理矛盾，照录（见文件头“原卷笔误/存疑”③）
\item 以下三个代数式：（1）$\dfrac yx+\dfrac xy$；（2）$|x|+\dfrac1{|x|}$；（3）$x^2+2+\dfrac1{x^2+2}$．其中的字母均表示实数，且代数式均有意义，那么代数式的最小值为 2 的个数是\mbox{（\quad）}

\keepwithprev
A. $0$\qquad B. $1$\qquad C. $2$\qquad D. $3$

\ifanswer
\solbegin
\step{（1）仅当 $x,y$ 同号时，$\dfrac yx,\dfrac xy>0$，由基本不等式得 $\dfrac yx+\dfrac xy\geqslant2$；}
\step{若 $x,y$ 异号，$\dfrac yx,\dfrac xy<0$，故 $\dfrac yx+\dfrac xy<0$，此时最小值不是 2；}
\step{（2）$|x|+\dfrac1{|x|}$ 中，$|x|>0$，由基本不等式 $|x|+\dfrac1{|x|}\geqslant2\sqrt{|x|\cdot\dfrac1{|x|}}=2$，}
\step{当且仅当 $|x|=1$，即 $x=\pm1$ 时取等号，故最小值为 2；}
\step{（3）令 $t=x^2+2$，则 $t\geqslant2$，函数化为 $t+\dfrac1t$，}
\step{对勾函数 $t+\dfrac1t$ 在 $[2,+\infty)$ 上严格增，最小值为 $2+\dfrac12=\dfrac52>2$，}
\step{等号成立条件 $t=1$ 无法取到，故最小值不是 2；}
\step{综上，最小值为 2 的代数式有 1 个，故选 $B$.}
\solend
\fi

\item 已知 $A$、$B$ 为非空实数集，$\Omega$ 为平面直角坐标系中的一些点构成的集合，集合 $C=\{y\mid$对任意 $x\in A$，有 $(x,y)\in\Omega\}$，集合 $D=\{x\mid$对任意 $y\in B$，有 $(x,y)\in\Omega\}$，对于下列两个命题：①若 $C\subseteq B$，则 $D\subseteq A$；②若 $C\supseteq B$，则 $D\subseteq A$；其中判断正确的是\mbox{（\quad）}

\keepwithprev
A. ①②都正确\qquad B. ①②都错误\qquad C. ①正确②错误\qquad D. ①错误②正确

\ifanswer
\solbegin
\step{设 $A=\{a_1,a_2,a_3,\cdots,a_n\},B=\{b_1,b_2,b_3,\cdots,b_n\},n>2$，}
\step{若 $\Omega=\{(c,b_1),(c,b_2),\cdots,(c,b_n)\},c\notin A$，}
\step{此时 $C=\varnothing\subseteq B,D=\{c\},D\subseteq A$ 错误，故①错误；}
% 原卷此处 N* 排正体（\mathrm{N}*），本稿统一 \N^*（见文件头“排版差异”④）；
% 全称量词“∀”照原卷符号排（见文件头“排版差异”⑫）
\step{若 $C\supseteq B$，则 $\forall(a_i,b_j)(1\leqslant i,j\leqslant n,i,j\in\N^*)\in\Omega$，则 $D\supseteq A$，}
\step{且 $A=\{a_1\},B=\{b_1,b_2\}$，}
\step{若 $\Omega=\{(a_1,b_1),(a_2,b_1),(a_1,b_2),(a_2,b_2),(c,b_1),(c,b_2)\},D$ 真包含 $A$，故②错误.}
\step{故选 $B$.}
\solend
\fi

\end{enumerate}

\bigsec{三、解答题}

\begin{enumerate}
\setcounter{enumi}{16}

\item 解下列关于 $x$ 的不等式组：$\begin{cases}|x-1|<2\\ \dfrac{x-2}{x-1}\geqslant0\end{cases}$.

\ifanswer
\solbegin
\step{解 $|x-1|<2$ 得 $-2<x-1<2$，解得 $-1<x<3$.}
\step{分式不等式 $\dfrac{x-2}{x-1}\geqslant0$ 等价于 $\begin{cases}(x-2)(x-1)\geqslant0\\ x-1\neq0\end{cases}$，解得 $x<1$ 或 $x\geqslant2$.}
\step{即 $\begin{cases}|x-1|<2\\ \dfrac{x-2}{x-1}\geqslant0\end{cases}$ 的解集为 $(-1,1)\cup[2,3)$.}
\solend
\fi

\ansspace[5]

\item 比较多项式 $2x^2-3x+1$ 与 $x^2-1$ 的大小.

\ifanswer
\solbegin
\step{作差得 $(2x^2-3x+1)-(x^2-1)=x^2-3x+2=(x-1)(x-2)$，}
\step{当 $x=1$ 或 $x=2$ 时，$(x-1)(x-2)=0$，即 $2x^2-3x+1=x^2-1$，}
\step{当 $x<1$ 或 $x>2$ 时，$(x-1)(x-2)>0$，即 $2x^2-3x+1>x^2-1$，}
\step{当 $1<x<2$ 时，$(x-1)(x-2)<0$，即 $2x^2-3x+1<x^2-1$；}
\step{综上，当 $x=1$ 或 $x=2$ 时，$2x^2-3x+1=x^2-1$；}
\step{当 $x<1$ 或 $x>2$ 时，$2x^2-3x+1>x^2-1$；}
\step{当 $1<x<2$ 时，$2x^2-3x+1<x^2-1$.}
\solend
\fi

\ansspace[6]

% 原卷 (1) 连用两个“若”（“若 a=4，若全集为实数集 R”），照录（见文件头“原卷笔误/存疑”①）
\item 已知集合 $A=\left\{x\mid\dfrac{(a+1)x}{x-1}<a\right\}$，集合 $B=[1,4]$，$a$ 为实常数；

（1）若 $a=4$，若全集为实数集 $\R$，求集合 $\overline{A}$（用区间形式表示）；

（2）若 $A\subseteq B$，求 $a$ 的取值范围.

\ifanswer
\solbegin
\stepm{(1)}{若 $a=4$，集合 $A=\left\{x\mid\dfrac{5x}{x-1}<4\right\}=\left\{x\mid\dfrac{5x}{x-1}-4<0\right\}=\left\{x\mid\dfrac{x+4}{x-1}<0\right\}$}
\step{$=(-4,1)$，因此 $\overline{A}=(-\infty,-4]\cup[1,+\infty)$；}
\stepm{(2)}{当 $-a<1$，即 $a>-1$ 时，$A=(-a,1)$，}
\step{要使 $A\subseteq[1,4]$，需 $-a\geqslant1$，即 $a\leqslant-1$，与 $a>-1$ 矛盾，舍去；}
\step{当 $-a>1$，即 $a<-1$ 时，$A=(1,-a)$，}
\step{要使 $A\subseteq[1,4]$，需 $-a\leqslant4$，即 $a\geqslant-4$，所以 $-4\leqslant a<-1$；}
\step{当 $-a=1$，即 $a=-1$ 时，$A=\varnothing$，满足 $A\subseteq B$；}
\step{综上，实数 $a$ 的取值范围为 $[-4,-1]$.}
\solend
\fi

\ansspace[7]

% 原卷题干“可以售出 8 万本、据市场调查”处用顿号，照录（见文件头“原卷笔误/存疑”②）
\item 某种杂志原以每本 25 元的价格销售，可以售出 8 万本、据市场调查，杂志的单价每提高 1 元，销售量就减少 0.2 万本，假定杂志的成本是每本 10 元（不计其他成本）.

（1）当杂志以每本 30 元定价销售时，求销售该杂志所获利润；

（2）在每本 25 元的价格的基础上提高定价，试确定杂志的定价在什么范围内可使得销售利润不低于 140 万元.

\ifanswer
\solbegin
\stepm{(1)}{当杂志以 30 元定价销售时，}
\step{销售该杂志所获利润为 $[8-(30-25)\times0.2]\times(30-10)=140$ 万元；}
\stepm{(2)}{在每本 25 元的价格的基础上提高定价，设杂志的定价为每本 $x$ 元，}
\step{由题意得 $[8-(x-25)\times0.2]\times(x-10)\geqslant140$，解得 $30\leqslant x\leqslant45$，}
\step{所以杂志的定价在每本不低于 30 元且不高于 45 元的范围内，}
\step{可使得销售利润不低于 140 万元.}
\solend
\fi

\ansspace[7]

% 原卷题干“定义：对于集合 A，则称集合 B=…”句式杂糅，照录（见文件头“原卷笔误/存疑”④）
\item 定义：对于集合 $A$，则称集合 $B=\left\{x\mid x=\dfrac{m+n}{1+mn},m,n\in A,1+mn\neq0\right\}$ 为集合 $A$ 的生成集（其中 $m$、$n$ 可以相等）.

（1）若集合 $A=\{-1,1\}$ 的生成集为 $B$，求集合 $B$；

（2）若集合 $A=\{a,2\}$ 的生成集为 $B$，$B$ 的子集个数为 4 个，求实数 $a$ 的值；

（3）若集合 $A=\{x\mid-1<x<1\}$ 的生成集为 $B$，求证：$B\subseteq A$.

\ifanswer
\solbegin
\stepm{(1)}{因为 $1+mn\neq0$，}
\step{当 $m=n=1$ 时，$x=\dfrac{m+n}{1+mn}=1$；当 $m=n=-1$ 时，$x=\dfrac{m+n}{1+mn}=-1$；}
\step{所以 $B=\{-1,1\}$；}
\stepm{(2)}{$B$ 的子集个数为 4 个，故 $B$ 集合的元素有 2 个，}
\step{当 $m=n=2$ 时，$x=\dfrac{m+n}{1+mn}=\dfrac45$，当 $m=n=a$ 时，$x=\dfrac{m+n}{1+mn}=\dfrac{2a}{1+a^2}$，}
\step{当 $m$ 与 $n$ 分别取 2 和 $a$ 时，$x=\dfrac{a+2}{1+2a}$，}
\stepm{①}{$\dfrac{2a}{1+a^2}=\dfrac45$，解得 $a=2$（舍）或 $a=\dfrac12$，}
\stepm{②}{$\dfrac{a+2}{1+2a}=\dfrac45$，解得 $a=2$（舍），}
\stepm{③}{$\dfrac{2a}{1+a^2}=\dfrac{a+2}{1+2a}$，解得 $a=2$（舍）或 $a=1$ 或 $a=-1$，}
\step{综上所述，$a=1$ 或 $a=-1$ 或 $a=\dfrac12$，即 $a\in\left\{-1,1,\dfrac12\right\}$；}
\stepm{(3)}{$\forall m$、$n\in A$，则 $\dfrac{m+n}{1+mn}+1=\dfrac{(m+1)(n+1)}{1+mn}>0$，故 $\dfrac{m+n}{1+mn}>-1$；}
\step{$\dfrac{m+n}{1+mn}-1=-\dfrac{(m-1)(n-1)}{1+mn}<0$，故 $\dfrac{m+n}{1+mn}<1$；}
\step{故 $B\subseteq\{x\mid-1<x<1\}$，即 $B\subseteq A$.}
\solend
\fi

\ansspace*[10]

\end{enumerate}

\end{document}
"
% 紧凑版：xelatex "\def\iscompact{1}% 高一71_中国中学_10月月考.tex
%   —— 高一上数学“2026-2027 年中国中学高一上 10 月月考”（讲评版）
% 试卷版：xelatex 高一71_中国中学_10月月考.tex
% 答案版：xelatex "\def\isanswer{1}\input{高一71_中国中学_10月月考.tex}"
% 紧凑版：xelatex "\def\iscompact{1}\input{高一71_中国中学_10月月考.tex}"
%
% 原始材料：微信公众号“魔都数学加油站”文章，连载“27学年高一71”；
% 原文标题“27学年高一71——2026-2027年上海市中国中学高一上10月月考”；
% 原文链接：https://mp.weixin.qq.com/s/UZrP53l3GIin040L1EkljA。
% 页面标注发布时间 2026-10-19 22:47（上海）（与 高一69、高一70 页面读出同一 ct 值）。
% 正文为 7 张 PNG 页（`原图/高一71_中国中学_10月月考/01..07.png`），全部 4762×6735，
% 均按页序读取。7 页全为试卷正文页（卷首标题在第 1 页页顶，卷末解答题第 21 题 (3)
% 解析止于第 7 页），无封面页、目录页、二维码或宣传页。粉色斜排“魔都数学加油站”
% 为卷面水印，与公众号页脚一并未录入。
%
% 卷首标题：照卷面印作“2026-2027 年中国中学高一上 10 月月考”（黑体加粗一行）。
% 原卷年份区间用半角连字符，本稿按全稿惯例排 en dash（同 高一68、高一69、高一70）；
% 原卷只有一行标题、无副标题。
%
% 篇幅与题号：原卷自第 3 题起编号（第 1、2 题不在本卷内）。填空题 3—12（10 题）、
% 选择题 13—16（4 题）、解答题 17—21（5 题），共 19 题。本稿沿用原节与原题号，
% 只改排为全稿统一的大节样式（原卷小节标题为“一、填空题”“二、选择题”“三、解答题”）。
%
% 讲评版：原卷每题下直接印【答案】或【解析】。**只印【答案】的 3 题**（第 3、4、5 题），
% 本稿只给 \ans{}；**第 13 题既无【答案】也无【解析】段，答案“B”直接印在题干括号内
% （“(  B  )”）**，按 高一37 第 13 题、高一63 第 14 题、高一67 第 15 题成例把括号留空、
% 答案另提为 \ans{$B$}。其余 15 题全部印【解析】（按全稿口径这类题不另提 \ans{}——
% 第 14—16 题的“故选 D／B／B”写在解析末句，亦照录在解析内）。
% 除第 13 题外，全卷选择题作答括号内均为空白（第 14、15、16 题）。
%
% 副标题：原卷未印副标题。本稿按“知识模块”口径标作“集合与常用逻辑用语、不等式”：
% 19 题中集合与常用逻辑用语 10 题（第 3、4、5、6、9、11、12、13、16、21 题；其中第 4 题
% 属反证法，一并计入常用逻辑用语），不等式 9 题（第 7、8、10、14、15、17、18、19、20 题；
% 第 19 题 (1) 属补集、(2) 属分式不等式，按主干计入不等式），两类均确有分量，故并标。
% 口径同 高一63、高一68、高一69、高一70。
%
% 与原卷的排版差异：
%   ① 原卷每题下直接印【答案】/【解析】，本稿按全稿惯例将答案与解析置于答案版，试卷版
%      保持干净；第 13 题印在题干括号内的答案“B”同样移入 \ans{}，见上；
%   ② 原卷未印副标题，本稿按“知识模块”口径补出；
%   ③ 原卷填空横线约四个字宽，本稿按全稿惯例统一排 \blankline[4.4em]；
%   ④ 原卷数集记号不统一：实数集排斜体 $R$（第 19 题 (1)），自然数集排斜体 $N$
%      （第 5 题题干、第 12 题题干与解析共三处），第 16 题解析又排正体 $\mathrm{N}^*$；
%      本稿按全稿惯例统一改用 $\R$、$\N$、$\N^*$（同 高一68／69／70 排版差异④）；
%      空集记号原卷作 $\varnothing$，照录；
%   ⑤ 原卷句末标点混用：第 15 题题干“（3）$x^2+2+\dfrac1{x^2+2}$**．**其中”、
%      第 17 题题干末“$\dfrac{x-2}{x-1}\geqslant0$**．**”两处放大确认（第 15 题为全角
%      句点、第 17 题为西文句点），其余各题末为西文句点，已逐处照录；题干内冒号——
%      第 5 题“$\{y\mid\cdots\}$**：**”、第 12 题““×”**：**”、第 17 题“不等式组**：**”、
%      第 21 题“定义**：**”——均为全角冒号，照录；数学式内的逗号（如第 12 题
%      “$x\in A,y\in B$”、第 16 题“$A=\{a_1,\cdots\},B=\{b_1,\cdots\},n>2$”）按原卷作
%      半角，中文行文里的逗号按原卷作全角，均未强改；
%   ⑥ 第 17—21 题的子问原卷印作“（1）…（2）…”（另起一段），本稿按 高一69／70 成例用
%      行内“（1）（2）（3）”段，解析中的子问标号一律用 \stepm{(1)}{…}；
%      第 21 题 (2) 解析内层的“①”“②”“③”照原卷作 \stepm{①}{…} 等；
%   ⑦ 每道解答题原卷未留作答空白，本稿按全稿惯例加 \ansspace（末题用 \ansspace*）；
%   ⑧ 选择题的作答括号原卷排半角“(　　)”（第 14、15、16 题括号内空白、第 13 题括号内
%      印 B），本稿按全稿惯例排 \mbox{（\quad）}；选项行排照原卷——第 13 题两行两条
%      （A B 一行、C D 一行），第 14、15、16 题一行四条；
%   ⑨ 第 8 题解析、第 9 题解析、第 17 题题干与解析中的多行不等式组，原卷均排作“{”
%      连接的多行式，本稿按全稿惯例排 cases；
%   ⑩ 第 10 题解析、第 17 题解析、第 19 题 (1) 解析的并集符号原卷排大号 $\bigcup$
%      （显示体），本稿按全稿惯例排行内 $\cup$；
%   ⑪ 第 8 题解析的判别式照原卷作空心“$\triangle$”（与第 7 题解析的“$\Delta$”不同形），
%      照录未改；
%   ⑫ 第 16 题解析第 4 行的“$\forall$”（全称量词）照原卷符号排，未改写为中文“任意”。
%
% 原卷笔误/存疑（均照录未改，正文对应处另有 % 注）：
%   ① 第 19 题 (1) 连用两个“若”——“若 $a=4$，若全集为实数集 $R$”，后一“若”应作
%      “当”或删；原卷如此，照录未改。
%   ② 第 20 题题干“可以售出 8 万本**、**据市场调查”——“万本”后用顿号，按文意应为
%      逗号或句号；原卷如此，照录未改。
%   ③ 第 15 题解析对代数式 (1) 只得出“$y/x+x/y\geqslant2$”（同号时），未讨论取等是否
%      可得——取 $x=y$ 时确能取到 2——却总结“最小值为 2 的代数式有 1 个，故选 $B$”；
%      按解析自身的推理 (1) 也应计入，结论“1 个”与推理矛盾（若计 (1) 应为 2 个、选 C）。
%      原卷如此，照录未改。
%   ④ 第 21 题题干“定义：对于集合 $A$，则称集合 $B=\cdots$”——“对于……”与“则称……”
%      句式杂糅（应作“对于集合 $A$，称……”或“若集合 $A$，则称……”）；原卷如此，照录未改。
%   ⑤ 第 13 题题干末“……的（　）条件”后无任何标点（其余各题题干末均有句点）；
%      原卷如此，照录未改。
%
% 行数核对：逐题数 \step/\stepm 条数与原卷【解析】印刷行数，15 题逐一相符，
% 合计 86 条 / 86 行（Q6：1；Q7：1；Q8：7；Q9：2；Q10：4；Q11：7；Q12：5；Q14：7；
% Q15：8；Q16：7；Q17：3；Q18：7；Q19：8；Q20：6；Q21：13）；另有 3 题只印【答案】
% （第 3、4、5 题）与第 13 题括号内答案，各提 \ans{}，不计入 \step 账。
%
% 插图：无题目配图或数学图形（7 页题面均无“如图/右图/图中”字样；页面另有粉色斜排水印，
% 见上）。
%
\documentclass[a4paper,11pt]{article}
\usepackage{common}

\begin{document}

\examtitlest[3]{2026--2027 年中国中学高一上 10 月月考}{集合与常用逻辑用语、不等式}

\bigsec{一、填空题}

\begin{enumerate}
\setcounter{enumi}{2}

\item 已知集合 $A=[1,+\infty)$，$B=[-1,3)$，则 $A\cap B=$\blankline[4.4em].

\ifanswer
\ans{$[1,3)$}
\fi

\item 用反证法证明命题“已知 $a$、$b$ 为正数，$n$ 为正整数，求证：如果 $a^n>b^n$，那么 $a>b$”时，应当作出的假设为：“假设\blankline[4.4em]”.

\ifanswer
\ans{$a\leqslant b$}
\fi

\item 用列举法表示集合 $\{y\mid y=x+1,x\in\N,x\leqslant1\}$：\blankline[4.4em].

\ifanswer
\ans{$\{1,2\}$}
\fi

\item 某班在一次测验中，有 36 人数学成绩不低于 80 分，有 20 人物理成绩不低于 80 分，有 15 人的数学、物理成绩都不低于 80 分，则这两科成绩中至少有一科不低于 80 分的人数为\blankline[4.4em]人.

\ifanswer
\solbegin
\step{至少有 $36+20-15=41$ 人.}
\solend
\fi

\item 若关于 $x$ 的不等式 $x^2-x+m<0$ 的解集是 $\varnothing$，则实数 $m$ 的取值范围是\blankline[4.4em].

\ifanswer
\solbegin
\step{由题意得 $\Delta=1-4m\leqslant0$，所以 $m\geqslant\dfrac14$.}
\solend
\fi

% 第 8 题解析的判别式原卷作空心“△”（与第 7 题解析的“Δ”不同形），照录（见文件头“排版差异”⑪）
\item 已知关于 $x$ 的方程 $x^2+mx+2m=0$ 的两个实数根分别为 $\alpha$，$\beta$，若 $\alpha^2+\beta^2=5$，则实数 $m=$\blankline[4.4em].

\ifanswer
\solbegin
\step{方程 $x^2+mx+2m=0$ 的两个实数根为 $\alpha$、$\beta$，}
\step{由根与系数的关系得 $\begin{cases}\alpha+\beta=-m\\ \alpha\beta=2m\end{cases}$，又因为 $\alpha^2+\beta^2=5$，}
\step{所以 $\alpha^2+\beta^2=(\alpha+\beta)^2-2\alpha\beta=m^2-4m=5$，整理得 $m^2-4m-5=0$，}
\step{解得 $m=5$ 或 $m=-1$，}
\step{当 $m=5$ 时，$\triangle=25-40=-15<0$，无实根，舍去；}
\step{当 $m=-1$ 时，$\triangle=1+8=9>0$，符合题意，}
\step{所以实数 $m=-1$.}
\solend
\fi

\item 已知集合 $A=\{x\mid6x+a>0\}$，若 $1\in A$，$0\notin A$，则实数 $a$ 的取值范围为\blankline[4.4em].

\ifanswer
\solbegin
\step{由 $A=\{x\mid6x+a>0\}$，$1\in A$，$0\notin A$，得 $\begin{cases}6+a>0\\ a\leqslant0\end{cases}$，解得 $-6<a\leqslant0$，}
\step{即实数 $a$ 的取值范围是 $(-6,0]$.}
\solend
\fi

\item 关于 $x$ 的不等式 $ax-b>0$ 的解集为 $(1,+\infty)$，则关于 $x$ 的不等式 $(ax+b)(x-2)>0$ 的解集为\blankline[4.4em].

\ifanswer
\solbegin
\step{由不等式 $ax-b>0$ 的解集为 $(1,+\infty)$，得 $a>0$ 且 $\dfrac ba=1$，即 $b=a$，}
\step{则 $(ax+b)(x-2)>0$，可化为 $a(x+1)(x-2)>0$.}
\step{即 $(x+1)(x-2)>0$，解得 $x<-1$ 或 $x>2$，}
\step{所以不等式的解集为 $(-\infty,-1)\cup(2,+\infty)$.}
\solend
\fi

\item 已知实数 $a_1$、$a_2$、$a_3$、$a_4$、$a_5$ 满足 $\{x\mid(x-a_1)(x-a_2)(x-a_3)(x-a_4)(x-a_5)=0\}=\{1,2,3,4\}$，则 $a_1+a_2+a_3+a_4+a_5$ 的最大值为\blankline[4.4em].

\ifanswer
\solbegin
\step{5 个实数 $a_1$、$a_2$、$a_3$、$a_4$、$a_5$ 的不同取值只有 1，2，3，4，}
\step{即这 5 个数均从 1，2，3，4 中选取.}
\step{集合有 4 个元素，说明 1，2，3，4 每个数至少出现一次，}
\step{因此恰好有一个数重复出现两次，其余三个数各出现一次.}
\step{要让 5 个数的总和最大，应让重复出现两次的数取集合中的最大值 4，}
\step{此时 5 个数为 1，2，3，4，4，总和为 $1+2+3+4+4=14$，}
\step{即所求最大值为 14.}
\solend
\fi

\item 定义集合运算“$\times$”：$A\times B=\{(x,y)\mid x\in A,y\in B\}$，称为 $A$、$B$ 两个集合的“卡氏积”，若 $A=\{x\mid x^2-|x|\leqslant0,x\in\N\}$，$B=\{1,2,3\}$，则 $(A\times B)\cap(B\times A)=$\blankline[4.4em].

\ifanswer
\solbegin
\step{因为 $x^2-|x|\leqslant0$，所以 $0\leqslant|x|\leqslant1$，又因为 $x\in\N$，}
\step{所以 $A=\{x\mid x^2-|x|\leqslant0,x\in\N\}=\{0,1\}$，又 $B=\{1,2,3\}$，}
\step{由“卡氏积”的定义得 $A\times B=\{(0,1),(0,2),(0,3),(1,1),(1,2),(1,3)\}$，}
\step{$B\times A=\{(1,0),(1,1),(2,0),(2,1),(3,0),(3,1)\}$，}
\step{所以 $(A\times B)\cap(B\times A)=\{(1,1)\}$.}
\solend
\fi

\end{enumerate}

\bigsec{二、选择题}

\begin{enumerate}
\setcounter{enumi}{12}

% 第 13 题答案“B”原卷印在题干括号内，且无【答案】/【解析】段；
% 按 高一37 第 13 题、高一63 第 14 题、高一67 第 15 题成例括号留空、答案另提 \ans{$B$}
% 原卷题干末“条件”后无标点，照录（见文件头“原卷笔误/存疑”⑤）
\item “$A\subseteq B$”是“$A\subset B$”的\mbox{（\quad）}条件

\keepwithprev
A. 充分不必要条件\qquad B. 必要不充分条件

C. 充要条件\qquad D. 既不充分也不必要条件

\ifanswer
\ans{$B$}
\fi

\item 已知 $a$、$b$、$c$ 均为实数，若 $a>b$ 且 $ab\neq0$，则下列不等关系中成立的是\mbox{（\quad）}

\keepwithprev
A. $ac^2>bc^2$\qquad B. $\dfrac1b>\dfrac1a$\qquad C. $|a|>|b|$\qquad D. $c-b>c-a$

\ifanswer
\solbegin
\step{对于 $A$，若 $a>b$，当 $c=0$ 时，有 $ac^2=bc^2=0$，故 $A$ 错误；}
\step{对于 $B$，$\dfrac1b-\dfrac1a=\dfrac{a-b}{ab}$，因为 $a>b$，所以 $a-b>0$，}
\step{当 $a$、$b$ 同号时，$ab>0$，则 $\dfrac1b-\dfrac1a>0$，即 $\dfrac1b>\dfrac1a$，}
\step{当 $a$、$b$ 异号时，$ab<0$，则 $\dfrac1b-\dfrac1a<0$，即 $\dfrac1b<\dfrac1a$，故 $B$ 错误；}
\step{对于 $C$，当 $a=1$，$b=-3$ 时，$|a|<|b|$，故 $C$ 错误；}
\step{对于 $D$，若 $a>b$，则 $-b>-a$，$c-b>c-a$，故 $D$ 正确.}
\step{故选 $D$.}
\solend
\fi

% 原卷解析对 (1) 未讨论取等是否可得（x=y 时确能取 2），却总结“有 1 个、故选 B”，
% 与推理矛盾，照录（见文件头“原卷笔误/存疑”③）
\item 以下三个代数式：（1）$\dfrac yx+\dfrac xy$；（2）$|x|+\dfrac1{|x|}$；（3）$x^2+2+\dfrac1{x^2+2}$．其中的字母均表示实数，且代数式均有意义，那么代数式的最小值为 2 的个数是\mbox{（\quad）}

\keepwithprev
A. $0$\qquad B. $1$\qquad C. $2$\qquad D. $3$

\ifanswer
\solbegin
\step{（1）仅当 $x,y$ 同号时，$\dfrac yx,\dfrac xy>0$，由基本不等式得 $\dfrac yx+\dfrac xy\geqslant2$；}
\step{若 $x,y$ 异号，$\dfrac yx,\dfrac xy<0$，故 $\dfrac yx+\dfrac xy<0$，此时最小值不是 2；}
\step{（2）$|x|+\dfrac1{|x|}$ 中，$|x|>0$，由基本不等式 $|x|+\dfrac1{|x|}\geqslant2\sqrt{|x|\cdot\dfrac1{|x|}}=2$，}
\step{当且仅当 $|x|=1$，即 $x=\pm1$ 时取等号，故最小值为 2；}
\step{（3）令 $t=x^2+2$，则 $t\geqslant2$，函数化为 $t+\dfrac1t$，}
\step{对勾函数 $t+\dfrac1t$ 在 $[2,+\infty)$ 上严格增，最小值为 $2+\dfrac12=\dfrac52>2$，}
\step{等号成立条件 $t=1$ 无法取到，故最小值不是 2；}
\step{综上，最小值为 2 的代数式有 1 个，故选 $B$.}
\solend
\fi

\item 已知 $A$、$B$ 为非空实数集，$\Omega$ 为平面直角坐标系中的一些点构成的集合，集合 $C=\{y\mid$对任意 $x\in A$，有 $(x,y)\in\Omega\}$，集合 $D=\{x\mid$对任意 $y\in B$，有 $(x,y)\in\Omega\}$，对于下列两个命题：①若 $C\subseteq B$，则 $D\subseteq A$；②若 $C\supseteq B$，则 $D\subseteq A$；其中判断正确的是\mbox{（\quad）}

\keepwithprev
A. ①②都正确\qquad B. ①②都错误\qquad C. ①正确②错误\qquad D. ①错误②正确

\ifanswer
\solbegin
\step{设 $A=\{a_1,a_2,a_3,\cdots,a_n\},B=\{b_1,b_2,b_3,\cdots,b_n\},n>2$，}
\step{若 $\Omega=\{(c,b_1),(c,b_2),\cdots,(c,b_n)\},c\notin A$，}
\step{此时 $C=\varnothing\subseteq B,D=\{c\},D\subseteq A$ 错误，故①错误；}
% 原卷此处 N* 排正体（\mathrm{N}*），本稿统一 \N^*（见文件头“排版差异”④）；
% 全称量词“∀”照原卷符号排（见文件头“排版差异”⑫）
\step{若 $C\supseteq B$，则 $\forall(a_i,b_j)(1\leqslant i,j\leqslant n,i,j\in\N^*)\in\Omega$，则 $D\supseteq A$，}
\step{且 $A=\{a_1\},B=\{b_1,b_2\}$，}
\step{若 $\Omega=\{(a_1,b_1),(a_2,b_1),(a_1,b_2),(a_2,b_2),(c,b_1),(c,b_2)\},D$ 真包含 $A$，故②错误.}
\step{故选 $B$.}
\solend
\fi

\end{enumerate}

\bigsec{三、解答题}

\begin{enumerate}
\setcounter{enumi}{16}

\item 解下列关于 $x$ 的不等式组：$\begin{cases}|x-1|<2\\ \dfrac{x-2}{x-1}\geqslant0\end{cases}$.

\ifanswer
\solbegin
\step{解 $|x-1|<2$ 得 $-2<x-1<2$，解得 $-1<x<3$.}
\step{分式不等式 $\dfrac{x-2}{x-1}\geqslant0$ 等价于 $\begin{cases}(x-2)(x-1)\geqslant0\\ x-1\neq0\end{cases}$，解得 $x<1$ 或 $x\geqslant2$.}
\step{即 $\begin{cases}|x-1|<2\\ \dfrac{x-2}{x-1}\geqslant0\end{cases}$ 的解集为 $(-1,1)\cup[2,3)$.}
\solend
\fi

\ansspace[5]

\item 比较多项式 $2x^2-3x+1$ 与 $x^2-1$ 的大小.

\ifanswer
\solbegin
\step{作差得 $(2x^2-3x+1)-(x^2-1)=x^2-3x+2=(x-1)(x-2)$，}
\step{当 $x=1$ 或 $x=2$ 时，$(x-1)(x-2)=0$，即 $2x^2-3x+1=x^2-1$，}
\step{当 $x<1$ 或 $x>2$ 时，$(x-1)(x-2)>0$，即 $2x^2-3x+1>x^2-1$，}
\step{当 $1<x<2$ 时，$(x-1)(x-2)<0$，即 $2x^2-3x+1<x^2-1$；}
\step{综上，当 $x=1$ 或 $x=2$ 时，$2x^2-3x+1=x^2-1$；}
\step{当 $x<1$ 或 $x>2$ 时，$2x^2-3x+1>x^2-1$；}
\step{当 $1<x<2$ 时，$2x^2-3x+1<x^2-1$.}
\solend
\fi

\ansspace[6]

% 原卷 (1) 连用两个“若”（“若 a=4，若全集为实数集 R”），照录（见文件头“原卷笔误/存疑”①）
\item 已知集合 $A=\left\{x\mid\dfrac{(a+1)x}{x-1}<a\right\}$，集合 $B=[1,4]$，$a$ 为实常数；

（1）若 $a=4$，若全集为实数集 $\R$，求集合 $\overline{A}$（用区间形式表示）；

（2）若 $A\subseteq B$，求 $a$ 的取值范围.

\ifanswer
\solbegin
\stepm{(1)}{若 $a=4$，集合 $A=\left\{x\mid\dfrac{5x}{x-1}<4\right\}=\left\{x\mid\dfrac{5x}{x-1}-4<0\right\}=\left\{x\mid\dfrac{x+4}{x-1}<0\right\}$}
\step{$=(-4,1)$，因此 $\overline{A}=(-\infty,-4]\cup[1,+\infty)$；}
\stepm{(2)}{当 $-a<1$，即 $a>-1$ 时，$A=(-a,1)$，}
\step{要使 $A\subseteq[1,4]$，需 $-a\geqslant1$，即 $a\leqslant-1$，与 $a>-1$ 矛盾，舍去；}
\step{当 $-a>1$，即 $a<-1$ 时，$A=(1,-a)$，}
\step{要使 $A\subseteq[1,4]$，需 $-a\leqslant4$，即 $a\geqslant-4$，所以 $-4\leqslant a<-1$；}
\step{当 $-a=1$，即 $a=-1$ 时，$A=\varnothing$，满足 $A\subseteq B$；}
\step{综上，实数 $a$ 的取值范围为 $[-4,-1]$.}
\solend
\fi

\ansspace[7]

% 原卷题干“可以售出 8 万本、据市场调查”处用顿号，照录（见文件头“原卷笔误/存疑”②）
\item 某种杂志原以每本 25 元的价格销售，可以售出 8 万本、据市场调查，杂志的单价每提高 1 元，销售量就减少 0.2 万本，假定杂志的成本是每本 10 元（不计其他成本）.

（1）当杂志以每本 30 元定价销售时，求销售该杂志所获利润；

（2）在每本 25 元的价格的基础上提高定价，试确定杂志的定价在什么范围内可使得销售利润不低于 140 万元.

\ifanswer
\solbegin
\stepm{(1)}{当杂志以 30 元定价销售时，}
\step{销售该杂志所获利润为 $[8-(30-25)\times0.2]\times(30-10)=140$ 万元；}
\stepm{(2)}{在每本 25 元的价格的基础上提高定价，设杂志的定价为每本 $x$ 元，}
\step{由题意得 $[8-(x-25)\times0.2]\times(x-10)\geqslant140$，解得 $30\leqslant x\leqslant45$，}
\step{所以杂志的定价在每本不低于 30 元且不高于 45 元的范围内，}
\step{可使得销售利润不低于 140 万元.}
\solend
\fi

\ansspace[7]

% 原卷题干“定义：对于集合 A，则称集合 B=…”句式杂糅，照录（见文件头“原卷笔误/存疑”④）
\item 定义：对于集合 $A$，则称集合 $B=\left\{x\mid x=\dfrac{m+n}{1+mn},m,n\in A,1+mn\neq0\right\}$ 为集合 $A$ 的生成集（其中 $m$、$n$ 可以相等）.

（1）若集合 $A=\{-1,1\}$ 的生成集为 $B$，求集合 $B$；

（2）若集合 $A=\{a,2\}$ 的生成集为 $B$，$B$ 的子集个数为 4 个，求实数 $a$ 的值；

（3）若集合 $A=\{x\mid-1<x<1\}$ 的生成集为 $B$，求证：$B\subseteq A$.

\ifanswer
\solbegin
\stepm{(1)}{因为 $1+mn\neq0$，}
\step{当 $m=n=1$ 时，$x=\dfrac{m+n}{1+mn}=1$；当 $m=n=-1$ 时，$x=\dfrac{m+n}{1+mn}=-1$；}
\step{所以 $B=\{-1,1\}$；}
\stepm{(2)}{$B$ 的子集个数为 4 个，故 $B$ 集合的元素有 2 个，}
\step{当 $m=n=2$ 时，$x=\dfrac{m+n}{1+mn}=\dfrac45$，当 $m=n=a$ 时，$x=\dfrac{m+n}{1+mn}=\dfrac{2a}{1+a^2}$，}
\step{当 $m$ 与 $n$ 分别取 2 和 $a$ 时，$x=\dfrac{a+2}{1+2a}$，}
\stepm{①}{$\dfrac{2a}{1+a^2}=\dfrac45$，解得 $a=2$（舍）或 $a=\dfrac12$，}
\stepm{②}{$\dfrac{a+2}{1+2a}=\dfrac45$，解得 $a=2$（舍），}
\stepm{③}{$\dfrac{2a}{1+a^2}=\dfrac{a+2}{1+2a}$，解得 $a=2$（舍）或 $a=1$ 或 $a=-1$，}
\step{综上所述，$a=1$ 或 $a=-1$ 或 $a=\dfrac12$，即 $a\in\left\{-1,1,\dfrac12\right\}$；}
\stepm{(3)}{$\forall m$、$n\in A$，则 $\dfrac{m+n}{1+mn}+1=\dfrac{(m+1)(n+1)}{1+mn}>0$，故 $\dfrac{m+n}{1+mn}>-1$；}
\step{$\dfrac{m+n}{1+mn}-1=-\dfrac{(m-1)(n-1)}{1+mn}<0$，故 $\dfrac{m+n}{1+mn}<1$；}
\step{故 $B\subseteq\{x\mid-1<x<1\}$，即 $B\subseteq A$.}
\solend
\fi

\ansspace*[10]

\end{enumerate}

\end{document}
"
%
% 原始材料：微信公众号“魔都数学加油站”文章，连载“27学年高一71”；
% 原文标题“27学年高一71——2026-2027年上海市中国中学高一上10月月考”；
% 原文链接：https://mp.weixin.qq.com/s/UZrP53l3GIin040L1EkljA。
% 页面标注发布时间 2026-10-19 22:47（上海）（与 高一69、高一70 页面读出同一 ct 值）。
% 正文为 7 张 PNG 页（`原图/高一71_中国中学_10月月考/01..07.png`），全部 4762×6735，
% 均按页序读取。7 页全为试卷正文页（卷首标题在第 1 页页顶，卷末解答题第 21 题 (3)
% 解析止于第 7 页），无封面页、目录页、二维码或宣传页。粉色斜排“魔都数学加油站”
% 为卷面水印，与公众号页脚一并未录入。
%
% 卷首标题：照卷面印作“2026-2027 年中国中学高一上 10 月月考”（黑体加粗一行）。
% 原卷年份区间用半角连字符，本稿按全稿惯例排 en dash（同 高一68、高一69、高一70）；
% 原卷只有一行标题、无副标题。
%
% 篇幅与题号：原卷自第 3 题起编号（第 1、2 题不在本卷内）。填空题 3—12（10 题）、
% 选择题 13—16（4 题）、解答题 17—21（5 题），共 19 题。本稿沿用原节与原题号，
% 只改排为全稿统一的大节样式（原卷小节标题为“一、填空题”“二、选择题”“三、解答题”）。
%
% 讲评版：原卷每题下直接印【答案】或【解析】。**只印【答案】的 3 题**（第 3、4、5 题），
% 本稿只给 \ans{}；**第 13 题既无【答案】也无【解析】段，答案“B”直接印在题干括号内
% （“(  B  )”）**，按 高一37 第 13 题、高一63 第 14 题、高一67 第 15 题成例把括号留空、
% 答案另提为 \ans{$B$}。其余 15 题全部印【解析】（按全稿口径这类题不另提 \ans{}——
% 第 14—16 题的“故选 D／B／B”写在解析末句，亦照录在解析内）。
% 除第 13 题外，全卷选择题作答括号内均为空白（第 14、15、16 题）。
%
% 副标题：原卷未印副标题。本稿按“知识模块”口径标作“集合与常用逻辑用语、不等式”：
% 19 题中集合与常用逻辑用语 10 题（第 3、4、5、6、9、11、12、13、16、21 题；其中第 4 题
% 属反证法，一并计入常用逻辑用语），不等式 9 题（第 7、8、10、14、15、17、18、19、20 题；
% 第 19 题 (1) 属补集、(2) 属分式不等式，按主干计入不等式），两类均确有分量，故并标。
% 口径同 高一63、高一68、高一69、高一70。
%
% 与原卷的排版差异：
%   ① 原卷每题下直接印【答案】/【解析】，本稿按全稿惯例将答案与解析置于答案版，试卷版
%      保持干净；第 13 题印在题干括号内的答案“B”同样移入 \ans{}，见上；
%   ② 原卷未印副标题，本稿按“知识模块”口径补出；
%   ③ 原卷填空横线约四个字宽，本稿按全稿惯例统一排 \blankline[4.4em]；
%   ④ 原卷数集记号不统一：实数集排斜体 $R$（第 19 题 (1)），自然数集排斜体 $N$
%      （第 5 题题干、第 12 题题干与解析共三处），第 16 题解析又排正体 $\mathrm{N}^*$；
%      本稿按全稿惯例统一改用 $\R$、$\N$、$\N^*$（同 高一68／69／70 排版差异④）；
%      空集记号原卷作 $\varnothing$，照录；
%   ⑤ 原卷句末标点混用：第 15 题题干“（3）$x^2+2+\dfrac1{x^2+2}$**．**其中”、
%      第 17 题题干末“$\dfrac{x-2}{x-1}\geqslant0$**．**”两处放大确认（第 15 题为全角
%      句点、第 17 题为西文句点），其余各题末为西文句点，已逐处照录；题干内冒号——
%      第 5 题“$\{y\mid\cdots\}$**：**”、第 12 题““×”**：**”、第 17 题“不等式组**：**”、
%      第 21 题“定义**：**”——均为全角冒号，照录；数学式内的逗号（如第 12 题
%      “$x\in A,y\in B$”、第 16 题“$A=\{a_1,\cdots\},B=\{b_1,\cdots\},n>2$”）按原卷作
%      半角，中文行文里的逗号按原卷作全角，均未强改；
%   ⑥ 第 17—21 题的子问原卷印作“（1）…（2）…”（另起一段），本稿按 高一69／70 成例用
%      行内“（1）（2）（3）”段，解析中的子问标号一律用 \stepm{(1)}{…}；
%      第 21 题 (2) 解析内层的“①”“②”“③”照原卷作 \stepm{①}{…} 等；
%   ⑦ 每道解答题原卷未留作答空白，本稿按全稿惯例加 \ansspace（末题用 \ansspace*）；
%   ⑧ 选择题的作答括号原卷排半角“(　　)”（第 14、15、16 题括号内空白、第 13 题括号内
%      印 B），本稿按全稿惯例排 \mbox{（\quad）}；选项行排照原卷——第 13 题两行两条
%      （A B 一行、C D 一行），第 14、15、16 题一行四条；
%   ⑨ 第 8 题解析、第 9 题解析、第 17 题题干与解析中的多行不等式组，原卷均排作“{”
%      连接的多行式，本稿按全稿惯例排 cases；
%   ⑩ 第 10 题解析、第 17 题解析、第 19 题 (1) 解析的并集符号原卷排大号 $\bigcup$
%      （显示体），本稿按全稿惯例排行内 $\cup$；
%   ⑪ 第 8 题解析的判别式照原卷作空心“$\triangle$”（与第 7 题解析的“$\Delta$”不同形），
%      照录未改；
%   ⑫ 第 16 题解析第 4 行的“$\forall$”（全称量词）照原卷符号排，未改写为中文“任意”。
%
% 原卷笔误/存疑（均照录未改，正文对应处另有 % 注）：
%   ① 第 19 题 (1) 连用两个“若”——“若 $a=4$，若全集为实数集 $R$”，后一“若”应作
%      “当”或删；原卷如此，照录未改。
%   ② 第 20 题题干“可以售出 8 万本**、**据市场调查”——“万本”后用顿号，按文意应为
%      逗号或句号；原卷如此，照录未改。
%   ③ 第 15 题解析对代数式 (1) 只得出“$y/x+x/y\geqslant2$”（同号时），未讨论取等是否
%      可得——取 $x=y$ 时确能取到 2——却总结“最小值为 2 的代数式有 1 个，故选 $B$”；
%      按解析自身的推理 (1) 也应计入，结论“1 个”与推理矛盾（若计 (1) 应为 2 个、选 C）。
%      原卷如此，照录未改。
%   ④ 第 21 题题干“定义：对于集合 $A$，则称集合 $B=\cdots$”——“对于……”与“则称……”
%      句式杂糅（应作“对于集合 $A$，称……”或“若集合 $A$，则称……”）；原卷如此，照录未改。
%   ⑤ 第 13 题题干末“……的（　）条件”后无任何标点（其余各题题干末均有句点）；
%      原卷如此，照录未改。
%
% 行数核对：逐题数 \step/\stepm 条数与原卷【解析】印刷行数，15 题逐一相符，
% 合计 86 条 / 86 行（Q6：1；Q7：1；Q8：7；Q9：2；Q10：4；Q11：7；Q12：5；Q14：7；
% Q15：8；Q16：7；Q17：3；Q18：7；Q19：8；Q20：6；Q21：13）；另有 3 题只印【答案】
% （第 3、4、5 题）与第 13 题括号内答案，各提 \ans{}，不计入 \step 账。
%
% 插图：无题目配图或数学图形（7 页题面均无“如图/右图/图中”字样；页面另有粉色斜排水印，
% 见上）。
%
\documentclass[a4paper,11pt]{article}
\usepackage{common}

\begin{document}

\examtitlest[3]{2026--2027 年中国中学高一上 10 月月考}{集合与常用逻辑用语、不等式}

\bigsec{一、填空题}

\begin{enumerate}
\setcounter{enumi}{2}

\item 已知集合 $A=[1,+\infty)$，$B=[-1,3)$，则 $A\cap B=$\blankline[4.4em].

\ifanswer
\ans{$[1,3)$}
\fi

\item 用反证法证明命题“已知 $a$、$b$ 为正数，$n$ 为正整数，求证：如果 $a^n>b^n$，那么 $a>b$”时，应当作出的假设为：“假设\blankline[4.4em]”.

\ifanswer
\ans{$a\leqslant b$}
\fi

\item 用列举法表示集合 $\{y\mid y=x+1,x\in\N,x\leqslant1\}$：\blankline[4.4em].

\ifanswer
\ans{$\{1,2\}$}
\fi

\item 某班在一次测验中，有 36 人数学成绩不低于 80 分，有 20 人物理成绩不低于 80 分，有 15 人的数学、物理成绩都不低于 80 分，则这两科成绩中至少有一科不低于 80 分的人数为\blankline[4.4em]人.

\ifanswer
\solbegin
\step{至少有 $36+20-15=41$ 人.}
\solend
\fi

\item 若关于 $x$ 的不等式 $x^2-x+m<0$ 的解集是 $\varnothing$，则实数 $m$ 的取值范围是\blankline[4.4em].

\ifanswer
\solbegin
\step{由题意得 $\Delta=1-4m\leqslant0$，所以 $m\geqslant\dfrac14$.}
\solend
\fi

% 第 8 题解析的判别式原卷作空心“△”（与第 7 题解析的“Δ”不同形），照录（见文件头“排版差异”⑪）
\item 已知关于 $x$ 的方程 $x^2+mx+2m=0$ 的两个实数根分别为 $\alpha$，$\beta$，若 $\alpha^2+\beta^2=5$，则实数 $m=$\blankline[4.4em].

\ifanswer
\solbegin
\step{方程 $x^2+mx+2m=0$ 的两个实数根为 $\alpha$、$\beta$，}
\step{由根与系数的关系得 $\begin{cases}\alpha+\beta=-m\\ \alpha\beta=2m\end{cases}$，又因为 $\alpha^2+\beta^2=5$，}
\step{所以 $\alpha^2+\beta^2=(\alpha+\beta)^2-2\alpha\beta=m^2-4m=5$，整理得 $m^2-4m-5=0$，}
\step{解得 $m=5$ 或 $m=-1$，}
\step{当 $m=5$ 时，$\triangle=25-40=-15<0$，无实根，舍去；}
\step{当 $m=-1$ 时，$\triangle=1+8=9>0$，符合题意，}
\step{所以实数 $m=-1$.}
\solend
\fi

\item 已知集合 $A=\{x\mid6x+a>0\}$，若 $1\in A$，$0\notin A$，则实数 $a$ 的取值范围为\blankline[4.4em].

\ifanswer
\solbegin
\step{由 $A=\{x\mid6x+a>0\}$，$1\in A$，$0\notin A$，得 $\begin{cases}6+a>0\\ a\leqslant0\end{cases}$，解得 $-6<a\leqslant0$，}
\step{即实数 $a$ 的取值范围是 $(-6,0]$.}
\solend
\fi

\item 关于 $x$ 的不等式 $ax-b>0$ 的解集为 $(1,+\infty)$，则关于 $x$ 的不等式 $(ax+b)(x-2)>0$ 的解集为\blankline[4.4em].

\ifanswer
\solbegin
\step{由不等式 $ax-b>0$ 的解集为 $(1,+\infty)$，得 $a>0$ 且 $\dfrac ba=1$，即 $b=a$，}
\step{则 $(ax+b)(x-2)>0$，可化为 $a(x+1)(x-2)>0$.}
\step{即 $(x+1)(x-2)>0$，解得 $x<-1$ 或 $x>2$，}
\step{所以不等式的解集为 $(-\infty,-1)\cup(2,+\infty)$.}
\solend
\fi

\item 已知实数 $a_1$、$a_2$、$a_3$、$a_4$、$a_5$ 满足 $\{x\mid(x-a_1)(x-a_2)(x-a_3)(x-a_4)(x-a_5)=0\}=\{1,2,3,4\}$，则 $a_1+a_2+a_3+a_4+a_5$ 的最大值为\blankline[4.4em].

\ifanswer
\solbegin
\step{5 个实数 $a_1$、$a_2$、$a_3$、$a_4$、$a_5$ 的不同取值只有 1，2，3，4，}
\step{即这 5 个数均从 1，2，3，4 中选取.}
\step{集合有 4 个元素，说明 1，2，3，4 每个数至少出现一次，}
\step{因此恰好有一个数重复出现两次，其余三个数各出现一次.}
\step{要让 5 个数的总和最大，应让重复出现两次的数取集合中的最大值 4，}
\step{此时 5 个数为 1，2，3，4，4，总和为 $1+2+3+4+4=14$，}
\step{即所求最大值为 14.}
\solend
\fi

\item 定义集合运算“$\times$”：$A\times B=\{(x,y)\mid x\in A,y\in B\}$，称为 $A$、$B$ 两个集合的“卡氏积”，若 $A=\{x\mid x^2-|x|\leqslant0,x\in\N\}$，$B=\{1,2,3\}$，则 $(A\times B)\cap(B\times A)=$\blankline[4.4em].

\ifanswer
\solbegin
\step{因为 $x^2-|x|\leqslant0$，所以 $0\leqslant|x|\leqslant1$，又因为 $x\in\N$，}
\step{所以 $A=\{x\mid x^2-|x|\leqslant0,x\in\N\}=\{0,1\}$，又 $B=\{1,2,3\}$，}
\step{由“卡氏积”的定义得 $A\times B=\{(0,1),(0,2),(0,3),(1,1),(1,2),(1,3)\}$，}
\step{$B\times A=\{(1,0),(1,1),(2,0),(2,1),(3,0),(3,1)\}$，}
\step{所以 $(A\times B)\cap(B\times A)=\{(1,1)\}$.}
\solend
\fi

\end{enumerate}

\bigsec{二、选择题}

\begin{enumerate}
\setcounter{enumi}{12}

% 第 13 题答案“B”原卷印在题干括号内，且无【答案】/【解析】段；
% 按 高一37 第 13 题、高一63 第 14 题、高一67 第 15 题成例括号留空、答案另提 \ans{$B$}
% 原卷题干末“条件”后无标点，照录（见文件头“原卷笔误/存疑”⑤）
\item “$A\subseteq B$”是“$A\subset B$”的\mbox{（\quad）}条件

\keepwithprev
A. 充分不必要条件\qquad B. 必要不充分条件

C. 充要条件\qquad D. 既不充分也不必要条件

\ifanswer
\ans{$B$}
\fi

\item 已知 $a$、$b$、$c$ 均为实数，若 $a>b$ 且 $ab\neq0$，则下列不等关系中成立的是\mbox{（\quad）}

\keepwithprev
A. $ac^2>bc^2$\qquad B. $\dfrac1b>\dfrac1a$\qquad C. $|a|>|b|$\qquad D. $c-b>c-a$

\ifanswer
\solbegin
\step{对于 $A$，若 $a>b$，当 $c=0$ 时，有 $ac^2=bc^2=0$，故 $A$ 错误；}
\step{对于 $B$，$\dfrac1b-\dfrac1a=\dfrac{a-b}{ab}$，因为 $a>b$，所以 $a-b>0$，}
\step{当 $a$、$b$ 同号时，$ab>0$，则 $\dfrac1b-\dfrac1a>0$，即 $\dfrac1b>\dfrac1a$，}
\step{当 $a$、$b$ 异号时，$ab<0$，则 $\dfrac1b-\dfrac1a<0$，即 $\dfrac1b<\dfrac1a$，故 $B$ 错误；}
\step{对于 $C$，当 $a=1$，$b=-3$ 时，$|a|<|b|$，故 $C$ 错误；}
\step{对于 $D$，若 $a>b$，则 $-b>-a$，$c-b>c-a$，故 $D$ 正确.}
\step{故选 $D$.}
\solend
\fi

% 原卷解析对 (1) 未讨论取等是否可得（x=y 时确能取 2），却总结“有 1 个、故选 B”，
% 与推理矛盾，照录（见文件头“原卷笔误/存疑”③）
\item 以下三个代数式：（1）$\dfrac yx+\dfrac xy$；（2）$|x|+\dfrac1{|x|}$；（3）$x^2+2+\dfrac1{x^2+2}$．其中的字母均表示实数，且代数式均有意义，那么代数式的最小值为 2 的个数是\mbox{（\quad）}

\keepwithprev
A. $0$\qquad B. $1$\qquad C. $2$\qquad D. $3$

\ifanswer
\solbegin
\step{（1）仅当 $x,y$ 同号时，$\dfrac yx,\dfrac xy>0$，由基本不等式得 $\dfrac yx+\dfrac xy\geqslant2$；}
\step{若 $x,y$ 异号，$\dfrac yx,\dfrac xy<0$，故 $\dfrac yx+\dfrac xy<0$，此时最小值不是 2；}
\step{（2）$|x|+\dfrac1{|x|}$ 中，$|x|>0$，由基本不等式 $|x|+\dfrac1{|x|}\geqslant2\sqrt{|x|\cdot\dfrac1{|x|}}=2$，}
\step{当且仅当 $|x|=1$，即 $x=\pm1$ 时取等号，故最小值为 2；}
\step{（3）令 $t=x^2+2$，则 $t\geqslant2$，函数化为 $t+\dfrac1t$，}
\step{对勾函数 $t+\dfrac1t$ 在 $[2,+\infty)$ 上严格增，最小值为 $2+\dfrac12=\dfrac52>2$，}
\step{等号成立条件 $t=1$ 无法取到，故最小值不是 2；}
\step{综上，最小值为 2 的代数式有 1 个，故选 $B$.}
\solend
\fi

\item 已知 $A$、$B$ 为非空实数集，$\Omega$ 为平面直角坐标系中的一些点构成的集合，集合 $C=\{y\mid$对任意 $x\in A$，有 $(x,y)\in\Omega\}$，集合 $D=\{x\mid$对任意 $y\in B$，有 $(x,y)\in\Omega\}$，对于下列两个命题：①若 $C\subseteq B$，则 $D\subseteq A$；②若 $C\supseteq B$，则 $D\subseteq A$；其中判断正确的是\mbox{（\quad）}

\keepwithprev
A. ①②都正确\qquad B. ①②都错误\qquad C. ①正确②错误\qquad D. ①错误②正确

\ifanswer
\solbegin
\step{设 $A=\{a_1,a_2,a_3,\cdots,a_n\},B=\{b_1,b_2,b_3,\cdots,b_n\},n>2$，}
\step{若 $\Omega=\{(c,b_1),(c,b_2),\cdots,(c,b_n)\},c\notin A$，}
\step{此时 $C=\varnothing\subseteq B,D=\{c\},D\subseteq A$ 错误，故①错误；}
% 原卷此处 N* 排正体（\mathrm{N}*），本稿统一 \N^*（见文件头“排版差异”④）；
% 全称量词“∀”照原卷符号排（见文件头“排版差异”⑫）
\step{若 $C\supseteq B$，则 $\forall(a_i,b_j)(1\leqslant i,j\leqslant n,i,j\in\N^*)\in\Omega$，则 $D\supseteq A$，}
\step{且 $A=\{a_1\},B=\{b_1,b_2\}$，}
\step{若 $\Omega=\{(a_1,b_1),(a_2,b_1),(a_1,b_2),(a_2,b_2),(c,b_1),(c,b_2)\},D$ 真包含 $A$，故②错误.}
\step{故选 $B$.}
\solend
\fi

\end{enumerate}

\bigsec{三、解答题}

\begin{enumerate}
\setcounter{enumi}{16}

\item 解下列关于 $x$ 的不等式组：$\begin{cases}|x-1|<2\\ \dfrac{x-2}{x-1}\geqslant0\end{cases}$.

\ifanswer
\solbegin
\step{解 $|x-1|<2$ 得 $-2<x-1<2$，解得 $-1<x<3$.}
\step{分式不等式 $\dfrac{x-2}{x-1}\geqslant0$ 等价于 $\begin{cases}(x-2)(x-1)\geqslant0\\ x-1\neq0\end{cases}$，解得 $x<1$ 或 $x\geqslant2$.}
\step{即 $\begin{cases}|x-1|<2\\ \dfrac{x-2}{x-1}\geqslant0\end{cases}$ 的解集为 $(-1,1)\cup[2,3)$.}
\solend
\fi

\ansspace[5]

\item 比较多项式 $2x^2-3x+1$ 与 $x^2-1$ 的大小.

\ifanswer
\solbegin
\step{作差得 $(2x^2-3x+1)-(x^2-1)=x^2-3x+2=(x-1)(x-2)$，}
\step{当 $x=1$ 或 $x=2$ 时，$(x-1)(x-2)=0$，即 $2x^2-3x+1=x^2-1$，}
\step{当 $x<1$ 或 $x>2$ 时，$(x-1)(x-2)>0$，即 $2x^2-3x+1>x^2-1$，}
\step{当 $1<x<2$ 时，$(x-1)(x-2)<0$，即 $2x^2-3x+1<x^2-1$；}
\step{综上，当 $x=1$ 或 $x=2$ 时，$2x^2-3x+1=x^2-1$；}
\step{当 $x<1$ 或 $x>2$ 时，$2x^2-3x+1>x^2-1$；}
\step{当 $1<x<2$ 时，$2x^2-3x+1<x^2-1$.}
\solend
\fi

\ansspace[6]

% 原卷 (1) 连用两个“若”（“若 a=4，若全集为实数集 R”），照录（见文件头“原卷笔误/存疑”①）
\item 已知集合 $A=\left\{x\mid\dfrac{(a+1)x}{x-1}<a\right\}$，集合 $B=[1,4]$，$a$ 为实常数；

（1）若 $a=4$，若全集为实数集 $\R$，求集合 $\overline{A}$（用区间形式表示）；

（2）若 $A\subseteq B$，求 $a$ 的取值范围.

\ifanswer
\solbegin
\stepm{(1)}{若 $a=4$，集合 $A=\left\{x\mid\dfrac{5x}{x-1}<4\right\}=\left\{x\mid\dfrac{5x}{x-1}-4<0\right\}=\left\{x\mid\dfrac{x+4}{x-1}<0\right\}$}
\step{$=(-4,1)$，因此 $\overline{A}=(-\infty,-4]\cup[1,+\infty)$；}
\stepm{(2)}{当 $-a<1$，即 $a>-1$ 时，$A=(-a,1)$，}
\step{要使 $A\subseteq[1,4]$，需 $-a\geqslant1$，即 $a\leqslant-1$，与 $a>-1$ 矛盾，舍去；}
\step{当 $-a>1$，即 $a<-1$ 时，$A=(1,-a)$，}
\step{要使 $A\subseteq[1,4]$，需 $-a\leqslant4$，即 $a\geqslant-4$，所以 $-4\leqslant a<-1$；}
\step{当 $-a=1$，即 $a=-1$ 时，$A=\varnothing$，满足 $A\subseteq B$；}
\step{综上，实数 $a$ 的取值范围为 $[-4,-1]$.}
\solend
\fi

\ansspace[7]

% 原卷题干“可以售出 8 万本、据市场调查”处用顿号，照录（见文件头“原卷笔误/存疑”②）
\item 某种杂志原以每本 25 元的价格销售，可以售出 8 万本、据市场调查，杂志的单价每提高 1 元，销售量就减少 0.2 万本，假定杂志的成本是每本 10 元（不计其他成本）.

（1）当杂志以每本 30 元定价销售时，求销售该杂志所获利润；

（2）在每本 25 元的价格的基础上提高定价，试确定杂志的定价在什么范围内可使得销售利润不低于 140 万元.

\ifanswer
\solbegin
\stepm{(1)}{当杂志以 30 元定价销售时，}
\step{销售该杂志所获利润为 $[8-(30-25)\times0.2]\times(30-10)=140$ 万元；}
\stepm{(2)}{在每本 25 元的价格的基础上提高定价，设杂志的定价为每本 $x$ 元，}
\step{由题意得 $[8-(x-25)\times0.2]\times(x-10)\geqslant140$，解得 $30\leqslant x\leqslant45$，}
\step{所以杂志的定价在每本不低于 30 元且不高于 45 元的范围内，}
\step{可使得销售利润不低于 140 万元.}
\solend
\fi

\ansspace[7]

% 原卷题干“定义：对于集合 A，则称集合 B=…”句式杂糅，照录（见文件头“原卷笔误/存疑”④）
\item 定义：对于集合 $A$，则称集合 $B=\left\{x\mid x=\dfrac{m+n}{1+mn},m,n\in A,1+mn\neq0\right\}$ 为集合 $A$ 的生成集（其中 $m$、$n$ 可以相等）.

（1）若集合 $A=\{-1,1\}$ 的生成集为 $B$，求集合 $B$；

（2）若集合 $A=\{a,2\}$ 的生成集为 $B$，$B$ 的子集个数为 4 个，求实数 $a$ 的值；

（3）若集合 $A=\{x\mid-1<x<1\}$ 的生成集为 $B$，求证：$B\subseteq A$.

\ifanswer
\solbegin
\stepm{(1)}{因为 $1+mn\neq0$，}
\step{当 $m=n=1$ 时，$x=\dfrac{m+n}{1+mn}=1$；当 $m=n=-1$ 时，$x=\dfrac{m+n}{1+mn}=-1$；}
\step{所以 $B=\{-1,1\}$；}
\stepm{(2)}{$B$ 的子集个数为 4 个，故 $B$ 集合的元素有 2 个，}
\step{当 $m=n=2$ 时，$x=\dfrac{m+n}{1+mn}=\dfrac45$，当 $m=n=a$ 时，$x=\dfrac{m+n}{1+mn}=\dfrac{2a}{1+a^2}$，}
\step{当 $m$ 与 $n$ 分别取 2 和 $a$ 时，$x=\dfrac{a+2}{1+2a}$，}
\stepm{①}{$\dfrac{2a}{1+a^2}=\dfrac45$，解得 $a=2$（舍）或 $a=\dfrac12$，}
\stepm{②}{$\dfrac{a+2}{1+2a}=\dfrac45$，解得 $a=2$（舍），}
\stepm{③}{$\dfrac{2a}{1+a^2}=\dfrac{a+2}{1+2a}$，解得 $a=2$（舍）或 $a=1$ 或 $a=-1$，}
\step{综上所述，$a=1$ 或 $a=-1$ 或 $a=\dfrac12$，即 $a\in\left\{-1,1,\dfrac12\right\}$；}
\stepm{(3)}{$\forall m$、$n\in A$，则 $\dfrac{m+n}{1+mn}+1=\dfrac{(m+1)(n+1)}{1+mn}>0$，故 $\dfrac{m+n}{1+mn}>-1$；}
\step{$\dfrac{m+n}{1+mn}-1=-\dfrac{(m-1)(n-1)}{1+mn}<0$，故 $\dfrac{m+n}{1+mn}<1$；}
\step{故 $B\subseteq\{x\mid-1<x<1\}$，即 $B\subseteq A$.}
\solend
\fi

\ansspace*[10]

\end{enumerate}

\end{document}
"
% 紧凑版：xelatex "\def\iscompact{1}% 高一71_中国中学_10月月考.tex
%   —— 高一上数学“2026-2027 年中国中学高一上 10 月月考”（讲评版）
% 试卷版：xelatex 高一71_中国中学_10月月考.tex
% 答案版：xelatex "\def\isanswer{1}% 高一71_中国中学_10月月考.tex
%   —— 高一上数学“2026-2027 年中国中学高一上 10 月月考”（讲评版）
% 试卷版：xelatex 高一71_中国中学_10月月考.tex
% 答案版：xelatex "\def\isanswer{1}\input{高一71_中国中学_10月月考.tex}"
% 紧凑版：xelatex "\def\iscompact{1}\input{高一71_中国中学_10月月考.tex}"
%
% 原始材料：微信公众号“魔都数学加油站”文章，连载“27学年高一71”；
% 原文标题“27学年高一71——2026-2027年上海市中国中学高一上10月月考”；
% 原文链接：https://mp.weixin.qq.com/s/UZrP53l3GIin040L1EkljA。
% 页面标注发布时间 2026-10-19 22:47（上海）（与 高一69、高一70 页面读出同一 ct 值）。
% 正文为 7 张 PNG 页（`原图/高一71_中国中学_10月月考/01..07.png`），全部 4762×6735，
% 均按页序读取。7 页全为试卷正文页（卷首标题在第 1 页页顶，卷末解答题第 21 题 (3)
% 解析止于第 7 页），无封面页、目录页、二维码或宣传页。粉色斜排“魔都数学加油站”
% 为卷面水印，与公众号页脚一并未录入。
%
% 卷首标题：照卷面印作“2026-2027 年中国中学高一上 10 月月考”（黑体加粗一行）。
% 原卷年份区间用半角连字符，本稿按全稿惯例排 en dash（同 高一68、高一69、高一70）；
% 原卷只有一行标题、无副标题。
%
% 篇幅与题号：原卷自第 3 题起编号（第 1、2 题不在本卷内）。填空题 3—12（10 题）、
% 选择题 13—16（4 题）、解答题 17—21（5 题），共 19 题。本稿沿用原节与原题号，
% 只改排为全稿统一的大节样式（原卷小节标题为“一、填空题”“二、选择题”“三、解答题”）。
%
% 讲评版：原卷每题下直接印【答案】或【解析】。**只印【答案】的 3 题**（第 3、4、5 题），
% 本稿只给 \ans{}；**第 13 题既无【答案】也无【解析】段，答案“B”直接印在题干括号内
% （“(  B  )”）**，按 高一37 第 13 题、高一63 第 14 题、高一67 第 15 题成例把括号留空、
% 答案另提为 \ans{$B$}。其余 15 题全部印【解析】（按全稿口径这类题不另提 \ans{}——
% 第 14—16 题的“故选 D／B／B”写在解析末句，亦照录在解析内）。
% 除第 13 题外，全卷选择题作答括号内均为空白（第 14、15、16 题）。
%
% 副标题：原卷未印副标题。本稿按“知识模块”口径标作“集合与常用逻辑用语、不等式”：
% 19 题中集合与常用逻辑用语 10 题（第 3、4、5、6、9、11、12、13、16、21 题；其中第 4 题
% 属反证法，一并计入常用逻辑用语），不等式 9 题（第 7、8、10、14、15、17、18、19、20 题；
% 第 19 题 (1) 属补集、(2) 属分式不等式，按主干计入不等式），两类均确有分量，故并标。
% 口径同 高一63、高一68、高一69、高一70。
%
% 与原卷的排版差异：
%   ① 原卷每题下直接印【答案】/【解析】，本稿按全稿惯例将答案与解析置于答案版，试卷版
%      保持干净；第 13 题印在题干括号内的答案“B”同样移入 \ans{}，见上；
%   ② 原卷未印副标题，本稿按“知识模块”口径补出；
%   ③ 原卷填空横线约四个字宽，本稿按全稿惯例统一排 \blankline[4.4em]；
%   ④ 原卷数集记号不统一：实数集排斜体 $R$（第 19 题 (1)），自然数集排斜体 $N$
%      （第 5 题题干、第 12 题题干与解析共三处），第 16 题解析又排正体 $\mathrm{N}^*$；
%      本稿按全稿惯例统一改用 $\R$、$\N$、$\N^*$（同 高一68／69／70 排版差异④）；
%      空集记号原卷作 $\varnothing$，照录；
%   ⑤ 原卷句末标点混用：第 15 题题干“（3）$x^2+2+\dfrac1{x^2+2}$**．**其中”、
%      第 17 题题干末“$\dfrac{x-2}{x-1}\geqslant0$**．**”两处放大确认（第 15 题为全角
%      句点、第 17 题为西文句点），其余各题末为西文句点，已逐处照录；题干内冒号——
%      第 5 题“$\{y\mid\cdots\}$**：**”、第 12 题““×”**：**”、第 17 题“不等式组**：**”、
%      第 21 题“定义**：**”——均为全角冒号，照录；数学式内的逗号（如第 12 题
%      “$x\in A,y\in B$”、第 16 题“$A=\{a_1,\cdots\},B=\{b_1,\cdots\},n>2$”）按原卷作
%      半角，中文行文里的逗号按原卷作全角，均未强改；
%   ⑥ 第 17—21 题的子问原卷印作“（1）…（2）…”（另起一段），本稿按 高一69／70 成例用
%      行内“（1）（2）（3）”段，解析中的子问标号一律用 \stepm{(1)}{…}；
%      第 21 题 (2) 解析内层的“①”“②”“③”照原卷作 \stepm{①}{…} 等；
%   ⑦ 每道解答题原卷未留作答空白，本稿按全稿惯例加 \ansspace（末题用 \ansspace*）；
%   ⑧ 选择题的作答括号原卷排半角“(　　)”（第 14、15、16 题括号内空白、第 13 题括号内
%      印 B），本稿按全稿惯例排 \mbox{（\quad）}；选项行排照原卷——第 13 题两行两条
%      （A B 一行、C D 一行），第 14、15、16 题一行四条；
%   ⑨ 第 8 题解析、第 9 题解析、第 17 题题干与解析中的多行不等式组，原卷均排作“{”
%      连接的多行式，本稿按全稿惯例排 cases；
%   ⑩ 第 10 题解析、第 17 题解析、第 19 题 (1) 解析的并集符号原卷排大号 $\bigcup$
%      （显示体），本稿按全稿惯例排行内 $\cup$；
%   ⑪ 第 8 题解析的判别式照原卷作空心“$\triangle$”（与第 7 题解析的“$\Delta$”不同形），
%      照录未改；
%   ⑫ 第 16 题解析第 4 行的“$\forall$”（全称量词）照原卷符号排，未改写为中文“任意”。
%
% 原卷笔误/存疑（均照录未改，正文对应处另有 % 注）：
%   ① 第 19 题 (1) 连用两个“若”——“若 $a=4$，若全集为实数集 $R$”，后一“若”应作
%      “当”或删；原卷如此，照录未改。
%   ② 第 20 题题干“可以售出 8 万本**、**据市场调查”——“万本”后用顿号，按文意应为
%      逗号或句号；原卷如此，照录未改。
%   ③ 第 15 题解析对代数式 (1) 只得出“$y/x+x/y\geqslant2$”（同号时），未讨论取等是否
%      可得——取 $x=y$ 时确能取到 2——却总结“最小值为 2 的代数式有 1 个，故选 $B$”；
%      按解析自身的推理 (1) 也应计入，结论“1 个”与推理矛盾（若计 (1) 应为 2 个、选 C）。
%      原卷如此，照录未改。
%   ④ 第 21 题题干“定义：对于集合 $A$，则称集合 $B=\cdots$”——“对于……”与“则称……”
%      句式杂糅（应作“对于集合 $A$，称……”或“若集合 $A$，则称……”）；原卷如此，照录未改。
%   ⑤ 第 13 题题干末“……的（　）条件”后无任何标点（其余各题题干末均有句点）；
%      原卷如此，照录未改。
%
% 行数核对：逐题数 \step/\stepm 条数与原卷【解析】印刷行数，15 题逐一相符，
% 合计 86 条 / 86 行（Q6：1；Q7：1；Q8：7；Q9：2；Q10：4；Q11：7；Q12：5；Q14：7；
% Q15：8；Q16：7；Q17：3；Q18：7；Q19：8；Q20：6；Q21：13）；另有 3 题只印【答案】
% （第 3、4、5 题）与第 13 题括号内答案，各提 \ans{}，不计入 \step 账。
%
% 插图：无题目配图或数学图形（7 页题面均无“如图/右图/图中”字样；页面另有粉色斜排水印，
% 见上）。
%
\documentclass[a4paper,11pt]{article}
\usepackage{common}

\begin{document}

\examtitlest[3]{2026--2027 年中国中学高一上 10 月月考}{集合与常用逻辑用语、不等式}

\bigsec{一、填空题}

\begin{enumerate}
\setcounter{enumi}{2}

\item 已知集合 $A=[1,+\infty)$，$B=[-1,3)$，则 $A\cap B=$\blankline[4.4em].

\ifanswer
\ans{$[1,3)$}
\fi

\item 用反证法证明命题“已知 $a$、$b$ 为正数，$n$ 为正整数，求证：如果 $a^n>b^n$，那么 $a>b$”时，应当作出的假设为：“假设\blankline[4.4em]”.

\ifanswer
\ans{$a\leqslant b$}
\fi

\item 用列举法表示集合 $\{y\mid y=x+1,x\in\N,x\leqslant1\}$：\blankline[4.4em].

\ifanswer
\ans{$\{1,2\}$}
\fi

\item 某班在一次测验中，有 36 人数学成绩不低于 80 分，有 20 人物理成绩不低于 80 分，有 15 人的数学、物理成绩都不低于 80 分，则这两科成绩中至少有一科不低于 80 分的人数为\blankline[4.4em]人.

\ifanswer
\solbegin
\step{至少有 $36+20-15=41$ 人.}
\solend
\fi

\item 若关于 $x$ 的不等式 $x^2-x+m<0$ 的解集是 $\varnothing$，则实数 $m$ 的取值范围是\blankline[4.4em].

\ifanswer
\solbegin
\step{由题意得 $\Delta=1-4m\leqslant0$，所以 $m\geqslant\dfrac14$.}
\solend
\fi

% 第 8 题解析的判别式原卷作空心“△”（与第 7 题解析的“Δ”不同形），照录（见文件头“排版差异”⑪）
\item 已知关于 $x$ 的方程 $x^2+mx+2m=0$ 的两个实数根分别为 $\alpha$，$\beta$，若 $\alpha^2+\beta^2=5$，则实数 $m=$\blankline[4.4em].

\ifanswer
\solbegin
\step{方程 $x^2+mx+2m=0$ 的两个实数根为 $\alpha$、$\beta$，}
\step{由根与系数的关系得 $\begin{cases}\alpha+\beta=-m\\ \alpha\beta=2m\end{cases}$，又因为 $\alpha^2+\beta^2=5$，}
\step{所以 $\alpha^2+\beta^2=(\alpha+\beta)^2-2\alpha\beta=m^2-4m=5$，整理得 $m^2-4m-5=0$，}
\step{解得 $m=5$ 或 $m=-1$，}
\step{当 $m=5$ 时，$\triangle=25-40=-15<0$，无实根，舍去；}
\step{当 $m=-1$ 时，$\triangle=1+8=9>0$，符合题意，}
\step{所以实数 $m=-1$.}
\solend
\fi

\item 已知集合 $A=\{x\mid6x+a>0\}$，若 $1\in A$，$0\notin A$，则实数 $a$ 的取值范围为\blankline[4.4em].

\ifanswer
\solbegin
\step{由 $A=\{x\mid6x+a>0\}$，$1\in A$，$0\notin A$，得 $\begin{cases}6+a>0\\ a\leqslant0\end{cases}$，解得 $-6<a\leqslant0$，}
\step{即实数 $a$ 的取值范围是 $(-6,0]$.}
\solend
\fi

\item 关于 $x$ 的不等式 $ax-b>0$ 的解集为 $(1,+\infty)$，则关于 $x$ 的不等式 $(ax+b)(x-2)>0$ 的解集为\blankline[4.4em].

\ifanswer
\solbegin
\step{由不等式 $ax-b>0$ 的解集为 $(1,+\infty)$，得 $a>0$ 且 $\dfrac ba=1$，即 $b=a$，}
\step{则 $(ax+b)(x-2)>0$，可化为 $a(x+1)(x-2)>0$.}
\step{即 $(x+1)(x-2)>0$，解得 $x<-1$ 或 $x>2$，}
\step{所以不等式的解集为 $(-\infty,-1)\cup(2,+\infty)$.}
\solend
\fi

\item 已知实数 $a_1$、$a_2$、$a_3$、$a_4$、$a_5$ 满足 $\{x\mid(x-a_1)(x-a_2)(x-a_3)(x-a_4)(x-a_5)=0\}=\{1,2,3,4\}$，则 $a_1+a_2+a_3+a_4+a_5$ 的最大值为\blankline[4.4em].

\ifanswer
\solbegin
\step{5 个实数 $a_1$、$a_2$、$a_3$、$a_4$、$a_5$ 的不同取值只有 1，2，3，4，}
\step{即这 5 个数均从 1，2，3，4 中选取.}
\step{集合有 4 个元素，说明 1，2，3，4 每个数至少出现一次，}
\step{因此恰好有一个数重复出现两次，其余三个数各出现一次.}
\step{要让 5 个数的总和最大，应让重复出现两次的数取集合中的最大值 4，}
\step{此时 5 个数为 1，2，3，4，4，总和为 $1+2+3+4+4=14$，}
\step{即所求最大值为 14.}
\solend
\fi

\item 定义集合运算“$\times$”：$A\times B=\{(x,y)\mid x\in A,y\in B\}$，称为 $A$、$B$ 两个集合的“卡氏积”，若 $A=\{x\mid x^2-|x|\leqslant0,x\in\N\}$，$B=\{1,2,3\}$，则 $(A\times B)\cap(B\times A)=$\blankline[4.4em].

\ifanswer
\solbegin
\step{因为 $x^2-|x|\leqslant0$，所以 $0\leqslant|x|\leqslant1$，又因为 $x\in\N$，}
\step{所以 $A=\{x\mid x^2-|x|\leqslant0,x\in\N\}=\{0,1\}$，又 $B=\{1,2,3\}$，}
\step{由“卡氏积”的定义得 $A\times B=\{(0,1),(0,2),(0,3),(1,1),(1,2),(1,3)\}$，}
\step{$B\times A=\{(1,0),(1,1),(2,0),(2,1),(3,0),(3,1)\}$，}
\step{所以 $(A\times B)\cap(B\times A)=\{(1,1)\}$.}
\solend
\fi

\end{enumerate}

\bigsec{二、选择题}

\begin{enumerate}
\setcounter{enumi}{12}

% 第 13 题答案“B”原卷印在题干括号内，且无【答案】/【解析】段；
% 按 高一37 第 13 题、高一63 第 14 题、高一67 第 15 题成例括号留空、答案另提 \ans{$B$}
% 原卷题干末“条件”后无标点，照录（见文件头“原卷笔误/存疑”⑤）
\item “$A\subseteq B$”是“$A\subset B$”的\mbox{（\quad）}条件

\keepwithprev
A. 充分不必要条件\qquad B. 必要不充分条件

C. 充要条件\qquad D. 既不充分也不必要条件

\ifanswer
\ans{$B$}
\fi

\item 已知 $a$、$b$、$c$ 均为实数，若 $a>b$ 且 $ab\neq0$，则下列不等关系中成立的是\mbox{（\quad）}

\keepwithprev
A. $ac^2>bc^2$\qquad B. $\dfrac1b>\dfrac1a$\qquad C. $|a|>|b|$\qquad D. $c-b>c-a$

\ifanswer
\solbegin
\step{对于 $A$，若 $a>b$，当 $c=0$ 时，有 $ac^2=bc^2=0$，故 $A$ 错误；}
\step{对于 $B$，$\dfrac1b-\dfrac1a=\dfrac{a-b}{ab}$，因为 $a>b$，所以 $a-b>0$，}
\step{当 $a$、$b$ 同号时，$ab>0$，则 $\dfrac1b-\dfrac1a>0$，即 $\dfrac1b>\dfrac1a$，}
\step{当 $a$、$b$ 异号时，$ab<0$，则 $\dfrac1b-\dfrac1a<0$，即 $\dfrac1b<\dfrac1a$，故 $B$ 错误；}
\step{对于 $C$，当 $a=1$，$b=-3$ 时，$|a|<|b|$，故 $C$ 错误；}
\step{对于 $D$，若 $a>b$，则 $-b>-a$，$c-b>c-a$，故 $D$ 正确.}
\step{故选 $D$.}
\solend
\fi

% 原卷解析对 (1) 未讨论取等是否可得（x=y 时确能取 2），却总结“有 1 个、故选 B”，
% 与推理矛盾，照录（见文件头“原卷笔误/存疑”③）
\item 以下三个代数式：（1）$\dfrac yx+\dfrac xy$；（2）$|x|+\dfrac1{|x|}$；（3）$x^2+2+\dfrac1{x^2+2}$．其中的字母均表示实数，且代数式均有意义，那么代数式的最小值为 2 的个数是\mbox{（\quad）}

\keepwithprev
A. $0$\qquad B. $1$\qquad C. $2$\qquad D. $3$

\ifanswer
\solbegin
\step{（1）仅当 $x,y$ 同号时，$\dfrac yx,\dfrac xy>0$，由基本不等式得 $\dfrac yx+\dfrac xy\geqslant2$；}
\step{若 $x,y$ 异号，$\dfrac yx,\dfrac xy<0$，故 $\dfrac yx+\dfrac xy<0$，此时最小值不是 2；}
\step{（2）$|x|+\dfrac1{|x|}$ 中，$|x|>0$，由基本不等式 $|x|+\dfrac1{|x|}\geqslant2\sqrt{|x|\cdot\dfrac1{|x|}}=2$，}
\step{当且仅当 $|x|=1$，即 $x=\pm1$ 时取等号，故最小值为 2；}
\step{（3）令 $t=x^2+2$，则 $t\geqslant2$，函数化为 $t+\dfrac1t$，}
\step{对勾函数 $t+\dfrac1t$ 在 $[2,+\infty)$ 上严格增，最小值为 $2+\dfrac12=\dfrac52>2$，}
\step{等号成立条件 $t=1$ 无法取到，故最小值不是 2；}
\step{综上，最小值为 2 的代数式有 1 个，故选 $B$.}
\solend
\fi

\item 已知 $A$、$B$ 为非空实数集，$\Omega$ 为平面直角坐标系中的一些点构成的集合，集合 $C=\{y\mid$对任意 $x\in A$，有 $(x,y)\in\Omega\}$，集合 $D=\{x\mid$对任意 $y\in B$，有 $(x,y)\in\Omega\}$，对于下列两个命题：①若 $C\subseteq B$，则 $D\subseteq A$；②若 $C\supseteq B$，则 $D\subseteq A$；其中判断正确的是\mbox{（\quad）}

\keepwithprev
A. ①②都正确\qquad B. ①②都错误\qquad C. ①正确②错误\qquad D. ①错误②正确

\ifanswer
\solbegin
\step{设 $A=\{a_1,a_2,a_3,\cdots,a_n\},B=\{b_1,b_2,b_3,\cdots,b_n\},n>2$，}
\step{若 $\Omega=\{(c,b_1),(c,b_2),\cdots,(c,b_n)\},c\notin A$，}
\step{此时 $C=\varnothing\subseteq B,D=\{c\},D\subseteq A$ 错误，故①错误；}
% 原卷此处 N* 排正体（\mathrm{N}*），本稿统一 \N^*（见文件头“排版差异”④）；
% 全称量词“∀”照原卷符号排（见文件头“排版差异”⑫）
\step{若 $C\supseteq B$，则 $\forall(a_i,b_j)(1\leqslant i,j\leqslant n,i,j\in\N^*)\in\Omega$，则 $D\supseteq A$，}
\step{且 $A=\{a_1\},B=\{b_1,b_2\}$，}
\step{若 $\Omega=\{(a_1,b_1),(a_2,b_1),(a_1,b_2),(a_2,b_2),(c,b_1),(c,b_2)\},D$ 真包含 $A$，故②错误.}
\step{故选 $B$.}
\solend
\fi

\end{enumerate}

\bigsec{三、解答题}

\begin{enumerate}
\setcounter{enumi}{16}

\item 解下列关于 $x$ 的不等式组：$\begin{cases}|x-1|<2\\ \dfrac{x-2}{x-1}\geqslant0\end{cases}$.

\ifanswer
\solbegin
\step{解 $|x-1|<2$ 得 $-2<x-1<2$，解得 $-1<x<3$.}
\step{分式不等式 $\dfrac{x-2}{x-1}\geqslant0$ 等价于 $\begin{cases}(x-2)(x-1)\geqslant0\\ x-1\neq0\end{cases}$，解得 $x<1$ 或 $x\geqslant2$.}
\step{即 $\begin{cases}|x-1|<2\\ \dfrac{x-2}{x-1}\geqslant0\end{cases}$ 的解集为 $(-1,1)\cup[2,3)$.}
\solend
\fi

\ansspace[5]

\item 比较多项式 $2x^2-3x+1$ 与 $x^2-1$ 的大小.

\ifanswer
\solbegin
\step{作差得 $(2x^2-3x+1)-(x^2-1)=x^2-3x+2=(x-1)(x-2)$，}
\step{当 $x=1$ 或 $x=2$ 时，$(x-1)(x-2)=0$，即 $2x^2-3x+1=x^2-1$，}
\step{当 $x<1$ 或 $x>2$ 时，$(x-1)(x-2)>0$，即 $2x^2-3x+1>x^2-1$，}
\step{当 $1<x<2$ 时，$(x-1)(x-2)<0$，即 $2x^2-3x+1<x^2-1$；}
\step{综上，当 $x=1$ 或 $x=2$ 时，$2x^2-3x+1=x^2-1$；}
\step{当 $x<1$ 或 $x>2$ 时，$2x^2-3x+1>x^2-1$；}
\step{当 $1<x<2$ 时，$2x^2-3x+1<x^2-1$.}
\solend
\fi

\ansspace[6]

% 原卷 (1) 连用两个“若”（“若 a=4，若全集为实数集 R”），照录（见文件头“原卷笔误/存疑”①）
\item 已知集合 $A=\left\{x\mid\dfrac{(a+1)x}{x-1}<a\right\}$，集合 $B=[1,4]$，$a$ 为实常数；

（1）若 $a=4$，若全集为实数集 $\R$，求集合 $\overline{A}$（用区间形式表示）；

（2）若 $A\subseteq B$，求 $a$ 的取值范围.

\ifanswer
\solbegin
\stepm{(1)}{若 $a=4$，集合 $A=\left\{x\mid\dfrac{5x}{x-1}<4\right\}=\left\{x\mid\dfrac{5x}{x-1}-4<0\right\}=\left\{x\mid\dfrac{x+4}{x-1}<0\right\}$}
\step{$=(-4,1)$，因此 $\overline{A}=(-\infty,-4]\cup[1,+\infty)$；}
\stepm{(2)}{当 $-a<1$，即 $a>-1$ 时，$A=(-a,1)$，}
\step{要使 $A\subseteq[1,4]$，需 $-a\geqslant1$，即 $a\leqslant-1$，与 $a>-1$ 矛盾，舍去；}
\step{当 $-a>1$，即 $a<-1$ 时，$A=(1,-a)$，}
\step{要使 $A\subseteq[1,4]$，需 $-a\leqslant4$，即 $a\geqslant-4$，所以 $-4\leqslant a<-1$；}
\step{当 $-a=1$，即 $a=-1$ 时，$A=\varnothing$，满足 $A\subseteq B$；}
\step{综上，实数 $a$ 的取值范围为 $[-4,-1]$.}
\solend
\fi

\ansspace[7]

% 原卷题干“可以售出 8 万本、据市场调查”处用顿号，照录（见文件头“原卷笔误/存疑”②）
\item 某种杂志原以每本 25 元的价格销售，可以售出 8 万本、据市场调查，杂志的单价每提高 1 元，销售量就减少 0.2 万本，假定杂志的成本是每本 10 元（不计其他成本）.

（1）当杂志以每本 30 元定价销售时，求销售该杂志所获利润；

（2）在每本 25 元的价格的基础上提高定价，试确定杂志的定价在什么范围内可使得销售利润不低于 140 万元.

\ifanswer
\solbegin
\stepm{(1)}{当杂志以 30 元定价销售时，}
\step{销售该杂志所获利润为 $[8-(30-25)\times0.2]\times(30-10)=140$ 万元；}
\stepm{(2)}{在每本 25 元的价格的基础上提高定价，设杂志的定价为每本 $x$ 元，}
\step{由题意得 $[8-(x-25)\times0.2]\times(x-10)\geqslant140$，解得 $30\leqslant x\leqslant45$，}
\step{所以杂志的定价在每本不低于 30 元且不高于 45 元的范围内，}
\step{可使得销售利润不低于 140 万元.}
\solend
\fi

\ansspace[7]

% 原卷题干“定义：对于集合 A，则称集合 B=…”句式杂糅，照录（见文件头“原卷笔误/存疑”④）
\item 定义：对于集合 $A$，则称集合 $B=\left\{x\mid x=\dfrac{m+n}{1+mn},m,n\in A,1+mn\neq0\right\}$ 为集合 $A$ 的生成集（其中 $m$、$n$ 可以相等）.

（1）若集合 $A=\{-1,1\}$ 的生成集为 $B$，求集合 $B$；

（2）若集合 $A=\{a,2\}$ 的生成集为 $B$，$B$ 的子集个数为 4 个，求实数 $a$ 的值；

（3）若集合 $A=\{x\mid-1<x<1\}$ 的生成集为 $B$，求证：$B\subseteq A$.

\ifanswer
\solbegin
\stepm{(1)}{因为 $1+mn\neq0$，}
\step{当 $m=n=1$ 时，$x=\dfrac{m+n}{1+mn}=1$；当 $m=n=-1$ 时，$x=\dfrac{m+n}{1+mn}=-1$；}
\step{所以 $B=\{-1,1\}$；}
\stepm{(2)}{$B$ 的子集个数为 4 个，故 $B$ 集合的元素有 2 个，}
\step{当 $m=n=2$ 时，$x=\dfrac{m+n}{1+mn}=\dfrac45$，当 $m=n=a$ 时，$x=\dfrac{m+n}{1+mn}=\dfrac{2a}{1+a^2}$，}
\step{当 $m$ 与 $n$ 分别取 2 和 $a$ 时，$x=\dfrac{a+2}{1+2a}$，}
\stepm{①}{$\dfrac{2a}{1+a^2}=\dfrac45$，解得 $a=2$（舍）或 $a=\dfrac12$，}
\stepm{②}{$\dfrac{a+2}{1+2a}=\dfrac45$，解得 $a=2$（舍），}
\stepm{③}{$\dfrac{2a}{1+a^2}=\dfrac{a+2}{1+2a}$，解得 $a=2$（舍）或 $a=1$ 或 $a=-1$，}
\step{综上所述，$a=1$ 或 $a=-1$ 或 $a=\dfrac12$，即 $a\in\left\{-1,1,\dfrac12\right\}$；}
\stepm{(3)}{$\forall m$、$n\in A$，则 $\dfrac{m+n}{1+mn}+1=\dfrac{(m+1)(n+1)}{1+mn}>0$，故 $\dfrac{m+n}{1+mn}>-1$；}
\step{$\dfrac{m+n}{1+mn}-1=-\dfrac{(m-1)(n-1)}{1+mn}<0$，故 $\dfrac{m+n}{1+mn}<1$；}
\step{故 $B\subseteq\{x\mid-1<x<1\}$，即 $B\subseteq A$.}
\solend
\fi

\ansspace*[10]

\end{enumerate}

\end{document}
"
% 紧凑版：xelatex "\def\iscompact{1}% 高一71_中国中学_10月月考.tex
%   —— 高一上数学“2026-2027 年中国中学高一上 10 月月考”（讲评版）
% 试卷版：xelatex 高一71_中国中学_10月月考.tex
% 答案版：xelatex "\def\isanswer{1}\input{高一71_中国中学_10月月考.tex}"
% 紧凑版：xelatex "\def\iscompact{1}\input{高一71_中国中学_10月月考.tex}"
%
% 原始材料：微信公众号“魔都数学加油站”文章，连载“27学年高一71”；
% 原文标题“27学年高一71——2026-2027年上海市中国中学高一上10月月考”；
% 原文链接：https://mp.weixin.qq.com/s/UZrP53l3GIin040L1EkljA。
% 页面标注发布时间 2026-10-19 22:47（上海）（与 高一69、高一70 页面读出同一 ct 值）。
% 正文为 7 张 PNG 页（`原图/高一71_中国中学_10月月考/01..07.png`），全部 4762×6735，
% 均按页序读取。7 页全为试卷正文页（卷首标题在第 1 页页顶，卷末解答题第 21 题 (3)
% 解析止于第 7 页），无封面页、目录页、二维码或宣传页。粉色斜排“魔都数学加油站”
% 为卷面水印，与公众号页脚一并未录入。
%
% 卷首标题：照卷面印作“2026-2027 年中国中学高一上 10 月月考”（黑体加粗一行）。
% 原卷年份区间用半角连字符，本稿按全稿惯例排 en dash（同 高一68、高一69、高一70）；
% 原卷只有一行标题、无副标题。
%
% 篇幅与题号：原卷自第 3 题起编号（第 1、2 题不在本卷内）。填空题 3—12（10 题）、
% 选择题 13—16（4 题）、解答题 17—21（5 题），共 19 题。本稿沿用原节与原题号，
% 只改排为全稿统一的大节样式（原卷小节标题为“一、填空题”“二、选择题”“三、解答题”）。
%
% 讲评版：原卷每题下直接印【答案】或【解析】。**只印【答案】的 3 题**（第 3、4、5 题），
% 本稿只给 \ans{}；**第 13 题既无【答案】也无【解析】段，答案“B”直接印在题干括号内
% （“(  B  )”）**，按 高一37 第 13 题、高一63 第 14 题、高一67 第 15 题成例把括号留空、
% 答案另提为 \ans{$B$}。其余 15 题全部印【解析】（按全稿口径这类题不另提 \ans{}——
% 第 14—16 题的“故选 D／B／B”写在解析末句，亦照录在解析内）。
% 除第 13 题外，全卷选择题作答括号内均为空白（第 14、15、16 题）。
%
% 副标题：原卷未印副标题。本稿按“知识模块”口径标作“集合与常用逻辑用语、不等式”：
% 19 题中集合与常用逻辑用语 10 题（第 3、4、5、6、9、11、12、13、16、21 题；其中第 4 题
% 属反证法，一并计入常用逻辑用语），不等式 9 题（第 7、8、10、14、15、17、18、19、20 题；
% 第 19 题 (1) 属补集、(2) 属分式不等式，按主干计入不等式），两类均确有分量，故并标。
% 口径同 高一63、高一68、高一69、高一70。
%
% 与原卷的排版差异：
%   ① 原卷每题下直接印【答案】/【解析】，本稿按全稿惯例将答案与解析置于答案版，试卷版
%      保持干净；第 13 题印在题干括号内的答案“B”同样移入 \ans{}，见上；
%   ② 原卷未印副标题，本稿按“知识模块”口径补出；
%   ③ 原卷填空横线约四个字宽，本稿按全稿惯例统一排 \blankline[4.4em]；
%   ④ 原卷数集记号不统一：实数集排斜体 $R$（第 19 题 (1)），自然数集排斜体 $N$
%      （第 5 题题干、第 12 题题干与解析共三处），第 16 题解析又排正体 $\mathrm{N}^*$；
%      本稿按全稿惯例统一改用 $\R$、$\N$、$\N^*$（同 高一68／69／70 排版差异④）；
%      空集记号原卷作 $\varnothing$，照录；
%   ⑤ 原卷句末标点混用：第 15 题题干“（3）$x^2+2+\dfrac1{x^2+2}$**．**其中”、
%      第 17 题题干末“$\dfrac{x-2}{x-1}\geqslant0$**．**”两处放大确认（第 15 题为全角
%      句点、第 17 题为西文句点），其余各题末为西文句点，已逐处照录；题干内冒号——
%      第 5 题“$\{y\mid\cdots\}$**：**”、第 12 题““×”**：**”、第 17 题“不等式组**：**”、
%      第 21 题“定义**：**”——均为全角冒号，照录；数学式内的逗号（如第 12 题
%      “$x\in A,y\in B$”、第 16 题“$A=\{a_1,\cdots\},B=\{b_1,\cdots\},n>2$”）按原卷作
%      半角，中文行文里的逗号按原卷作全角，均未强改；
%   ⑥ 第 17—21 题的子问原卷印作“（1）…（2）…”（另起一段），本稿按 高一69／70 成例用
%      行内“（1）（2）（3）”段，解析中的子问标号一律用 \stepm{(1)}{…}；
%      第 21 题 (2) 解析内层的“①”“②”“③”照原卷作 \stepm{①}{…} 等；
%   ⑦ 每道解答题原卷未留作答空白，本稿按全稿惯例加 \ansspace（末题用 \ansspace*）；
%   ⑧ 选择题的作答括号原卷排半角“(　　)”（第 14、15、16 题括号内空白、第 13 题括号内
%      印 B），本稿按全稿惯例排 \mbox{（\quad）}；选项行排照原卷——第 13 题两行两条
%      （A B 一行、C D 一行），第 14、15、16 题一行四条；
%   ⑨ 第 8 题解析、第 9 题解析、第 17 题题干与解析中的多行不等式组，原卷均排作“{”
%      连接的多行式，本稿按全稿惯例排 cases；
%   ⑩ 第 10 题解析、第 17 题解析、第 19 题 (1) 解析的并集符号原卷排大号 $\bigcup$
%      （显示体），本稿按全稿惯例排行内 $\cup$；
%   ⑪ 第 8 题解析的判别式照原卷作空心“$\triangle$”（与第 7 题解析的“$\Delta$”不同形），
%      照录未改；
%   ⑫ 第 16 题解析第 4 行的“$\forall$”（全称量词）照原卷符号排，未改写为中文“任意”。
%
% 原卷笔误/存疑（均照录未改，正文对应处另有 % 注）：
%   ① 第 19 题 (1) 连用两个“若”——“若 $a=4$，若全集为实数集 $R$”，后一“若”应作
%      “当”或删；原卷如此，照录未改。
%   ② 第 20 题题干“可以售出 8 万本**、**据市场调查”——“万本”后用顿号，按文意应为
%      逗号或句号；原卷如此，照录未改。
%   ③ 第 15 题解析对代数式 (1) 只得出“$y/x+x/y\geqslant2$”（同号时），未讨论取等是否
%      可得——取 $x=y$ 时确能取到 2——却总结“最小值为 2 的代数式有 1 个，故选 $B$”；
%      按解析自身的推理 (1) 也应计入，结论“1 个”与推理矛盾（若计 (1) 应为 2 个、选 C）。
%      原卷如此，照录未改。
%   ④ 第 21 题题干“定义：对于集合 $A$，则称集合 $B=\cdots$”——“对于……”与“则称……”
%      句式杂糅（应作“对于集合 $A$，称……”或“若集合 $A$，则称……”）；原卷如此，照录未改。
%   ⑤ 第 13 题题干末“……的（　）条件”后无任何标点（其余各题题干末均有句点）；
%      原卷如此，照录未改。
%
% 行数核对：逐题数 \step/\stepm 条数与原卷【解析】印刷行数，15 题逐一相符，
% 合计 86 条 / 86 行（Q6：1；Q7：1；Q8：7；Q9：2；Q10：4；Q11：7；Q12：5；Q14：7；
% Q15：8；Q16：7；Q17：3；Q18：7；Q19：8；Q20：6；Q21：13）；另有 3 题只印【答案】
% （第 3、4、5 题）与第 13 题括号内答案，各提 \ans{}，不计入 \step 账。
%
% 插图：无题目配图或数学图形（7 页题面均无“如图/右图/图中”字样；页面另有粉色斜排水印，
% 见上）。
%
\documentclass[a4paper,11pt]{article}
\usepackage{common}

\begin{document}

\examtitlest[3]{2026--2027 年中国中学高一上 10 月月考}{集合与常用逻辑用语、不等式}

\bigsec{一、填空题}

\begin{enumerate}
\setcounter{enumi}{2}

\item 已知集合 $A=[1,+\infty)$，$B=[-1,3)$，则 $A\cap B=$\blankline[4.4em].

\ifanswer
\ans{$[1,3)$}
\fi

\item 用反证法证明命题“已知 $a$、$b$ 为正数，$n$ 为正整数，求证：如果 $a^n>b^n$，那么 $a>b$”时，应当作出的假设为：“假设\blankline[4.4em]”.

\ifanswer
\ans{$a\leqslant b$}
\fi

\item 用列举法表示集合 $\{y\mid y=x+1,x\in\N,x\leqslant1\}$：\blankline[4.4em].

\ifanswer
\ans{$\{1,2\}$}
\fi

\item 某班在一次测验中，有 36 人数学成绩不低于 80 分，有 20 人物理成绩不低于 80 分，有 15 人的数学、物理成绩都不低于 80 分，则这两科成绩中至少有一科不低于 80 分的人数为\blankline[4.4em]人.

\ifanswer
\solbegin
\step{至少有 $36+20-15=41$ 人.}
\solend
\fi

\item 若关于 $x$ 的不等式 $x^2-x+m<0$ 的解集是 $\varnothing$，则实数 $m$ 的取值范围是\blankline[4.4em].

\ifanswer
\solbegin
\step{由题意得 $\Delta=1-4m\leqslant0$，所以 $m\geqslant\dfrac14$.}
\solend
\fi

% 第 8 题解析的判别式原卷作空心“△”（与第 7 题解析的“Δ”不同形），照录（见文件头“排版差异”⑪）
\item 已知关于 $x$ 的方程 $x^2+mx+2m=0$ 的两个实数根分别为 $\alpha$，$\beta$，若 $\alpha^2+\beta^2=5$，则实数 $m=$\blankline[4.4em].

\ifanswer
\solbegin
\step{方程 $x^2+mx+2m=0$ 的两个实数根为 $\alpha$、$\beta$，}
\step{由根与系数的关系得 $\begin{cases}\alpha+\beta=-m\\ \alpha\beta=2m\end{cases}$，又因为 $\alpha^2+\beta^2=5$，}
\step{所以 $\alpha^2+\beta^2=(\alpha+\beta)^2-2\alpha\beta=m^2-4m=5$，整理得 $m^2-4m-5=0$，}
\step{解得 $m=5$ 或 $m=-1$，}
\step{当 $m=5$ 时，$\triangle=25-40=-15<0$，无实根，舍去；}
\step{当 $m=-1$ 时，$\triangle=1+8=9>0$，符合题意，}
\step{所以实数 $m=-1$.}
\solend
\fi

\item 已知集合 $A=\{x\mid6x+a>0\}$，若 $1\in A$，$0\notin A$，则实数 $a$ 的取值范围为\blankline[4.4em].

\ifanswer
\solbegin
\step{由 $A=\{x\mid6x+a>0\}$，$1\in A$，$0\notin A$，得 $\begin{cases}6+a>0\\ a\leqslant0\end{cases}$，解得 $-6<a\leqslant0$，}
\step{即实数 $a$ 的取值范围是 $(-6,0]$.}
\solend
\fi

\item 关于 $x$ 的不等式 $ax-b>0$ 的解集为 $(1,+\infty)$，则关于 $x$ 的不等式 $(ax+b)(x-2)>0$ 的解集为\blankline[4.4em].

\ifanswer
\solbegin
\step{由不等式 $ax-b>0$ 的解集为 $(1,+\infty)$，得 $a>0$ 且 $\dfrac ba=1$，即 $b=a$，}
\step{则 $(ax+b)(x-2)>0$，可化为 $a(x+1)(x-2)>0$.}
\step{即 $(x+1)(x-2)>0$，解得 $x<-1$ 或 $x>2$，}
\step{所以不等式的解集为 $(-\infty,-1)\cup(2,+\infty)$.}
\solend
\fi

\item 已知实数 $a_1$、$a_2$、$a_3$、$a_4$、$a_5$ 满足 $\{x\mid(x-a_1)(x-a_2)(x-a_3)(x-a_4)(x-a_5)=0\}=\{1,2,3,4\}$，则 $a_1+a_2+a_3+a_4+a_5$ 的最大值为\blankline[4.4em].

\ifanswer
\solbegin
\step{5 个实数 $a_1$、$a_2$、$a_3$、$a_4$、$a_5$ 的不同取值只有 1，2，3，4，}
\step{即这 5 个数均从 1，2，3，4 中选取.}
\step{集合有 4 个元素，说明 1，2，3，4 每个数至少出现一次，}
\step{因此恰好有一个数重复出现两次，其余三个数各出现一次.}
\step{要让 5 个数的总和最大，应让重复出现两次的数取集合中的最大值 4，}
\step{此时 5 个数为 1，2，3，4，4，总和为 $1+2+3+4+4=14$，}
\step{即所求最大值为 14.}
\solend
\fi

\item 定义集合运算“$\times$”：$A\times B=\{(x,y)\mid x\in A,y\in B\}$，称为 $A$、$B$ 两个集合的“卡氏积”，若 $A=\{x\mid x^2-|x|\leqslant0,x\in\N\}$，$B=\{1,2,3\}$，则 $(A\times B)\cap(B\times A)=$\blankline[4.4em].

\ifanswer
\solbegin
\step{因为 $x^2-|x|\leqslant0$，所以 $0\leqslant|x|\leqslant1$，又因为 $x\in\N$，}
\step{所以 $A=\{x\mid x^2-|x|\leqslant0,x\in\N\}=\{0,1\}$，又 $B=\{1,2,3\}$，}
\step{由“卡氏积”的定义得 $A\times B=\{(0,1),(0,2),(0,3),(1,1),(1,2),(1,3)\}$，}
\step{$B\times A=\{(1,0),(1,1),(2,0),(2,1),(3,0),(3,1)\}$，}
\step{所以 $(A\times B)\cap(B\times A)=\{(1,1)\}$.}
\solend
\fi

\end{enumerate}

\bigsec{二、选择题}

\begin{enumerate}
\setcounter{enumi}{12}

% 第 13 题答案“B”原卷印在题干括号内，且无【答案】/【解析】段；
% 按 高一37 第 13 题、高一63 第 14 题、高一67 第 15 题成例括号留空、答案另提 \ans{$B$}
% 原卷题干末“条件”后无标点，照录（见文件头“原卷笔误/存疑”⑤）
\item “$A\subseteq B$”是“$A\subset B$”的\mbox{（\quad）}条件

\keepwithprev
A. 充分不必要条件\qquad B. 必要不充分条件

C. 充要条件\qquad D. 既不充分也不必要条件

\ifanswer
\ans{$B$}
\fi

\item 已知 $a$、$b$、$c$ 均为实数，若 $a>b$ 且 $ab\neq0$，则下列不等关系中成立的是\mbox{（\quad）}

\keepwithprev
A. $ac^2>bc^2$\qquad B. $\dfrac1b>\dfrac1a$\qquad C. $|a|>|b|$\qquad D. $c-b>c-a$

\ifanswer
\solbegin
\step{对于 $A$，若 $a>b$，当 $c=0$ 时，有 $ac^2=bc^2=0$，故 $A$ 错误；}
\step{对于 $B$，$\dfrac1b-\dfrac1a=\dfrac{a-b}{ab}$，因为 $a>b$，所以 $a-b>0$，}
\step{当 $a$、$b$ 同号时，$ab>0$，则 $\dfrac1b-\dfrac1a>0$，即 $\dfrac1b>\dfrac1a$，}
\step{当 $a$、$b$ 异号时，$ab<0$，则 $\dfrac1b-\dfrac1a<0$，即 $\dfrac1b<\dfrac1a$，故 $B$ 错误；}
\step{对于 $C$，当 $a=1$，$b=-3$ 时，$|a|<|b|$，故 $C$ 错误；}
\step{对于 $D$，若 $a>b$，则 $-b>-a$，$c-b>c-a$，故 $D$ 正确.}
\step{故选 $D$.}
\solend
\fi

% 原卷解析对 (1) 未讨论取等是否可得（x=y 时确能取 2），却总结“有 1 个、故选 B”，
% 与推理矛盾，照录（见文件头“原卷笔误/存疑”③）
\item 以下三个代数式：（1）$\dfrac yx+\dfrac xy$；（2）$|x|+\dfrac1{|x|}$；（3）$x^2+2+\dfrac1{x^2+2}$．其中的字母均表示实数，且代数式均有意义，那么代数式的最小值为 2 的个数是\mbox{（\quad）}

\keepwithprev
A. $0$\qquad B. $1$\qquad C. $2$\qquad D. $3$

\ifanswer
\solbegin
\step{（1）仅当 $x,y$ 同号时，$\dfrac yx,\dfrac xy>0$，由基本不等式得 $\dfrac yx+\dfrac xy\geqslant2$；}
\step{若 $x,y$ 异号，$\dfrac yx,\dfrac xy<0$，故 $\dfrac yx+\dfrac xy<0$，此时最小值不是 2；}
\step{（2）$|x|+\dfrac1{|x|}$ 中，$|x|>0$，由基本不等式 $|x|+\dfrac1{|x|}\geqslant2\sqrt{|x|\cdot\dfrac1{|x|}}=2$，}
\step{当且仅当 $|x|=1$，即 $x=\pm1$ 时取等号，故最小值为 2；}
\step{（3）令 $t=x^2+2$，则 $t\geqslant2$，函数化为 $t+\dfrac1t$，}
\step{对勾函数 $t+\dfrac1t$ 在 $[2,+\infty)$ 上严格增，最小值为 $2+\dfrac12=\dfrac52>2$，}
\step{等号成立条件 $t=1$ 无法取到，故最小值不是 2；}
\step{综上，最小值为 2 的代数式有 1 个，故选 $B$.}
\solend
\fi

\item 已知 $A$、$B$ 为非空实数集，$\Omega$ 为平面直角坐标系中的一些点构成的集合，集合 $C=\{y\mid$对任意 $x\in A$，有 $(x,y)\in\Omega\}$，集合 $D=\{x\mid$对任意 $y\in B$，有 $(x,y)\in\Omega\}$，对于下列两个命题：①若 $C\subseteq B$，则 $D\subseteq A$；②若 $C\supseteq B$，则 $D\subseteq A$；其中判断正确的是\mbox{（\quad）}

\keepwithprev
A. ①②都正确\qquad B. ①②都错误\qquad C. ①正确②错误\qquad D. ①错误②正确

\ifanswer
\solbegin
\step{设 $A=\{a_1,a_2,a_3,\cdots,a_n\},B=\{b_1,b_2,b_3,\cdots,b_n\},n>2$，}
\step{若 $\Omega=\{(c,b_1),(c,b_2),\cdots,(c,b_n)\},c\notin A$，}
\step{此时 $C=\varnothing\subseteq B,D=\{c\},D\subseteq A$ 错误，故①错误；}
% 原卷此处 N* 排正体（\mathrm{N}*），本稿统一 \N^*（见文件头“排版差异”④）；
% 全称量词“∀”照原卷符号排（见文件头“排版差异”⑫）
\step{若 $C\supseteq B$，则 $\forall(a_i,b_j)(1\leqslant i,j\leqslant n,i,j\in\N^*)\in\Omega$，则 $D\supseteq A$，}
\step{且 $A=\{a_1\},B=\{b_1,b_2\}$，}
\step{若 $\Omega=\{(a_1,b_1),(a_2,b_1),(a_1,b_2),(a_2,b_2),(c,b_1),(c,b_2)\},D$ 真包含 $A$，故②错误.}
\step{故选 $B$.}
\solend
\fi

\end{enumerate}

\bigsec{三、解答题}

\begin{enumerate}
\setcounter{enumi}{16}

\item 解下列关于 $x$ 的不等式组：$\begin{cases}|x-1|<2\\ \dfrac{x-2}{x-1}\geqslant0\end{cases}$.

\ifanswer
\solbegin
\step{解 $|x-1|<2$ 得 $-2<x-1<2$，解得 $-1<x<3$.}
\step{分式不等式 $\dfrac{x-2}{x-1}\geqslant0$ 等价于 $\begin{cases}(x-2)(x-1)\geqslant0\\ x-1\neq0\end{cases}$，解得 $x<1$ 或 $x\geqslant2$.}
\step{即 $\begin{cases}|x-1|<2\\ \dfrac{x-2}{x-1}\geqslant0\end{cases}$ 的解集为 $(-1,1)\cup[2,3)$.}
\solend
\fi

\ansspace[5]

\item 比较多项式 $2x^2-3x+1$ 与 $x^2-1$ 的大小.

\ifanswer
\solbegin
\step{作差得 $(2x^2-3x+1)-(x^2-1)=x^2-3x+2=(x-1)(x-2)$，}
\step{当 $x=1$ 或 $x=2$ 时，$(x-1)(x-2)=0$，即 $2x^2-3x+1=x^2-1$，}
\step{当 $x<1$ 或 $x>2$ 时，$(x-1)(x-2)>0$，即 $2x^2-3x+1>x^2-1$，}
\step{当 $1<x<2$ 时，$(x-1)(x-2)<0$，即 $2x^2-3x+1<x^2-1$；}
\step{综上，当 $x=1$ 或 $x=2$ 时，$2x^2-3x+1=x^2-1$；}
\step{当 $x<1$ 或 $x>2$ 时，$2x^2-3x+1>x^2-1$；}
\step{当 $1<x<2$ 时，$2x^2-3x+1<x^2-1$.}
\solend
\fi

\ansspace[6]

% 原卷 (1) 连用两个“若”（“若 a=4，若全集为实数集 R”），照录（见文件头“原卷笔误/存疑”①）
\item 已知集合 $A=\left\{x\mid\dfrac{(a+1)x}{x-1}<a\right\}$，集合 $B=[1,4]$，$a$ 为实常数；

（1）若 $a=4$，若全集为实数集 $\R$，求集合 $\overline{A}$（用区间形式表示）；

（2）若 $A\subseteq B$，求 $a$ 的取值范围.

\ifanswer
\solbegin
\stepm{(1)}{若 $a=4$，集合 $A=\left\{x\mid\dfrac{5x}{x-1}<4\right\}=\left\{x\mid\dfrac{5x}{x-1}-4<0\right\}=\left\{x\mid\dfrac{x+4}{x-1}<0\right\}$}
\step{$=(-4,1)$，因此 $\overline{A}=(-\infty,-4]\cup[1,+\infty)$；}
\stepm{(2)}{当 $-a<1$，即 $a>-1$ 时，$A=(-a,1)$，}
\step{要使 $A\subseteq[1,4]$，需 $-a\geqslant1$，即 $a\leqslant-1$，与 $a>-1$ 矛盾，舍去；}
\step{当 $-a>1$，即 $a<-1$ 时，$A=(1,-a)$，}
\step{要使 $A\subseteq[1,4]$，需 $-a\leqslant4$，即 $a\geqslant-4$，所以 $-4\leqslant a<-1$；}
\step{当 $-a=1$，即 $a=-1$ 时，$A=\varnothing$，满足 $A\subseteq B$；}
\step{综上，实数 $a$ 的取值范围为 $[-4,-1]$.}
\solend
\fi

\ansspace[7]

% 原卷题干“可以售出 8 万本、据市场调查”处用顿号，照录（见文件头“原卷笔误/存疑”②）
\item 某种杂志原以每本 25 元的价格销售，可以售出 8 万本、据市场调查，杂志的单价每提高 1 元，销售量就减少 0.2 万本，假定杂志的成本是每本 10 元（不计其他成本）.

（1）当杂志以每本 30 元定价销售时，求销售该杂志所获利润；

（2）在每本 25 元的价格的基础上提高定价，试确定杂志的定价在什么范围内可使得销售利润不低于 140 万元.

\ifanswer
\solbegin
\stepm{(1)}{当杂志以 30 元定价销售时，}
\step{销售该杂志所获利润为 $[8-(30-25)\times0.2]\times(30-10)=140$ 万元；}
\stepm{(2)}{在每本 25 元的价格的基础上提高定价，设杂志的定价为每本 $x$ 元，}
\step{由题意得 $[8-(x-25)\times0.2]\times(x-10)\geqslant140$，解得 $30\leqslant x\leqslant45$，}
\step{所以杂志的定价在每本不低于 30 元且不高于 45 元的范围内，}
\step{可使得销售利润不低于 140 万元.}
\solend
\fi

\ansspace[7]

% 原卷题干“定义：对于集合 A，则称集合 B=…”句式杂糅，照录（见文件头“原卷笔误/存疑”④）
\item 定义：对于集合 $A$，则称集合 $B=\left\{x\mid x=\dfrac{m+n}{1+mn},m,n\in A,1+mn\neq0\right\}$ 为集合 $A$ 的生成集（其中 $m$、$n$ 可以相等）.

（1）若集合 $A=\{-1,1\}$ 的生成集为 $B$，求集合 $B$；

（2）若集合 $A=\{a,2\}$ 的生成集为 $B$，$B$ 的子集个数为 4 个，求实数 $a$ 的值；

（3）若集合 $A=\{x\mid-1<x<1\}$ 的生成集为 $B$，求证：$B\subseteq A$.

\ifanswer
\solbegin
\stepm{(1)}{因为 $1+mn\neq0$，}
\step{当 $m=n=1$ 时，$x=\dfrac{m+n}{1+mn}=1$；当 $m=n=-1$ 时，$x=\dfrac{m+n}{1+mn}=-1$；}
\step{所以 $B=\{-1,1\}$；}
\stepm{(2)}{$B$ 的子集个数为 4 个，故 $B$ 集合的元素有 2 个，}
\step{当 $m=n=2$ 时，$x=\dfrac{m+n}{1+mn}=\dfrac45$，当 $m=n=a$ 时，$x=\dfrac{m+n}{1+mn}=\dfrac{2a}{1+a^2}$，}
\step{当 $m$ 与 $n$ 分别取 2 和 $a$ 时，$x=\dfrac{a+2}{1+2a}$，}
\stepm{①}{$\dfrac{2a}{1+a^2}=\dfrac45$，解得 $a=2$（舍）或 $a=\dfrac12$，}
\stepm{②}{$\dfrac{a+2}{1+2a}=\dfrac45$，解得 $a=2$（舍），}
\stepm{③}{$\dfrac{2a}{1+a^2}=\dfrac{a+2}{1+2a}$，解得 $a=2$（舍）或 $a=1$ 或 $a=-1$，}
\step{综上所述，$a=1$ 或 $a=-1$ 或 $a=\dfrac12$，即 $a\in\left\{-1,1,\dfrac12\right\}$；}
\stepm{(3)}{$\forall m$、$n\in A$，则 $\dfrac{m+n}{1+mn}+1=\dfrac{(m+1)(n+1)}{1+mn}>0$，故 $\dfrac{m+n}{1+mn}>-1$；}
\step{$\dfrac{m+n}{1+mn}-1=-\dfrac{(m-1)(n-1)}{1+mn}<0$，故 $\dfrac{m+n}{1+mn}<1$；}
\step{故 $B\subseteq\{x\mid-1<x<1\}$，即 $B\subseteq A$.}
\solend
\fi

\ansspace*[10]

\end{enumerate}

\end{document}
"
%
% 原始材料：微信公众号“魔都数学加油站”文章，连载“27学年高一71”；
% 原文标题“27学年高一71——2026-2027年上海市中国中学高一上10月月考”；
% 原文链接：https://mp.weixin.qq.com/s/UZrP53l3GIin040L1EkljA。
% 页面标注发布时间 2026-10-19 22:47（上海）（与 高一69、高一70 页面读出同一 ct 值）。
% 正文为 7 张 PNG 页（`原图/高一71_中国中学_10月月考/01..07.png`），全部 4762×6735，
% 均按页序读取。7 页全为试卷正文页（卷首标题在第 1 页页顶，卷末解答题第 21 题 (3)
% 解析止于第 7 页），无封面页、目录页、二维码或宣传页。粉色斜排“魔都数学加油站”
% 为卷面水印，与公众号页脚一并未录入。
%
% 卷首标题：照卷面印作“2026-2027 年中国中学高一上 10 月月考”（黑体加粗一行）。
% 原卷年份区间用半角连字符，本稿按全稿惯例排 en dash（同 高一68、高一69、高一70）；
% 原卷只有一行标题、无副标题。
%
% 篇幅与题号：原卷自第 3 题起编号（第 1、2 题不在本卷内）。填空题 3—12（10 题）、
% 选择题 13—16（4 题）、解答题 17—21（5 题），共 19 题。本稿沿用原节与原题号，
% 只改排为全稿统一的大节样式（原卷小节标题为“一、填空题”“二、选择题”“三、解答题”）。
%
% 讲评版：原卷每题下直接印【答案】或【解析】。**只印【答案】的 3 题**（第 3、4、5 题），
% 本稿只给 \ans{}；**第 13 题既无【答案】也无【解析】段，答案“B”直接印在题干括号内
% （“(  B  )”）**，按 高一37 第 13 题、高一63 第 14 题、高一67 第 15 题成例把括号留空、
% 答案另提为 \ans{$B$}。其余 15 题全部印【解析】（按全稿口径这类题不另提 \ans{}——
% 第 14—16 题的“故选 D／B／B”写在解析末句，亦照录在解析内）。
% 除第 13 题外，全卷选择题作答括号内均为空白（第 14、15、16 题）。
%
% 副标题：原卷未印副标题。本稿按“知识模块”口径标作“集合与常用逻辑用语、不等式”：
% 19 题中集合与常用逻辑用语 10 题（第 3、4、5、6、9、11、12、13、16、21 题；其中第 4 题
% 属反证法，一并计入常用逻辑用语），不等式 9 题（第 7、8、10、14、15、17、18、19、20 题；
% 第 19 题 (1) 属补集、(2) 属分式不等式，按主干计入不等式），两类均确有分量，故并标。
% 口径同 高一63、高一68、高一69、高一70。
%
% 与原卷的排版差异：
%   ① 原卷每题下直接印【答案】/【解析】，本稿按全稿惯例将答案与解析置于答案版，试卷版
%      保持干净；第 13 题印在题干括号内的答案“B”同样移入 \ans{}，见上；
%   ② 原卷未印副标题，本稿按“知识模块”口径补出；
%   ③ 原卷填空横线约四个字宽，本稿按全稿惯例统一排 \blankline[4.4em]；
%   ④ 原卷数集记号不统一：实数集排斜体 $R$（第 19 题 (1)），自然数集排斜体 $N$
%      （第 5 题题干、第 12 题题干与解析共三处），第 16 题解析又排正体 $\mathrm{N}^*$；
%      本稿按全稿惯例统一改用 $\R$、$\N$、$\N^*$（同 高一68／69／70 排版差异④）；
%      空集记号原卷作 $\varnothing$，照录；
%   ⑤ 原卷句末标点混用：第 15 题题干“（3）$x^2+2+\dfrac1{x^2+2}$**．**其中”、
%      第 17 题题干末“$\dfrac{x-2}{x-1}\geqslant0$**．**”两处放大确认（第 15 题为全角
%      句点、第 17 题为西文句点），其余各题末为西文句点，已逐处照录；题干内冒号——
%      第 5 题“$\{y\mid\cdots\}$**：**”、第 12 题““×”**：**”、第 17 题“不等式组**：**”、
%      第 21 题“定义**：**”——均为全角冒号，照录；数学式内的逗号（如第 12 题
%      “$x\in A,y\in B$”、第 16 题“$A=\{a_1,\cdots\},B=\{b_1,\cdots\},n>2$”）按原卷作
%      半角，中文行文里的逗号按原卷作全角，均未强改；
%   ⑥ 第 17—21 题的子问原卷印作“（1）…（2）…”（另起一段），本稿按 高一69／70 成例用
%      行内“（1）（2）（3）”段，解析中的子问标号一律用 \stepm{(1)}{…}；
%      第 21 题 (2) 解析内层的“①”“②”“③”照原卷作 \stepm{①}{…} 等；
%   ⑦ 每道解答题原卷未留作答空白，本稿按全稿惯例加 \ansspace（末题用 \ansspace*）；
%   ⑧ 选择题的作答括号原卷排半角“(　　)”（第 14、15、16 题括号内空白、第 13 题括号内
%      印 B），本稿按全稿惯例排 \mbox{（\quad）}；选项行排照原卷——第 13 题两行两条
%      （A B 一行、C D 一行），第 14、15、16 题一行四条；
%   ⑨ 第 8 题解析、第 9 题解析、第 17 题题干与解析中的多行不等式组，原卷均排作“{”
%      连接的多行式，本稿按全稿惯例排 cases；
%   ⑩ 第 10 题解析、第 17 题解析、第 19 题 (1) 解析的并集符号原卷排大号 $\bigcup$
%      （显示体），本稿按全稿惯例排行内 $\cup$；
%   ⑪ 第 8 题解析的判别式照原卷作空心“$\triangle$”（与第 7 题解析的“$\Delta$”不同形），
%      照录未改；
%   ⑫ 第 16 题解析第 4 行的“$\forall$”（全称量词）照原卷符号排，未改写为中文“任意”。
%
% 原卷笔误/存疑（均照录未改，正文对应处另有 % 注）：
%   ① 第 19 题 (1) 连用两个“若”——“若 $a=4$，若全集为实数集 $R$”，后一“若”应作
%      “当”或删；原卷如此，照录未改。
%   ② 第 20 题题干“可以售出 8 万本**、**据市场调查”——“万本”后用顿号，按文意应为
%      逗号或句号；原卷如此，照录未改。
%   ③ 第 15 题解析对代数式 (1) 只得出“$y/x+x/y\geqslant2$”（同号时），未讨论取等是否
%      可得——取 $x=y$ 时确能取到 2——却总结“最小值为 2 的代数式有 1 个，故选 $B$”；
%      按解析自身的推理 (1) 也应计入，结论“1 个”与推理矛盾（若计 (1) 应为 2 个、选 C）。
%      原卷如此，照录未改。
%   ④ 第 21 题题干“定义：对于集合 $A$，则称集合 $B=\cdots$”——“对于……”与“则称……”
%      句式杂糅（应作“对于集合 $A$，称……”或“若集合 $A$，则称……”）；原卷如此，照录未改。
%   ⑤ 第 13 题题干末“……的（　）条件”后无任何标点（其余各题题干末均有句点）；
%      原卷如此，照录未改。
%
% 行数核对：逐题数 \step/\stepm 条数与原卷【解析】印刷行数，15 题逐一相符，
% 合计 86 条 / 86 行（Q6：1；Q7：1；Q8：7；Q9：2；Q10：4；Q11：7；Q12：5；Q14：7；
% Q15：8；Q16：7；Q17：3；Q18：7；Q19：8；Q20：6；Q21：13）；另有 3 题只印【答案】
% （第 3、4、5 题）与第 13 题括号内答案，各提 \ans{}，不计入 \step 账。
%
% 插图：无题目配图或数学图形（7 页题面均无“如图/右图/图中”字样；页面另有粉色斜排水印，
% 见上）。
%
\documentclass[a4paper,11pt]{article}
\usepackage{common}

\begin{document}

\examtitlest[3]{2026--2027 年中国中学高一上 10 月月考}{集合与常用逻辑用语、不等式}

\bigsec{一、填空题}

\begin{enumerate}
\setcounter{enumi}{2}

\item 已知集合 $A=[1,+\infty)$，$B=[-1,3)$，则 $A\cap B=$\blankline[4.4em].

\ifanswer
\ans{$[1,3)$}
\fi

\item 用反证法证明命题“已知 $a$、$b$ 为正数，$n$ 为正整数，求证：如果 $a^n>b^n$，那么 $a>b$”时，应当作出的假设为：“假设\blankline[4.4em]”.

\ifanswer
\ans{$a\leqslant b$}
\fi

\item 用列举法表示集合 $\{y\mid y=x+1,x\in\N,x\leqslant1\}$：\blankline[4.4em].

\ifanswer
\ans{$\{1,2\}$}
\fi

\item 某班在一次测验中，有 36 人数学成绩不低于 80 分，有 20 人物理成绩不低于 80 分，有 15 人的数学、物理成绩都不低于 80 分，则这两科成绩中至少有一科不低于 80 分的人数为\blankline[4.4em]人.

\ifanswer
\solbegin
\step{至少有 $36+20-15=41$ 人.}
\solend
\fi

\item 若关于 $x$ 的不等式 $x^2-x+m<0$ 的解集是 $\varnothing$，则实数 $m$ 的取值范围是\blankline[4.4em].

\ifanswer
\solbegin
\step{由题意得 $\Delta=1-4m\leqslant0$，所以 $m\geqslant\dfrac14$.}
\solend
\fi

% 第 8 题解析的判别式原卷作空心“△”（与第 7 题解析的“Δ”不同形），照录（见文件头“排版差异”⑪）
\item 已知关于 $x$ 的方程 $x^2+mx+2m=0$ 的两个实数根分别为 $\alpha$，$\beta$，若 $\alpha^2+\beta^2=5$，则实数 $m=$\blankline[4.4em].

\ifanswer
\solbegin
\step{方程 $x^2+mx+2m=0$ 的两个实数根为 $\alpha$、$\beta$，}
\step{由根与系数的关系得 $\begin{cases}\alpha+\beta=-m\\ \alpha\beta=2m\end{cases}$，又因为 $\alpha^2+\beta^2=5$，}
\step{所以 $\alpha^2+\beta^2=(\alpha+\beta)^2-2\alpha\beta=m^2-4m=5$，整理得 $m^2-4m-5=0$，}
\step{解得 $m=5$ 或 $m=-1$，}
\step{当 $m=5$ 时，$\triangle=25-40=-15<0$，无实根，舍去；}
\step{当 $m=-1$ 时，$\triangle=1+8=9>0$，符合题意，}
\step{所以实数 $m=-1$.}
\solend
\fi

\item 已知集合 $A=\{x\mid6x+a>0\}$，若 $1\in A$，$0\notin A$，则实数 $a$ 的取值范围为\blankline[4.4em].

\ifanswer
\solbegin
\step{由 $A=\{x\mid6x+a>0\}$，$1\in A$，$0\notin A$，得 $\begin{cases}6+a>0\\ a\leqslant0\end{cases}$，解得 $-6<a\leqslant0$，}
\step{即实数 $a$ 的取值范围是 $(-6,0]$.}
\solend
\fi

\item 关于 $x$ 的不等式 $ax-b>0$ 的解集为 $(1,+\infty)$，则关于 $x$ 的不等式 $(ax+b)(x-2)>0$ 的解集为\blankline[4.4em].

\ifanswer
\solbegin
\step{由不等式 $ax-b>0$ 的解集为 $(1,+\infty)$，得 $a>0$ 且 $\dfrac ba=1$，即 $b=a$，}
\step{则 $(ax+b)(x-2)>0$，可化为 $a(x+1)(x-2)>0$.}
\step{即 $(x+1)(x-2)>0$，解得 $x<-1$ 或 $x>2$，}
\step{所以不等式的解集为 $(-\infty,-1)\cup(2,+\infty)$.}
\solend
\fi

\item 已知实数 $a_1$、$a_2$、$a_3$、$a_4$、$a_5$ 满足 $\{x\mid(x-a_1)(x-a_2)(x-a_3)(x-a_4)(x-a_5)=0\}=\{1,2,3,4\}$，则 $a_1+a_2+a_3+a_4+a_5$ 的最大值为\blankline[4.4em].

\ifanswer
\solbegin
\step{5 个实数 $a_1$、$a_2$、$a_3$、$a_4$、$a_5$ 的不同取值只有 1，2，3，4，}
\step{即这 5 个数均从 1，2，3，4 中选取.}
\step{集合有 4 个元素，说明 1，2，3，4 每个数至少出现一次，}
\step{因此恰好有一个数重复出现两次，其余三个数各出现一次.}
\step{要让 5 个数的总和最大，应让重复出现两次的数取集合中的最大值 4，}
\step{此时 5 个数为 1，2，3，4，4，总和为 $1+2+3+4+4=14$，}
\step{即所求最大值为 14.}
\solend
\fi

\item 定义集合运算“$\times$”：$A\times B=\{(x,y)\mid x\in A,y\in B\}$，称为 $A$、$B$ 两个集合的“卡氏积”，若 $A=\{x\mid x^2-|x|\leqslant0,x\in\N\}$，$B=\{1,2,3\}$，则 $(A\times B)\cap(B\times A)=$\blankline[4.4em].

\ifanswer
\solbegin
\step{因为 $x^2-|x|\leqslant0$，所以 $0\leqslant|x|\leqslant1$，又因为 $x\in\N$，}
\step{所以 $A=\{x\mid x^2-|x|\leqslant0,x\in\N\}=\{0,1\}$，又 $B=\{1,2,3\}$，}
\step{由“卡氏积”的定义得 $A\times B=\{(0,1),(0,2),(0,3),(1,1),(1,2),(1,3)\}$，}
\step{$B\times A=\{(1,0),(1,1),(2,0),(2,1),(3,0),(3,1)\}$，}
\step{所以 $(A\times B)\cap(B\times A)=\{(1,1)\}$.}
\solend
\fi

\end{enumerate}

\bigsec{二、选择题}

\begin{enumerate}
\setcounter{enumi}{12}

% 第 13 题答案“B”原卷印在题干括号内，且无【答案】/【解析】段；
% 按 高一37 第 13 题、高一63 第 14 题、高一67 第 15 题成例括号留空、答案另提 \ans{$B$}
% 原卷题干末“条件”后无标点，照录（见文件头“原卷笔误/存疑”⑤）
\item “$A\subseteq B$”是“$A\subset B$”的\mbox{（\quad）}条件

\keepwithprev
A. 充分不必要条件\qquad B. 必要不充分条件

C. 充要条件\qquad D. 既不充分也不必要条件

\ifanswer
\ans{$B$}
\fi

\item 已知 $a$、$b$、$c$ 均为实数，若 $a>b$ 且 $ab\neq0$，则下列不等关系中成立的是\mbox{（\quad）}

\keepwithprev
A. $ac^2>bc^2$\qquad B. $\dfrac1b>\dfrac1a$\qquad C. $|a|>|b|$\qquad D. $c-b>c-a$

\ifanswer
\solbegin
\step{对于 $A$，若 $a>b$，当 $c=0$ 时，有 $ac^2=bc^2=0$，故 $A$ 错误；}
\step{对于 $B$，$\dfrac1b-\dfrac1a=\dfrac{a-b}{ab}$，因为 $a>b$，所以 $a-b>0$，}
\step{当 $a$、$b$ 同号时，$ab>0$，则 $\dfrac1b-\dfrac1a>0$，即 $\dfrac1b>\dfrac1a$，}
\step{当 $a$、$b$ 异号时，$ab<0$，则 $\dfrac1b-\dfrac1a<0$，即 $\dfrac1b<\dfrac1a$，故 $B$ 错误；}
\step{对于 $C$，当 $a=1$，$b=-3$ 时，$|a|<|b|$，故 $C$ 错误；}
\step{对于 $D$，若 $a>b$，则 $-b>-a$，$c-b>c-a$，故 $D$ 正确.}
\step{故选 $D$.}
\solend
\fi

% 原卷解析对 (1) 未讨论取等是否可得（x=y 时确能取 2），却总结“有 1 个、故选 B”，
% 与推理矛盾，照录（见文件头“原卷笔误/存疑”③）
\item 以下三个代数式：（1）$\dfrac yx+\dfrac xy$；（2）$|x|+\dfrac1{|x|}$；（3）$x^2+2+\dfrac1{x^2+2}$．其中的字母均表示实数，且代数式均有意义，那么代数式的最小值为 2 的个数是\mbox{（\quad）}

\keepwithprev
A. $0$\qquad B. $1$\qquad C. $2$\qquad D. $3$

\ifanswer
\solbegin
\step{（1）仅当 $x,y$ 同号时，$\dfrac yx,\dfrac xy>0$，由基本不等式得 $\dfrac yx+\dfrac xy\geqslant2$；}
\step{若 $x,y$ 异号，$\dfrac yx,\dfrac xy<0$，故 $\dfrac yx+\dfrac xy<0$，此时最小值不是 2；}
\step{（2）$|x|+\dfrac1{|x|}$ 中，$|x|>0$，由基本不等式 $|x|+\dfrac1{|x|}\geqslant2\sqrt{|x|\cdot\dfrac1{|x|}}=2$，}
\step{当且仅当 $|x|=1$，即 $x=\pm1$ 时取等号，故最小值为 2；}
\step{（3）令 $t=x^2+2$，则 $t\geqslant2$，函数化为 $t+\dfrac1t$，}
\step{对勾函数 $t+\dfrac1t$ 在 $[2,+\infty)$ 上严格增，最小值为 $2+\dfrac12=\dfrac52>2$，}
\step{等号成立条件 $t=1$ 无法取到，故最小值不是 2；}
\step{综上，最小值为 2 的代数式有 1 个，故选 $B$.}
\solend
\fi

\item 已知 $A$、$B$ 为非空实数集，$\Omega$ 为平面直角坐标系中的一些点构成的集合，集合 $C=\{y\mid$对任意 $x\in A$，有 $(x,y)\in\Omega\}$，集合 $D=\{x\mid$对任意 $y\in B$，有 $(x,y)\in\Omega\}$，对于下列两个命题：①若 $C\subseteq B$，则 $D\subseteq A$；②若 $C\supseteq B$，则 $D\subseteq A$；其中判断正确的是\mbox{（\quad）}

\keepwithprev
A. ①②都正确\qquad B. ①②都错误\qquad C. ①正确②错误\qquad D. ①错误②正确

\ifanswer
\solbegin
\step{设 $A=\{a_1,a_2,a_3,\cdots,a_n\},B=\{b_1,b_2,b_3,\cdots,b_n\},n>2$，}
\step{若 $\Omega=\{(c,b_1),(c,b_2),\cdots,(c,b_n)\},c\notin A$，}
\step{此时 $C=\varnothing\subseteq B,D=\{c\},D\subseteq A$ 错误，故①错误；}
% 原卷此处 N* 排正体（\mathrm{N}*），本稿统一 \N^*（见文件头“排版差异”④）；
% 全称量词“∀”照原卷符号排（见文件头“排版差异”⑫）
\step{若 $C\supseteq B$，则 $\forall(a_i,b_j)(1\leqslant i,j\leqslant n,i,j\in\N^*)\in\Omega$，则 $D\supseteq A$，}
\step{且 $A=\{a_1\},B=\{b_1,b_2\}$，}
\step{若 $\Omega=\{(a_1,b_1),(a_2,b_1),(a_1,b_2),(a_2,b_2),(c,b_1),(c,b_2)\},D$ 真包含 $A$，故②错误.}
\step{故选 $B$.}
\solend
\fi

\end{enumerate}

\bigsec{三、解答题}

\begin{enumerate}
\setcounter{enumi}{16}

\item 解下列关于 $x$ 的不等式组：$\begin{cases}|x-1|<2\\ \dfrac{x-2}{x-1}\geqslant0\end{cases}$.

\ifanswer
\solbegin
\step{解 $|x-1|<2$ 得 $-2<x-1<2$，解得 $-1<x<3$.}
\step{分式不等式 $\dfrac{x-2}{x-1}\geqslant0$ 等价于 $\begin{cases}(x-2)(x-1)\geqslant0\\ x-1\neq0\end{cases}$，解得 $x<1$ 或 $x\geqslant2$.}
\step{即 $\begin{cases}|x-1|<2\\ \dfrac{x-2}{x-1}\geqslant0\end{cases}$ 的解集为 $(-1,1)\cup[2,3)$.}
\solend
\fi

\ansspace[5]

\item 比较多项式 $2x^2-3x+1$ 与 $x^2-1$ 的大小.

\ifanswer
\solbegin
\step{作差得 $(2x^2-3x+1)-(x^2-1)=x^2-3x+2=(x-1)(x-2)$，}
\step{当 $x=1$ 或 $x=2$ 时，$(x-1)(x-2)=0$，即 $2x^2-3x+1=x^2-1$，}
\step{当 $x<1$ 或 $x>2$ 时，$(x-1)(x-2)>0$，即 $2x^2-3x+1>x^2-1$，}
\step{当 $1<x<2$ 时，$(x-1)(x-2)<0$，即 $2x^2-3x+1<x^2-1$；}
\step{综上，当 $x=1$ 或 $x=2$ 时，$2x^2-3x+1=x^2-1$；}
\step{当 $x<1$ 或 $x>2$ 时，$2x^2-3x+1>x^2-1$；}
\step{当 $1<x<2$ 时，$2x^2-3x+1<x^2-1$.}
\solend
\fi

\ansspace[6]

% 原卷 (1) 连用两个“若”（“若 a=4，若全集为实数集 R”），照录（见文件头“原卷笔误/存疑”①）
\item 已知集合 $A=\left\{x\mid\dfrac{(a+1)x}{x-1}<a\right\}$，集合 $B=[1,4]$，$a$ 为实常数；

（1）若 $a=4$，若全集为实数集 $\R$，求集合 $\overline{A}$（用区间形式表示）；

（2）若 $A\subseteq B$，求 $a$ 的取值范围.

\ifanswer
\solbegin
\stepm{(1)}{若 $a=4$，集合 $A=\left\{x\mid\dfrac{5x}{x-1}<4\right\}=\left\{x\mid\dfrac{5x}{x-1}-4<0\right\}=\left\{x\mid\dfrac{x+4}{x-1}<0\right\}$}
\step{$=(-4,1)$，因此 $\overline{A}=(-\infty,-4]\cup[1,+\infty)$；}
\stepm{(2)}{当 $-a<1$，即 $a>-1$ 时，$A=(-a,1)$，}
\step{要使 $A\subseteq[1,4]$，需 $-a\geqslant1$，即 $a\leqslant-1$，与 $a>-1$ 矛盾，舍去；}
\step{当 $-a>1$，即 $a<-1$ 时，$A=(1,-a)$，}
\step{要使 $A\subseteq[1,4]$，需 $-a\leqslant4$，即 $a\geqslant-4$，所以 $-4\leqslant a<-1$；}
\step{当 $-a=1$，即 $a=-1$ 时，$A=\varnothing$，满足 $A\subseteq B$；}
\step{综上，实数 $a$ 的取值范围为 $[-4,-1]$.}
\solend
\fi

\ansspace[7]

% 原卷题干“可以售出 8 万本、据市场调查”处用顿号，照录（见文件头“原卷笔误/存疑”②）
\item 某种杂志原以每本 25 元的价格销售，可以售出 8 万本、据市场调查，杂志的单价每提高 1 元，销售量就减少 0.2 万本，假定杂志的成本是每本 10 元（不计其他成本）.

（1）当杂志以每本 30 元定价销售时，求销售该杂志所获利润；

（2）在每本 25 元的价格的基础上提高定价，试确定杂志的定价在什么范围内可使得销售利润不低于 140 万元.

\ifanswer
\solbegin
\stepm{(1)}{当杂志以 30 元定价销售时，}
\step{销售该杂志所获利润为 $[8-(30-25)\times0.2]\times(30-10)=140$ 万元；}
\stepm{(2)}{在每本 25 元的价格的基础上提高定价，设杂志的定价为每本 $x$ 元，}
\step{由题意得 $[8-(x-25)\times0.2]\times(x-10)\geqslant140$，解得 $30\leqslant x\leqslant45$，}
\step{所以杂志的定价在每本不低于 30 元且不高于 45 元的范围内，}
\step{可使得销售利润不低于 140 万元.}
\solend
\fi

\ansspace[7]

% 原卷题干“定义：对于集合 A，则称集合 B=…”句式杂糅，照录（见文件头“原卷笔误/存疑”④）
\item 定义：对于集合 $A$，则称集合 $B=\left\{x\mid x=\dfrac{m+n}{1+mn},m,n\in A,1+mn\neq0\right\}$ 为集合 $A$ 的生成集（其中 $m$、$n$ 可以相等）.

（1）若集合 $A=\{-1,1\}$ 的生成集为 $B$，求集合 $B$；

（2）若集合 $A=\{a,2\}$ 的生成集为 $B$，$B$ 的子集个数为 4 个，求实数 $a$ 的值；

（3）若集合 $A=\{x\mid-1<x<1\}$ 的生成集为 $B$，求证：$B\subseteq A$.

\ifanswer
\solbegin
\stepm{(1)}{因为 $1+mn\neq0$，}
\step{当 $m=n=1$ 时，$x=\dfrac{m+n}{1+mn}=1$；当 $m=n=-1$ 时，$x=\dfrac{m+n}{1+mn}=-1$；}
\step{所以 $B=\{-1,1\}$；}
\stepm{(2)}{$B$ 的子集个数为 4 个，故 $B$ 集合的元素有 2 个，}
\step{当 $m=n=2$ 时，$x=\dfrac{m+n}{1+mn}=\dfrac45$，当 $m=n=a$ 时，$x=\dfrac{m+n}{1+mn}=\dfrac{2a}{1+a^2}$，}
\step{当 $m$ 与 $n$ 分别取 2 和 $a$ 时，$x=\dfrac{a+2}{1+2a}$，}
\stepm{①}{$\dfrac{2a}{1+a^2}=\dfrac45$，解得 $a=2$（舍）或 $a=\dfrac12$，}
\stepm{②}{$\dfrac{a+2}{1+2a}=\dfrac45$，解得 $a=2$（舍），}
\stepm{③}{$\dfrac{2a}{1+a^2}=\dfrac{a+2}{1+2a}$，解得 $a=2$（舍）或 $a=1$ 或 $a=-1$，}
\step{综上所述，$a=1$ 或 $a=-1$ 或 $a=\dfrac12$，即 $a\in\left\{-1,1,\dfrac12\right\}$；}
\stepm{(3)}{$\forall m$、$n\in A$，则 $\dfrac{m+n}{1+mn}+1=\dfrac{(m+1)(n+1)}{1+mn}>0$，故 $\dfrac{m+n}{1+mn}>-1$；}
\step{$\dfrac{m+n}{1+mn}-1=-\dfrac{(m-1)(n-1)}{1+mn}<0$，故 $\dfrac{m+n}{1+mn}<1$；}
\step{故 $B\subseteq\{x\mid-1<x<1\}$，即 $B\subseteq A$.}
\solend
\fi

\ansspace*[10]

\end{enumerate}

\end{document}
"
%
% 原始材料：微信公众号“魔都数学加油站”文章，连载“27学年高一71”；
% 原文标题“27学年高一71——2026-2027年上海市中国中学高一上10月月考”；
% 原文链接：https://mp.weixin.qq.com/s/UZrP53l3GIin040L1EkljA。
% 页面标注发布时间 2026-10-19 22:47（上海）（与 高一69、高一70 页面读出同一 ct 值）。
% 正文为 7 张 PNG 页（`原图/高一71_中国中学_10月月考/01..07.png`），全部 4762×6735，
% 均按页序读取。7 页全为试卷正文页（卷首标题在第 1 页页顶，卷末解答题第 21 题 (3)
% 解析止于第 7 页），无封面页、目录页、二维码或宣传页。粉色斜排“魔都数学加油站”
% 为卷面水印，与公众号页脚一并未录入。
%
% 卷首标题：照卷面印作“2026-2027 年中国中学高一上 10 月月考”（黑体加粗一行）。
% 原卷年份区间用半角连字符，本稿按全稿惯例排 en dash（同 高一68、高一69、高一70）；
% 原卷只有一行标题、无副标题。
%
% 篇幅与题号：原卷自第 3 题起编号（第 1、2 题不在本卷内）。填空题 3—12（10 题）、
% 选择题 13—16（4 题）、解答题 17—21（5 题），共 19 题。本稿沿用原节与原题号，
% 只改排为全稿统一的大节样式（原卷小节标题为“一、填空题”“二、选择题”“三、解答题”）。
%
% 讲评版：原卷每题下直接印【答案】或【解析】。**只印【答案】的 3 题**（第 3、4、5 题），
% 本稿只给 \ans{}；**第 13 题既无【答案】也无【解析】段，答案“B”直接印在题干括号内
% （“(  B  )”）**，按 高一37 第 13 题、高一63 第 14 题、高一67 第 15 题成例把括号留空、
% 答案另提为 \ans{$B$}。其余 15 题全部印【解析】（按全稿口径这类题不另提 \ans{}——
% 第 14—16 题的“故选 D／B／B”写在解析末句，亦照录在解析内）。
% 除第 13 题外，全卷选择题作答括号内均为空白（第 14、15、16 题）。
%
% 副标题：原卷未印副标题。本稿按“知识模块”口径标作“集合与常用逻辑用语、不等式”：
% 19 题中集合与常用逻辑用语 10 题（第 3、4、5、6、9、11、12、13、16、21 题；其中第 4 题
% 属反证法，一并计入常用逻辑用语），不等式 9 题（第 7、8、10、14、15、17、18、19、20 题；
% 第 19 题 (1) 属补集、(2) 属分式不等式，按主干计入不等式），两类均确有分量，故并标。
% 口径同 高一63、高一68、高一69、高一70。
%
% 与原卷的排版差异：
%   ① 原卷每题下直接印【答案】/【解析】，本稿按全稿惯例将答案与解析置于答案版，试卷版
%      保持干净；第 13 题印在题干括号内的答案“B”同样移入 \ans{}，见上；
%   ② 原卷未印副标题，本稿按“知识模块”口径补出；
%   ③ 原卷填空横线约四个字宽，本稿按全稿惯例统一排 \blankline[4.4em]；
%   ④ 原卷数集记号不统一：实数集排斜体 $R$（第 19 题 (1)），自然数集排斜体 $N$
%      （第 5 题题干、第 12 题题干与解析共三处），第 16 题解析又排正体 $\mathrm{N}^*$；
%      本稿按全稿惯例统一改用 $\R$、$\N$、$\N^*$（同 高一68／69／70 排版差异④）；
%      空集记号原卷作 $\varnothing$，照录；
%   ⑤ 原卷句末标点混用：第 15 题题干“（3）$x^2+2+\dfrac1{x^2+2}$**．**其中”、
%      第 17 题题干末“$\dfrac{x-2}{x-1}\geqslant0$**．**”两处放大确认（第 15 题为全角
%      句点、第 17 题为西文句点），其余各题末为西文句点，已逐处照录；题干内冒号——
%      第 5 题“$\{y\mid\cdots\}$**：**”、第 12 题““×”**：**”、第 17 题“不等式组**：**”、
%      第 21 题“定义**：**”——均为全角冒号，照录；数学式内的逗号（如第 12 题
%      “$x\in A,y\in B$”、第 16 题“$A=\{a_1,\cdots\},B=\{b_1,\cdots\},n>2$”）按原卷作
%      半角，中文行文里的逗号按原卷作全角，均未强改；
%   ⑥ 第 17—21 题的子问原卷印作“（1）…（2）…”（另起一段），本稿按 高一69／70 成例用
%      行内“（1）（2）（3）”段，解析中的子问标号一律用 \stepm{(1)}{…}；
%      第 21 题 (2) 解析内层的“①”“②”“③”照原卷作 \stepm{①}{…} 等；
%   ⑦ 每道解答题原卷未留作答空白，本稿按全稿惯例加 \ansspace（末题用 \ansspace*）；
%   ⑧ 选择题的作答括号原卷排半角“(　　)”（第 14、15、16 题括号内空白、第 13 题括号内
%      印 B），本稿按全稿惯例排 \mbox{（\quad）}；选项行排照原卷——第 13 题两行两条
%      （A B 一行、C D 一行），第 14、15、16 题一行四条；
%   ⑨ 第 8 题解析、第 9 题解析、第 17 题题干与解析中的多行不等式组，原卷均排作“{”
%      连接的多行式，本稿按全稿惯例排 cases；
%   ⑩ 第 10 题解析、第 17 题解析、第 19 题 (1) 解析的并集符号原卷排大号 $\bigcup$
%      （显示体），本稿按全稿惯例排行内 $\cup$；
%   ⑪ 第 8 题解析的判别式照原卷作空心“$\triangle$”（与第 7 题解析的“$\Delta$”不同形），
%      照录未改；
%   ⑫ 第 16 题解析第 4 行的“$\forall$”（全称量词）照原卷符号排，未改写为中文“任意”。
%
% 原卷笔误/存疑（均照录未改，正文对应处另有 % 注）：
%   ① 第 19 题 (1) 连用两个“若”——“若 $a=4$，若全集为实数集 $R$”，后一“若”应作
%      “当”或删；原卷如此，照录未改。
%   ② 第 20 题题干“可以售出 8 万本**、**据市场调查”——“万本”后用顿号，按文意应为
%      逗号或句号；原卷如此，照录未改。
%   ③ 第 15 题解析对代数式 (1) 只得出“$y/x+x/y\geqslant2$”（同号时），未讨论取等是否
%      可得——取 $x=y$ 时确能取到 2——却总结“最小值为 2 的代数式有 1 个，故选 $B$”；
%      按解析自身的推理 (1) 也应计入，结论“1 个”与推理矛盾（若计 (1) 应为 2 个、选 C）。
%      原卷如此，照录未改。
%   ④ 第 21 题题干“定义：对于集合 $A$，则称集合 $B=\cdots$”——“对于……”与“则称……”
%      句式杂糅（应作“对于集合 $A$，称……”或“若集合 $A$，则称……”）；原卷如此，照录未改。
%   ⑤ 第 13 题题干末“……的（　）条件”后无任何标点（其余各题题干末均有句点）；
%      原卷如此，照录未改。
%
% 行数核对：逐题数 \step/\stepm 条数与原卷【解析】印刷行数，15 题逐一相符，
% 合计 86 条 / 86 行（Q6：1；Q7：1；Q8：7；Q9：2；Q10：4；Q11：7；Q12：5；Q14：7；
% Q15：8；Q16：7；Q17：3；Q18：7；Q19：8；Q20：6；Q21：13）；另有 3 题只印【答案】
% （第 3、4、5 题）与第 13 题括号内答案，各提 \ans{}，不计入 \step 账。
%
% 插图：无题目配图或数学图形（7 页题面均无“如图/右图/图中”字样；页面另有粉色斜排水印，
% 见上）。
%
\documentclass[a4paper,11pt]{article}
\usepackage{common}

\begin{document}

\examtitlest[3]{2026--2027 年中国中学高一上 10 月月考}{集合与常用逻辑用语、不等式}

\bigsec{一、填空题}

\begin{enumerate}
\setcounter{enumi}{2}

\item 已知集合 $A=[1,+\infty)$，$B=[-1,3)$，则 $A\cap B=$\blankline[4.4em].

\ifanswer
\ans{$[1,3)$}
\fi

\item 用反证法证明命题“已知 $a$、$b$ 为正数，$n$ 为正整数，求证：如果 $a^n>b^n$，那么 $a>b$”时，应当作出的假设为：“假设\blankline[4.4em]”.

\ifanswer
\ans{$a\leqslant b$}
\fi

\item 用列举法表示集合 $\{y\mid y=x+1,x\in\N,x\leqslant1\}$：\blankline[4.4em].

\ifanswer
\ans{$\{1,2\}$}
\fi

\item 某班在一次测验中，有 36 人数学成绩不低于 80 分，有 20 人物理成绩不低于 80 分，有 15 人的数学、物理成绩都不低于 80 分，则这两科成绩中至少有一科不低于 80 分的人数为\blankline[4.4em]人.

\ifanswer
\solbegin
\step{至少有 $36+20-15=41$ 人.}
\solend
\fi

\item 若关于 $x$ 的不等式 $x^2-x+m<0$ 的解集是 $\varnothing$，则实数 $m$ 的取值范围是\blankline[4.4em].

\ifanswer
\solbegin
\step{由题意得 $\Delta=1-4m\leqslant0$，所以 $m\geqslant\dfrac14$.}
\solend
\fi

% 第 8 题解析的判别式原卷作空心“△”（与第 7 题解析的“Δ”不同形），照录（见文件头“排版差异”⑪）
\item 已知关于 $x$ 的方程 $x^2+mx+2m=0$ 的两个实数根分别为 $\alpha$，$\beta$，若 $\alpha^2+\beta^2=5$，则实数 $m=$\blankline[4.4em].

\ifanswer
\solbegin
\step{方程 $x^2+mx+2m=0$ 的两个实数根为 $\alpha$、$\beta$，}
\step{由根与系数的关系得 $\begin{cases}\alpha+\beta=-m\\ \alpha\beta=2m\end{cases}$，又因为 $\alpha^2+\beta^2=5$，}
\step{所以 $\alpha^2+\beta^2=(\alpha+\beta)^2-2\alpha\beta=m^2-4m=5$，整理得 $m^2-4m-5=0$，}
\step{解得 $m=5$ 或 $m=-1$，}
\step{当 $m=5$ 时，$\triangle=25-40=-15<0$，无实根，舍去；}
\step{当 $m=-1$ 时，$\triangle=1+8=9>0$，符合题意，}
\step{所以实数 $m=-1$.}
\solend
\fi

\item 已知集合 $A=\{x\mid6x+a>0\}$，若 $1\in A$，$0\notin A$，则实数 $a$ 的取值范围为\blankline[4.4em].

\ifanswer
\solbegin
\step{由 $A=\{x\mid6x+a>0\}$，$1\in A$，$0\notin A$，得 $\begin{cases}6+a>0\\ a\leqslant0\end{cases}$，解得 $-6<a\leqslant0$，}
\step{即实数 $a$ 的取值范围是 $(-6,0]$.}
\solend
\fi

\item 关于 $x$ 的不等式 $ax-b>0$ 的解集为 $(1,+\infty)$，则关于 $x$ 的不等式 $(ax+b)(x-2)>0$ 的解集为\blankline[4.4em].

\ifanswer
\solbegin
\step{由不等式 $ax-b>0$ 的解集为 $(1,+\infty)$，得 $a>0$ 且 $\dfrac ba=1$，即 $b=a$，}
\step{则 $(ax+b)(x-2)>0$，可化为 $a(x+1)(x-2)>0$.}
\step{即 $(x+1)(x-2)>0$，解得 $x<-1$ 或 $x>2$，}
\step{所以不等式的解集为 $(-\infty,-1)\cup(2,+\infty)$.}
\solend
\fi

\item 已知实数 $a_1$、$a_2$、$a_3$、$a_4$、$a_5$ 满足 $\{x\mid(x-a_1)(x-a_2)(x-a_3)(x-a_4)(x-a_5)=0\}=\{1,2,3,4\}$，则 $a_1+a_2+a_3+a_4+a_5$ 的最大值为\blankline[4.4em].

\ifanswer
\solbegin
\step{5 个实数 $a_1$、$a_2$、$a_3$、$a_4$、$a_5$ 的不同取值只有 1，2，3，4，}
\step{即这 5 个数均从 1，2，3，4 中选取.}
\step{集合有 4 个元素，说明 1，2，3，4 每个数至少出现一次，}
\step{因此恰好有一个数重复出现两次，其余三个数各出现一次.}
\step{要让 5 个数的总和最大，应让重复出现两次的数取集合中的最大值 4，}
\step{此时 5 个数为 1，2，3，4，4，总和为 $1+2+3+4+4=14$，}
\step{即所求最大值为 14.}
\solend
\fi

\item 定义集合运算“$\times$”：$A\times B=\{(x,y)\mid x\in A,y\in B\}$，称为 $A$、$B$ 两个集合的“卡氏积”，若 $A=\{x\mid x^2-|x|\leqslant0,x\in\N\}$，$B=\{1,2,3\}$，则 $(A\times B)\cap(B\times A)=$\blankline[4.4em].

\ifanswer
\solbegin
\step{因为 $x^2-|x|\leqslant0$，所以 $0\leqslant|x|\leqslant1$，又因为 $x\in\N$，}
\step{所以 $A=\{x\mid x^2-|x|\leqslant0,x\in\N\}=\{0,1\}$，又 $B=\{1,2,3\}$，}
\step{由“卡氏积”的定义得 $A\times B=\{(0,1),(0,2),(0,3),(1,1),(1,2),(1,3)\}$，}
\step{$B\times A=\{(1,0),(1,1),(2,0),(2,1),(3,0),(3,1)\}$，}
\step{所以 $(A\times B)\cap(B\times A)=\{(1,1)\}$.}
\solend
\fi

\end{enumerate}

\bigsec{二、选择题}

\begin{enumerate}
\setcounter{enumi}{12}

% 第 13 题答案“B”原卷印在题干括号内，且无【答案】/【解析】段；
% 按 高一37 第 13 题、高一63 第 14 题、高一67 第 15 题成例括号留空、答案另提 \ans{$B$}
% 原卷题干末“条件”后无标点，照录（见文件头“原卷笔误/存疑”⑤）
\item “$A\subseteq B$”是“$A\subset B$”的\mbox{（\quad）}条件

\keepwithprev
A. 充分不必要条件\qquad B. 必要不充分条件

C. 充要条件\qquad D. 既不充分也不必要条件

\ifanswer
\ans{$B$}
\fi

\item 已知 $a$、$b$、$c$ 均为实数，若 $a>b$ 且 $ab\neq0$，则下列不等关系中成立的是\mbox{（\quad）}

\keepwithprev
A. $ac^2>bc^2$\qquad B. $\dfrac1b>\dfrac1a$\qquad C. $|a|>|b|$\qquad D. $c-b>c-a$

\ifanswer
\solbegin
\step{对于 $A$，若 $a>b$，当 $c=0$ 时，有 $ac^2=bc^2=0$，故 $A$ 错误；}
\step{对于 $B$，$\dfrac1b-\dfrac1a=\dfrac{a-b}{ab}$，因为 $a>b$，所以 $a-b>0$，}
\step{当 $a$、$b$ 同号时，$ab>0$，则 $\dfrac1b-\dfrac1a>0$，即 $\dfrac1b>\dfrac1a$，}
\step{当 $a$、$b$ 异号时，$ab<0$，则 $\dfrac1b-\dfrac1a<0$，即 $\dfrac1b<\dfrac1a$，故 $B$ 错误；}
\step{对于 $C$，当 $a=1$，$b=-3$ 时，$|a|<|b|$，故 $C$ 错误；}
\step{对于 $D$，若 $a>b$，则 $-b>-a$，$c-b>c-a$，故 $D$ 正确.}
\step{故选 $D$.}
\solend
\fi

% 原卷解析对 (1) 未讨论取等是否可得（x=y 时确能取 2），却总结“有 1 个、故选 B”，
% 与推理矛盾，照录（见文件头“原卷笔误/存疑”③）
\item 以下三个代数式：（1）$\dfrac yx+\dfrac xy$；（2）$|x|+\dfrac1{|x|}$；（3）$x^2+2+\dfrac1{x^2+2}$．其中的字母均表示实数，且代数式均有意义，那么代数式的最小值为 2 的个数是\mbox{（\quad）}

\keepwithprev
A. $0$\qquad B. $1$\qquad C. $2$\qquad D. $3$

\ifanswer
\solbegin
\step{（1）仅当 $x,y$ 同号时，$\dfrac yx,\dfrac xy>0$，由基本不等式得 $\dfrac yx+\dfrac xy\geqslant2$；}
\step{若 $x,y$ 异号，$\dfrac yx,\dfrac xy<0$，故 $\dfrac yx+\dfrac xy<0$，此时最小值不是 2；}
\step{（2）$|x|+\dfrac1{|x|}$ 中，$|x|>0$，由基本不等式 $|x|+\dfrac1{|x|}\geqslant2\sqrt{|x|\cdot\dfrac1{|x|}}=2$，}
\step{当且仅当 $|x|=1$，即 $x=\pm1$ 时取等号，故最小值为 2；}
\step{（3）令 $t=x^2+2$，则 $t\geqslant2$，函数化为 $t+\dfrac1t$，}
\step{对勾函数 $t+\dfrac1t$ 在 $[2,+\infty)$ 上严格增，最小值为 $2+\dfrac12=\dfrac52>2$，}
\step{等号成立条件 $t=1$ 无法取到，故最小值不是 2；}
\step{综上，最小值为 2 的代数式有 1 个，故选 $B$.}
\solend
\fi

\item 已知 $A$、$B$ 为非空实数集，$\Omega$ 为平面直角坐标系中的一些点构成的集合，集合 $C=\{y\mid$对任意 $x\in A$，有 $(x,y)\in\Omega\}$，集合 $D=\{x\mid$对任意 $y\in B$，有 $(x,y)\in\Omega\}$，对于下列两个命题：①若 $C\subseteq B$，则 $D\subseteq A$；②若 $C\supseteq B$，则 $D\subseteq A$；其中判断正确的是\mbox{（\quad）}

\keepwithprev
A. ①②都正确\qquad B. ①②都错误\qquad C. ①正确②错误\qquad D. ①错误②正确

\ifanswer
\solbegin
\step{设 $A=\{a_1,a_2,a_3,\cdots,a_n\},B=\{b_1,b_2,b_3,\cdots,b_n\},n>2$，}
\step{若 $\Omega=\{(c,b_1),(c,b_2),\cdots,(c,b_n)\},c\notin A$，}
\step{此时 $C=\varnothing\subseteq B,D=\{c\},D\subseteq A$ 错误，故①错误；}
% 原卷此处 N* 排正体（\mathrm{N}*），本稿统一 \N^*（见文件头“排版差异”④）；
% 全称量词“∀”照原卷符号排（见文件头“排版差异”⑫）
\step{若 $C\supseteq B$，则 $\forall(a_i,b_j)(1\leqslant i,j\leqslant n,i,j\in\N^*)\in\Omega$，则 $D\supseteq A$，}
\step{且 $A=\{a_1\},B=\{b_1,b_2\}$，}
\step{若 $\Omega=\{(a_1,b_1),(a_2,b_1),(a_1,b_2),(a_2,b_2),(c,b_1),(c,b_2)\},D$ 真包含 $A$，故②错误.}
\step{故选 $B$.}
\solend
\fi

\end{enumerate}

\bigsec{三、解答题}

\begin{enumerate}
\setcounter{enumi}{16}

\item 解下列关于 $x$ 的不等式组：$\begin{cases}|x-1|<2\\ \dfrac{x-2}{x-1}\geqslant0\end{cases}$.

\ifanswer
\solbegin
\step{解 $|x-1|<2$ 得 $-2<x-1<2$，解得 $-1<x<3$.}
\step{分式不等式 $\dfrac{x-2}{x-1}\geqslant0$ 等价于 $\begin{cases}(x-2)(x-1)\geqslant0\\ x-1\neq0\end{cases}$，解得 $x<1$ 或 $x\geqslant2$.}
\step{即 $\begin{cases}|x-1|<2\\ \dfrac{x-2}{x-1}\geqslant0\end{cases}$ 的解集为 $(-1,1)\cup[2,3)$.}
\solend
\fi

\ansspace[5]

\item 比较多项式 $2x^2-3x+1$ 与 $x^2-1$ 的大小.

\ifanswer
\solbegin
\step{作差得 $(2x^2-3x+1)-(x^2-1)=x^2-3x+2=(x-1)(x-2)$，}
\step{当 $x=1$ 或 $x=2$ 时，$(x-1)(x-2)=0$，即 $2x^2-3x+1=x^2-1$，}
\step{当 $x<1$ 或 $x>2$ 时，$(x-1)(x-2)>0$，即 $2x^2-3x+1>x^2-1$，}
\step{当 $1<x<2$ 时，$(x-1)(x-2)<0$，即 $2x^2-3x+1<x^2-1$；}
\step{综上，当 $x=1$ 或 $x=2$ 时，$2x^2-3x+1=x^2-1$；}
\step{当 $x<1$ 或 $x>2$ 时，$2x^2-3x+1>x^2-1$；}
\step{当 $1<x<2$ 时，$2x^2-3x+1<x^2-1$.}
\solend
\fi

\ansspace[6]

% 原卷 (1) 连用两个“若”（“若 a=4，若全集为实数集 R”），照录（见文件头“原卷笔误/存疑”①）
\item 已知集合 $A=\left\{x\mid\dfrac{(a+1)x}{x-1}<a\right\}$，集合 $B=[1,4]$，$a$ 为实常数；

（1）若 $a=4$，若全集为实数集 $\R$，求集合 $\overline{A}$（用区间形式表示）；

（2）若 $A\subseteq B$，求 $a$ 的取值范围.

\ifanswer
\solbegin
\stepm{(1)}{若 $a=4$，集合 $A=\left\{x\mid\dfrac{5x}{x-1}<4\right\}=\left\{x\mid\dfrac{5x}{x-1}-4<0\right\}=\left\{x\mid\dfrac{x+4}{x-1}<0\right\}$}
\step{$=(-4,1)$，因此 $\overline{A}=(-\infty,-4]\cup[1,+\infty)$；}
\stepm{(2)}{当 $-a<1$，即 $a>-1$ 时，$A=(-a,1)$，}
\step{要使 $A\subseteq[1,4]$，需 $-a\geqslant1$，即 $a\leqslant-1$，与 $a>-1$ 矛盾，舍去；}
\step{当 $-a>1$，即 $a<-1$ 时，$A=(1,-a)$，}
\step{要使 $A\subseteq[1,4]$，需 $-a\leqslant4$，即 $a\geqslant-4$，所以 $-4\leqslant a<-1$；}
\step{当 $-a=1$，即 $a=-1$ 时，$A=\varnothing$，满足 $A\subseteq B$；}
\step{综上，实数 $a$ 的取值范围为 $[-4,-1]$.}
\solend
\fi

\ansspace[7]

% 原卷题干“可以售出 8 万本、据市场调查”处用顿号，照录（见文件头“原卷笔误/存疑”②）
\item 某种杂志原以每本 25 元的价格销售，可以售出 8 万本、据市场调查，杂志的单价每提高 1 元，销售量就减少 0.2 万本，假定杂志的成本是每本 10 元（不计其他成本）.

（1）当杂志以每本 30 元定价销售时，求销售该杂志所获利润；

（2）在每本 25 元的价格的基础上提高定价，试确定杂志的定价在什么范围内可使得销售利润不低于 140 万元.

\ifanswer
\solbegin
\stepm{(1)}{当杂志以 30 元定价销售时，}
\step{销售该杂志所获利润为 $[8-(30-25)\times0.2]\times(30-10)=140$ 万元；}
\stepm{(2)}{在每本 25 元的价格的基础上提高定价，设杂志的定价为每本 $x$ 元，}
\step{由题意得 $[8-(x-25)\times0.2]\times(x-10)\geqslant140$，解得 $30\leqslant x\leqslant45$，}
\step{所以杂志的定价在每本不低于 30 元且不高于 45 元的范围内，}
\step{可使得销售利润不低于 140 万元.}
\solend
\fi

\ansspace[7]

% 原卷题干“定义：对于集合 A，则称集合 B=…”句式杂糅，照录（见文件头“原卷笔误/存疑”④）
\item 定义：对于集合 $A$，则称集合 $B=\left\{x\mid x=\dfrac{m+n}{1+mn},m,n\in A,1+mn\neq0\right\}$ 为集合 $A$ 的生成集（其中 $m$、$n$ 可以相等）.

（1）若集合 $A=\{-1,1\}$ 的生成集为 $B$，求集合 $B$；

（2）若集合 $A=\{a,2\}$ 的生成集为 $B$，$B$ 的子集个数为 4 个，求实数 $a$ 的值；

（3）若集合 $A=\{x\mid-1<x<1\}$ 的生成集为 $B$，求证：$B\subseteq A$.

\ifanswer
\solbegin
\stepm{(1)}{因为 $1+mn\neq0$，}
\step{当 $m=n=1$ 时，$x=\dfrac{m+n}{1+mn}=1$；当 $m=n=-1$ 时，$x=\dfrac{m+n}{1+mn}=-1$；}
\step{所以 $B=\{-1,1\}$；}
\stepm{(2)}{$B$ 的子集个数为 4 个，故 $B$ 集合的元素有 2 个，}
\step{当 $m=n=2$ 时，$x=\dfrac{m+n}{1+mn}=\dfrac45$，当 $m=n=a$ 时，$x=\dfrac{m+n}{1+mn}=\dfrac{2a}{1+a^2}$，}
\step{当 $m$ 与 $n$ 分别取 2 和 $a$ 时，$x=\dfrac{a+2}{1+2a}$，}
\stepm{①}{$\dfrac{2a}{1+a^2}=\dfrac45$，解得 $a=2$（舍）或 $a=\dfrac12$，}
\stepm{②}{$\dfrac{a+2}{1+2a}=\dfrac45$，解得 $a=2$（舍），}
\stepm{③}{$\dfrac{2a}{1+a^2}=\dfrac{a+2}{1+2a}$，解得 $a=2$（舍）或 $a=1$ 或 $a=-1$，}
\step{综上所述，$a=1$ 或 $a=-1$ 或 $a=\dfrac12$，即 $a\in\left\{-1,1,\dfrac12\right\}$；}
\stepm{(3)}{$\forall m$、$n\in A$，则 $\dfrac{m+n}{1+mn}+1=\dfrac{(m+1)(n+1)}{1+mn}>0$，故 $\dfrac{m+n}{1+mn}>-1$；}
\step{$\dfrac{m+n}{1+mn}-1=-\dfrac{(m-1)(n-1)}{1+mn}<0$，故 $\dfrac{m+n}{1+mn}<1$；}
\step{故 $B\subseteq\{x\mid-1<x<1\}$，即 $B\subseteq A$.}
\solend
\fi

\ansspace*[10]

\end{enumerate}

\end{document}
"
%
% 原始材料：微信公众号“魔都数学加油站”文章，连载“27学年高一71”；
% 原文标题“27学年高一71——2026-2027年上海市中国中学高一上10月月考”；
% 原文链接：https://mp.weixin.qq.com/s/UZrP53l3GIin040L1EkljA。
% 页面标注发布时间 2026-10-19 22:47（上海）（与 高一69、高一70 页面读出同一 ct 值）。
% 正文为 7 张 PNG 页（`原图/高一71_中国中学_10月月考/01..07.png`），全部 4762×6735，
% 均按页序读取。7 页全为试卷正文页（卷首标题在第 1 页页顶，卷末解答题第 21 题 (3)
% 解析止于第 7 页），无封面页、目录页、二维码或宣传页。粉色斜排“魔都数学加油站”
% 为卷面水印，与公众号页脚一并未录入。
%
% 卷首标题：照卷面印作“2026-2027 年中国中学高一上 10 月月考”（黑体加粗一行）。
% 原卷年份区间用半角连字符，本稿按全稿惯例排 en dash（同 高一68、高一69、高一70）；
% 原卷只有一行标题、无副标题。
%
% 篇幅与题号：原卷自第 3 题起编号（第 1、2 题不在本卷内）。填空题 3—12（10 题）、
% 选择题 13—16（4 题）、解答题 17—21（5 题），共 19 题。本稿沿用原节与原题号，
% 只改排为全稿统一的大节样式（原卷小节标题为“一、填空题”“二、选择题”“三、解答题”）。
%
% 讲评版：原卷每题下直接印【答案】或【解析】。**只印【答案】的 3 题**（第 3、4、5 题），
% 本稿只给 \ans{}；**第 13 题既无【答案】也无【解析】段，答案“B”直接印在题干括号内
% （“(  B  )”）**，按 高一37 第 13 题、高一63 第 14 题、高一67 第 15 题成例把括号留空、
% 答案另提为 \ans{$B$}。其余 15 题全部印【解析】（按全稿口径这类题不另提 \ans{}——
% 第 14—16 题的“故选 D／B／B”写在解析末句，亦照录在解析内）。
% 除第 13 题外，全卷选择题作答括号内均为空白（第 14、15、16 题）。
%
% 副标题：原卷未印副标题。本稿按“知识模块”口径标作“集合与常用逻辑用语、不等式”：
% 19 题中集合与常用逻辑用语 10 题（第 3、4、5、6、9、11、12、13、16、21 题；其中第 4 题
% 属反证法，一并计入常用逻辑用语），不等式 9 题（第 7、8、10、14、15、17、18、19、20 题；
% 第 19 题 (1) 属补集、(2) 属分式不等式，按主干计入不等式），两类均确有分量，故并标。
% 口径同 高一63、高一68、高一69、高一70。
%
% 与原卷的排版差异：
%   ① 原卷每题下直接印【答案】/【解析】，本稿按全稿惯例将答案与解析置于答案版，试卷版
%      保持干净；第 13 题印在题干括号内的答案“B”同样移入 \ans{}，见上；
%   ② 原卷未印副标题，本稿按“知识模块”口径补出；
%   ③ 原卷填空横线约四个字宽，本稿按全稿惯例统一排 \blankline[4.4em]；
%   ④ 原卷数集记号不统一：实数集排斜体 $R$（第 19 题 (1)），自然数集排斜体 $N$
%      （第 5 题题干、第 12 题题干与解析共三处），第 16 题解析又排正体 $\mathrm{N}^*$；
%      本稿按全稿惯例统一改用 $\R$、$\N$、$\N^*$（同 高一68／69／70 排版差异④）；
%      空集记号原卷作 $\varnothing$，照录；
%   ⑤ 原卷句末标点混用：第 15 题题干“（3）$x^2+2+\dfrac1{x^2+2}$**．**其中”、
%      第 17 题题干末“$\dfrac{x-2}{x-1}\geqslant0$**．**”两处放大确认（第 15 题为全角
%      句点、第 17 题为西文句点），其余各题末为西文句点，已逐处照录；题干内冒号——
%      第 5 题“$\{y\mid\cdots\}$**：**”、第 12 题““×”**：**”、第 17 题“不等式组**：**”、
%      第 21 题“定义**：**”——均为全角冒号，照录；数学式内的逗号（如第 12 题
%      “$x\in A,y\in B$”、第 16 题“$A=\{a_1,\cdots\},B=\{b_1,\cdots\},n>2$”）按原卷作
%      半角，中文行文里的逗号按原卷作全角，均未强改；
%   ⑥ 第 17—21 题的子问原卷印作“（1）…（2）…”（另起一段），本稿按 高一69／70 成例用
%      行内“（1）（2）（3）”段，解析中的子问标号一律用 \stepm{(1)}{…}；
%      第 21 题 (2) 解析内层的“①”“②”“③”照原卷作 \stepm{①}{…} 等；
%   ⑦ 每道解答题原卷未留作答空白，本稿按全稿惯例加 \ansspace（末题用 \ansspace*）；
%   ⑧ 选择题的作答括号原卷排半角“(　　)”（第 14、15、16 题括号内空白、第 13 题括号内
%      印 B），本稿按全稿惯例排 \mbox{（\quad）}；选项行排照原卷——第 13 题两行两条
%      （A B 一行、C D 一行），第 14、15、16 题一行四条；
%   ⑨ 第 8 题解析、第 9 题解析、第 17 题题干与解析中的多行不等式组，原卷均排作“{”
%      连接的多行式，本稿按全稿惯例排 cases；
%   ⑩ 第 10 题解析、第 17 题解析、第 19 题 (1) 解析的并集符号原卷排大号 $\bigcup$
%      （显示体），本稿按全稿惯例排行内 $\cup$；
%   ⑪ 第 8 题解析的判别式照原卷作空心“$\triangle$”（与第 7 题解析的“$\Delta$”不同形），
%      照录未改；
%   ⑫ 第 16 题解析第 4 行的“$\forall$”（全称量词）照原卷符号排，未改写为中文“任意”。
%
% 原卷笔误/存疑（均照录未改，正文对应处另有 % 注）：
%   ① 第 19 题 (1) 连用两个“若”——“若 $a=4$，若全集为实数集 $R$”，后一“若”应作
%      “当”或删；原卷如此，照录未改。
%   ② 第 20 题题干“可以售出 8 万本**、**据市场调查”——“万本”后用顿号，按文意应为
%      逗号或句号；原卷如此，照录未改。
%   ③ 第 15 题解析对代数式 (1) 只得出“$y/x+x/y\geqslant2$”（同号时），未讨论取等是否
%      可得——取 $x=y$ 时确能取到 2——却总结“最小值为 2 的代数式有 1 个，故选 $B$”；
%      按解析自身的推理 (1) 也应计入，结论“1 个”与推理矛盾（若计 (1) 应为 2 个、选 C）。
%      原卷如此，照录未改。
%   ④ 第 21 题题干“定义：对于集合 $A$，则称集合 $B=\cdots$”——“对于……”与“则称……”
%      句式杂糅（应作“对于集合 $A$，称……”或“若集合 $A$，则称……”）；原卷如此，照录未改。
%   ⑤ 第 13 题题干末“……的（　）条件”后无任何标点（其余各题题干末均有句点）；
%      原卷如此，照录未改。
%
% 行数核对：逐题数 \step/\stepm 条数与原卷【解析】印刷行数，15 题逐一相符，
% 合计 86 条 / 86 行（Q6：1；Q7：1；Q8：7；Q9：2；Q10：4；Q11：7；Q12：5；Q14：7；
% Q15：8；Q16：7；Q17：3；Q18：7；Q19：8；Q20：6；Q21：13）；另有 3 题只印【答案】
% （第 3、4、5 题）与第 13 题括号内答案，各提 \ans{}，不计入 \step 账。
%
% 插图：无题目配图或数学图形（7 页题面均无“如图/右图/图中”字样；页面另有粉色斜排水印，
% 见上）。
%
\documentclass[a4paper,11pt]{article}
\usepackage{common}

\begin{document}

\examtitlest[3]{2026--2027 年中国中学高一上 10 月月考}{集合与常用逻辑用语、不等式}

\bigsec{一、填空题}

\begin{enumerate}
\setcounter{enumi}{2}

\item 已知集合 $A=[1,+\infty)$，$B=[-1,3)$，则 $A\cap B=$\blankline[4.4em].

\ifanswer
\ans{$[1,3)$}
\fi

\item 用反证法证明命题“已知 $a$、$b$ 为正数，$n$ 为正整数，求证：如果 $a^n>b^n$，那么 $a>b$”时，应当作出的假设为：“假设\blankline[4.4em]”.

\ifanswer
\ans{$a\leqslant b$}
\fi

\item 用列举法表示集合 $\{y\mid y=x+1,x\in\N,x\leqslant1\}$：\blankline[4.4em].

\ifanswer
\ans{$\{1,2\}$}
\fi

\item 某班在一次测验中，有 36 人数学成绩不低于 80 分，有 20 人物理成绩不低于 80 分，有 15 人的数学、物理成绩都不低于 80 分，则这两科成绩中至少有一科不低于 80 分的人数为\blankline[4.4em]人.

\ifanswer
\solbegin
\step{至少有 $36+20-15=41$ 人.}
\solend
\fi

\item 若关于 $x$ 的不等式 $x^2-x+m<0$ 的解集是 $\varnothing$，则实数 $m$ 的取值范围是\blankline[4.4em].

\ifanswer
\solbegin
\step{由题意得 $\Delta=1-4m\leqslant0$，所以 $m\geqslant\dfrac14$.}
\solend
\fi

% 第 8 题解析的判别式原卷作空心“△”（与第 7 题解析的“Δ”不同形），照录（见文件头“排版差异”⑪）
\item 已知关于 $x$ 的方程 $x^2+mx+2m=0$ 的两个实数根分别为 $\alpha$，$\beta$，若 $\alpha^2+\beta^2=5$，则实数 $m=$\blankline[4.4em].

\ifanswer
\solbegin
\step{方程 $x^2+mx+2m=0$ 的两个实数根为 $\alpha$、$\beta$，}
\step{由根与系数的关系得 $\begin{cases}\alpha+\beta=-m\\ \alpha\beta=2m\end{cases}$，又因为 $\alpha^2+\beta^2=5$，}
\step{所以 $\alpha^2+\beta^2=(\alpha+\beta)^2-2\alpha\beta=m^2-4m=5$，整理得 $m^2-4m-5=0$，}
\step{解得 $m=5$ 或 $m=-1$，}
\step{当 $m=5$ 时，$\triangle=25-40=-15<0$，无实根，舍去；}
\step{当 $m=-1$ 时，$\triangle=1+8=9>0$，符合题意，}
\step{所以实数 $m=-1$.}
\solend
\fi

\item 已知集合 $A=\{x\mid6x+a>0\}$，若 $1\in A$，$0\notin A$，则实数 $a$ 的取值范围为\blankline[4.4em].

\ifanswer
\solbegin
\step{由 $A=\{x\mid6x+a>0\}$，$1\in A$，$0\notin A$，得 $\begin{cases}6+a>0\\ a\leqslant0\end{cases}$，解得 $-6<a\leqslant0$，}
\step{即实数 $a$ 的取值范围是 $(-6,0]$.}
\solend
\fi

\item 关于 $x$ 的不等式 $ax-b>0$ 的解集为 $(1,+\infty)$，则关于 $x$ 的不等式 $(ax+b)(x-2)>0$ 的解集为\blankline[4.4em].

\ifanswer
\solbegin
\step{由不等式 $ax-b>0$ 的解集为 $(1,+\infty)$，得 $a>0$ 且 $\dfrac ba=1$，即 $b=a$，}
\step{则 $(ax+b)(x-2)>0$，可化为 $a(x+1)(x-2)>0$.}
\step{即 $(x+1)(x-2)>0$，解得 $x<-1$ 或 $x>2$，}
\step{所以不等式的解集为 $(-\infty,-1)\cup(2,+\infty)$.}
\solend
\fi

\item 已知实数 $a_1$、$a_2$、$a_3$、$a_4$、$a_5$ 满足 $\{x\mid(x-a_1)(x-a_2)(x-a_3)(x-a_4)(x-a_5)=0\}=\{1,2,3,4\}$，则 $a_1+a_2+a_3+a_4+a_5$ 的最大值为\blankline[4.4em].

\ifanswer
\solbegin
\step{5 个实数 $a_1$、$a_2$、$a_3$、$a_4$、$a_5$ 的不同取值只有 1，2，3，4，}
\step{即这 5 个数均从 1，2，3，4 中选取.}
\step{集合有 4 个元素，说明 1，2，3，4 每个数至少出现一次，}
\step{因此恰好有一个数重复出现两次，其余三个数各出现一次.}
\step{要让 5 个数的总和最大，应让重复出现两次的数取集合中的最大值 4，}
\step{此时 5 个数为 1，2，3，4，4，总和为 $1+2+3+4+4=14$，}
\step{即所求最大值为 14.}
\solend
\fi

\item 定义集合运算“$\times$”：$A\times B=\{(x,y)\mid x\in A,y\in B\}$，称为 $A$、$B$ 两个集合的“卡氏积”，若 $A=\{x\mid x^2-|x|\leqslant0,x\in\N\}$，$B=\{1,2,3\}$，则 $(A\times B)\cap(B\times A)=$\blankline[4.4em].

\ifanswer
\solbegin
\step{因为 $x^2-|x|\leqslant0$，所以 $0\leqslant|x|\leqslant1$，又因为 $x\in\N$，}
\step{所以 $A=\{x\mid x^2-|x|\leqslant0,x\in\N\}=\{0,1\}$，又 $B=\{1,2,3\}$，}
\step{由“卡氏积”的定义得 $A\times B=\{(0,1),(0,2),(0,3),(1,1),(1,2),(1,3)\}$，}
\step{$B\times A=\{(1,0),(1,1),(2,0),(2,1),(3,0),(3,1)\}$，}
\step{所以 $(A\times B)\cap(B\times A)=\{(1,1)\}$.}
\solend
\fi

\end{enumerate}

\bigsec{二、选择题}

\begin{enumerate}
\setcounter{enumi}{12}

% 第 13 题答案“B”原卷印在题干括号内，且无【答案】/【解析】段；
% 按 高一37 第 13 题、高一63 第 14 题、高一67 第 15 题成例括号留空、答案另提 \ans{$B$}
% 原卷题干末“条件”后无标点，照录（见文件头“原卷笔误/存疑”⑤）
\item “$A\subseteq B$”是“$A\subset B$”的\mbox{（\quad）}条件

\keepwithprev
A. 充分不必要条件\qquad B. 必要不充分条件

C. 充要条件\qquad D. 既不充分也不必要条件

\ifanswer
\ans{$B$}
\fi

\item 已知 $a$、$b$、$c$ 均为实数，若 $a>b$ 且 $ab\neq0$，则下列不等关系中成立的是\mbox{（\quad）}

\keepwithprev
A. $ac^2>bc^2$\qquad B. $\dfrac1b>\dfrac1a$\qquad C. $|a|>|b|$\qquad D. $c-b>c-a$

\ifanswer
\solbegin
\step{对于 $A$，若 $a>b$，当 $c=0$ 时，有 $ac^2=bc^2=0$，故 $A$ 错误；}
\step{对于 $B$，$\dfrac1b-\dfrac1a=\dfrac{a-b}{ab}$，因为 $a>b$，所以 $a-b>0$，}
\step{当 $a$、$b$ 同号时，$ab>0$，则 $\dfrac1b-\dfrac1a>0$，即 $\dfrac1b>\dfrac1a$，}
\step{当 $a$、$b$ 异号时，$ab<0$，则 $\dfrac1b-\dfrac1a<0$，即 $\dfrac1b<\dfrac1a$，故 $B$ 错误；}
\step{对于 $C$，当 $a=1$，$b=-3$ 时，$|a|<|b|$，故 $C$ 错误；}
\step{对于 $D$，若 $a>b$，则 $-b>-a$，$c-b>c-a$，故 $D$ 正确.}
\step{故选 $D$.}
\solend
\fi

% 原卷解析对 (1) 未讨论取等是否可得（x=y 时确能取 2），却总结“有 1 个、故选 B”，
% 与推理矛盾，照录（见文件头“原卷笔误/存疑”③）
\item 以下三个代数式：（1）$\dfrac yx+\dfrac xy$；（2）$|x|+\dfrac1{|x|}$；（3）$x^2+2+\dfrac1{x^2+2}$．其中的字母均表示实数，且代数式均有意义，那么代数式的最小值为 2 的个数是\mbox{（\quad）}

\keepwithprev
A. $0$\qquad B. $1$\qquad C. $2$\qquad D. $3$

\ifanswer
\solbegin
\step{（1）仅当 $x,y$ 同号时，$\dfrac yx,\dfrac xy>0$，由基本不等式得 $\dfrac yx+\dfrac xy\geqslant2$；}
\step{若 $x,y$ 异号，$\dfrac yx,\dfrac xy<0$，故 $\dfrac yx+\dfrac xy<0$，此时最小值不是 2；}
\step{（2）$|x|+\dfrac1{|x|}$ 中，$|x|>0$，由基本不等式 $|x|+\dfrac1{|x|}\geqslant2\sqrt{|x|\cdot\dfrac1{|x|}}=2$，}
\step{当且仅当 $|x|=1$，即 $x=\pm1$ 时取等号，故最小值为 2；}
\step{（3）令 $t=x^2+2$，则 $t\geqslant2$，函数化为 $t+\dfrac1t$，}
\step{对勾函数 $t+\dfrac1t$ 在 $[2,+\infty)$ 上严格增，最小值为 $2+\dfrac12=\dfrac52>2$，}
\step{等号成立条件 $t=1$ 无法取到，故最小值不是 2；}
\step{综上，最小值为 2 的代数式有 1 个，故选 $B$.}
\solend
\fi

\item 已知 $A$、$B$ 为非空实数集，$\Omega$ 为平面直角坐标系中的一些点构成的集合，集合 $C=\{y\mid$对任意 $x\in A$，有 $(x,y)\in\Omega\}$，集合 $D=\{x\mid$对任意 $y\in B$，有 $(x,y)\in\Omega\}$，对于下列两个命题：①若 $C\subseteq B$，则 $D\subseteq A$；②若 $C\supseteq B$，则 $D\subseteq A$；其中判断正确的是\mbox{（\quad）}

\keepwithprev
A. ①②都正确\qquad B. ①②都错误\qquad C. ①正确②错误\qquad D. ①错误②正确

\ifanswer
\solbegin
\step{设 $A=\{a_1,a_2,a_3,\cdots,a_n\},B=\{b_1,b_2,b_3,\cdots,b_n\},n>2$，}
\step{若 $\Omega=\{(c,b_1),(c,b_2),\cdots,(c,b_n)\},c\notin A$，}
\step{此时 $C=\varnothing\subseteq B,D=\{c\},D\subseteq A$ 错误，故①错误；}
% 原卷此处 N* 排正体（\mathrm{N}*），本稿统一 \N^*（见文件头“排版差异”④）；
% 全称量词“∀”照原卷符号排（见文件头“排版差异”⑫）
\step{若 $C\supseteq B$，则 $\forall(a_i,b_j)(1\leqslant i,j\leqslant n,i,j\in\N^*)\in\Omega$，则 $D\supseteq A$，}
\step{且 $A=\{a_1\},B=\{b_1,b_2\}$，}
\step{若 $\Omega=\{(a_1,b_1),(a_2,b_1),(a_1,b_2),(a_2,b_2),(c,b_1),(c,b_2)\},D$ 真包含 $A$，故②错误.}
\step{故选 $B$.}
\solend
\fi

\end{enumerate}

\bigsec{三、解答题}

\begin{enumerate}
\setcounter{enumi}{16}

\item 解下列关于 $x$ 的不等式组：$\begin{cases}|x-1|<2\\ \dfrac{x-2}{x-1}\geqslant0\end{cases}$.

\ifanswer
\solbegin
\step{解 $|x-1|<2$ 得 $-2<x-1<2$，解得 $-1<x<3$.}
\step{分式不等式 $\dfrac{x-2}{x-1}\geqslant0$ 等价于 $\begin{cases}(x-2)(x-1)\geqslant0\\ x-1\neq0\end{cases}$，解得 $x<1$ 或 $x\geqslant2$.}
\step{即 $\begin{cases}|x-1|<2\\ \dfrac{x-2}{x-1}\geqslant0\end{cases}$ 的解集为 $(-1,1)\cup[2,3)$.}
\solend
\fi

\ansspace[5]

\item 比较多项式 $2x^2-3x+1$ 与 $x^2-1$ 的大小.

\ifanswer
\solbegin
\step{作差得 $(2x^2-3x+1)-(x^2-1)=x^2-3x+2=(x-1)(x-2)$，}
\step{当 $x=1$ 或 $x=2$ 时，$(x-1)(x-2)=0$，即 $2x^2-3x+1=x^2-1$，}
\step{当 $x<1$ 或 $x>2$ 时，$(x-1)(x-2)>0$，即 $2x^2-3x+1>x^2-1$，}
\step{当 $1<x<2$ 时，$(x-1)(x-2)<0$，即 $2x^2-3x+1<x^2-1$；}
\step{综上，当 $x=1$ 或 $x=2$ 时，$2x^2-3x+1=x^2-1$；}
\step{当 $x<1$ 或 $x>2$ 时，$2x^2-3x+1>x^2-1$；}
\step{当 $1<x<2$ 时，$2x^2-3x+1<x^2-1$.}
\solend
\fi

\ansspace[6]

% 原卷 (1) 连用两个“若”（“若 a=4，若全集为实数集 R”），照录（见文件头“原卷笔误/存疑”①）
\item 已知集合 $A=\left\{x\mid\dfrac{(a+1)x}{x-1}<a\right\}$，集合 $B=[1,4]$，$a$ 为实常数；

（1）若 $a=4$，若全集为实数集 $\R$，求集合 $\overline{A}$（用区间形式表示）；

（2）若 $A\subseteq B$，求 $a$ 的取值范围.

\ifanswer
\solbegin
\stepm{(1)}{若 $a=4$，集合 $A=\left\{x\mid\dfrac{5x}{x-1}<4\right\}=\left\{x\mid\dfrac{5x}{x-1}-4<0\right\}=\left\{x\mid\dfrac{x+4}{x-1}<0\right\}$}
\step{$=(-4,1)$，因此 $\overline{A}=(-\infty,-4]\cup[1,+\infty)$；}
\stepm{(2)}{当 $-a<1$，即 $a>-1$ 时，$A=(-a,1)$，}
\step{要使 $A\subseteq[1,4]$，需 $-a\geqslant1$，即 $a\leqslant-1$，与 $a>-1$ 矛盾，舍去；}
\step{当 $-a>1$，即 $a<-1$ 时，$A=(1,-a)$，}
\step{要使 $A\subseteq[1,4]$，需 $-a\leqslant4$，即 $a\geqslant-4$，所以 $-4\leqslant a<-1$；}
\step{当 $-a=1$，即 $a=-1$ 时，$A=\varnothing$，满足 $A\subseteq B$；}
\step{综上，实数 $a$ 的取值范围为 $[-4,-1]$.}
\solend
\fi

\ansspace[7]

% 原卷题干“可以售出 8 万本、据市场调查”处用顿号，照录（见文件头“原卷笔误/存疑”②）
\item 某种杂志原以每本 25 元的价格销售，可以售出 8 万本、据市场调查，杂志的单价每提高 1 元，销售量就减少 0.2 万本，假定杂志的成本是每本 10 元（不计其他成本）.

（1）当杂志以每本 30 元定价销售时，求销售该杂志所获利润；

（2）在每本 25 元的价格的基础上提高定价，试确定杂志的定价在什么范围内可使得销售利润不低于 140 万元.

\ifanswer
\solbegin
\stepm{(1)}{当杂志以 30 元定价销售时，}
\step{销售该杂志所获利润为 $[8-(30-25)\times0.2]\times(30-10)=140$ 万元；}
\stepm{(2)}{在每本 25 元的价格的基础上提高定价，设杂志的定价为每本 $x$ 元，}
\step{由题意得 $[8-(x-25)\times0.2]\times(x-10)\geqslant140$，解得 $30\leqslant x\leqslant45$，}
\step{所以杂志的定价在每本不低于 30 元且不高于 45 元的范围内，}
\step{可使得销售利润不低于 140 万元.}
\solend
\fi

\ansspace[7]

% 原卷题干“定义：对于集合 A，则称集合 B=…”句式杂糅，照录（见文件头“原卷笔误/存疑”④）
\item 定义：对于集合 $A$，则称集合 $B=\left\{x\mid x=\dfrac{m+n}{1+mn},m,n\in A,1+mn\neq0\right\}$ 为集合 $A$ 的生成集（其中 $m$、$n$ 可以相等）.

（1）若集合 $A=\{-1,1\}$ 的生成集为 $B$，求集合 $B$；

（2）若集合 $A=\{a,2\}$ 的生成集为 $B$，$B$ 的子集个数为 4 个，求实数 $a$ 的值；

（3）若集合 $A=\{x\mid-1<x<1\}$ 的生成集为 $B$，求证：$B\subseteq A$.

\ifanswer
\solbegin
\stepm{(1)}{因为 $1+mn\neq0$，}
\step{当 $m=n=1$ 时，$x=\dfrac{m+n}{1+mn}=1$；当 $m=n=-1$ 时，$x=\dfrac{m+n}{1+mn}=-1$；}
\step{所以 $B=\{-1,1\}$；}
\stepm{(2)}{$B$ 的子集个数为 4 个，故 $B$ 集合的元素有 2 个，}
\step{当 $m=n=2$ 时，$x=\dfrac{m+n}{1+mn}=\dfrac45$，当 $m=n=a$ 时，$x=\dfrac{m+n}{1+mn}=\dfrac{2a}{1+a^2}$，}
\step{当 $m$ 与 $n$ 分别取 2 和 $a$ 时，$x=\dfrac{a+2}{1+2a}$，}
\stepm{①}{$\dfrac{2a}{1+a^2}=\dfrac45$，解得 $a=2$（舍）或 $a=\dfrac12$，}
\stepm{②}{$\dfrac{a+2}{1+2a}=\dfrac45$，解得 $a=2$（舍），}
\stepm{③}{$\dfrac{2a}{1+a^2}=\dfrac{a+2}{1+2a}$，解得 $a=2$（舍）或 $a=1$ 或 $a=-1$，}
\step{综上所述，$a=1$ 或 $a=-1$ 或 $a=\dfrac12$，即 $a\in\left\{-1,1,\dfrac12\right\}$；}
\stepm{(3)}{$\forall m$、$n\in A$，则 $\dfrac{m+n}{1+mn}+1=\dfrac{(m+1)(n+1)}{1+mn}>0$，故 $\dfrac{m+n}{1+mn}>-1$；}
\step{$\dfrac{m+n}{1+mn}-1=-\dfrac{(m-1)(n-1)}{1+mn}<0$，故 $\dfrac{m+n}{1+mn}<1$；}
\step{故 $B\subseteq\{x\mid-1<x<1\}$，即 $B\subseteq A$.}
\solend
\fi

\ansspace*[10]

\end{enumerate}

\end{document}
